\documentclass[11pt,a4paper]{article}

\usepackage[margin=0.82in]{geometry}
\usepackage{fontspec}
\usepackage{xeCJK}
\xeCJKsetup{CJKspace=true}
\setmainfont{DejaVu Serif}
\setsansfont{DejaVu Sans}
\setmonofont{DejaVu Sans Mono}
\setCJKmainfont{Noto Sans CJK KR}
\setCJKsansfont{Noto Sans CJK KR}
\setCJKmonofont{Noto Sans Mono CJK KR}
\usepackage{amsmath,amssymb,mathtools,bm}
\usepackage{booktabs}
\usepackage{longtable}
\usepackage{tabularx}
\usepackage{array}
\usepackage{enumitem}
\usepackage{xcolor}
\usepackage{hyperref}
\usepackage{listings}
\usepackage{titlesec}
\usepackage{fancyhdr}
\usepackage{microtype}
\usepackage{ragged2e}
\usepackage{graphicx}
\usepackage{float}
\usepackage{setspace}
\usepackage{multirow}

\hypersetup{
  colorlinks=true,
  linkcolor=blue!55!black,
  urlcolor=blue!55!black,
  citecolor=blue!55!black,
  pdftitle={Topology-Distilled, Error-Triggered Global-Local Wind Dynamics for Static 3D Gaussian Thin Surfaces},
  pdfauthor={Wind3DGS idea sketch}
}

\definecolor{notegray}{RGB}{245,246,248}
\definecolor{darkblue}{RGB}{30,58,138}
\definecolor{darkgreen}{RGB}{22,101,52}
\definecolor{darkred}{RGB}{153,27,27}
\definecolor{darkorange}{RGB}{154,52,18}
\definecolor{darkpurple}{RGB}{107,33,168}
\definecolor{softblue}{RGB}{239,246,255}
\definecolor{softgreen}{RGB}{240,253,244}
\definecolor{softred}{RGB}{254,242,242}

\lstset{
  basicstyle=\ttfamily\small,
  breaklines=true,
  columns=fullflexible,
  frame=single,
  backgroundcolor=\color{notegray},
  xleftmargin=0.5em,
  xrightmargin=0.5em,
  framerule=0pt
}

\setlength{\headheight}{14pt}
\setlength{\tabcolsep}{4pt}
\emergencystretch=3em
\allowdisplaybreaks
\pagestyle{fancy}
\fancyhf{}
\lhead{Topology-Distilled Error-Triggered Wind Dynamics}
\rhead{2026-08-13}
\cfoot{\thepage}

\titleformat{\section}{\Large\bfseries\color{darkblue}}{\thesection}{0.8em}{}
\titleformat{\subsection}{\large\bfseries}{\thesubsection}{0.8em}{}
\titleformat{\subsubsection}{\normalsize\bfseries}{\thesubsubsection}{0.8em}{}

\makeatletter
\renewcommand{\@pnumwidth}{2.4em}
\renewcommand{\@tocrmarg}{3.0em}
\renewcommand*\l@subsection{\@dottedtocline{2}{1.5em}{3.3em}}
\renewcommand*\l@subsubsection{\@dottedtocline{3}{4.8em}{4.2em}}
\makeatother

\newcommand{\R}{\mathbb{R}}
\newcommand{\eps}{\varepsilon}
\newcommand{\diag}{\operatorname{diag}}
\newcommand{\orth}{\operatorname{orth}}
\newcommand{\cof}{\operatorname{cof}}
\newcommand{\clamp}{\operatorname{clamp}}
\newcommand{\sigmoid}{\operatorname{sigmoid}}
\newcommand{\argmin}{\operatorname*{arg\,min}}
\newcommand{\ours}{\textcolor{darkgreen}{\textbf{[제안]}}}
\newcommand{\physics}{\textcolor{darkblue}{\textbf{[물리식]}}}
\newcommand{\learned}{\textcolor{darkpurple}{\textbf{[학습]}}}
\newcommand{\deterministic}{\textcolor{darkgreen}{\textbf{[고정 연산]}}}
\newcommand{\constrained}{\textcolor{darkorange}{\textbf{[제약 풀이]}}}
\newcommand{\trainonly}{\textcolor{darkorange}{\textbf{[학습 전용]}}}
\newcommand{\runtime}{\textcolor{darkgreen}{\textbf{[Runtime]}}}
\newcommand{\optional}{\textcolor{darkblue}{\textbf{[선택]}}}
\newcommand{\risk}{\textcolor{darkred}{\textbf{[위험]}}}
\newcommand{\killgate}{\textcolor{darkred}{\textbf{[중단 기준]}}}
\newcommand{\baseline}{\textcolor{darkpurple}{\textbf{[비교]}}}
\newcommand{\reviewbox}[4]{%
  \par\medskip\noindent
  \fcolorbox{#1}{#2}{%
    \parbox{\dimexpr\linewidth-2\fboxsep-2\fboxrule\relax}{%
      \textbf{V2 리뷰 판정 --- #3.} #4}}
  \par\medskip}
\newcommand{\reviewaccept}[1]{\reviewbox{darkgreen}{softgreen}{수용}{#1}}
\newcommand{\reviewpartial}[1]{\reviewbox{darkorange}{notegray}{부분 수용}{#1}}
\newcommand{\reviewrebut}[1]{\reviewbox{darkred}{softred}{반박/현행 유지}{#1}}
\newcommand{\reviewdefer}[1]{\reviewbox{darkorange}{notegray}{증거 전 보류}{#1}}
\newcommand{\reviewvthreebox}[4]{%
  \par\medskip\noindent
  \begingroup
  \setlength{\fboxsep}{4pt}%
  \fcolorbox{#1}{#2}{\parbox{0.955\linewidth}{%
      \textbf{V3 리뷰 판정 --- #3.} #4}}%
  \endgroup\par\medskip}
\newcommand{\reviewvthreeaccept}[1]{\reviewvthreebox{darkgreen}{softgreen}{수용}{#1}}
\newcommand{\reviewvthreepartial}[1]{\reviewvthreebox{darkorange}{notegray}{부분 수용}{#1}}
\newcommand{\reviewvthreerebut}[1]{\reviewvthreebox{darkred}{softred}{반박/현행 유지}{#1}}
\newcommand{\reviewvthreedefer}[1]{\reviewvthreebox{darkorange}{notegray}{증거 전 보류}{#1}}
\newcommand{\reviewvfourbox}[4]{%
  \par\medskip\noindent
  \begingroup
  \setlength{\fboxsep}{4pt}%
  \fcolorbox{#1}{#2}{\parbox{0.955\linewidth}{%
      \textbf{V4 리뷰 판정 --- #3.} #4}}%
  \endgroup\par\medskip}
\newcommand{\reviewvfouraccept}[1]{\reviewvfourbox{darkgreen}{softgreen}{수용}{#1}}
\newcommand{\reviewvfourpartial}[1]{\reviewvfourbox{darkorange}{notegray}{부분 수용}{#1}}
\newcommand{\reviewvfourrebut}[1]{\reviewvfourbox{darkred}{softred}{반박/현행 유지}{#1}}
\newcommand{\reviewvfourdefer}[1]{\reviewvfourbox{darkorange}{notegray}{감수/보류}{#1}}
\newcommand{\reviewvfivebox}[4]{%
  \par\medskip\noindent
  \begingroup
  \setlength{\fboxsep}{4pt}%
  \fcolorbox{#1}{#2}{\parbox{0.955\linewidth}{%
      \textbf{V5 리뷰 판정 --- #3.} #4}}%
  \endgroup\par\medskip}
\newcommand{\reviewvfiveaccept}[1]{\reviewvfivebox{darkgreen}{softgreen}{수용}{#1}}
\newcommand{\reviewvfivepartial}[1]{\reviewvfivebox{darkorange}{notegray}{부분 수용}{#1}}
\newcommand{\reviewvfiverebut}[1]{\reviewvfivebox{darkred}{softred}{반박/현행 유지}{#1}}
\newcommand{\reviewvfivedefer}[1]{\reviewvfivebox{darkorange}{notegray}{감수/보류}{#1}}
\newcolumntype{Y}{>{\RaggedRight\arraybackslash}X}
\newcolumntype{L}[1]{>{\RaggedRight\arraybackslash}p{#1}}

\title{\textbf{Topology-Distilled, Error-Triggered Global--Local Wind Dynamics}\\
\textbf{for Static 3D Gaussian Thin Surfaces}\\
\large 독립 아이디어 스케치: 물리식 기반 Global predictor와 선택적 complementary Local dynamics}
\author{Wind3DGS 연구 설계 문서}
\date{2026년 8월 13일}

\begin{document}
\maketitle

\begin{center}
\fbox{\parbox{0.95\linewidth}{
\textbf{한 문장 주장.}
독립적으로 재구성된 static 3DGS thin-surface asset에서 training-only mesh teacher로
material topology와 축약 물리 연산자를 distill하고,
runtime에는 현재 변형 표면의 공기력을 object-level reduced dynamics로 적분한 뒤,
고정 cardinality의 atomic gating unit으로 선택한 영역에서만 global modal complement의 Local physics를 계산한다.
Learned future-benefit gate와 missing-force prediction은 각각 독립 검증을 통과한 조건부 확장이다.\\[0.45em]
\textbf{Global 원리.}
현재 위치, 속도, 법선, 면적과 spatial wind에서 quasi-steady aerodynamic force를 계산하고,
V0에서는 모든 physical anchor의 force를 직접 Global generalized force로 투영한 뒤,
항상 켜진 mass--stiffness--damping system을 적분한다.\\[0.45em]
\textbf{Local 원리.}
Patch identity를 유지하는 fixed-superset complement basis에서
fixed-cardinality Top-\(K\)로 선택한 영역만 Local excitation과 선택적 missing-force network를 실행한다.
Global mode와 겹치지 않는 하나의 coupled local physical system이 그 힘을 적분한다.
겹치는 patch는 force와 torque를 보존하도록 먼저 조립하며,
수정된 형상의 aerodynamic load 차이는 다시 Global로 되돌린다.\\[0.45em]
\textbf{조건부 모듈.}
Aerodynamic cubature, learned missing-force network, learned future-benefit gate는
각각 profiling 또는 oracle upper bound가 필요성을 보인 뒤에만 mainline에 남긴다.
V0 selector는 비용 예측 없이 fixed-cardinality atomic-unit ranking을 사용한다.\\[0.45em]
\textbf{Runtime 비사용 항목.}
명시적인 target triangle-mesh reconstruction, full MPM/FEM,
Navier--Stokes/CFD, frame-wise Gaussian 재학습은 사용하지 않는다.
RSH는 mainline에서 제거하고 별도의 공력 표현 비교 항목으로만 남긴다.\\[0.45em]
\textbf{초기 범위.}
정량 core는 고정 또는 compliant attachment를 가진 flat/shallow-rest, single-layer,
asset-uniform isotropic thin surface이며 canonical V0 material은 \(\nu=0\)으로 제한한다:
cantilever strip, flag, pennant와 curtain strip.
Hanging cloth는 secondary stress test, paper-like sheet와 compliant-stem leaf/petal은
qualitative/transfer asset이다. Nonzero Poisson ratio, anisotropy, pre-curved shell coupling,
영구 crease와 plasticity는 첫 논문의 정량 claim이 아니다.
Self-contact, tearing, detached flight, full tree, two-way fluid field update는 초기 논문 범위가 아니다.
}}
\end{center}

\tableofcontents
\newpage

\section{문서 목적과 오늘 확정한 결정}

이 문서는 기존 RSH 중심 문서의 revision이 아니다.
2026년 8월 13일 논의에서 새로 정리한 연구 방향을 독립적인 문제 정의, 방법, 학습 프로토콜,
비교 실험, 선행연구 경계로 처음부터 다시 작성한 아이디어 스케치다.
기존 파일은 수정하지 않으며, 본 문서는 다음 여섯 결정을 출발점으로 삼는다.

\begin{enumerate}[leftmargin=1.7em]
  \item \textbf{Global motion은 축약 물리식으로 푼다.}
  물체 전체의 관성, 복원, 감쇠, attachment, 큰 굽힘과 저주파 진동은 항상 켜진 reduced structural dynamics가 담당한다.
  \item \textbf{Local motion도 물리식으로 적분한다.}
  V0의 작은 local mass--stiffness--damping system은 analytic excitation을 적분한다.
  조건부 F가 통과하면 network는 위치가 아니라 base가 놓친 missing force만 예측한다.
  \item \textbf{Selector는 고정 cardinality 선택으로 실제 sparse execution을 만든다.}
  V0는 이미 계산한 force와 reduced operator를 재사용하는 cheap dynamic-compliance indicator를 쓰고,
  조건부 G가 통과하면
  global-only와 oracle-local counterfactual rollout의 미래 오차 감소량을 예측한다.
  \item \textbf{중복 제거는 두 채널로 분리한다.}
  Force label에서는 base가 이미 설명한 항을 빼고,
  motion space에서는 local basis를 global modal complement로 제한한다.
  실제 local reaction과 corrected-geometry aerodynamic delta는 제거하지 않고 Global에 feedback한다.
  \item \textbf{Patch를 각각 적분한 뒤 평균하지 않는다.}
  Overlap된 patch force를 먼저 하나의 보존적인 local force로 조립하고,
  그 뒤 global complement에 정의된 단일 local state를 적분한다.
  \item \textbf{V0와 후보 연구를 분리한다.}
  V0는 full-anchor analytic aero, patch-indexed fixed Local basis,
  fixed-cardinality selection과 Oracle/physics-only Local로 시작한다.
  Cubature, learned missing force, learned gate와 dynamic basis는 각각 선행 kill test를 통과한 뒤에만 추가한다.
\end{enumerate}

\subsection{``Error-triggered'' 용어 계약}

이 문서의 \emph{error-triggered}는 certified a posteriori error estimator를 뜻하지 않는다.
Training/development의 Oracle future-benefit은 실제 omitted-complement error reduction을 측정하지만,
V0 runtime의 cached dynamic compliance는 그 필요성을 저비용으로 순위화하는
complement-response proxy다. 따라서 calibration은 Oracle regret과 activation delay로 검증하며,
strict tolerance 보장이나 posterior-error bound를 주장하지 않는다.
Learned gate는 이 명칭의 필요조건도 충분조건도 아니다.
Deterministic cached-compliance라도 Oracle benefit과의 rank correlation, regret,
activation delay/high-benefit miss와 strongest heuristic 대비 동일 \(K_A\)/same-p95 quality--latency 개선을
held-out object/gust에서 반복해서 보이면
error relevance를 방어할 수 있다. 반대로 learned gate라도 그 증거가 없으면 명칭을 정당화하지 못한다.
최종 결과가 compliance response만 지지하면 \emph{Compliance-Guided},
error relevance 자체가 약하면 \emph{Selective Complementary}로 제목을 낮춘다.
Every-frame V0 score는 residual이 아니므로 \emph{Residual-Guided}로 바꾸지 않는다.

\subsection{핵심 연구 질문}

\begin{enumerate}[leftmargin=1.7em]
  \item Static 3DGS와 setup metadata가 주어진 target에서 explicit triangle-mesh reconstruction 없이
  thin-surface material connectivity, area/mass, bending/stretching operator와 reduced bases를 구성할 수 있는가?
  \item Full-anchor current-state aerodynamic projection의 실제 비용은 얼마이며,
  profiling 후에만 추가하는 hyper-reduction이 generalized load 정확도를 유지하면서 wall-clock를 줄이는가?
  \item Global reduced physics가 놓치는 local flutter와 fold를
  선택적 complementary Local physics로 복원할 수 있으며,
  physics-only Local이 부족할 때 missing-force network가 추가 이득을 주는가?
  \item Error-triggered Local이 always-on Local과 유사한 품질을 유지하면서
  실제 GPU kernel 실행과 wall-clock cost를 의미 있게 줄이는가?
  \item Local correction이 단순 visual displacement가 아니라
  corrected aerodynamic geometry와 structural coupling을 통해 Global state에 기계적으로 영향을 주는가?
\end{enumerate}

\section{문제 정의}

\subsection{Runtime 입력}

Static appearance만으로 metric material과 attachment를 식별할 수 있다고 가정하지 않는다.
Target asset setup의 입력은 다음으로 고정한다.

\begin{equation}
  \mathcal I_{\mathrm{setup}}
  =
  \left(
    G^0,\;
    s_{\mathrm{metric}},\;
    \vartheta,\;
    \mathcal C
  \right).
\end{equation}

\begin{itemize}[leftmargin=1.5em]
  \item \(G^0=\{(\mu_j^0,\Sigma_j^0,\alpha_j,c_j)\}_{j=1}^{N_G}\):
  독립적으로 재구성된 canonical static 3DGS.
  \item \(s_{\mathrm{metric}}\): known dimension, scale marker 또는 명시적인 nondimensional preset.
  \item \(\vartheta\): density, thickness, stretching/bending stiffness, damping과 aerodynamic coefficient를 포함하는
  user-specified material metadata.
  \item \(\mathcal C\): setup 시 사용자가 지정하는 hard/soft attachment layout과 compliance.
\end{itemize}

Setup은 versioned object package

\begin{equation}
  \mathcal P
  =
  \mathcal P(G^0;s_{\mathrm{metric}},\vartheta,\mathcal C)
\end{equation}

를 만든다. Attachment layout 또는 basis를 바꾸는 material parameter를 수정하면
\(K,\Phi,\{\Psi_r\}\)를 다시 만들거나 refactorize해야 한다.
첫 논문에서는 이를 frame 중 편집 기능으로 주장하지 않고 asset setup event로 취급한다.
고정된 support DOF의 prescribed motion만 같은 package에서 runtime control로 허용한다.

Frame runtime 입력은

\begin{equation}
  \mathcal I_{\mathrm{run}}^t
  =
  \left(
    \mathcal P,\;
    W(\mathbf x,t),\;
    u_{\mathcal C}(t),\;
    \Delta t
  \right)
\end{equation}

이다. 여기서 \(W(\mathbf x,t)\)는 prescribed spatial wind이고,
\(u_{\mathcal C}(t)\)는 이미 정해진 support layout의 선택적 handle motion이다.
Target triangle mesh, target FEM/MPM particles와 source camera trajectory는 runtime 필수 입력이 아니다.

\subsection{Runtime 출력}

각 frame의 출력은 다음이다.

\begin{equation}
  \mathcal O_t
  =
  \left(
    G^t,\;
    q^t,\dot q^t,\;
    z^t,\dot z^t,\;
    \mathcal A_t,\;
    \mathcal D_t,\;
    \mathfrak d_t
  \right),
\end{equation}

여기서 \(q\)는 object-wide Global state,
\(z\)는 patch identity가 고정된 Global-complement Local state,
\(\mathcal A_t\)와 \(\mathcal D_t\)는 각각 active/decay patch 또는 overlap group,
\(\mathfrak d_t\)는 force, gate, conservation, latency diagnostics다.
최종 Gaussian \(G^t\)는 standard 3DGS renderer에 바로 전달된다.

\subsection{Training-only privileged input}

Training asset에서는 source triangle mesh와 high-resolution teacher를 사용할 수 있다.

\begin{equation}
  \mathcal I_{\mathrm{train}}
  =
  \left(
    \mathcal M,\;
    \mathcal T^{\mathrm{shell/MPM}},\;
    \mathcal T^{\mathrm{fluid}}\;(\text{optional}),\;
    G^0_{\mathrm{independent}}
  \right).
\end{equation}

Training 3DGS는 source vertices에 Gaussian을 그대로 부착한 oracle/debug variant와 분리하고,
rendered multi-view images에서 독립적으로 reconstruct한 variant를 포함한다.

\section{범위, 가정, 비주장}

\subsection{초기 물리 범위}

\begin{itemize}[leftmargin=1.5em]
  \item 정량 core는 flat/shallow-rest이고 single-layer 또는 locally single-layer로 볼 수 있는 thin surfaces다.
  V0 material family는 asset-uniform isotropic \(\nu=0\)이며,
  nonzero \(\nu\), anisotropy와 pre-curved shell은 secondary/qualitative 범위다.
  \item Hard 또는 compliant attachment가 있는 object.
  \item Moderate deformation과 flutter; 심한 tearing과 topology change 제외.
  \item Prescribed one-way wind: object deformation은 aerodynamic load를 바꾸지만,
  object가 surrounding wind field 자체를 다시 푸는 full two-way CFD/FSI는 제외.
  \item Contact와 self-contact는 첫 논문에서 제외하거나 teacher-only stress test로 제한.
\end{itemize}

\subsection{명시적으로 주장하지 않는 것}

\begin{itemize}[leftmargin=1.5em]
  \item 최초의 modal wind simulation.
  \item 최초의 Gaussian surface aerodynamic force.
  \item 최초의 ``mesh-free'' reduced 3DGS dynamics.
  \item 최초의 hierarchical/global--local Gaussian simulator.
  \item 최초의 learned residual physics.
  \item 최초의 adaptive local/full-space simulation.
  \item Full CFD 또는 engineering-grade aeroelastic prediction.
  \item Appearance SH 또는 affine Gaussian transport 자체의 신규성.
\end{itemize}

\subsection{V0 core와 조건부 확장}

전체 시스템을 한 번에 구현하지 않는다. 다음 계층은 문서의 우선순위 계약이다.

\begin{longtable}{L{0.18\linewidth}L{0.35\linewidth}L{0.39\linewidth}}
\toprule
계층 & 포함 항목 & 진입 또는 유지 조건 \\
\midrule
\endfirsthead
\toprule
계층 & 포함 항목 & 진입 또는 유지 조건 \\
\midrule
\endhead
V0 core
& Full-anchor current-state analytic aero, always-on Global ROM,
patch-indexed fixed-superset Local basis, Oracle/cached-compliance Top-\(K\),
physics-only Local, conservative assembly/complement, one corrector, affine transport
& 먼저 구현하고 same-anchor sparse solver 및 Oracle Local upper bound로 반증한다. \\
\addlinespace
조건부 H: aero hyper-reduction
& Cubature/sample selection과 hardware-aware aero backend
& Full-anchor aero+projection이 target hardware에서 dynamics 병목이고,
held-out/localized gust에서 generalized-load Pareto를 개선할 때만 유지한다. \\
\addlinespace
조건부 F: learned missing force
& Reduced Local physics가 놓친 nodal force defect prediction
& Full-resolution Local과 physics-only Local 사이에 반복 가능한 gap이 있고,
유일한 nodal label 또는 검증된 lift가 있을 때만 학습한다. \\
\addlinespace
조건부 G: learned gate
& Counterfactual future-benefit predictor와 uncertainty
& Oracle Top-\(K\)와 cached-compliance Top-\(K\) 사이에 충분한 gap이 있고,
같은 \(K_A\)와 measured p95에서 learned selector가 이를 줄일 때만 유지한다. \\
\addlinespace
후속 최적화
& Per-patch hardware cost model, dynamic basis dimension, state remapping
& V0 quality--latency와 profiler 결과가 추가 복잡성을 정당화할 때만 검토한다. \\
\bottomrule
\end{longtable}

조건부 모듈이 no-go 판정을 받는 것은 미완성이 아니라
그 모듈을 제거한 더 작은 방법으로 정상 종료하는 연구 결정이다.

\subsection{Revision 정합성과 source-of-truth 판정}

이 sketch를 method/formula의 canonical source로 둔다.
동기화 이전 snapshot을 기준으로 제기된
``전역 Local orthogonalization'', ``cost-normalized gate mainline'',
``필수 cubature/learned F/G'', ``rank-2 map에 대한 3D cofactor'',
``corrector의 전체 structural reaction 재적용'' 충돌은 현재 revision에는 해당하지 않는다.
현재 계약은 각각 patch-local basis와 explicit cross block,
raw positive benefit 또는 cached-compliance fixed Top-\(K\),
조건부 H/F/G, typed tangent/full-shell map,
updated-minus-consumed structural delta다.
따라서 그 다섯 항목을 현재 설계의 source-of-truth conflict로 보는 판정은 수용하지 않는다.
다만 sketch와 구현 checklist가 이후 다시 어긋날 가능성은 실제 위험이므로,
구현 전 equation label, tensor shape, owner, activation 소비 지점을 자동 대조한다.

\section{핵심 상태와 물리량 소유권}

\subsection{Global--Local 상태 분해}

Sparse material anchors의 rest position을 \(x^0\in\R^{3K}\)라 한다.
Runtime deformation은 다음처럼 분해한다.

\begin{equation}
  u^t
  =
  u_{\mathrm p}^t+\Phi q^t
  +
  \Psi_{\mathrm{sup}}z_{\mathrm{sup}}^t
  =u_{\mathrm p}^t+\Phi q^t+\Psi_{\mathcal S_t}z_{\mathcal S_t}^t,
  \qquad
  x^t=x^0+u^t,
  \qquad
  v^t=\dot u_{\mathrm p}^t+\Phi\dot q^t+\Psi_{\mathrm{sup}}\dot z_{\mathrm{sup}}^t.
  \label{eq:state-decomposition}
\end{equation}

\(\Phi\in\R^{3K\times r_g}\)는 object-wide low-frequency basis다.
\(u_{\mathrm p}^t\)는 prescribed hard-support motion의 particular lift이며
zero/static support이면 0이다.
Local은 하나의 전역 QR/SVD로 섞은 basis가 아니라
patch identity가 고정된 block basis

\begin{equation}
  \Psi_{\mathcal S_t}
  =
  [\Psi_r]_{r\in\mathcal S_t},
  \qquad
  z_{\mathcal S_t}^t=G_{\mathcal S_t}z_{\mathrm{sup}}^t,
  \qquad
  \mathcal S_t=\mathcal A_t\cup\mathcal D_t
\end{equation}

로 둔다. 각 \(\Psi_r\)는 patch 내부 column identity를 보존하면서
Global span만 제거하고 다음 조건을 만족한다.

\begin{equation}
  \Phi^\top M\Psi_r=0,
  \qquad \forall r.
  \label{eq:mass-orthogonality}
\end{equation}

식~\eqref{eq:mass-orthogonality}의 의미는 Local이 Global이 이미 표현한
mass-weighted motion을 다시 만들지 않는다는 것이다.
서로 다른 patch 사이의 \(M\)-orthogonality는 강제하지 않으며,
그 coupling은 \(M_{rs},C_{rs},K_{rs}\) block으로 명시적으로 보존한다.
V0는 모든 \(\Psi_r\)의 coordinate slot을 setup에서 고정하고
active/decay index의 sub-block만 선택하여 푼다.
즉 저장 state는 항상 fixed-size \(z_{\mathrm{sup}}\)이고,
\(G_{\mathcal S_t}\)가 solve에 필요한 block만 gather/scatter한다.
따라서 active set 변화가 basis vector를 다시 혼합하거나 state remapping을 요구하지 않는다.

\subsection{물리량 owner 계약}

\begin{longtable}{L{0.22\linewidth}L{0.27\linewidth}L{0.43\linewidth}}
\toprule
물리량 & 유일한 owner & 계약 \\
\midrule
\endfirsthead
\toprule
물리량 & 유일한 owner & 계약 \\
\midrule
\endhead
Mass/inertia & Global/Local reduced physical system & Network가 inertia를 다시 예측하지 않는다. \\
\addlinespace
동결한 structural restoring & Fixed \(K_g,K_\ell,K_{g\ell}\)와 projective target ledger & Core asset 전체에서 one-shot objective normal-curvature projective/implicit midpoint를 사용한다. Target은 RHS에 한 번만 들어가며, selector와 무관하게 기존 Local energy는 물리적으로 감쇠한다. \\
\addlinespace
Damping & \(C_g,C_\ell,C_{g\ell}\) & Network residual과 중복 학습하지 않는다. \\
\addlinespace
Gravity & explicit external force & Teacher residual에서 반드시 제거한다. \\
\addlinespace
Attachment / support & constraint 또는 compliant force 중 하나 & 같은 attachment를 force와 projection으로 동시에 적용하지 않는다. \\
\addlinespace
Quasi-steady aero & current-surface analytic law & Relative wind, current normal/area로 매 frame 계산한다. \\
\addlinespace
Missing local force & 조건부 F learned residual; V0에서는 channel 없음 & Base analytic/structural force를 제거한 teacher defect만 학습한다. \\
\addlinespace
Global nonlinear defect & First-paper method에는 별도 owner 없음 & Learned generalized-force/state prediction은 비교 baseline 또는 후속 연구다. \\
\addlinespace
Overlap reconciliation & constrained force assembly & Network가 patch seam을 임의로 숨기게 두지 않는다. \\
\addlinespace
Gaussian transport & final total anchor deformation & Global과 Local이 Gaussian을 각각 두 번 움직이지 않는다. \\
\bottomrule
\end{longtable}

\subsection{Selector가 끄는 것과 끄지 않는 것}

Selector는 Global restoring, Global damping, gravity, attachment를 절대 끄지 않는다.
Selector가 제어하는 것은 다음이다.

\begin{itemize}[leftmargin=1.5em]
  \item dormant patch의 Local enrichment 시작 여부,
  \item 조건부 F가 유지된 경우 active patch의 missing-force network 실행 여부,
  \item active aerodynamic correction sample 평가 여부.
\end{itemize}

한 번 local energy를 가진 patch는 selection에서 빠져도 즉시 \(z,\dot z\)를 0으로 reset하거나 freeze하지 않는다.
새 excitation과 조건부 expert를 끄고 작은 local physical system으로 decay를 계속한다.
Group energy, cross load와 selector score가 충분히 작을 때만 dormant state로 반환한다.

\section{Contribution hierarchy와 신규성 경계}

\subsection{권장 contribution}

\begin{enumerate}[leftmargin=1.7em]
  \item \textbf{Topology-distilled reduced physical scaffold.}
  독립적으로 재구성된 static Gaussian thin surface에서
  explicit target triangle-mesh reconstruction 없이 sparse point/edge/local-relation scaffold,
  area/mass와 Global/patch-indexed Local operators를 구성한다.
  Material과 attachment는 appearance에서 자동 식별한다고 주장하지 않고 setup metadata로 받는다.
  \item \textbf{Error-triggered complementary Local dynamics.}
  Global ROM이 놓치는 high-frequency wind response를 patch identity가 유지되는 complement physics로 선택적으로 복원한다.
  Learned future benefit와 learned missing force는 각각 Oracle/physics baseline을 넘을 때만 이 기여의 실현 방식으로 남긴다.
  \item \textbf{Closed-loop geometry-to-load update와 직접 렌더링 가능한 transport.}
  Patch force를 적분 전에 보존적으로 조립하고, corrected-minus-predictor aerodynamic/structural delta만 Global에 되먹인 뒤,
  최종 anchor deformation을 Gaussian mean/covariance에 한 번 운반한다.
  이는 독립된 새 transport 연산이라기보다 앞의 두 algorithmic contribution을 완결하는 system contribution이다.
  Affine transport 자체가 새 연산이라는 주장은 하지 않는다.
\end{enumerate}

Aerodynamic hyper-reduction은 main contribution이 아니라 profiling 기반 go/no-go 후보 모듈이다.
Full-anchor analytic evaluation이 충분히 싸면 이를 제거하고 정확한 generalized-load projection을 사용한다.
또한 고정 physical-anchor/reduced-coordinate budget에서 structural solve 비용이
render Gaussian 수와 분리되는지는 중요한 scaling evidence이지만 독립 algorithmic novelty는 아니다.
Gaussian transport와 rendering 비용까지 \(N_G\)와 무관하다고 주장하지 않는다.

\reviewvfivepartial{Relation-local MLS, objective normal-curvature projective target와
implicit midpoint를 별도의 headline novelty로 올리지 말라는 지적은 수용한다.
이 backend의 기술적 복잡성 자체는 신규성 증거가 아니며 MC1/MC2와 독립된 네 번째 기여로
세지 않는다. 다만 단순 구현 세부로도 격하하지 않는다. 이는 topology-distilled scaffold를
실현하고 Global--Local 비교에 같은 물리 reference를 제공하는
\emph{deterministic objective enabling backend}다. MC1의 operator/eigenspace/action 보존을
판정하는 수식, fixture와 ablation은 유지하되 headline claim은 MC1/MC2/SYS에 둔다.}

\subsection{개별 식, exact chain과 검증 가능한 기여의 구분}

\begin{longtable}{L{0.27\linewidth}L{0.18\linewidth}L{0.47\linewidth}}
\toprule
연산 & 개별 신규성 & 정확한 경계 \\
\midrule
\endfirsthead
\toprule
연산 & 개별 신규성 & 정확한 경계 \\
\midrule
\endhead
\(x=x^0+\Phi q\), reduced ODE & 낮음 & 고전 subspace dynamics와 동일 계열이다 \cite{barbic2005stvk}. \\
\addlinespace
Wind force를 modal force로 projection & 낮음 & Wind Projection Basis가 직접적인 선행이다 \cite{diener2009windbasis}. \\
\addlinespace
Current Gaussian normal/area aero & 낮음 & Gaussian Swaying이 직접적인 선행이다 \cite{yan2026gaussianswaying}. \\
\addlinespace
Static GS reduced deformation & 낮음--중간 & DynamicTree, FreeForm, PhysSkin과 가깝다 \cite{li2026dynamictree,xiang2026freeform,lei2026physskin}. \\
\addlinespace
Physics + learned residual & 낮음 & PGRD와 Learned Contact Corrections가 직접 경계를 만든다 \cite{patel2026pgrd,romero2021contact}. \\
\addlinespace
Adaptive local enrichment & 낮음 & Subspace Condensation과 Trading Spaces가 선행한다 \cite{teng2015condensation,trusty2024tradingspaces}. \\
\addlinespace
Global complement motion & 낮음 & Complementary dynamics 계열이 선행한다 \cite{benchekroun2023complementary}. \\
\addlinespace
제안 V0 chain & 차별화 단서이나 독립 신규성 증거는 아님 & Static GS thin surface, training-only topology/operator distillation, full-anchor current-state load, fixed-cardinality selective complementary Local physics, 보존적 force assembly와 corrected-minus-predictor feedback의 연결은 조사 범위에서 발견하지 못했다. 그러나 MC1의 operator/eigenspace 보존과 MC2의 matched quality--latency 우위를 각각 증명해야 한다. Learned gate, missing-force와 cubature는 독립적인 조건부 실현 방식이다. \\
\bottomrule
\end{longtable}

\subsection{Obvious composition을 넘기 위한 조건}

Gaussian Swaying의 force law를 FreeForm 또는 DynamicTree의 reduced solver에 연결하고
일반 residual network를 더하는 것만으로는 강한 novelty가 되기 어렵다.
``동일한 조합을 찾지 못했다''는 문장은 novelty search의 경계일 뿐 algorithmic contribution의 증명이 아니다.
각 contribution은 자신이 담당하는 claim과 ablation을 독립적으로 통과해야 하며,
한 claim의 실패를 다른 claim의 자동 실패로 확대하지 않는다.

\begin{enumerate}[leftmargin=1.7em]
  \item \textbf{MC1:} 단순 Euclidean kNN이 아닌 oriented material topology와
  setup-conditioned thin-shell operator distillation. Independent GS의 seed, density, hole,
  floater와 close-layer 변화에서도 area/mass, low spectrum과 rollout이 보존되어야 한다.
  \item \textbf{MC2:} Patch identity를 보존하는 complement basis, coupled single-system Local update와
  Oracle/cached-compliance fixed-cardinality sparse execution.
  \item \textbf{SYS:} Joint conservation--complement assembly, local structural coupling,
  corrected-aero delta와 coherent Gaussian transport를 포함하는 안정적인 closed-loop update와
  structural-solve/rendering-resolution scaling evidence.
\end{enumerate}

\section{세 단계 전체 파이프라인}

\begin{longtable}{L{0.17\linewidth}L{0.28\linewidth}L{0.47\linewidth}}
\toprule
Phase & 입력 & 출력 \\
\midrule
\endfirsthead
\toprule
Phase & 입력 & 출력 \\
\midrule
\endhead
Teacher generation & Source mesh, material/attachment, wind program, render cameras & High-resolution structural trajectories, optional fluid loads, multi-view images \\
\addlinespace
Distillation / training & Teacher trajectories, independently reconstructed static GS & Target-runtime object package; 조건부 gate/missing-force model \\
\addlinespace
Target runtime & Versioned object package, prescribed wind, fixed-layout handle motion & Global/Local state, deformed Gaussians, diagnostics and rendered frames \\
\bottomrule
\end{longtable}

\subsection{Target-runtime object package}

한 target asset의 setup 단계는 다음 package를 출력한다.

\begin{equation}
\begin{aligned}
  \mathcal P(G^0;s_{\mathrm{metric}},\vartheta,\mathcal C)=\{&
  x^0,\mathcal E,\mathcal H,\mathcal Q_{\mathrm{op}},\mathcal R,\;
  A^0,n^0,M,\{c_{k,\mathrm{abs}}^c\}_{k,c},\;
  \vartheta,\mathcal C,\\
  &K_{\mathrm{lin}}^0,C_{\mathrm{lin}}^0,K^0,C^0,N_{\mathrm h},M_{\mathrm f},\\
  &K_{\mathrm{lin,f}},C_{\mathrm{lin,f}},K_{\mathrm f},C_{\mathrm f},\\
  &\mathcal U_{\mathrm p},\mathcal U_{\mathrm c},
  \widetilde\Phi_{\mathrm f},\Phi,\\
  &\{B_r,\widetilde\Psi_{\mathrm f,r},\Psi_r,\mathcal I_r^{\mathrm{own}},
  \rho_{r,\mathrm{tail}}\}_{r=1}^{P},\;
  \sigma_{\mathrm{struct}},\;
  \mathcal S_{\mathrm{full}},\mathcal H_{\mathrm{aero}}^{\mathrm{optional}},\;
  \mathcal B_{\mathrm{GS}}
  \}.
\end{aligned}
\end{equation}

\(\mathcal E\)는 material adjacency,
\(\mathcal H\)는 bending/hinge-like 관계,
\(\mathcal Q_{\mathrm{op}}\)는 식~\eqref{eq:relation-mls-map}의
metric MLS stencil, normalized relation-area ownership과 constitutive channel을 포함하는
versioned operator relation 집합,
\(N_{\mathrm h}\)는 hard-attachment free-DOF null-space map,
\(K_{\mathrm{lin}}^0,C_{\mathrm{lin}}^0\)은 mode construction과 small-strain calibration에 쓰는
rest-tangent operator이고, \(K^0,C^0\)는 canonical normal-curvature-projective solve의
anchor-space fixed Hessian,
\((K_{\mathrm{lin,f}},C_{\mathrm{lin,f}})\)는 rest-linear free-space reference,
\((M_{\mathrm f},K_{\mathrm f},C_{\mathrm f})\)는 canonical free-space solve operator,
\(\mathcal U_{\mathrm p}\)는 prescribed support motion에서
\((u_{\mathrm p},\dot u_{\mathrm p},\ddot u_{\mathrm p})\)를 만드는 typed lift,
\(\mathcal U_{\mathrm c}\)는 같은 runtime handle control에서 compliant target
\((d_{\mathrm c},\dot d_{\mathrm c})\)를 만드는 별도 typed map,
\(\sigma_{\mathrm{struct}}\)는 식~\eqref{eq:v0-structural-freeze-signature}의
backend/integrator/timestep/operator signature,
\(\mathcal R\)은 overlapping core--halo patches,
\(\rho_{r,\mathrm{tail}}\)은 식~\eqref{eq:patch-complement-tail}에서 정의하는
package-validation locality diagnostic,
\(\mathcal S_{\mathrm{full}}\)은 모든 physical anchor의 aerodynamic support,
\(\mathcal H_{\mathrm{aero}}^{\mathrm{optional}}\)은 profiling을 통과했을 때만 포함하는
cubature sample/weight backend,
\(\mathcal B_{\mathrm{GS}}\)는 anchor-to-Gaussian binding이다.

\section{Teacher와 학습 데이터 생성}

\subsection{Teacher-S: 넓은 구조 상태를 위한 teacher}

주 데이터는 high-resolution nonlinear thin-shell FEM, cloth FEM,
또는 검증된 dense MPM에서 만든다.
Teacher-S는 다음 label을 제공한다.

\begin{itemize}[leftmargin=1.5em]
  \item mesh positions, velocities, accelerations,
  \item strain, curvature, bending/stretching energy,
  \item internal/external/reaction force breakdown,
  \item surface normal, area, deformation gradient,
  \item attachment reaction,
  \item optional Gaussian mean/covariance target.
\end{itemize}

Teacher-S가 사용하는 aerodynamic base가 runtime과 같은 quasi-steady law라면,
learned missing force는 basis truncation, reduced operator error,
nonlinear shell response와 topology/material distillation error만 학습한다고 주장한다.

\subsection{Teacher-FSI: 제한된 unsteady aerodynamic teacher}

Wake, vortex shedding, history-dependent flutter를 주장하려면
작은 asset subset에 CFD, LBM, 또는 실제 계측 기반 teacher가 추가로 필요하다.
이 경우에만 missing-force label에 다음 항이 들어갈 수 있다.

\begin{equation}
  \Delta f_{\mathrm{unsteady}}^*
  =
  f_{\mathrm{FSI/measurement}}^*
  -
  f_{\mathrm{quasi\mbox{-}steady}}^*.
\end{equation}

\risk Teacher-FSI 없이 \emph{wake를 학습한다} 또는
\emph{vortex shedding을 복원한다}고 주장하면 안 된다.
Teacher-S만 사용하는 MVP에서는 해당 항을 basis/structure correction으로 제한한다.

\subsection{Factorized data sweep}

\begin{equation}
  \mathcal D
  =
  \mathcal D_{\mathrm{shape}}
  \times\mathcal D_{\mathrm{material}}
  \times\mathcal D_{\mathrm{attachment}}
  \times\mathcal D_{\mathrm{wind}}
  \times\mathcal D_{\mathrm{GS}}
  \times\mathcal D_{\mathrm{initial}}.
\end{equation}

\begin{itemize}[leftmargin=1.5em]
  \item Shape: flat/shallow flag, strip, pennant와 curtain; paper/leaf는 secondary qualitative transfer.
  \item Material: V0 isotropic density, thickness, Young's modulus와 damping.
  Canonical constant-Hessian projective backend는 \(\nu=0\)의 corotated small-strain limit로 고정한다.
  Nonzero Poisson ratio와 anisotropy는 state-dependent/extra projective constraint가 필요한
  후속 또는 qualitative stress test로만 둔다.
  \item Attachment: one edge, two points, partial compliant support, stem support.
  \item Wind direction/magnitude: yaw, elevation, weak-to-strong flutter range.
  \item Temporal wind: steady, sinusoidal gust, chirp, abrupt start/stop.
  \item Spatial wind: traveling gust, localized gust, nonuniform macro zones.
  \item GS quality: density, view count, covariance noise, opacity perturbation,
  holes, missing backside, floaters, reconstruction seed.
\end{itemize}

\subsection{Split protocol}

Frame-random split은 동일 trajectory의 이웃 frame을 train/test에 섞으므로 금지한다.
최소 split 단위는 trajectory이며, 주요 일반화 결과는 object-disjoint split으로 보고한다.

\begin{enumerate}[leftmargin=1.7em]
  \item Seen object / unseen wind.
  \item Seen object / unseen V0 scalar material parameters
  \((\rho,h,E,\tau_{\mathrm m},\tau_{\mathrm b})\) with \(\nu=0\).
  \item Seen object / unseen attachment.
  \item Unseen geometry within the same category.
  \item Unseen object category.
  \item Independent GS reconstruction of a seen source shape.
  \item Unseen object + declared-V0-family material + wind combined OOD.
\end{enumerate}

\section{Topology와 reduced operator distillation}

\subsection{Surface-aware anchors}

Raw Gaussian count를 physical importance로 사용하지 않는다.
Opacity, covariance surface-likeness, multi-view contribution,
local isolation, boundary/curvature evidence를 함께 사용하여
surface-reliable Gaussian support를 만들고 \(K\ll N_G\)개의 material anchors를 선택한다.

각 anchor는 다음 attribute를 갖는다.

\begin{equation}
  a_k^0
  =
  \left(
  x_k^0,n_k^0,A_k^0,m_k,
  \vartheta_k,c_k,\eta_k
  \right),
\end{equation}

여기서 \(c_k\)는 attachment/compliance,
\(\eta_k\)는 reconstruction reliability다.
면적과 질량은 Gaussian별 중복 support를 그대로 합하지 않고
partition-of-unity로 총량을 보존한다.

\subsection{Oriented material topology}

Topology network는 Euclidean proximity만 출력하지 않는다.
후보 edge \((i,j)\)에 대해 다음을 예측한다.

\begin{equation}
  (\widehat y_{ij}^{\mathrm{mat}},
   \widehat \ell_{ij}^0,
   \widehat t_{ij},
   \widehat\eta_{ij}).
\end{equation}

Accepted operator relation \(\alpha\)의 각 incidence \(k\in\mathcal I_\alpha\)에는
\((\widehat a_{\alpha k}^{\mathrm m},\widehat a_{\alpha k}^{\mathrm b},
\widehat v_{\alpha k}^{\mathrm m},\widehat v_{\alpha k}^{\mathrm b},
\widehat\eta_{\alpha k})\)를 출력한다.
별도의 local relation도 이 incidence schema를 사용하며 bending/hinge-like coupling,
orientation consistency, disconnected-layer rejection을 나타낸다.
Training label은 mesh geodesic/material adjacency에서 오지만,
runtime graph는 explicit target triangle face connectivity를 요구하지 않는다.
\(\widehat a^{\mathrm m},\widehat a^{\mathrm b}\)는 뒤의
식~\eqref{eq:relation-area-partition}에서만 소비하는
dimensionless relation-area logit이고, \(\widehat\eta\)는 reconstruction confidence다.
\(\widehat v^{c}\)는 area share와 별도로 supervision/calibration하는 channel별 relation validity다.
이 값을 physical stiffness로 직접 해석하지 않는다.
Metric stiffness, thickness와 density는 setup 입력 \(\vartheta\)에서만 조립한다.
Attachment layout도 \(G^0\) appearance에서 자동 예측하는 mainline 항목이 아니라
setup 입력 \(\mathcal C\)를 anchor에 transfer한 값이다.

\subsection{Relation-local metric stencil}

\reviewvthreepartial{V3 리뷰가 지적한 대로
\(K_{\mathrm{stretch}},K_{\mathrm{shear}},K_{\mathrm{bend}}\)의 실제 조립식과
material scaling은 MC1의 핵심 계약이므로 수용한다.
다만 \(K=B^\top D B\)와
\(w=\operatorname{softplus}(\widehat w)\)만으로는 algebraic PSD만 보장한다.
그 식만으로는 \(B\)의 rigid-motion nullspace와 polynomial reproduction,
finite-rotation objectivity, anchor valence/density 변화에 대한 resolution consistency,
geometry weight와 metric material의 분리가 보장되지 않는다.
따라서 V0는 아래 metric MLS stencil, normalized area partition,
user-specified constitutive matrix와 typed attachment assembly를 함께 고정한다.
Finite-rotation objectivity는 rest operator 자체의 성질로 과장하지 않고
뒤의 structural-backend freeze에서 별도로 판정한다.}

Operator relation 집합을 \(\mathcal Q_{\mathrm{op}}\)라 한다.
각 \(\alpha\in\mathcal Q_{\mathrm{op}}\)는

\begin{equation}
  \alpha=
  (\mathcal I_\alpha,S_\alpha,T_\alpha^0,
   \omega_\alpha,\ell_\alpha),
\end{equation}

를 저장한다.
\(\mathcal I_\alpha\)는 oriented point/edge/local-relation graph에서 얻은 local anchor support,
\(S_\alpha\in\{0,1\}^{3n_\alpha\times3K}\)는 support extraction,
\(T_\alpha^0=[t_{\alpha1}^0,t_{\alpha2}^0,n_\alpha^0]\in SO(3)\)는 rest local frame이다.
이 relation은 global triangle face, watertightness 또는 manifold connectivity를 요구하지 않는다.
V0 stencil은 flat 또는 shallow-rest strip/flag/pennant/curtain을 위한 plate operator다.
Curved shell의 covariant membrane--curvature coupling은 포함하지 않으므로
pre-curved leaf/petal은 qualitative stress test로만 보고, 정량 core claim에는 사용하지 않는다.

\begin{equation}
\begin{aligned}
  \bar x_\alpha^0
  &=\sum_{i\in\mathcal I_\alpha}\omega_{\alpha i}x_i^0,
  &\sum_i\omega_{\alpha i}&=1,\\
  \ell_\alpha^2
  &=\sum_{i\in\mathcal I_\alpha}\omega_{\alpha i}
  \left\|
  [t_{\alpha1}^0,t_{\alpha2}^0]^\top
  (x_i^0-\bar x_\alpha^0)
  \right\|_2^2 .
\end{aligned}
\end{equation}

Metric tangent coordinate와 quadratic polynomial basis를

\begin{equation}
\begin{aligned}
  \zeta_{\alpha i}
  &=\ell_\alpha^{-1}
  [t_{\alpha1}^0,t_{\alpha2}^0]^\top
  (x_i^0-\bar x_\alpha^0),\\
  p(\zeta)
  &=
  [1,\zeta_1,\zeta_2,
    \tfrac12\zeta_1^2,\zeta_1\zeta_2,
    \tfrac12\zeta_2^2]^\top
\end{aligned}
\end{equation}

로 둔다. \(\mathsf V_\alpha=[p(\zeta_{\alpha i})^\top]_i\)와
positive dimensionless fit weight \(\mathsf W_\alpha^{\mathrm{fit}}\)에서

\begin{equation}
  \mathsf G_\alpha
  =
  (\mathsf V_\alpha^\top
   \mathsf W_\alpha^{\mathrm{fit}}\mathsf V_\alpha)^{-1}
  \mathsf V_\alpha^\top\mathsf W_\alpha^{\mathrm{fit}}
  \label{eq:relation-mls-map}
\end{equation}

를 만든다.
모든 \(\omega_{\alpha i}>0\), \(\ell_\alpha>0\),
\(\operatorname{rank}(\mathsf V_\alpha)=6\)과 package condition-number tolerance를
통과하지 못하면 support를 넓히거나 relation을 거부한다.
Regularization으로 rank failure를 숨기지 않는다.
\(e_j^{(6)}\)를 \(\R^6\)의 표준기저라 하면 derivative row는

\begin{equation}
\begin{aligned}
  L_{\alpha1}&=\ell_\alpha^{-1}(e_2^{(6)})^\top\mathsf G_\alpha,
  &L_{\alpha2}&=\ell_\alpha^{-1}(e_3^{(6)})^\top\mathsf G_\alpha,\\
  H_{\alpha11}&=\ell_\alpha^{-2}(e_4^{(6)})^\top\mathsf G_\alpha,
  &H_{\alpha12}&=\ell_\alpha^{-2}(e_5^{(6)})^\top\mathsf G_\alpha,
  &H_{\alpha22}&=\ell_\alpha^{-2}(e_6^{(6)})^\top\mathsf G_\alpha .
\end{aligned}
\end{equation}

Local rigid displacement를 construction으로 제거하기 위해

\begin{equation}
  Z_\alpha^0
  =
  \begin{bmatrix}
    I_3&-[x_{\alpha1}^0-\bar x_\alpha^0]_\times\\
    \vdots&\vdots\\
    I_3&-[x_{\alpha n_\alpha}^0-\bar x_\alpha^0]_\times
  \end{bmatrix},
\end{equation}

\begin{equation}
  P_\alpha^{\mathrm{rb}}
  =I-Z_\alpha^0
  \bigl((Z_\alpha^0)^\top
        W_\alpha^{\mathrm{dof}}Z_\alpha^0\bigr)^+
  (Z_\alpha^0)^\top W_\alpha^{\mathrm{dof}},
  \qquad
  W_\alpha^{\mathrm{dof}}
  =\diag_i(\omega_{\alpha i}I_3)
\end{equation}

로 둔다.
\(\mathcal T_\alpha=I_{n_\alpha}\otimes(T_\alpha^0)^\top\),
\(e_a^{(3)}\)를 \(\R^3\)의 표준기저,
\(J_{\alpha a}=I_{n_\alpha}\otimes(e_a^{(3)})^\top\)라 하면
global anchor displacement에 작용하는 rest-linear relation operator는

\begin{equation}
\begin{aligned}
  B_{\alpha,\mathrm{s}}
  &=
  \begin{bmatrix}
    L_{\alpha1}J_{\alpha1}\mathcal T_\alpha\\
    L_{\alpha2}J_{\alpha2}\mathcal T_\alpha
  \end{bmatrix}
  P_\alpha^{\mathrm{rb}}S_\alpha,\\
  B_{\alpha,\mathrm{sh}}
  &=
  \left(
    L_{\alpha2}J_{\alpha1}\mathcal T_\alpha+
    L_{\alpha1}J_{\alpha2}\mathcal T_\alpha
  \right)P_\alpha^{\mathrm{rb}}S_\alpha,\\
  B_{\alpha,\mathrm{b}}
  &=
  \begin{bmatrix}
    H_{\alpha11}J_{\alpha3}\mathcal T_\alpha\\
    H_{\alpha22}J_{\alpha3}\mathcal T_\alpha\\
    2H_{\alpha12}J_{\alpha3}\mathcal T_\alpha
  \end{bmatrix}
  P_\alpha^{\mathrm{rb}}S_\alpha.
  \label{eq:exact-relation-strain-operators}
\end{aligned}
\end{equation}

Relation-local factor는
\(B_{\alpha,c}=\widehat B_{\alpha,c}S_\alpha\),
\(c\in\{\mathrm m,\mathrm b\}\)로 저장하여 teacher probe와 package hard check에서 재사용한다.
Isotropic membrane channel은 tangent-axis 회전에 따라 normal strain과 shear가 섞여도
동일 energy를 내도록

\begin{equation}
  B_{\alpha,\mathrm m}
  =
  \begin{bmatrix}
    B_{\alpha,\mathrm s}\\
    B_{\alpha,\mathrm{sh}}
  \end{bmatrix}
  \label{eq:relation-membrane-operator}
\end{equation}

로 묶는다. Stretch와 shear에 서로 다른 relation area를 주지 않는다.

따라서 translation과 infinitesimal rigid rotation에 대해
\(B_{\alpha,c}Z^0=0\),
\(c\in\{\mathrm m,\mathrm b\}\)가 construction으로 성립한다.

\subsection{Normalized relation area와 material separation}

Network의 dimensionless geometry logit은 physical stiffness로 직접 사용하지 않는다.
Channel \(c\in\{\mathrm m,\mathrm b\}\)와 anchor \(k\)를 소유할 수 있는
relation 집합을 \(\mathcal Q_c(k)\)라 하면

\begin{equation}
  q_{\alpha k}^{c,\mathrm{val}}
  :=\sigmoid(\widehat v_{\alpha k}^{c}),
  \qquad
  p_{\alpha k}^{c,\mathrm{val}}
  :=
  \operatorname{Cal}_{\mathrm{rel},c}
  \left(
    q_{\alpha k}^{c,\mathrm{val}},
    \widehat\eta_{\alpha k},
    \kappa(\mathsf V_\alpha^\top
            \mathsf W_\alpha^{\mathrm{fit}}\mathsf V_\alpha)
  \right)
  \in[0,1]
  \label{eq:relation-validity-confidence}
\end{equation}

로 relation-validity probability를 별도 출력하고 development split에서 calibration한다.
\(\widehat a_{\alpha k}^{c}\)는 normalized area share만 정하며,
\(\widehat v_{\alpha k}^{c}\)의 raw probability \(q^{c,\mathrm{val}}\)는 teacher relation validity로
학습하고, \(\operatorname{Cal}_{\mathrm{rel},c}\)은 그 뒤 development split에서 fit/freeze한다.
그러면

\begin{equation}
\begin{aligned}
  \widetilde\pi_{\alpha k}^{c}
  &=
  \mathbf 1[\alpha\in\mathcal Q_c(k)]
  \,p_{\alpha k}^{c,\mathrm{val}}
  \,\operatorname{softplus}(\widehat a_{\alpha k}^{c}),\\
  \pi_{\alpha k}^{c}
  &=
  \frac{\widetilde\pi_{\alpha k}^{c}}
  {\sum_{\beta\in\mathcal Q_c(k)}
   \widetilde\pi_{\beta k}^{c}},
  &a_\alpha^c&=\sum_k\pi_{\alpha k}^cA_k^0.
  \label{eq:relation-area-partition}
\end{aligned}
\end{equation}

\reviewvfourpartial{모든 candidate가 낮은 confidence여도 normalized share가
anchor area 전체를 분배한다는 실패 모드는 수용한다. 그러나
\(\sum_\alpha\widetilde\pi_{\alpha k}^{c}\)를 그대로 absolute confidence로 쓰는 처방은
채택하지 않는다. Area logit은 normalized share를 통해서만 식별되므로 공통 logit shift,
candidate 수와 valence에 따라 raw sum이 임의로 달라질 수 있다. 따라서 area ownership과
별도로 supervision/calibration한 validity probability와 count-aware package confidence를 쓴다.
이 confidence를 Young's modulus나 stiffness에 곱해 불확실성을 재료 연화로 숨기지 않는다.}

Anchor별 absolute evidence와 package confidence는
\(\kappa_\alpha:=\kappa(\mathsf V_\alpha^\top
\mathsf W_\alpha^{\mathrm{fit}}\mathsf V_\alpha)\)라 두고

\begin{equation}
\begin{aligned}
  n_k^c&:=|\mathcal Q_c(k)|,
  &e_{k,\Sigma}^{c}&:=\sum_{\alpha\in\mathcal Q_c(k)}
    p_{\alpha k}^{c,\mathrm{val}},\\
  \bar e_k^c&:=\frac{e_{k,\Sigma}^{c}}{\max(1,n_k^c)},
  &p_{k,\max}^c&:=\max_{\alpha\in\mathcal Q_c(k)}
    p_{\alpha k}^{c,\mathrm{val}},\\
  c_{k,\mathrm{abs}}^c
  &:={}
  \operatorname{Cal}_{\mathrm{anchor},c}
  \left(
    p_{k,\max}^c,\bar e_k^c,
    \log(1+e_{k,\Sigma}^{c}),n_k^c,
    \{\kappa_\alpha\}_{\alpha\in\mathcal Q_c(k)}
  \right)
  &&\in[0,1].
  \label{eq:relation-absolute-confidence}
\end{aligned}
\end{equation}

Anchor-level calibration target은 teacher-valid relation이 단순히 후보에 존재하는지가 아니라
operator를 실제로 구성할 수 있는지를 나타낸다. Channel별 minimum accepted rank를
\(r_{\min,c}\), condition threshold를 \(\kappa_{\max,c}\)라 두고

\begin{equation}
  y_{k,c}^{\mathrm{abs},*}
  :=
  \mathbf 1\!\left[
  \exists\alpha\in\mathcal Q_c(k):
  y_{\alpha k}^{c,\mathrm{val},*}=1,
  \ \operatorname{rank}(\mathsf V_\alpha)\ge r_{\min,c},
  \ \kappa_\alpha\le\kappa_{\max,c},
  \ \mathrm{orientOK}_\alpha=1
  \right].
  \label{eq:anchor-absolute-confidence-target}
\end{equation}

\(\operatorname{Cal}_{\mathrm{anchor},c}\)은 development split에서 이 binary target에 대한
NLL/Brier를 최소화해 fit한 뒤 freeze한다. Held-out에서는 Brier, ECE, AUPRC와
\(c_{\min,c}\)에서의 accepted coverage/false-accept rate를 함께 보고한다.
Development split에서 고정한 calibration을 held-out/test에서 다시 맞추지 않는다.
\(n_k^c=0\)이면 \(p_{k,\max}^c=0\), \(c_{k,\mathrm{abs}}^c=0\)으로 둔다.
\(c_{k,\mathrm{abs}}^c<c_{\min,c}\)이면 support 확대와 MLS rank/conditioning 재검사를 먼저
수행하고, 여전히 실패하면 user hint, conservative fallback 또는 package rejection을 반환한다.
Normalized share가 존재한다는 이유만으로 package를 통과시키지 않는다.

\reviewvfiveaccept{Absolute confidence는 operator 품질의 대체 점수가 아니라
setup package를 support 확대, user hint, fallback 또는 rejection으로 보내는 triage 신호라는
지적을 수용한다. Calibration split은 object-disjoint로 만들고, development에서 고정한
anchor-density 및 candidate-valence 구간별 Brier/ECE, coverage와 false-accept rate를
held-out/test에 보고한다. Test를 보고 calibrator, threshold 또는 bin boundary를 다시 맞추지 않는다.}

Confidence triage를 통과해도 package acceptance가 자동으로 주어지지 않는다.
최종 acceptance는 total/local area와 mass 보존, relation rank/conditioning,
affine membrane 및 quadratic normal-curvature reproduction, rest-linear pair의 첫 8--16개
eigenfrequency와 modal principal angle, common-probe rollout과 stability hard check의
conjunction으로 판정한다. 하나라도 실패하면 높은 \(c_{k,\mathrm{abs}}^c\)로 상쇄하지 않는다.

Package는 denominator가 0인 anchor를 거부하고

\begin{equation}
  \sum_{\alpha\in\mathcal Q_{\mathrm{op}}}a_\alpha^c
  =\sum_kA_k^0,
  \qquad
  c\in\{\mathrm m,\mathrm b\}
\end{equation}

를 검사한다.
Softplus는 positivity를 위한 내부 parameterization일 뿐이며,
normalization과 area conservation이 resolution-scaling 계약이다.

\subsection{Mass와 rest-linear constitutive reference}

첫 논문 V0의 material family는 user metadata로 지정한
asset-uniform isotropic thin surface로 고정한다.
Setup은 \(\vartheta=(\rho,h,E,\tau_{\mathrm m},\tau_{\mathrm b})\)를 제공하며,
\(\rho,h,E>0\), \(\tau_c\ge0\)를 요구한다.
Canonical constant-Hessian projective backend는 \(\nu=0\)의 corotated
small-strain limit만 주장한다. Nonzero Poisson ratio에는 area-like projective constraint나
state-dependent co-rotated tangent가 추가로 필요하므로 V0 material sweep에 넣지 않는다.
Heterogeneous 또는 anisotropic constitutive field도 후속 범위이며,
V0 network가 relation별 modulus나 arbitrary SPD matrix를 예측하지 않는다.
\(\rho\)는 volumetric density다.
만약 preset이 areal density \(\bar\rho\)를 제공하면 별도 type을 사용하고
같은 package에서 두 density convention을 섞지 않는다.

\begin{equation}
  m_k=\rho hA_k^0
  \quad\text{또는}\quad
  m_k=\bar\rho A_k^0,
  \qquad
  M=\diag(m_1I_3,\ldots,m_KI_3).
  \label{eq:exact-anchor-mass}
\end{equation}

모든 relation은 같은 선언 material을 사용한다. \(\nu=0\) plane-stress의
rest-linear calibration matrix는

\begin{equation}
\begin{aligned}
  D_{\alpha,\mathrm m}^{\mathrm{lin}}
  &=a_\alpha^{\mathrm m}Eh
  \begin{bmatrix}
    1&0&0\\
    0&1&0\\
    0&0&1/2
  \end{bmatrix},\\
  D_{\alpha,\mathrm{b}}^{\mathrm{lin}}
  &=a_\alpha^{\mathrm{b}}\frac{Eh^3}{12}
  \begin{bmatrix}
    1&0&0\\
    0&1&0\\
    0&0&1/2
  \end{bmatrix}.
  \label{eq:relation-constitutive-matrices}
\end{aligned}
\end{equation}

\begin{equation}
\begin{aligned}
  K_{\mathrm{membrane}}^{\mathrm{lin}}
  &=\sum_\alpha B_{\alpha,\mathrm m}^\top
    D_{\alpha,\mathrm m}^{\mathrm{lin}}B_{\alpha,\mathrm m},\\
  K_{\mathrm{bend}}^{\mathrm{lin}}
  &=\sum_\alpha B_{\alpha,\mathrm{b}}^\top
    D_{\alpha,\mathrm{b}}^{\mathrm{lin}}B_{\alpha,\mathrm{b}},\\
  K_{\mathrm{lin}}^0
  &=K_{\mathrm{membrane}}^{\mathrm{lin}}
    +K_{\mathrm{bend}}^{\mathrm{lin}}.
  \label{eq:exact-positive-stiffness-assembly}
\end{aligned}
\end{equation}

Rest-linear Kelvin--Voigt reference는
\(D_{\alpha,c}^{\eta,\mathrm{lin}}
=\tau_cD_{\alpha,c}^{\mathrm{lin}}\succeq0\)로 두고

\begin{equation}
\begin{aligned}
  C_{\mathrm{lin}}^0
  &=\sum_{\alpha,c}B_{\alpha,c}^\top
    D_{\alpha,c}^{\eta,\mathrm{lin}}B_{\alpha,c}.
  \label{eq:exact-positive-damping-assembly}
\end{aligned}
\end{equation}

이 \(B^\top DB\) operator는 topology/metric/material distillation의
small-strain calibration, modal basis와 rest-linear baseline에 사용한다.
이를 finite-rotation runtime energy에 그대로 넣고 회전된 strain target만 RHS에 두면
순수 rigid state는 맞출 수 있어도 이미 변형된 state에 rigid rotation을 중첩했을 때
material-frame objectivity가 깨질 수 있다. 따라서 다음 world-vector normal-curvature projective operator와 구분한다.

\subsection{Objective normal-curvature projective runtime operator}

\reviewvthreepartial{V3 리뷰가 concrete corotational backend를 고정하라는 요구는 수용한다.
그러나 rest-frame scalar \(B\)에 회전 target만 더하는 fixed-Hessian 식은
strained state에 superposed rigid rotation을 가했을 때 stretch와 shear channel을 섞을 수 있으므로
그대로 채택하지 않는다. Canonical runtime은 scalar MLS coefficient를 world-vector에 동일하게 작용시키는
left-rotation-equivariant projective constraint를 사용한다. 이 선택은 fixed Hessian과 exact
superposed-rotation objectivity를 함께 보존하지만 arbitrary Poisson response를 무료로 주지는 않는다.
따라서 V0는 \(\nu=0\)으로 범위를 좁히고 nonzero \(\nu\)는 결과 없이 주장하지 않는다.}

\reviewvfouraccept{V4 리뷰의 Critical A는 수용한다.
현재 \(\mathsf A_{\alpha,\mathrm b}\)와 block-isotropic metric을 그대로 energy에 쓰면
normal curvature뿐 아니라 spatially varying in-plane deformation의 tangential second derivative까지
bending으로 벌점화한다. 이는 rest-linear \(B_{\alpha,\mathrm b}\)가 local normal component만
사용하는 계약과 다르고, flutter/wrinkle spectrum을 불필요하게 평활화할 수 있다.
따라서 vector Hessian descriptor와 fixed global-step Hessian은 유지하되 실제 bending stiffness와
damping residual은 predictor normal에 대한 projective normal-curvature target으로 정의한다.
Frozen target을 한 번만 쓰는 final state에서 tangential residual이 정확히 0이라고 주장하지 않고
그 predictor lag를 splitting error로 측정한다.}

Relation마다 scalar MLS derivative row를 각 world component에 동일하게 적용하는
vector constraint를

\begin{equation}
\begin{aligned}
  \mathsf A_{\alpha,\mathrm m}
  &=
  \left(
  \begin{bmatrix}L_{\alpha1}\\L_{\alpha2}\end{bmatrix}
  \otimes I_3\right)S_\alpha,\\
  \mathsf A_{\alpha,\mathrm b}
  &=
  \left(
  \begin{bmatrix}H_{\alpha11}\\H_{\alpha22}\\
  \sqrt2H_{\alpha12}\end{bmatrix}
  \otimes I_3\right)S_\alpha,
  \qquad
  d_{\alpha,c}^0:=\mathsf A_{\alpha,c}x^0 .
  \label{eq:objective-vector-projective-constraints}
\end{aligned}
\end{equation}

로 둔다. \(m_{\mathrm m}=2,m_{\mathrm b}=3\)라 하면
\(\mathsf A_{\alpha,c}\in\R^{3m_c\times3K}\)이고,
MLS polynomial reproduction 때문에 translation을 소거한다.
모든 coefficient row가 scalar이고 world component에는 \(I_3\)로 작용하므로
상수 \(Q\in SO(3)\), translation \(a\)에 대해

\begin{equation}
  \mathsf A_{\alpha,c}
  \bigl((I_K\otimes Q)x+\mathbf1_K\otimes a\bigr)
  =(I_{m_c}\otimes Q)\mathsf A_{\alpha,c}x.
  \label{eq:vector-constraint-equivariance}
\end{equation}

Positive block-isotropic metric과 fixed projective Hessian은

\begin{equation}
\begin{aligned}
  \Omega_{\alpha,\mathrm m}
  &=a_\alpha^{\mathrm m}Eh\,I_6,
  &\Omega_{\alpha,\mathrm b}
  &=a_\alpha^{\mathrm b}\frac{Eh^3}{12}I_9,\\
  \Omega_{\alpha,c}^{\eta}
  &=\tau_c\Omega_{\alpha,c},
  &K_{\mathrm{proj}}^0
  &=\sum_{\alpha,c}\mathsf A_{\alpha,c}^\top
    \Omega_{\alpha,c}\mathsf A_{\alpha,c},\\
  &&C_{\mathrm{proj}}^0
  &=\sum_{\alpha,c}\mathsf A_{\alpha,c}^\top
    \Omega_{\alpha,c}^{\eta}\mathsf A_{\alpha,c}.
  \label{eq:fixed-vector-projective-hessian}
\end{aligned}
\end{equation}

으로 조립한다. \(\Omega=\widehat\Omega\otimes I_3\)가 모든 global rotation과 commute하므로
뒤에서 회전 target을 local step으로 갱신해도 global Hessian은 고정된다.
\(K_{\mathrm{lin}}^0\)은 mode/infinitesimal reference이고
\(K_{\mathrm{proj}}^0\)은 local--global solve matrix이므로 같은 matrix라고 주장하지 않는다.
특히 \(K_{\mathrm{proj}}^0\)의 bending block 자체를 physical thin-shell tangent라고 부르지 않는다.
이는 아래 normal-curvature local projection의 constant global-step Hessian이며,
tangential Hessian component는 frame target의 RHS에서 상쇄한다.
\(\nu=0\)의 planar quadratic in-plane, quadratic normal-curvature,
shallow-rest curvature와 superposed-rotation probe에서 operator action, 첫 8--16개 modal subspace와
nonlinear rollout이 사전 등록 tolerance를 통과할 때만 이 projective approximation을 사용한다.
실패하면 \(K_{\mathrm{proj}}^0\)를 teacher residual로 덮지 않고 state-dependent co-rotated tangent나
추가 projective constraint로 structural contract를 다시 설계한다.

Mass-proportional \(\alpha_MM\)은 rigid velocity도 감쇠하므로 internal damping으로 숨기지 않는다.
같은 이유로 target-rate가 없는 Rayleigh \(\alpha_KK\) damping도 canonical V0에 넣지 않는다.
필요하면 명시적인 world-drag channel \(C_{\mathrm{world}}\succeq0\)로 분리하고,
free-object objectivity test에서는 \(C_{\mathrm{world}}=0\)으로 둔다.
Anisotropic material은 후속 확장이며, 첫 논문에서 network가 arbitrary constitutive matrix를 예측하지 않는다.

\subsection{Hard/compliant attachment와 free-DOF operator}

Hard row와 compliant row는
\(\mathcal C_{\mathrm h}\cap\mathcal C_{\mathrm c}=\varnothing\)으로 분리한다.
Compliant attachment만

\begin{equation}
  K_{\mathrm{attach}}
  =A_{\mathrm c}^\top D_{\mathrm c}A_{\mathrm c},
  \qquad
  C_{\mathrm{attach}}
  =A_{\mathrm c}^\top V_{\mathrm c}A_{\mathrm c},
  \qquad
  D_{\mathrm c},V_{\mathrm c}\succeq0
\end{equation}

로 조립한다. Moving compliant target은 world-zero spring으로 숨기지 않고

\begin{equation}
\begin{aligned}
  E_{\mathrm c}(u,t)
  &=\frac12\|A_{\mathrm c}u-d_{\mathrm c}(t)\|_{D_{\mathrm c}}^2,\\
  f_{\mathrm c}(u,\dot u,t)
  &=-A_{\mathrm c}^\top D_{\mathrm c}
    (A_{\mathrm c}u-d_{\mathrm c})
    -A_{\mathrm c}^\top V_{\mathrm c}
    (A_{\mathrm c}\dot u-\dot d_{\mathrm c})
\end{aligned}
\end{equation}

로 정의한다. 따라서 inhomogeneous target RHS는
\(h_{\mathrm c}^t=A_{\mathrm c}^\top
(D_{\mathrm c}d_{\mathrm c}^t+V_{\mathrm c}\dot d_{\mathrm c}^t)\)이고,
같은 row의 \(K_{\mathrm{attach}},C_{\mathrm{attach}}\)와 정확히 한 번 짝을 이룬다.
Package의 \(\mathcal U_{\mathrm c}\)는 fixed-layout runtime handle control
\(u_{\mathcal C}(t)\)에서 이 \(d_{\mathrm c}(t),\dot d_{\mathrm c}(t)\)를 만든다.
Moving compliant control을 허용한 asset은 이 map을 package signature에 포함하며,
map이 없는 asset에서 임의의 moving target을 받지 않는다.
Hard attachment는 penalty stiffness를 더하지 않고
\(A_{\mathrm h}u=d_{\mathrm h}(t)\)를 만족하도록 제거한다.
V0의 hard row는 whole-anchor 또는 명시된 Cartesian component를 고르는
coordinate-selection constraint로 제한한다. 따라서 Boolean free injection \(N_{\mathrm h}\)에 대해

\begin{equation}
  A_{\mathrm h}A_{\mathrm h}^\top=I,\qquad
  A_{\mathrm h}N_{\mathrm h}=0,\qquad
  A_{\mathrm h}^\top A_{\mathrm h}
  +N_{\mathrm h}N_{\mathrm h}^\top=I
\end{equation}

가 성립한다. Typed lift와 dual free-force mask는

\begin{equation}
  u_{\mathrm p}(t)=A_{\mathrm h}^\top d_{\mathrm h}(t),
  \qquad
  u=u_{\mathrm p}(t)+N_{\mathrm h}y,
  \qquad
  P_{\mathrm{free}}:=N_{\mathrm h}N_{\mathrm h}^\top
  \label{eq:coordinate-hard-support-contract}
\end{equation}

로 고정한다. Diagonal anchor mass에서 \(N_{\mathrm h}^\top Mu_{\mathrm p}=0\)이고
\(P_{\mathrm{free}}\)는 displacement와 force covector 모두에 일관된 binary mask다.
일반 multipoint constraint를 쓰려면 null-space force coordinate
\(N_{\mathrm h}^\top f\)와 mass-dual lift를 별도로 정의해야 하므로 V0 범위에 넣지 않는다.

\begin{equation}
\begin{aligned}
  \bar K_{\mathrm{lin}}&=K_{\mathrm{lin}}^0+K_{\mathrm{attach}},
  &\bar K&=K_{\mathrm{proj}}^0+K_{\mathrm{attach}},\\
  \bar C_{\mathrm{lin}}&=C_{\mathrm{lin}}^0+C_{\mathrm{attach}}+C_{\mathrm{world}},
  &\bar C&=C_{\mathrm{proj}}^0+C_{\mathrm{attach}}+C_{\mathrm{world}},\\
  M_{\mathrm f}&=N_{\mathrm h}^\top MN_{\mathrm h},
  &K_{\mathrm{lin,f}}&=N_{\mathrm h}^\top\bar K_{\mathrm{lin}}N_{\mathrm h},\\
  C_{\mathrm{lin,f}}&=N_{\mathrm h}^\top\bar C_{\mathrm{lin}}N_{\mathrm h},\\
  K_{\mathrm f}&=N_{\mathrm h}^\top\bar KN_{\mathrm h},
  &C_{\mathrm f}&=N_{\mathrm h}^\top\bar CN_{\mathrm h}.
  \label{eq:hard-attachment-free-operators}
\end{aligned}
\end{equation}

Anchor-space fixed Hessian은 \(K^0:=\bar K\), \(C^0:=\bar C\)로 typed한다.
Prescribed hard motion이 만드는 operator-owned lift는

\begin{equation}
  f_{\mathrm p}^t
  :=-M\ddot u_{\mathrm p}^t-C^0\dot u_{\mathrm p}^t-K^0u_{\mathrm p}^t.
  \label{eq:hard-support-lift-force}
\end{equation}

로 두며 external/attachment force와 별도 lineage로 기록한다.

Eigenproblem과 mode/POD calibration은 \((M_{\mathrm f},K_{\mathrm{lin,f}})\)에서 만들고,
rest-linear reference solve는
\((M_{\mathrm f},K_{\mathrm{lin,f}},C_{\mathrm{lin,f}})\),
runtime reduced solve block은 \((M_{\mathrm f},K_{\mathrm f},C_{\mathrm f})\)에서 만든다.
Free-coordinate basis를 \(\widetilde\Phi_{\mathrm f},\widetilde\Psi_{\mathrm f}\)라 하면
anchor reconstruction과 force virtual work에는

\begin{equation}
  \Phi=N_{\mathrm h}\widetilde\Phi_{\mathrm f},
  \qquad
  \Psi=N_{\mathrm h}\widetilde\Psi_{\mathrm f},
  \qquad
  u=u_{\mathrm p}(t)+\Phi q+\Psi z
  \label{eq:hard-attachment-basis-lift}
\end{equation}

를 저장한다.
이후 reduced \(M,K,C\) block 표기는 free operator와 free basis의 projection을 줄여 쓴 것이며,
anchor-space \(\Phi,\Psi\)를 free matrix와 직접 곱하지 않는다.
식~\eqref{eq:hard-support-lift-force}와 hard reaction은 force ledger에 별도로 기록하며,
같은 attachment row를 penalty와 hard projection으로 동시에 적용하지 않는다.

\subsection{Operator hard checks와 dimensional scaling}

\(Z^0\)를 full-space rigid-motion basis라 하면 package는

\begin{equation}
\begin{gathered}
  M_{\mathrm f}\succ0,
  \qquad K_{\mathrm{lin,f}},C_{\mathrm{lin,f}},K_{\mathrm f},C_{\mathrm f}\succeq0,
  \qquad K_{\mathrm f}=K_{\mathrm f}^\top,
  \quad C_{\mathrm f}=C_{\mathrm f}^\top,\\
  K_{\mathrm{lin}}^0Z^0=0,
  \qquad C_{\mathrm{lin}}^0Z^0=0,
  \qquad
  \dim\ker K_{\mathrm{lin,f}}=d_{\mathrm{rigid}}^{\mathrm{expected}}
  \label{eq:operator-construction-hard-checks}
\end{gathered}
\end{equation}

를 검사한다.
Attachment가 모든 rigid mode를 제거한다고 가정하지 않고 실제 support layout에서
\(d_{\mathrm{rigid}}^{\mathrm{expected}}\)를 계산한다.
\(K_{\mathrm{proj}}^0\)는 local rotation target을 고정했을 때의 global-step Hessian이므로
그 matrix 단독에 rotational nullspace를 요구하지 않는다.
Canonical finite-motion hard check는 Hessian kernel 대신
식~\eqref{eq:projective-structural-energy} 전체의 rigid-motion zero energy,
strained-state superposed-rotation energy invariance와 force equivariance를 사용한다.
Volumetric-density branch에서 attachment를 끄거나 그 length/stiffness도 같은 물리 similarity로
scaling한 fixture에 한해
uniform similarity \(x\mapsto sx\), \(A\mapsto s^2A\), \(h\mapsto sh\)에서는

\begin{equation}
\begin{gathered}
  \ell\mapsto s\ell,
  \quad M\mapsto s^3M,
  \quad (B_{\mathrm m},\mathsf A_{\mathrm m})\mapsto
  s^{-1}(B_{\mathrm m},\mathsf A_{\mathrm m}),\\
  (B_{\mathrm b},\mathsf A_{\mathrm b})\mapsto
  s^{-2}(B_{\mathrm b},\mathsf A_{\mathrm b}),\qquad
  (K_{\mathrm{lin}}^0,K_{\mathrm{proj}}^0)\mapsto
  s(K_{\mathrm{lin}}^0,K_{\mathrm{proj}}^0),
  \qquad \omega_j\mapsto s^{-1}\omega_j
  \label{eq:operator-dimensional-scaling}
\end{gathered}
\end{equation}

을 unit test로 확인한다.
Fixed areal-density preset이면 \(M\mapsto s^2M\)이므로 별도 fixture와 expected law를 사용한다.
고정 compliant spring을 섞어 structural \(K\mapsto sK\) test를 모호하게 만들지 않는다.
Fixed geometry에서 anchor density만 바꿀 때에는 total area/mass invariance,
affine membrane patch, quadratic bending patch와 첫 8--16개 eigenfrequency convergence를 검사한다.
Network가 dense arbitrary \(M,C,K\)를 직접 출력하게 두지 않는다.

\subsection{V0 structural backend와 integrator freeze}

\reviewvthreepartial{V3 리뷰가 world-space rest-linear stiffness와 미확정 integrator를
구현 전에 하나로 고정하라는 진단은 수용한다.
다만 generic Newton corotation과 backward Euler를 측정 없이 정답으로 확정하지 않는다.
State-dependent tangent는 active-set factor를 매 frame 바꾸고,
backward Euler는 MC2가 복원하려는 high-frequency response와 PSD를 수치적으로 감쇠한다.
Canonical V0는 아래의
\emph{one-shot objective normal-curvature projective frozen target + implicit midpoint}로
고정한다. Rest-linear와 backward Euler는 각각 correctness와 robustness ablation일 뿐이며,
core asset별로 backend를 바꾸지 않는다. Step 0에서 이 canonical 조합이 사전 등록 tolerance를
통과하지 못하면 조용히 다른 backend를 고르는 대신 구현 계약을 재설계한다.}

식~\eqref{eq:objective-vector-projective-constraints}의
world-vector relation descriptor를 사용한다.
이전 최종 상태의 causal half-step geometry/frozen-velocity predictor
\((x^{\mathrm{pred},t-1/2},v^{\mathrm{pred},t-1/2})\)에서 predictor descriptor를

\begin{equation}
\begin{aligned}
  \bar h_{\alpha,c}^t
  &:=\mathsf A_{\alpha,c}x^{\mathrm{pred},t-1/2},\\
  \bar g_{\alpha,c}^t
  &:=\mathsf A_{\alpha,c}v^{\mathrm{pred},t-1/2}
  \label{eq:projective-predictor-descriptors}
\end{aligned}
\end{equation}

로 둔다. Bending Hessian까지 rotation fit에 넣으면 curvature를 줄이는 방향으로 normal이 기울 수 있으므로,
V0 relation frame은 oriented membrane/tangent descriptor에서만

\begin{equation}
  R_\alpha^t
  =\argmin_{R\in SO(3)}
  \left\|
  \bar h_{\alpha,\mathrm m}^t
  -(I_{m_{\mathrm m}}\otimes R)d_{\alpha,\mathrm m}^0
  \right\|_{\Omega_{\alpha,\mathrm m}}^2
  \label{eq:projective-rotation-fit}
\end{equation}

로 한 번 fit하고 한 frame 동안 고정한다.
Weighted det-positive polar/SVD를 사용하며, point-offset Procrustes는 rank 판정, 초기화 또는
동률 tie-break에만 쓰고 다른 목적함수의 rotation으로 대체하지 않는다.
정상적인 single-layer plate relation의 covariance는 rank 2이므로
full rank 3을 요구하지 않는다. Rest/current tangent rank가 2 이상이고 oriented normal의
sign consistency를 만족하면 허용하며, tangent rank가 2 미만이면 support를 넓히거나 relation을 거부한다.
Rotation derivative를 global solve에서 취하지 않는다.

\(m_{\mathrm m}=2,m_{\mathrm b}=3\)와
\(\mathsf R_{\alpha,c}^t=I_{m_c}\otimes R_\alpha^t\)를 쓰고,
normal-curvature projector를

\begin{equation}
\begin{aligned}
  \bar n_\alpha^t&:=R_\alpha^tn_\alpha^0,\\
  \mathsf P_{\alpha,n}^t
  &:={}I_{m_{\mathrm b}}\otimes
  (\bar n_\alpha^t(\bar n_\alpha^t)^\top),\\
  \mathsf P_{\alpha,t}^t&:=I_{3m_{\mathrm b}}-\mathsf P_{\alpha,n}^t,\\
  \mathsf P_{\alpha,n}^0
  &:={}I_{m_{\mathrm b}}\otimes(n_\alpha^0(n_\alpha^0)^\top),\\
  d_{\alpha,\mathrm b}^{0,n}
  &:=\mathsf P_{\alpha,n}^0d_{\alpha,\mathrm b}^0
  \label{eq:bending-normal-projectors}
\end{aligned}
\end{equation}

로 둔다. Flat rest에서는 quadratic reproduction에 의해
\(d_{\alpha,\mathrm b}^0\simeq d_{\alpha,\mathrm b}^{0,n}\simeq0\)이어야 하며 이를 hard check한다.
Shallow rest에서는 rotated rest-normal curvature \(d_{\alpha,\mathrm b}^{0,n}\)만 target으로 쓰고
rest descriptor의 tangential connection component를 curvature로 회전시키지 않는다.
World angular-velocity matrix는 positive tangent-kinematic metric의 rigid-spin fit으로 구한다.

\begin{equation}
\begin{aligned}
  \Xi_\alpha^t
  &:=
  \argmin_{\Xi\in\mathfrak{so}(3)}
  \left\|
  \bar g_{\alpha,\mathrm m}^t
  -(I_{m_{\mathrm m}}\otimes\Xi)\bar h_{\alpha,\mathrm m}^t
  \right\|_{\Omega_{\alpha,\mathrm m}}^2,\\
  \mathsf X_{\alpha,c}^t
  &:={}I_{m_c}\otimes\Xi_\alpha^t,\\
  y_{\alpha,\mathrm m}^{x,t}
  &:={}\mathsf R_{\alpha,\mathrm m}^td_{\alpha,\mathrm m}^0,\\
  y_{\alpha,\mathrm b}^{x,t}
  &:={}\mathsf P_{\alpha,t}^t\bar h_{\alpha,\mathrm b}^t
  +\mathsf R_{\alpha,\mathrm b}^td_{\alpha,\mathrm b}^{0,n},\\
  p_{\alpha,c}^{u,t}
  &:={}y_{\alpha,c}^{x,t}-d_{\alpha,c}^0,\\
  \bar r_{\alpha,\mathrm b}^{v,t}
  &:={}\bar g_{\alpha,\mathrm b}^t
  -\mathsf X_{\alpha,\mathrm b}^t\bar h_{\alpha,\mathrm b}^t,\\
  p_{\alpha,\mathrm m}^{v,t}
  &:={}\mathsf X_{\alpha,\mathrm m}^t\bar h_{\alpha,\mathrm m}^t,\\
  p_{\alpha,\mathrm b}^{v,t}
  &:={}\mathsf X_{\alpha,\mathrm b}^t\bar h_{\alpha,\mathrm b}^t
  +\mathsf P_{\alpha,t}^t\bar r_{\alpha,\mathrm b}^{v,t}
  \label{eq:projective-relation-targets}
\end{aligned}
\end{equation}

로 정의한다. 이는 relation마다 최대 \(3\times3\) normal equation이다.
Tangent rank가 2이면 rigid spin이 정의되며, 모든 \(\tau_c=0\)이면 force에 기여하지 않으므로
계산을 생략하고 \(\Xi_\alpha^t=0\)으로 둔다. Rank/finite-value가 실패하면 support를 넓히거나 relation을 거부한다.
따라서 rotation history나 matrix-log branch state가 필요하지 않으며,
\(p^{v,t}\)는 단순한 유한차분 \(\dot R\,d^0\)가 아니라 current relation descriptor의
best-fit rigid spin을 제거한다.
At predictor에서 bending stiffness residual은
\(\mathsf P_{\alpha,n}^t\bar h_{\alpha,\mathrm b}^t-
\mathsf R_{\alpha,\mathrm b}^td_{\alpha,\mathrm b}^{0,n}\),
damping residual은
\(\mathsf P_{\alpha,n}^t(\bar g_{\alpha,\mathrm b}^t-
\mathsf X_{\alpha,\mathrm b}^t\bar h_{\alpha,\mathrm b}^t)\)이므로
flat rest에서는 normal curvature와 그 relative rate만 남는다.

Frame-frozen structural energy와 relative-rate dissipation은

\begin{equation}
\begin{aligned}
  E_{\mathrm{str}}^t(u)
  =&\;\frac12\sum_{\alpha,c}
  \|\mathsf A_{\alpha,c}(x^0+u)
    -y_{\alpha,c}^{x,t}\|_{\Omega_{\alpha,c}}^2
  +\frac12\|A_{\mathrm c}u-d_{\mathrm c}^t\|_{D_{\mathrm c}}^2,\\
  \mathcal R_{\mathrm{str}}^t(\dot u)
  =&\;\frac12\sum_{\alpha,c}
  \|\mathsf A_{\alpha,c}\dot u-p_{\alpha,c}^{v,t}\|_{\Omega_{\alpha,c}^{\eta}}^2
  +\frac12\|A_{\mathrm c}\dot u-\dot d_{\mathrm c}^t\|_{V_{\mathrm c}}^2
  +\frac12\dot u^\top C_{\mathrm{world}}\dot u .
  \label{eq:projective-structural-energy}
\end{aligned}
\end{equation}

Frozen predictor target에 대한 arbitrary solve state의 bending residual은

\begin{equation}
\begin{aligned}
  r_{\alpha,\mathrm b}^{x,t}(x)
  ={}&\mathsf P_{\alpha,n}^t\mathsf A_{\alpha,\mathrm b}x
  -\mathsf R_{\alpha,\mathrm b}^td_{\alpha,\mathrm b}^{0,n}\\
  &+\mathsf P_{\alpha,t}^t
  (\mathsf A_{\alpha,\mathrm b}x-\bar h_{\alpha,\mathrm b}^t),\\
  r_{\alpha,\mathrm b}^{v,t}(v)
  ={}&\mathsf P_{\alpha,n}^t
  (\mathsf A_{\alpha,\mathrm b}v
   -\mathsf X_{\alpha,\mathrm b}^t\bar h_{\alpha,\mathrm b}^t)\\
  &+\mathsf P_{\alpha,t}^t
  (\mathsf A_{\alpha,\mathrm b}v-\bar g_{\alpha,\mathrm b}^t).
  \label{eq:normal-curvature-one-shot-residual}
\end{aligned}
\end{equation}

따라서 tangential term은 predictor state 또는 converged local--global projection에서 사라지지만,
one-shot final state에서는 predictor lag로 남을 수 있다. 이를 physical bending으로 숨기지 않고
splitting error로 측정한다. Tolerance를 넘으면 two-step/수렴 projection을 추가하거나
state-dependent projected tangent로 backend를 재설계한다.

따라서 Hessian \(K^0,C^0\)는 setup 이후 고정되고 target만

\begin{equation}
\begin{aligned}
  h_K^t&=\sum_{\alpha,c}\mathsf A_{\alpha,c}^\top
  \Omega_{\alpha,c}p_{\alpha,c}^{u,t}
  +A_{\mathrm c}^\top D_{\mathrm c}d_{\mathrm c}^t,\\
  h_C^t&=\sum_{\alpha,c}\mathsf A_{\alpha,c}^\top
  \Omega_{\alpha,c}^{\eta}p_{\alpha,c}^{v,t}
  +A_{\mathrm c}^\top V_{\mathrm c}\dot d_{\mathrm c}^t,\\
  M\ddot u+C^0\dot u+K^0u&=f_{\mathrm{ext}}^t+h_K^t+h_C^t
  \label{eq:canonical-structural-force-ledger}
\end{aligned}
\end{equation}

처럼 RHS로 바뀐다. Hard lift를 대입하면

\begin{equation}
  M_{\mathrm f}\ddot y+C_{\mathrm f}\dot y+K_{\mathrm f}y
  =N_{\mathrm h}^\top
  (f_{\mathrm{ext}}^t+h_K^t+h_C^t+f_{\mathrm p}^t).
  \label{eq:free-structural-force-ledger}
\end{equation}

이 식을 Global, Local과 corrector의 유일한 structural source of truth로 둔다.
\(h_K^t,h_C^t,f_{\mathrm p}^t\)는 unique relation/support graph에서 한 번 계산하는
operator-owned channel이며 patch proposal 또는 joint Schur에서 다시 조립하지 않는다.
Free object에서 relation derivative row의 translation nullspace는 frame-frozen solve에서도
net internal force를 0으로 만든다. Tangent rotation/spin fit의 stationarity는 membrane channel의
predictor torque를 제거한다. Normal-curvature bending target은 global rotation에 equivariant하지만,
one-shot global step은 frozen normal/target의 derivative를 취하지 않으므로 predictor와 final state의
bending torque residual을 별도로 측정한다. 사전 등록 tolerance를 넘으면 local--global iteration을
추가하거나 backend를 재설계하며, frame objectivity와 exact angular-momentum conservation을 같은
주장으로 취급하지 않는다.
Attachments와 world drag를 끈 상태에서 식~\eqref{eq:vector-constraint-equivariance}와
block-isotropic \(\Omega\) 때문에 상수 \(Q\in SO(3)\), translation \(a\)에 대해
\[
  \bar h'=(I_{m_c}\otimes Q)\bar h,
  \qquad
  \bar n'=Q\bar n,
  \qquad
  \mathsf P_n'=(I_{m_{\mathrm b}}\otimes Q)\mathsf P_n
  (I_{m_{\mathrm b}}\otimes Q^\top)
\]
이고 \(y^{x,t},p^{v,t}\)도 같은 block rotation으로 변환된다. 따라서 frame-frozen energy는
invariant이고 internal force도 \(f_{\mathrm{int}}(Qx+a;Q\mathcal Y)=
(I_K\otimes Q)f_{\mathrm{int}}(x;\mathcal Y)\)를 만족한다.
따라서 pure rigid motion뿐 아니라 \emph{strained state + superposed rigid rotation}을
hard test로 검사한다. 시간에 따라 변하는 superposed rigid motion \(Q(t)\)에 대해서도

\begin{equation}
\begin{aligned}
  \bar h'&=(I_{m_c}\otimes Q)\bar h,\\
  \bar g'&=(I_{m_c}\otimes Q)\bar g
    +(I_{m_c}\otimes\dot Q)\bar h,\\
  \Xi'&=\dot Q Q^\top+Q\Xi Q^\top,\\
  \bar g'-(I_{m_c}\otimes\Xi')\bar h'
  &=(I_{m_c}\otimes Q)
    [\bar g-(I_{m_c}\otimes\Xi)\bar h]
  \label{eq:relative-rate-objectivity}
\end{aligned}
\end{equation}

를 검사한다. 따라서 co-rotated rate target은 free rigid spin을 internal damping으로 감쇠시키지
않아야 하며, 이를 energy objectivity와 별도의 force/dissipation hard test로 둔다.
Rest-linear reference는 \(K_{\mathrm{lin}}^0,C_{\mathrm{lin}}^0\)를 쓰는 별도 backend이며
projective target을 0으로 만든 것과 같다고 주장하지 않는다.

\subsection{Rest-linear와 locally refit projective action closure gate}

\reviewvfourpartial{\(K_{\mathrm{lin}}\)에서 만든 Global basis가 projective runtime과
충분히 일치해야 한다는 진단은 수용한다. 다만 비교 대상을 local target을 고정한
global-step Hessian \(K_{\mathrm{proj}}^0\) 자체로 두지는 않는다. 그 matrix는 constant solve
Hessian이므로 rotational nullspace를 가질 필요가 없다. 비교 대상은 tangent rotation과
normal-curvature target을 다시 fit한 diagnostic energy, 실제 held-target force action과
그 rest action tangent다. Refit energy의 gradient와 held-target force를 같은 것으로 두지 않는다.}

Zero prescribed motion과 zero compliant-target offset에서 \(u=N_{\mathrm h}y\)라 둔다.
주어진 full state \(x\)에서 식~\eqref{eq:projective-rotation-fit}의
\(R_\alpha^\star(x)\)와 식~\eqref{eq:bending-normal-projectors}--
\eqref{eq:projective-relation-targets}의 target을 즉시 다시 계산한 map을
\(\mathcal Y_{\alpha,c}(x)\)라 한다. 구체적으로

\begin{equation}
\begin{aligned}
  \mathcal Y_{\alpha,\mathrm m}(x)
  &:={}(I_{m_{\mathrm m}}\otimes R_\alpha^\star(x))d_{\alpha,\mathrm m}^0,\\
  \mathcal Y_{\alpha,\mathrm b}(x)
  &:={}\mathsf P_{\alpha,t}^\star(x)\mathsf A_{\alpha,\mathrm b}x
  +(I_{m_{\mathrm b}}\otimes R_\alpha^\star(x))d_{\alpha,\mathrm b}^{0,n}.
  \label{eq:refit-projective-target-map}
\end{aligned}
\end{equation}

그러면

\begin{equation}
\begin{aligned}
  \mathcal E_{\mathrm{lin}}^{\mathrm f}(y)
  &:={}\frac12y^\top K_{\mathrm{lin,f}}y,\\
  \mathcal E_{\mathrm{proj,refit}}^{\mathrm f}(y)
  &:={}\frac12\sum_{\alpha,c}
  \left\|
    \mathsf A_{\alpha,c}(x^0+N_{\mathrm h}y)
    -\mathcal Y_{\alpha,c}(x^0+N_{\mathrm h}y)
  \right\|_{\Omega_{\alpha,c}}^2
  +\frac12\|A_{\mathrm c}N_{\mathrm h}y\|_{D_{\mathrm c}}^2,\\
  f_{\mathrm{lin}}(y)&:=-K_{\mathrm{lin,f}}y,\\
  f_{\mathrm{proj,held}}(y)
  &:={}-N_{\mathrm h}^\top
  \sum_{\alpha,c}\mathsf A_{\alpha,c}^\top\Omega_{\alpha,c}
  \left[
    \mathsf A_{\alpha,c}(x^0+N_{\mathrm h}y)
    -\mathcal Y_{\alpha,c}(x^0+N_{\mathrm h}y)
  \right]
  -N_{\mathrm h}^\top A_{\mathrm c}^\top D_{\mathrm c}A_{\mathrm c}N_{\mathrm h}y,\\
  K_{\mathrm{proj,act,f}}^0
  &:=-\left.\frac{\partial f_{\mathrm{proj,held}}}{\partial y}\right|_{y=0},
  \qquad
  K_{\mathrm{proj,sym,f}}^0
  :=\frac12(K_{\mathrm{proj,act,f}}^0+K_{\mathrm{proj,act,f}}^{0\top}).
  \label{eq:refit-projective-action-closure}
\end{aligned}
\end{equation}

\(f_{\mathrm{proj,held}}\)는 각 probe에서 target을 refit한 뒤 그 target derivative를 global force에서
취하지 않는 canonical runtime action이다. \(\mathcal Y_{\alpha,\mathrm b}\)는 full bending energy의
stationary minimizer가 아니므로 envelope theorem으로
\(-\nabla\mathcal E_{\mathrm{proj,refit}}\)와 동일시하지 않는다.
Action tangent는 target-refit finite difference 또는 검증된 analytic differentiation으로 계산하며
raw \(K_{\mathrm{proj}}^0\)로 대체하지 않는다.
Rigid-null probe를 제외한 small-strain planar in-plane/normal-curvature probe
\(\{y_j\}_{j=1}^{N_{\mathrm{probe}}}\)에 대해

\begin{equation}
\begin{aligned}
  \epsilon_E
  &:={}\max_j
  \frac{|\mathcal E_{\mathrm{proj,refit}}^{\mathrm f}(y_j)
         -\mathcal E_{\mathrm{lin}}^{\mathrm f}(y_j)|}
       {\mathcal E_{\mathrm{lin}}^{\mathrm f}(y_j)+\eps_E},\\
  \epsilon_f
  &:={}\max_j
  \frac{\|f_{\mathrm{proj,held}}(y_j)-f_{\mathrm{lin}}(y_j)\|_{M_{\mathrm f}^{-1}}}
       {\|f_{\mathrm{lin}}(y_j)\|_{M_{\mathrm f}^{-1}}+\eps_f}.
  \label{eq:klin-projective-action-error}
\end{aligned}
\end{equation}

Action tangent의 nonconservative splitting component와 symmetric stability를

\begin{equation}
  \epsilon_{\mathrm{skew}}
  :=\frac{\|K_{\mathrm{proj,act,f}}^0-K_{\mathrm{proj,act,f}}^{0\top}\|_F}
           {\|K_{\mathrm{proj,act,f}}^0\|_F+\eps_K},
  \qquad
  \lambda_{\min}(K_{\mathrm{proj,sym,f}}^0)\ge-\eps_{\mathrm{PSD}}
  \label{eq:projective-action-tangent-check}
\end{equation}

로 검사한다. 첫 non-rigid \(r\)개 \(M_{\mathrm f}\)-orthonormal mode를
\((K_{\mathrm{lin,f}},M_{\mathrm f})\)와
\((K_{\mathrm{proj,sym,f}}^0,M_{\mathrm f})\)에서 각각
\(V_{\mathrm{lin},r},V_{\mathrm{proj},r}\)라 하고
\(P_{\bullet,r}^{M}:=V_{\bullet,r}V_{\bullet,r}^\top M_{\mathrm f}\)라 두면

\begin{equation}
\begin{aligned}
  \epsilon_\angle
  &:={}\sin\theta_{\max}^{M_{\mathrm f}}
  (\operatorname{span}V_{\mathrm{lin},r},
   \operatorname{span}V_{\mathrm{proj},r}),\\
  \epsilon_{\mathrm{probe}}
  &:={}\max_j
  \frac{\|(P_{\mathrm{lin},r}^{M}-P_{\mathrm{proj},r}^{M})y_j\|_{M_{\mathrm f}}}
       {\|y_j\|_{M_{\mathrm f}}+\eps_P}.
  \label{eq:klin-projective-subspace-error}
\end{aligned}
\end{equation}

Projected descriptor-rate damping도 같은 closure를 별도로 통과해야 한다.
Zero target velocity의 rest state에서 small velocity probe \(w\)에 대해
식~\eqref{eq:projective-relation-targets}의 rigid-spin fit과 rate target을 다시 계산한 값을
\(p_{\alpha,c}^{v,\star}(w)\)라 두고

\begin{equation}
\begin{aligned}
  f_{\mathrm{damp,held}}(w)
  &:={}-N_{\mathrm h}^\top\Bigg[
  \sum_{\alpha,c}\mathsf A_{\alpha,c}^\top\Omega_{\alpha,c}^{\eta}
  \left(\mathsf A_{\alpha,c}N_{\mathrm h}w-p_{\alpha,c}^{v,\star}(w)\right)\\
  &\hspace{8.0em}
  +A_{\mathrm c}^\top V_{\mathrm c}A_{\mathrm c}N_{\mathrm h}w
  +C_{\mathrm{world}}N_{\mathrm h}w\Bigg],\\
  C_{\mathrm{proj,act,f}}^0
  &:=-\left.\frac{\partial f_{\mathrm{damp,held}}}{\partial w}\right|_{w=0},
  \qquad
  C_{\mathrm{proj,sym,f}}^0
  :=\frac12(C_{\mathrm{proj,act,f}}^0+C_{\mathrm{proj,act,f}}^{0\top}),\\
  \epsilon_C
  &:={}\max_j
  \frac{\|f_{\mathrm{damp,held}}(w_j)+C_{\mathrm{lin,f}}w_j\|_{M_{\mathrm f}^{-1}}}
       {\|C_{\mathrm{lin,f}}w_j\|_{M_{\mathrm f}^{-1}}+\eps_C},\\
  \epsilon_{C,\mathrm{skew}}
  &:={}\frac{\|C_{\mathrm{proj,act,f}}^0-C_{\mathrm{proj,act,f}}^{0\top}\|_F}
       {\|C_{\mathrm{proj,act,f}}^0\|_F+\eps_{C,K}}.
  \label{eq:rest-linear-projective-damping-closure}
\end{aligned}
\end{equation}

여기서 \(w_j\)는 rigid velocity를 제외한 membrane/normal-curvature rate probe다.
\(C_{\mathrm{proj,act,f}}^0\)는 raw fixed \(C_{\mathrm{proj}}^0\)가 아니라
rate target을 refit한 실제 held damping action의 tangent다.

V0 package는 development split에서 probe amplitude와 threshold를 고정한 뒤

\begin{equation}
\begin{gathered}
  \epsilon_E\le\tau_E,\qquad
  \epsilon_f\le\tau_f,\qquad
  \epsilon_{\mathrm{skew}}\le\tau_{\mathrm{skew}},\qquad
  \lambda_{\min}(K_{\mathrm{proj,sym,f}}^0)\ge-\eps_{\mathrm{PSD}},\\
  \epsilon_\angle\le\tau_\angle,\qquad
  \epsilon_{\mathrm{probe}}\le\tau_P,\qquad
  \epsilon_C\le\tau_C,\qquad
  \epsilon_{C,\mathrm{skew}}\le\tau_{C,\mathrm{skew}},\qquad
  \lambda_{\min}(C_{\mathrm{proj,sym,f}}^0)\ge-\eps_{C,\mathrm{PSD}}
\end{gathered}
  \label{eq:klin-projective-hard-gate}
\end{equation}

를 모두 통과해야 한다. 실패하면 learned residual로 차이를 덮지 않고 projective constraint,
local projection 또는 structural backend를 재설계한다.
여기서 \(\eps_{\mathrm{PSD}}\)와 \(\eps_{C,\mathrm{PSD}}\)는 각각 stiffness와 damping
단위를 가진 typed tolerance로 저장하거나, 명시한 reference scale로 무차원화한 뒤
asset 간 비교에 사용한다.

\reviewvfiverebut{Closure 실패 뒤의 fallback을 구체화하라는 진단은 수용한다.
그러나 모든 실패에 ``\(K_{\mathrm{proj,sym}}^0\) basis 재구성 \(\rightarrow\)
two-step/converged projection \(\rightarrow\) state-dependent tangent''를 무조건 적용하는
고정 순서는 채택하지 않는다. Skew/PSD, one-shot splitting과 basis mismatch는 서로 다른
실패이며, 유효하지 않은 action tangent의 eigenbasis는 fallback이 될 수 없다.}

식~\eqref{eq:klin-projective-hard-gate}는 두 판정을 합친 canonical acceptance gate다.
첫째, action-validity 판정은 rigid/objectivity, held-force action, stiffness/damping tangent의
skew와 PSD를 검사한다. 둘째, rest-linear-agreement 판정은 energy/force probe,
modal angle과 common-probe subspace가 rest-linear reference와 합의하는지 검사한다.
Canonical V0 package는 두 판정을 모두 통과해야 한다. 이 구분은 실패 원인을 찾기 위한 것이며,
한쪽 실패를 다른쪽 basis 선택으로 지우기 위한 우회로가 아니다. Artifact에 따라 다음
conditional triage를 고정한다.

\begin{enumerate}[leftmargin=1.7em]
  \item Skew/PSD, rigid/objectivity 또는 held-force action 자체가 실패하면 basis를 바꾸지 않고
  relation projector, projective target 또는 structural backend를 먼저 재설계한다.
  \item Action-validity는 통과하지만 rest-linear modal angle/common-probe가 실패하면
  canonical package와 MC1의 rest-linear operator/eigenspace agreement 주장은 실패로 판정한다.
  \((K_{\mathrm{proj,sym,f}}^0,M_{\mathrm f})\) Global basis는 별도의 versioned pivot experiment로만
  재구성할 수 있으며, 기존 hard gate를 통과한 것으로 간주하지 않는다. 이 branch를 채택하려면
  \(\Phi\), 모든 \(\Psi_r\), complement/cross block과 factor cache를 다시 만들고,
  Gate A의 증거를 runtime held-action/subspace fidelity와 downstream rollout으로 재정의하며
  rest-linear agreement claim의 철회를 명시한다.
  \item Rest-action closure는 통과하지만 one-shot 대비 two-step/converged reference의
  torque, phase, PSD 또는 timestep error만 지배적이면 더 수렴한 local projection을 runtime 후보로
  올려 정확도--비용을 다시 측정한다.
  \item 원인별 조치 뒤에도 실패하면 versioned state-dependent co-rotated tangent backend로
  전환하고 tangent assembly, factor invalidation과 p50/p95/p99를 새 계약으로 재검증한다.
\end{enumerate}

따라서 이는 무조건적인 실행 순서가 아니라 failure-mode별 decision tree다.

Canonical integrator는 implicit midpoint다. 한 frame의 aerodynamic force,
projective rotation/rate target, compliant target와 hard-support lift는
식~\eqref{eq:causal-midpoint-geometry-predictor}의 causal half-step predictor 및
같은 \(t-\tfrac12\) control sample에서 한 번 평가해
\(F_y^{t-1/2}:=F_{y,\mathrm{held}}^t\)로 고정한다.
Global--Local structural cross load만은 causal predicted-midpoint 값을 Global predictor가 먼저 소비하고,
Local update 뒤 actual endpoint 평균으로 만든 midpoint 값과의 차이를 Module 11이 한 번 보정한다.
따라서 이 coupling channel을 predictor에서 끝까지 고정했다고 표현하지 않는다.
Global \(q\) 또는 active Local \(z_{\mathcal S}\)를 일반화해 \(y\)라 쓰면

\begin{equation}
\begin{aligned}
  A_y^{\mathrm{MP}}
  &:=\frac{4}{\Delta t^2}M_y+\frac{2}{\Delta t}C_y+K_y,\\
  A_y^{\mathrm{MP}}\Delta y^t
  &=2F_{y,\mathrm{held}}^t-2K_yy^{t-1}
    +\frac{4}{\Delta t}M_y\dot y^{t-1},\\
  \dot y^t&=\frac{2}{\Delta t}\Delta y^t-\dot y^{t-1}.
  \label{eq:midpoint-effective-update}
\end{aligned}
\end{equation}

Backward Euler robustness reference는 같은 held load로

\begin{equation}
\begin{aligned}
  A_y^{\mathrm{BE}}
  &:=\frac{1}{\Delta t^2}M_y+\frac{1}{\Delta t}C_y+K_y,\\
  A_y^{\mathrm{BE}}\Delta y^t
  &=F_{y,\mathrm{held}}^t-K_yy^{t-1}
    +\frac{1}{\Delta t}M_y\dot y^{t-1},\\
  \dot y^t&=\Delta y^t/\Delta t.
  \label{eq:backward-euler-reference-update}
\end{aligned}
\end{equation}

를 사용한다. Undamped linear mode에서

\begin{equation}
  |g_{\mathrm{BE}}(\omega\Delta t)|
  =\frac{1}{\sqrt{1+(\omega\Delta t)^2}}<1,
  \qquad |g_{\mathrm{MP}}(\omega\Delta t)|=1
  \label{eq:integrator-modal-amplification}
\end{equation}

이다. Corrected endpoint aerodynamic delta를 midpoint-held load의 delta로 한 번 적용하는 것은
partitioned first-order correction이며, monolithic 또는 두 번 반복한 reference와 오차를 보고한다.
Canonical manifest는

\begin{equation}
\begin{aligned}
  \sigma_{\mathrm{struct}}:=\bigl(&
  \texttt{objective\_normal\_curvature\_PD},
  \texttt{implicit\_midpoint},\\
  &\Delta t,
  \texttt{operator\_and\_target\_version}\bigr)
  \label{eq:v0-structural-freeze-signature}
\end{aligned}
\end{equation}

를 모든 core asset에서 공유한다. Fixed Hessian이므로 factor-cache exact reuse는
operator/target version, \(\Delta t\)와 active-set signature가 같을 때 허용된다.
이 backend도 fixed reduced basis의 표현력 부족을 해결하지 않으므로
Global/POD rank sweep과 common-probe projection error를 별도로 통과해야 한다.

\subsection{Global basis}

Global basis는 generalized eigenproblem 또는 teacher POD와 physics mode의 hybrid에서 얻는다.

\begin{equation}
  K_{\mathrm{lin,f}}\widetilde\phi_j
  =\lambda_jM_{\mathrm f}\widetilde\phi_j,
  \qquad
  \widetilde\Phi_{\mathrm f}
  =[\widetilde\phi_1,\ldots,\widetilde\phi_{r_g}],
  \qquad
  \Phi=N_{\mathrm h}\widetilde\Phi_{\mathrm f},
  \qquad
  \Phi^\top M\Phi=I.
\end{equation}

Low-frequency bending, sway, twist와 attachment-driven motion을 우선한다.
Teacher POD를 사용할 때도 \(M\)-orthonormalization과 rest/attachment consistency를 적용한다.

\subsection{Patch Local basis와 complement}

Patch \(r\)의 raw local modes를 \(B_r\)라 하고,
\(R_r\)을 global anchor에서 patch anchor를 추출하는 operator,
\(W_r\)을 core--halo partition weight라 한다.
Patch raw basis를 global anchor space에 lift하되 patch ID를 유지한다.
먼저 hard-null space의 \(M\)-orthogonal projector

\begin{equation}
  P_{\mathrm h}^M
  =N_{\mathrm h}M_{\mathrm f}^{-1}N_{\mathrm h}^\top M
\end{equation}

를 적용한다.

\begin{equation}
  U_r
  =
  P_{\mathrm h}^M R_r^\top W_rB_r.
\end{equation}

따라서 \(A_{\mathrm h}U_r=0\)이며, Global complement가 hard constraint를 다시 깨뜨리지 않는다.

Global span 제거 operator는

\begin{equation}
  P_\perp
  =
  I-\Phi(\Phi^\top M\Phi)^{-1}\Phi^\top M.
\end{equation}

각 patch에서 Global span만 제거한다.

\begin{equation}
  \overline B_r=P_\perp U_r.
\end{equation}

Patch Gram eigendecomposition과 rank-revealing truncation 뒤
patch 내부 retained column끼리만 정규화한다.

\begin{equation}
  \overline B_r^\top M\overline B_r
  =V_r\Lambda_rV_r^\top,
  \qquad
  \Psi_r
  =
  \overline B_rV_{r,+}\Lambda_{r,+}^{-1/2},
  \qquad
  \Psi_r^\top M\Psi_r=I,
  \qquad
  \Phi^\top M\Psi_r=0.
  \label{eq:patch-local-basis}
\end{equation}

Free-coordinate patch basis는
\(\widetilde\Psi_{\mathrm f,r}
=M_{\mathrm f}^{-1}N_{\mathrm h}^\top M\Psi_r\)로 저장하며
\(\Psi_r=N_{\mathrm h}\widetilde\Psi_{\mathrm f,r}\)를 hard check한다.

여기서 \(+\)는 \(\lambda_i>\tau_{\mathrm{rank}}\)인 retained eigenspace만 뜻한다.

서로 다른 patch를 한 번에 \(M\)-QR/SVD하면 tip과 attachment basis가 섞여
gate owner, patch energy와 실제 compute unit이 사라질 수 있으므로 금지한다.
다른 patch 사이에는 \(\Psi_r^\top M\Psi_s=0\)를 요구하지 않는다.
여기서 locality는 patch ID, source support와 compute ownership의 locality다.
\(P_\perp\)가 낮은 진폭의 Global-mode tail을 만들 수 있으므로 엄밀한 compact support를 주장하지 않는다.

\reviewaccept{Global complement가 patch identity를 유지하더라도 선언된 core--halo support 밖에
작은 motion tail을 만들 수 있다는 지적을 수용한다. Compact support를 문구로만 부정하는 데서 그치지 않고,
package validation에서 support 밖의 mass-weighted tail을 정량화한다.}

Patch \(r\)의 선언된 owner support에 속하는 free vector DOF를
\(\mathcal I_r^{\mathrm{own}}\)이라 하고, 그 unique DOF를 full anchor space로 lift하는 matrix를
\(E_r\in\R^{3K\times n_r^{\mathrm{own}}}\)라 한다. Consistent mass에서도 support decomposition이
명확하도록 ordinary binary mask가 아니라 다음 \(M\)-orthogonal projector를 사용한다.

\begin{equation}
  \Pi_r^M
  =
  E_r(E_r^\top M E_r)^{-1}E_r^\top M,
  \qquad
  \|X\|_{M,F}^2
  :=
  \operatorname{tr}(X^\top M X).
\end{equation}

Retained rank가 0보다 큰 patch에 대해 complement tail ratio를

\begin{equation}
  \rho_{r,\mathrm{tail}}
  =
  \frac{
    \|(I-\Pi_r^M)\Psi_r\|_{M,F}^2
  }{
    \|\Psi_r\|_{M,F}^2+\eps_{\mathrm{tail}}
  },
  \qquad
  0\le \rho_{r,\mathrm{tail}}\le 1
  \label{eq:patch-complement-tail}
\end{equation}

로 저장한다. Lumped diagonal \(M\)에서는 \(\Pi_r^M\)가 support mask로 환원된다.
Patch를 하나의 atomic group \(\mathcal G\)로 병합한 경우에는 support union과
\(\Psi_{\mathcal G}\)로 같은 static 값을 \(\rho_{\mathcal G,\mathrm{tail}}\)로 다시 계산한다.
Rollout에서 실제로 excitation된 motion의 tail도

\begin{equation}
  \rho_{\mathcal G,\mathrm{roll}}
  =
  \frac{
    \sum_t\|(I-\Pi_{\mathcal G}^M)\Psi_{\mathcal G}z_{\mathcal G}^t\|_M^2
  }{
    \sum_t\|\Psi_{\mathcal G}z_{\mathcal G}^t\|_M^2+\eps
  }
\end{equation}

로 보조 보고한다.

\reviewvthreeaccept{\(\tau_{\mathrm{tail}}\)을 asset과 load에 무관한 보편 상수로 두지 않는다.
Tail 자체의 크기만으로 basis를 기각하지 않고, 동일 asset family와 wind/load program에서
downstream trajectory와 spectrum에 미치는 영향을 development split에서 calibration한 뒤
held-out split에서 검증한다. 관계가 약하거나 비단조이면 hard threshold를 사용하지 않고
tail distribution과 downstream curve만 보고한다.}

Package decision의 단위 \(\mathcal U\)는 singleton patch 또는 사전 병합한 atomic group이다.
확장 halo로 만든 동일 retained-rank reference basis를
\(\Psi_{\mathcal U}^{\mathrm{wide}}\)라 하고,
asset family \(a\), wind/load program \(w\)의 rollout degradation을

\begin{equation}
  \Delta e_{\mathrm{tail}}^{a,w}
  =
  \left[
  e_{\mathrm{roll}}^{a,w}(\Psi_{\mathcal U})
  -e_{\mathrm{roll}}^{a,w}(\Psi_{\mathcal U}^{\mathrm{wide}})
  \right]_+
  \label{eq:tail-downstream-degradation}
\end{equation}

로 정의한다.
Development split에서
\((\rho_{\mathcal U,\mathrm{tail}},\rho_{\mathcal U,\mathrm{roll}},
\Delta e_{\mathrm{tail}}^{a,w},e_{\mathrm{PSD}}^{a,w},
\kappa(M_{\mathcal SS}))\)의 scatter와 conditional quantile을 보고한다.
Static package decision에는 같은 통계량으로 calibration한 asset-family threshold만 사용한다.
사용 가능한 경우에만

\begin{equation}
\begin{aligned}
  \tau_{\mathrm{tail,static}}^{a}
  =\sup\biggl\{\tau\in[0,1]:\;&
  \max_{w\in\mathcal W_{\mathrm{dev}}}
  \operatorname{Quantile}_{0.95}^{\mathrm{dev}}
  \left(
  \Delta e_{\mathrm{tail}}^{a,w}
  \mid \rho_{\mathcal U,\mathrm{tail}}\le\tau
  \right)
  \le\epsilon_{\mathrm{tail}}^{a},\\
  &
  \min_{w\in\mathcal W_{\mathrm{dev}}}
  n_{\mathrm{dev}}^{a,w}
  (\rho_{\mathcal U,\mathrm{tail}}\le\tau)
  \ge n_{\min}
  \biggr\}
  \label{eq:tail-calibrated-threshold}
\end{aligned}
\end{equation}

를 development-only diagnostic threshold로 고정한다.
Feasible set이 비어 있거나 conditional sample이 \(n_{\min}\)보다 작으면 threshold는
\emph{undefined}로 기록하고 package를 자동 accept/reject하지 않는다.
모든 sample이 통과해도 domain을 \([0,1]\)로 제한하므로 threshold는 최대 1이다.
Held-out wind에서는 threshold를 다시 맞추지 않고 false accept/reject와 downstream error를 보고한다.
\(\rho_{\mathcal U,\mathrm{roll}}\)은 load-dependent 진단으로만 사용하고
static threshold를 적용하지 않는다.

\reviewpartial{식~\eqref{eq:tail-calibrated-threshold}이 정의되고
\(\rho_{\mathcal U,\mathrm{tail}}>\tau_{\mathrm{tail,static}}^{a}\)이면 support/halo 확대,
overlap-group 병합 또는 support-constrained basis를 검토한다. 그러나 기존 \(\Psi_r\)를 사후 hard-mask하는
처방은 \(\Phi^\top M\Psi_r=0\)를 다시 깨뜨릴 수 있으므로 채택하지 않는다.
Tail이 실제 conditioning, patch ownership 또는 rollout을 해칠 때만 basis construction을 바꾸고,
정확한 Global complement 계약은 유지한다.}

강한 overlap으로 full superset Gram matrix가 rank deficient하거나 ill-conditioned하면
overlap-connected patch를 하나의 gating group으로 묶거나
patch ownership metadata를 보존하는 block-aware pivoting으로 중복 column을 제거한다.

V0는

\begin{equation}
  \Psi_{\mathrm{sup}}
  =
  [\Psi_1,\ldots,\Psi_P]
\end{equation}

와 각 patch의 고정 coordinate range를 object package에 저장한다.
Runtime에는 \(\mathcal S_t=\mathcal A_t\cup\mathcal D_t\)의 principal block만
gather--solve--scatter한다. 이는 dynamic basis recomputation이 아니라 고정 slot 선택이므로
기존 coordinate를 재투영하는 state remapping이 필요 없다.
Patch force는 여전히 먼저 shared anchor space에 조립하고,
선택된 모든 block을 포함한 하나의 coupled Local system에서 한 번만 적분한다.

\subsection{Reduced block operators}

\begin{equation}
  M_{gg}=\widetilde\Phi_{\mathrm f}^\top
  M_{\mathrm f}\widetilde\Phi_{\mathrm f},
  \qquad
  M_{g r}=\widetilde\Phi_{\mathrm f}^\top
  M_{\mathrm f}\widetilde\Psi_{\mathrm f,r}=0,
\end{equation}

\begin{equation}
  K_{gg}=\widetilde\Phi_{\mathrm f}^\top
  K_{\mathrm f}\widetilde\Phi_{\mathrm f},
  \qquad
  C_{gg}=\widetilde\Phi_{\mathrm f}^\top
  C_{\mathrm f}\widetilde\Phi_{\mathrm f}.
\end{equation}

\begin{equation}
  M_{rs}=\widetilde\Psi_{\mathrm f,r}^\top
  M_{\mathrm f}\widetilde\Psi_{\mathrm f,s},
  \quad
  K_{rs}^{0}=\widetilde\Psi_{\mathrm f,r}^\top
  K_{\mathrm f}\widetilde\Psi_{\mathrm f,s},
  \quad
  C_{rs}=\widetilde\Psi_{\mathrm f,r}^\top
  C_{\mathrm f}\widetilde\Psi_{\mathrm f,s}.
  \label{eq:patch-cross-blocks}
\end{equation}

활성/decay 집합에 대해

\begin{equation}
  M_{\ell\ell}^{\mathcal S_t}=[M_{rs}]_{r,s\in\mathcal S_t},
  \quad
  K_{\ell\ell}^{0,\mathcal S_t}=[K_{rs}^{0}]_{r,s\in\mathcal S_t},
  \quad
  C_{\ell\ell}^{\mathcal S_t}=[C_{rs}]_{r,s\in\mathcal S_t}
\end{equation}

를 구성한다. Canonical objective normal-curvature projective backend는 fixed Hessian \(K^0\)와
frame target RHS를 사용하므로 이 block도 setup 이후 고정된다.
\(M_{rs},K_{rs}^0,C_{rs}\)를 버리고 patch별 diagonal block만 푸는 것은
overlap coupling을 무시하므로 mainline이 아니다.
Pre-factorized diagonal/group block은 full cross-block matvec를 유지하는 iterative backend의
preconditioner로는 사용할 수 있다.
Reachable한 모든 active/decay group에서 \(M_{\ell\ell}^{\mathcal S_t}\)의 rank와 condition number,
block symmetry/reciprocity를 검사한다.

\(K_{gr}^{0}=\widetilde\Phi_{\mathrm f}^\top
K_{\mathrm f}\widetilde\Psi_{\mathrm f,r}\)와
\(C_{gr}=\widetilde\Phi_{\mathrm f}^\top
C_{\mathrm f}\widetilde\Psi_{\mathrm f,r}\)가 0이 아니면
Local state의 structural reaction이 Global에 전달된다.
Symmetric operator에서는
\(K_{rg}^{0}=(K_{gr}^{0})^\top\),
\(C_{rg}=C_{gr}^\top\)이고,
\(K_{g,\mathrm{sup}}^{0}=[K_{g1}^{0},\ldots,K_{gP}^{0}]\),
\(C_{g,\mathrm{sup}}=[C_{g1},\ldots,C_{gP}]\)를 package에 저장한다.
Rest-linear eigenbasis에서 \(K_{\mathrm{lin},gr}=0\)이어도 canonical
normal-curvature projective \(K_{gr}^{0}\)가 0이라는 뜻은 아니다.
Runtime reaction은 저장된 canonical cross block의 실제 값만 사용하며 0인 항을 억지로 만들지 않는다.
Projective rotation target은 식~\eqref{eq:canonical-structural-force-ledger}의
operator-owned RHS로 별도 투영하며 cross block에 중복 삽입하지 않는다.

\subsection{Gaussian binding}

각 Gaussian은 소수 anchor의 affine MLS support에 고정된다.

\begin{equation}
  \mathcal B_j
  =
  \{(k,w_{jk},\xi_{jk}^0)\}_{k\in\mathcal N_j}.
\end{equation}

Binding은 identity, translation, rigid rotation,
in-plane affine deformation을 tolerance 내에서 재현해야 한다.
Setup 후 frame마다 neighbor를 다시 고르지 않는다.

\section{Oracle selector label과 조건부 F/G supervision contract}

\subsection{Teacher state projection}

Teacher anchor displacement \(u^{*,t}\)에서 Global coordinate를 얻는다.

\begin{equation}
  q^{*,t}
  =
  (\Phi^\top M\Phi)^{-1}\Phi^\top M
  (u^{*,t}-u_{\mathrm p}^t).
\end{equation}

Global reconstruction을 제거한 residual을 Local complement에 투영한다.

\begin{equation}
  z_{\mathrm{sup}}^{*,t}
  =
  (\Psi_{\mathrm{sup}}^\top M\Psi_{\mathrm{sup}})^{-1}
  \Psi_{\mathrm{sup}}^\top M
  (u^{*,t}-u_{\mathrm p}^t-\Phi q^{*,t}).
\end{equation}

이 projection은 full fixed-superset Gram matrix가 validation을 통과한 asset에서만 사용한다.
Rank/conditioning test가 실패하면 강한 overlap group 안에서 중복 mode를 제거한 뒤 package를 다시 만든다.

\subsection{Missing-force label}

Known local physics가 teacher \(z^*\)를 만들기 위해 필요로 하는 force를 inverse dynamics로 계산한다.

\begin{equation}
\begin{aligned}
  \Delta Q_{\ell}^{*,t}
  =&\;
  M_{\ell\ell}\ddot z^{*,t}
  +C_{\ell\ell}\dot z^{*,t}
  +K_{\ell\ell}^{0}z^{*,t}\\
  &+C_{\ell g}\dot q^{*,t}
  +K_{\ell g}^{0}q^{*,t}
  -\Psi_{\mathrm{sup}}^\top
  \left(
    f_{\mathrm{aero}}^{\mathrm{base},*,t}
    +f_{\mathrm{gravity}}^t
    +f_{\mathrm{other}}^t
    +h_K^t+h_C^t+f_{\mathrm p}^t
  \right).
  \label{eq:missing-force-label}
\end{aligned}
\end{equation}

식~\eqref{eq:missing-force-label}은 teacher total acceleration을 그대로 residual로 쓰지 않는다.
Mass, stiffness, damping, gravity, analytic aero와 attachment operator contribution 및
typed target/lift \((h_K,h_C,f_{\mathrm p})\)가 이미 설명한 부분을 모두 한 번씩 제거한다.
Teacher가 force breakdown을 직접 기록하면 이를 우선 사용한다.
Position만 있다면 smoothed differentiation과 differentiable rollout loss를 같이 사용한다.

\subsection{Generalized force와 patch anchor-force label의 연결}

식~\eqref{eq:missing-force-label}은 Local generalized force를 정의하지만,
Module 7 network와 Module 10 assembly는 patch anchor-force field를 요구한다.
두 공간 사이의 null space를 임의로 network에 떠넘기지 않고 다음 label contract를 사용한다.

Teacher mesh displacement와 anchor displacement의 interpolation을

\begin{equation}
  u_{\mathcal M}
  =
  \mathcal B_{A\rightarrow\mathcal M}u_A
\end{equation}

라고 하면 virtual work를 보존하는 mesh-to-anchor force transfer는

\begin{equation}
  T_{\mathcal M\rightarrow A}
  =
  \mathcal B_{A\rightarrow\mathcal M}^{\top}
\end{equation}

이다. Teacher가 nodal force breakdown을 제공하는 기본 경로에서는 channel별로

\begin{equation}
  \Delta f_{A,c}^{*,t}
  =
  T_{\mathcal M\rightarrow A}f_{\mathcal M,c}^{\mathrm{teacher},t}
  -f_{A,c}^{\mathrm{owned},t},
  \qquad
  c\in\{\mathrm{missing\ aero},\mathrm{missing\ structural}\}
  \label{eq:teacher-anchor-force-label}
\end{equation}

을 만든다. 여기서 \(f_{A,c}^{\mathrm{owned}}\)는 analytic aero, reduced structural base처럼
force ledger가 이미 소비한 동일 channel 항이다.
Hard coordinate-selection layout에서는 Local excitation과 support-row teacher force를 먼저

\begin{equation}
\begin{aligned}
  \Delta f_{A,c}^{*,\mathrm{free},t}
  &:={}P_{\mathrm{free}}\Delta f_{A,c}^{*,t},\\
  \rho_{h,c}^{*,\mathrm{teacher},t}
  &:={}A_{\mathrm h}\Delta f_{A,c}^{*,t}
  \label{eq:teacher-hard-free-force-split}
\end{aligned}
\end{equation}

로 분리한다. Missing-structural channel의 \(\rho_h^{*,\mathrm{teacher}}\)는
object-to-support intrinsic companion reference이고, missing-aero hard-row component는
support-reaction-only 진단값이다. 어느 쪽도 patch Local excitation head의 hard-row target으로 넣지 않는다.
Global virtual-work component와 그 minimum-\(M^{-1}\)-norm force representative는

\begin{equation}
  \Delta Q_{g,c}^{*,t}=\Phi^\top\Delta f_{A,c}^{*,\mathrm{free},t},
  \qquad
  \Delta f_{g,c}^{*,t}
  =M\Phi(\Phi^\top M\Phi)^{-1}\Delta Q_{g,c}^{*,t}.
\end{equation}

Force-label metric과 optional smoother의 source of truth는

\begin{equation}
  W_f\in\mathbb S_{++}^{3K},
  \qquad
  W_{f,t}:=E_{f,t}^\top W_fE_{f,t}\in\mathbb S_{++}^{n_t},
  \qquad
  L_f\in\R^{m_L\times3K}
  \label{eq:force-label-metric-smoother}
\end{equation}

로 둔다. \(W_f\)는 package-versioned anchor block-diagonal force metric이고,
식~\eqref{eq:support-restricted-operators}의 free support restriction을 runtime과 teacher가 함께 쓴다.
\(L_f\)는 accepted material-relation graph의 owner-aware force-gradient/Laplacian이며,
hard row, attachment boundary 또는 close-layer rejected relation을 가로질러 smoothing하지 않는다.
Boundary에서는 명시한 one-sided stencil만 사용하고 graph/stencil version을 label metadata에 저장한다.

먼저 \(\bar f_{\ell,c}^{*,t}=\Delta f_{A,c}^{*,\mathrm{free},t}-\Delta f_{g,c}^{*,t}\)와
그 Local remainder의 보존 target
\(b_{\ell,c}^{*,t}=A_{\mathrm{cons},c}^t\bar f_{\ell,c}^{*,t}\)를 만든다.
그 다음 complement와 wrench/constraint를 순차 projection하지 않고 하나의 joint KKT로 푼다.

\begin{equation}
\begin{aligned}
  \Delta f_{\ell,c}^{*,t}
  =\argmin_f\;&
  \frac12\|f-\bar f_{\ell,c}^{*,t}\|_{W_f}^2
  +\frac{\lambda_{f,\mathrm{smooth}}}{2}\|L_ff\|_2^2\\
  \mathrm{s.t.}\;&
  \Phi^\top f=0,\\
  &A_{\mathrm{cons},c}^tf=b_{\ell,c}^{*,t}.
  \label{eq:teacher-local-force-label}
\end{aligned}
\end{equation}

이 joint solve는 Global span 제거 뒤 force/torque 보정을 다시 적용해
첫 제약을 깨뜨리는 순차 projection을 금지한다.
이 식은 free object와 external free-force label의 shortcut이다.
Hard missing-structural label은 support assignment 뒤 runtime과 같은 augmented variable을 사용한다.

\begin{equation}
\begin{aligned}
  \zeta_{\mathrm{struct}}^{*,t}
  =
  \begin{bmatrix}y_{\mathrm{struct}}^{*,t}\\\rho_{h,\mathrm{struct}}^{*,t}\end{bmatrix}
  =\argmin_{y,\rho}\;&
  \frac12\left\|
  y-E_{f,t}^\top\bar f_{\ell,\mathrm{struct}}^{*,t}
  \right\|_{W_{f,t}}^2
  +\frac12\|\rho\|_{W_h}^2\\
  \mathrm{s.t.}\;&
  B_{\mathrm{struct}}^t
  \begin{bmatrix}y\\\rho\end{bmatrix}=0,\\
  \Delta f_{\ell,\mathrm{struct}}^{*,\mathrm{free},t}
  &:={}E_{f,t}y_{\mathrm{struct}}^{*,t},\qquad
  r_{o\rightarrow h,\mathrm{struct}}^{*,t}
  :=A_{\mathrm h}^\top\rho_{h,\mathrm{struct}}^{*,t}.
  \label{eq:teacher-hard-structural-augmented-label}
\end{aligned}
\end{equation}

여기서 \(B_{\mathrm{struct}}^t\)는 식~\eqref{eq:structural-augmented-constraints}와
동일한 current-geometry wrench/complement operator이고, metric과 reference
\([E_{f,t}^\top\bar f_{\ell,\mathrm{struct}}^*;0]\)도 runtime과 정확히 같다.
따라서 smoothing 또는 teacher reaction을 objective reference에 추가해 inference와 다른
label projection을 만들지 않는다. Raw teacher reaction
\(\rho_{h,\mathrm{struct}}^{*,\mathrm{teacher}}\)는 canonical projection output과의
reaction diagnostic에만 사용한다.

\begin{equation}
  \epsilon_{\rho,\mathrm{teacher}}^t
  :=
  \frac{
    \|\rho_{h,\mathrm{struct}}^{*,t}
    -\rho_{h,\mathrm{struct}}^{*,\mathrm{teacher},t}\|_2
  }{
    \|\rho_{h,\mathrm{struct}}^{*,\mathrm{teacher},t}\|_2+\eps
  }.
  \label{eq:teacher-reaction-projection-diagnostic}
\end{equation}

Local head는 raw free-support remainder만 예측하고, canonical
\((y_{\mathrm{struct}}^*,\rho_{h,\mathrm{struct}}^*)\)는 post-assembly
supervision/diagnostic에 사용한다.

Nodal-breakdown 경로의 channel별 Local generalized target과 split consistency는

\begin{equation}
\begin{aligned}
  \Delta Q_{\ell,c}^{*,t}
  &:={}\Psi_{\mathrm{sup}}^\top
  \Delta f_{\ell,c}^{*,\mathrm{asm},t},\\
  \epsilon_{\mathrm{split}}^t
  &:={}
  \frac{\left\|\sum_c\Delta Q_{\ell,c}^{*,t}-\Delta Q_\ell^{*,t}\right\|_2}
  {\|\Delta Q_\ell^{*,t}\|_2+\eps}
  \le\tau_{\mathrm{split}},
  \label{eq:channel-generalized-force-split}
\end{aligned}
\end{equation}

으로 고정한다. 여기서 assembled target은 아래 식~\eqref{eq:patch-force-label}의
channel-typed 정의를 사용한다. 이 consistency가 실패하면 channel label을 만들지 않고
conditional F를 거부한다.
First-paper method는 \(\Delta Q_{g,c}^{*}\)를 진단값으로만 기록한다.
별도 Global missing-force head는 범위 밖의 후속 연구다.

Generalized force 하나만으로 external resultant/reaction은 복원되지 않는다.
따라서 다음 lifting은 teacher가 channel별
\(\Delta Q_{\ell,c}^{*,\mathrm{input},t}\)와 별도 measured/teacher conservation target
\(b_{\mathrm{cons},c}^{*,t}\)를 모두 제공하는 후속 optional 경로로만 둔다.

\begin{equation}
\begin{aligned}
  \Delta f_{\ell,c}^{*,\mathrm{lift},t}
  =\argmin_f\;&
  \frac12\|f\|_{W_f}^2
  +\frac{\lambda_{f,\mathrm{smooth}}}{2}\|L_ff\|_2^2\\
  \mathrm{s.t.}\;&
  \Psi_{\mathrm{sup}}^\top f=\Delta Q_{\ell,c}^{*,\mathrm{input},t},\\
  &\Phi^\top f=0,\\
  &A_{\mathrm{cons},c}^tf=b_{\mathrm{cons},c}^{*,t}.
  \label{eq:generalized-force-lifting}
\end{aligned}
\end{equation}

식~\eqref{eq:generalized-force-lifting}은 free/no-hard shortcut이며,
\(b_{\mathrm{cons},c}^{*,t}\)가 없는 ``generalized-force only'' sample에는 적용하지 않는다.
Hard missing-structural generalized fallback은 raw support reference와 teacher reaction까지 있어야
식~\eqref{eq:teacher-hard-structural-augmented-label}의 동일 objective를 재현할 수 있으므로
first-paper에서는 no-go다. 이 auxiliary target들이 모두 있는 후속 branch에서만 generalized constraint를
\(\Psi_{\mathrm{sup}}^\top E_{f,t}y
=\Delta Q_{\ell,\mathrm{struct}}^{*,\mathrm{input}}\)로 바꾸고,
\([y;\rho_h]\)에 식~\eqref{eq:structural-augmented-constraints}를 동시에 적용한다.

External channel의 원래 full wrench와 \(\Phi^\top f=0\)가 충돌하면
Global component와 그 wrench를 먼저 \(\Delta Q_g^*\)로 route하고,
남은 Local target에 대해서만 보존 제약을 건다.
KKT rank/feasibility 또는 generalized consistency

\begin{equation}
  \epsilon_Q^{\mathrm{lift}}
  =
  \frac{
    \|\Psi_{\mathrm{sup}}^\top\Delta f_{\ell,c}^{*,\mathrm{lift},t}
    -\Delta Q_{\ell,c}^{*,\mathrm{input},t}\|_2
  }{
    \|\Delta Q_{\ell,c}^{*,\mathrm{input},t}\|_2+\eps
  }
\end{equation}

가 tolerance를 넘으면 label을 생성하지 않거나 low-confidence로 표시한다.
신뢰할 수 있는 nodal channel breakdown 또는 식~\eqref{eq:channel-generalized-force-split}의
channel별 generalized target이 없으면 conditional F를 no-go로 판정한다.
Missing-aero는 external remainder이고 missing-structural은 augmented reaction constraint를 쓰므로,
constraint type이 정의되지 않은 single total-defect head로 조용히 합치지 않는다.

Patch weights는 다음 partition-of-unity 계약을 따른다.

\begin{equation}
\begin{aligned}
\Delta f_{\ell,c}^{*,\mathrm{raw},t}
&:=
\begin{cases}
  \Delta f_{\ell,c}^{*,t},
  &\text{external 또는 free-object shortcut},\\
  E_{f,t}E_{f,t}^\top\bar f_{\ell,\mathrm{struct}}^{*,t},
  &\text{hard missing-structural channel},
\end{cases}\\
\Delta f_{\ell,c}^{*,\mathrm{asm},t}
&:=
\begin{cases}
  \Delta f_{\ell,c}^{*,t},
  &\text{external 또는 free-object shortcut},\\
  E_{f,t}y_{\mathrm{struct}}^{*,t},
  &\text{hard missing-structural channel},
\end{cases}\\
\sum_{r=1}^{P}R_r^\top W_rR_r
&=I,
\qquad
\Delta f_{r,c}^{*,t}=R_r\Delta f_{\ell,c}^{*,\mathrm{raw},t}.
\end{aligned}
  \label{eq:patch-force-label}
\end{equation}

따라서 모든 patch가 활성이고 \(\gamma_r=1\)이면

\begin{equation}
  \sum_rR_r^\top W_r\Delta f_{r,c}^{*,t}
  =
  \Delta f_{\ell,c}^{*,\mathrm{raw},t}.
\end{equation}

Label에 \(W_r\)를 다시 곱하거나 나누지 않는다.
동일 weight는 Module 10 assembly에서 정확히 한 번 적용한다.

\subsection{Gate benefit label}

Gate label은 단순 strain threshold가 아니다.
Patch ID의 disjoint cover를 이루는 고정 atomic gating-unit 집합을

\begin{equation}
  \mathfrak G=\{\mathcal G_1,\ldots,\mathcal G_{P_G}\},
  \qquad
  \bigcup_{\mathcal G\in\mathfrak G}\mathcal G=\{1,\ldots,P\}
\end{equation}

로 setup에서 정한다. 기본 unit은 singleton patch이고,
strong overlap 또는 큰 cross block을 가진 patch만 하나의 unit으로 병합한다.
모든 Oracle/runtime/learned selector는 동일한 \(\mathfrak G\)를 사용한다.
Horizon \(H\) 동안 Global-only rollout과 unit \(\mathcal G\)의
oracle Local correction을 함께 켠 counterfactual rollout을 비교한다.

\begin{equation}
  b_{\mathcal G}^t
  =
  [\mathcal E_H^{G}(\mathcal G,t)
  -\mathcal E_H^{G+L_{\mathcal G}}(\mathcal G,t)]_+.
  \label{eq:benefit-label}
\end{equation}

\(\mathcal E_H\)는 future position, velocity, normal, curvature,
strain, spectrum error의 weighted sum이다.
식~\eqref{eq:benefit-label}은 \emph{Local이 필요해 보이는가}가 아니라
\emph{Local을 켜면 미래 오차가 실제로 얼마나 줄어드는가}를 정답으로 만든다.
V0는 cost predictor를 학습하지 않고 \(K_A\) cardinality를 바꾸며 quality--latency curve를 그린다.

Patch benefit은 일반적으로 가산적이지 않다.

\begin{equation}
  B(\{r,s\})\ne B(\{r\})+B(\{s\}),
  \qquad
  C(\{r,s\})\ne C(\{r\})+C(\{s\}).
\end{equation}

따라서 singleton-patch benefit의 합을 selected-set quality나 GPU wall-clock로 해석하지 않는다.
강하게 겹치거나 cross block이 큰 patch는 고정 group으로 묶고,
작은 fixture에서는 group/set oracle과 greedy Top-\(K\) regret을 함께 측정한다.

모든 patch counterfactual이 너무 비싸면 다음 근사를 사용한다.

\begin{enumerate}[leftmargin=1.7em]
  \item Teacher complement energy가 큰 후보 patch만 선별한다.
  \item 강하게 겹치거나 cross block이 큰 patch를 group으로 묶는다.
  \item Short horizon \(H\)로 benefit을 계산하고,
  일부 sequence에서 long-horizon calibration을 수행한다.
\end{enumerate}

\subsection{Runtime input leakage 방지}

Gate input은 runtime에서 관측 가능한 값만 사용한다.

\begin{equation}
  \chi_{\mathcal G}^t
  =
  \left(
  q^t,\dot q^t,\;
  z_{\mathcal N(\mathcal G)}^t,\dot z_{\mathcal N(\mathcal G)}^t,\;
  \operatorname{pool}_{r\in\mathcal G}
  (\epsilon_r^t,\kappa_r^t,W_r^t,\vartheta_r,\mathcal C_r,h_r^t)
  \right).
\end{equation}

Teacher future state, teacher residual, teacher force를 gate input으로 넣지 않는다.
그 값들은 label 생성에만 사용한다.

\section{Runtime 14단계 개요}

\begin{longtable}{L{0.055\linewidth}L{0.20\linewidth}L{0.18\linewidth}L{0.24\linewidth}L{0.24\linewidth}}
\toprule
ID & 단계 & Solver 종류 & 입력 & 출력 \\
\midrule
\endfirsthead
\toprule
ID & 단계 & Solver 종류 & 입력 & 출력 \\
\midrule
\endhead
1 & 이전 상태 읽기 & \deterministic & \(q,\dot q,z,\dot z\), selector memory & 현재 frame 초기 상태 \\
2 & Causal half-step 표면 복원 & \deterministic & Previous state, bases, binding & Predicted position, frozen velocity, normal, area \\
3 & 상대풍/공기력 & \physics & Surface state, prescribed wind & Sample aerodynamic force \\
4 & Global load 구성 & \physics+\optional & Full-anchor force, \(\Phi\), optional cubature & Generalized load \(Q_g\) \\
5 & Global predictor & \physics & \(Q_g,M_g,C_g,K_g\) & \(\tilde q,\dot{\tilde q}\) \\
6 & Local cardinality 선택 & \physics+\learned & Global predictor, patch features/history & Score와 active set \\
7 & Local force proposal & \physics+\learned & Active patch state, wind, material & Analytic/base Local force + optional missing-force channels \\
10 & Overlap/support route와 conditional joint solve & \constrained & Base + optional defect proposals, overlap/support graph & Support-restricted complement force와 reaction \\
9 & Complement contract/projection & \deterministic & Route/constraint-verified force, \(\Psi_{\mathcal S}\) & Local generalized force/diagnostics \\
8 & Local physical update & \physics & Assembled local force, \(M_\ell,C_\ell,K_\ell\) & \(z,\dot z\) \\
11 & Global feedback/corrector & \physics & Cross reaction, corrected aero delta & Corrected \(q,\dot q\) \\
12 & 최종 anchor/reaction 진단 & \deterministic & Corrected \(q,z\), acceleration, bases/force ledger & Anchor state/deformation gradient, total reaction과 \(r_{\mathrm{ROM}}\) \\
13 & Gaussian transport/render & \deterministic & Anchor state, GS binding & \(G^t\), rendered image \\
14 & 상태/hysteresis 갱신 & \deterministic & Energy, score, optional uncertainty & Next state, dormant/active/decay set \\
\bottomrule
\end{longtable}

\paragraph{순서에 대한 구현 주의.}
위 표의 행 순서는 실제 data dependency이고 첫 열은 실행 번호가 아니라 Module ID다.
7번 patch force proposal을 10번에서 조립하고
support-restricted Global routing과 필요한 channel의 conservation--complement joint solve를 마친 뒤,
9번에서 contract/generalized projection을 검사하고 8번의 단일 Local update를 수행한다.
즉 \emph{patch별 적분 후 overlap 평균}은 금지한다.
상세 module은 실제 계산 순서와 데이터 의존성을 명시한다.

\reviewrebut{V2 리뷰의 축약 V0 도식처럼 physics-only complementary Local solve 뒤에
conservative overlap/complement assembly를 두는 순서는 채택하지 않는다.
Overlap ownership, force/torque constraint와 Global-span 제거가 끝나기 전에는 단일 Local generalized load가
정의되지 않는다. Patch별 solve를 먼저 수행하면 shared-anchor mass와 inter-patch \(M/K/C\) cross-block
소유권이 다시 모호해진다. 표의 번호는 module 식별자일 뿐 실행 dependency 순서가 아니다.}

\reviewvthreerebut{V3 리뷰 \S8의 축약 pipeline도
physics-only complementary Local 뒤에 unique-anchor assembly/routing을 두어
같은 data dependency를 다시 역전시킨다.
Assembly와 support route 또는 conservation--Global-complement joint solve 전에는 Local generalized load가
정의되지 않으므로 이 순서는 채택하지 않는다.
위 표와 식~\eqref{eq:canonical-runtime-dependency}가 canonical 실행 순서다.}

Canonical runtime dependency는

\begin{equation}
\begin{aligned}
\text{patch force proposal}
&\rightarrow
\text{unique-anchor assembly}\\
&\rightarrow
\text{support-restricted Global routing/feasibility}\\
&\rightarrow
\left\{
\begin{array}{l}
\text{external: explicit feasible remainder}\\
\text{structural/reaction: augmented joint solve}
\end{array}
\right.\\
&\rightarrow
\text{generalized projection}\\
&\rightarrow
\text{single coupled Local solve}
\rightarrow
\text{delta Global corrector}
\end{aligned}
\label{eq:canonical-runtime-dependency}
\end{equation}

이며 pipeline figure, profiler trace와 구현 API는 반드시 이 순서를 사용한다.

\section{Module 1: 이전 Global/Local 상태 읽기}

\subsection*{푸는 질문}

\emph{현재 바람만 보고 새 위치를 만드는 것이 아니라,
이전 프레임의 관성, 진동 phase, local 잔류 에너지를 어떻게 이어갈 것인가?}

\subsection*{Solver 분류}

\deterministic State gather와 buffer validation이다. Network와 새로운 물리 solve는 없다.

\subsection*{입력}

\begin{equation}
  s^{t-1}
  =
  \left(
  q^{t-1},\dot q^{t-1},
  z^{t-1},\dot z^{t-1},
  \mathcal A_{t-1},
  \mathcal D_{t-1},
  h^{t-1}
  \right).
\end{equation}

\(\mathcal A\)는 Local excitation과 solve를 허용하는 active patch,
\(\mathcal D\)는 새 excitation은 꺼졌지만 local energy를 물리적으로 감쇠시키는 decay patch,
\(h\)는 selector hysteresis와 선택적인 aerodynamic memory다.
조건부 F가 살아 있을 때만 \(\mathcal A\)에서 missing-force network도 실행한다.

\subsection*{적용 연산}

\begin{itemize}[leftmargin=1.5em]
  \item State dimension, active basis index, timestep, asset package version을 검사한다.
  \item First-paper runtime에서는 material과 attachment layout을 바꾸지 않는다.
  Setup package version이 달라지면 새 asset으로 재초기화하고, 고정 layout의 prescribed handle motion만 frame input으로 읽는다.
  \item Wind 변화만으로 topology, bases, Gaussian binding을 다시 만들지 않는다.
  \item Frame \(t\)의 causal half-step control sample
  \(u_{\mathcal C}(t-\tfrac12)\), \(W(\cdot,t-\tfrac12)\)와 endpoint control을 읽는다.
  \(\mathcal U_{\mathrm p}\)로
  \(u_{\mathrm p},\dot u_{\mathrm p},\ddot u_{\mathrm p}\)를,
  compliant handle이 있으면 \(\mathcal U_{\mathrm c}\)로
  \(d_{\mathrm c},\dot d_{\mathrm c}\)를 같은 sampling time에서 만든다.
\end{itemize}

\subsection*{출력과 다음 단계}

검증된 \(s^{t-1}\)가 Module 2의 causal half-step load-surface reconstruction으로 넘어간다.
State를 저장하기 때문에 resonance, wind-stop recovery, damping envelope가 프레임 사이에 유지된다.

\subsection*{필수 검증}

\begin{itemize}[leftmargin=1.5em]
  \item Zero-input state copy가 bitwise 또는 tolerance 내에서 일관적인가?
  \item Patch가 active에서 decay로 바뀌어도 \(z,\dot z\)가 reset되지 않는가?
  \item Asset/control 변경 시 오래된 operator와 새 state가 혼용되지 않는가?
\end{itemize}

\section{Module 2: 현재 anchor와 surface state 복원}

\subsection*{푸는 질문}

\emph{작은 Global/Local coordinate를 실제 공간의 위치, 속도, 법선, 면적으로 어떻게 바꿀 것인가?}

\subsection*{Solver 분류}

\deterministic Modal lifting, affine MLS, deformation-gradient geometry update다.

\subsection*{입력}

\begin{itemize}[leftmargin=1.5em]
  \item \(q^{t-1},\dot q^{t-1},z^{t-1},\dot z^{t-1}\),
  \item prescribed hard-support midpoint lift
  \(u_{\mathrm p}^{t-1/2},\dot u_{\mathrm p}^{t-1/2},
  \ddot u_{\mathrm p}^{t-1/2}\),
  \item rest anchors \(x^0\), Global basis \(\Phi\), current Local basis \(\Psi_{\mathcal S}\),
  \item aerodynamic samples와 anchor support,
  \item rest normal \(n_s^0\), rest area \(A_s^0\).
\end{itemize}

\subsection*{Causal half-step geometry와 frozen-velocity load predictor}

\begin{equation}
\begin{aligned}
  q^{\mathrm{pred},t-1/2}
  &=q^{t-1}+\frac{\Delta t}{2}\dot q^{t-1},\\
  z_{\mathrm{sup}}^{\mathrm{pred},t-1/2}
  &=z_{\mathrm{sup}}^{t-1}
    +\frac{\Delta t}{2}\dot z_{\mathrm{sup}}^{t-1},\\
  x^{\mathrm{pred},t-1/2}
  &=x^0+u_{\mathrm p}^{t-1/2}
  +\Phi q^{\mathrm{pred},t-1/2}
  +\Psi_{\mathrm{sup}}z_{\mathrm{sup}}^{\mathrm{pred},t-1/2}.
\end{aligned}
  \label{eq:causal-midpoint-geometry-predictor}
\end{equation}

\begin{equation}
  v^{\mathrm{pred},t-1/2}
  =\dot u_{\mathrm p}^{t-1/2}
  +\Phi\dot q^{t-1}+\Psi_{\mathrm{sup}}\dot z_{\mathrm{sup}}^{t-1}.
\end{equation}

이 predictor를 위해 auxiliary structural solve를 하지 않는다.
이전 state의 position/velocity와 prescribed midpoint control만 쓰는 causal half-step extrapolation이다.
Position은 half-step으로 외삽하지만 acceleration state를 저장하지 않으므로 velocity는
이전 endpoint 값으로 hold한다. 따라서 structural homogeneous update는 implicit midpoint여도
exogenous aerodynamic load와 one-pass partitioned coupling의 시간 정확도를 자동으로 2차라고 주장하지 않는다.
필요한 reference에서는 이전 acceleration을 포함한 true midpoint predictor 또는 fixed-point iteration과
비교해 phase/PSD 및 timestep convergence를 보고한다.
아래 geometry 식의 \(x^t,v^t\)는 generic evaluation argument이며,
Module 2--4의 첫 pass에서는 위
\(x^{\mathrm{pred},t-1/2},v^{\mathrm{pred},t-1/2}\)를 넣는다.

\subsection*{Tangential Jacobian과 area-normal}

Rank-2 surface map과 \(3\times3\) cofactor를 혼용하지 않는다.
Sample \(s\)의 oriented rest tangent frame을

\begin{equation}
  T_s^0=[t_{s,1}^0,t_{s,2}^0]\in\R^{3\times2},
  \qquad
  (T_s^0)^\top T_s^0=I_2,
  \qquad
  t_{s,1}^0\times t_{s,2}^0=n_s^0
\end{equation}

로 둔다. Rest neighbor coordinate와 current tangential Jacobian은

\begin{equation}
  \xi_{sk}^0=(T_s^0)^\top(x_k^0-x_s^0)\in\R^2,
\end{equation}

\begin{equation}
  J_s^t
  =
  \argmin_{J\in\R^{3\times2}}
  \sum_{k\in\mathcal N_s}\omega_{sk}
  \left\|(x_k^t-x_s^t)-J\xi_{sk}^0\right\|_2^2
  \label{eq:tangent-jacobian}
\end{equation}

로 맞춘다. Noisy/curved rest neighborhood에서 \(J_s^0\)가 정확히 \(T_s^0\)와 같다고
가정하지 않는다. Raw cross product를
\(\bar{\mathbf a}_s^t=(J_s^te_1)\times(J_s^te_2)\)라 하면
oriented area-normal vector는 rest fit으로 정규화해 계산한다.

\begin{equation}
  \mathbf a_s^t
  =
  A_s^0
  \frac{\bar{\mathbf a}_s^t}
  {\max(\|\bar{\mathbf a}_s^0\|_2,\eps)},
  \qquad
  A_s^t=\|\mathbf a_s^t\|_2,
  \qquad
  n_s^t=\frac{\mathbf a_s^t}{\max(A_s^t,\eps)}.
  \label{eq:area-normal}
\end{equation}

\(\bar n_s^0=\bar{\mathbf a}_s^0/\max(\|\bar{\mathbf a}_s^0\|_2,\eps)\)로 둔다.
Degenerate rest fit 또는 stretch에서는 \(J_s\)의 rank/singular value, area와 orientation continuity를 검사한다.

Gaussian covariance에는 \(3\times2\) \(J_s\)를 직접 사용할 수 없으므로
별도의 full-rank shell transport map을 정의한다.

\begin{equation}
  F_{s,\parallel}^t=J_s^t(J_s^0)^\dagger,
  \qquad
  F_{s,\mathrm{shell}}^t
  =
  F_{s,\parallel}^t
  +\lambda_{n,s}^t n_s^t(\bar n_s^0)^\top.
  \label{eq:full-shell-map}
\end{equation}

V0는 thickness-preserving \(\lambda_{n,s}^t=1\)을 사용한다.
Rank-2 \(J_s^0\)와 rigid motion \(J_s^t=RJ_s^0\)이면
\(F_{s,\mathrm{shell}}^0=I\), \(F_{s,\mathrm{shell}}^t=R\)가 된다.
Tangent stretch policy와 normal/thickness extension은 typed field로 분리하며,
rest identity와 rigid rotation reproduction을 반드시 검증한다.

\subsection*{출력과 다음 단계}

\begin{equation}
  \mathcal G_{\mathrm{surf}}^{\mathrm{pred},t-1/2}
  =
  \{x_s^{\mathrm{pred},t-1/2},v_s^{\mathrm{pred},t-1/2},
  J_s^{\mathrm{pred},t-1/2},F_{s,\mathrm{shell}}^{\mathrm{pred},t-1/2},
  n_s^{\mathrm{pred},t-1/2},A_s^{\mathrm{pred},t-1/2}\}_{s\in\mathcal S_{\mathrm{full}}}.
\end{equation}

이 surface state가 Module 3의 relative-wind/aerodynamic force 계산으로 넘어간다.
Area-normal은 \(J_s\)에서, Module 13 covariance transport는
명시적인 \(F_{s,\mathrm{shell}}\) 또는 Gaussian support에서 fitting한 full-rank map에서 얻는다.

\subsection*{필수 검증}

\begin{itemize}[leftmargin=1.5em]
  \item Identity와 rigid rotation에서 area가 보존되고 normal이 정확히 회전하는가?
  \item In-plane scale에서 area가 예상 비율로 변하는가?
  \item Aero normal과 rendering/transport normal이 동일한 convention을 쓰는가?
  \item Normal sign이 frame 사이에 갑자기 뒤집히지 않는가?
\end{itemize}

\section{Module 3: Relative wind와 quasi-steady surface force}

\subsection*{푸는 질문}

\emph{현재 위치에서 움직이고 있는 표면이 실제로 느끼는 바람과
그 바람이 만드는 기본 external force는 무엇인가?}

\subsection*{Solver 분류}

\physics Prescribed wind interpolation과 quasi-steady aerodynamic law다.
Runtime에 Navier--Stokes, LBM, full CFD를 풀지 않는다.

\subsection*{입력}

\begin{itemize}[leftmargin=1.5em]
  \item Module 2의 causal half-step geometry와 frozen velocity:
  \(
  \bigl(x_s^{\mathrm{pred},t-1/2},v_s^{\mathrm{pred},t-1/2}\bigr)
  \) 및
  \(
  \bigl(n_s^{\mathrm{pred},t-1/2},A_s^{\mathrm{pred},t-1/2}\bigr)
  \),
  \item prescribed wind field \(W(\mathbf x,t-\tfrac12)\),
  \item density \(\rho_{\mathrm{air}}\),
  \item drag/friction/lift coefficients \(\theta_s^{\mathrm{aero}}\),
  \item optional exposure/reliability \(e_s\).
\end{itemize}

\subsection*{Relative wind}

\begin{equation}
  r_s^{\mathrm{pred},t-1/2}
  =
  W(x_s^{\mathrm{pred},t-1/2},t-\tfrac12)
  -v_s^{\mathrm{pred},t-1/2},
  \qquad
  \widehat r_s^{\mathrm{pred},t-1/2}
  =
  \frac{r_s^{\mathrm{pred},t-1/2}}
  {\|r_s^{\mathrm{pred},t-1/2}\|_2+\eps}.
  \label{eq:relative-wind}
\end{equation}

표면이 바람과 같은 방향으로 움직이면 실제로 느끼는 속도가 줄고,
반대 방향으로 움직이면 증가한다.
Absolute wind만 쓰는 baseline은 이 효과를 제거하는 ablation이다.

\subsection*{Surface aerodynamic law}

구체적인 drag/lift convention은 teacher와 runtime에서 동일하게 고정한다.
문서 수준에서는 다음 generic calibrated map으로 쓴다.

\begin{equation}
  f_{a,s}^t
  =
  \frac12\rho_{\mathrm{air}}
  A_s^{\mathrm{pred},t-1/2}
  \|r_s^{\mathrm{pred},t-1/2}\|_2^2
  e_s^{\mathrm{pred},t-1/2}
  \mathcal C_{\mathrm{aero}}
  \left(
    \widehat r_s^{\mathrm{pred},t-1/2},
    n_s^{\mathrm{pred},t-1/2},
    \theta_s^{\mathrm{aero}}
  \right).
  \label{eq:surface-aero}
\end{equation}

\(f_{a,s}^t\)의 superscript \(t\)는 frame-\(t\) integrator에 hold되는 load를 뜻하며,
평가 시점은 \(t-\tfrac12\)이다.
\(\mathcal C_{\mathrm{aero}}\)는 normal pressure/drag,
tangential friction, optional lift를 반환한다.
MVP에서는 two-sided thin-sheet law를 사용하고,
front/back 계수가 다르면 oriented normal label을 유지한다.

\subsection*{수식이 풀지 않는 것}

식~\eqref{eq:surface-aero}은 local quasi-steady force다.
다음 현상은 직접 풀지 않는다.

\begin{itemize}[leftmargin=1.5em]
  \item object가 wind field에 만드는 wake,
  \item vortex shedding의 fluid state,
  \item long aerodynamic memory,
  \item dense self-shadowing과 full pressure solve.
\end{itemize}

이 현상을 주장하려면 Teacher-FSI가 필요하고,
그 차이만 Module 7의 missing-force target에 포함한다.

\subsection*{출력과 다음 단계}

Full-anchor force set \(\{f_{a,s}^t\}\)가 Module 4의 Global load construction으로 넘어간다.
Active Local correction 이후에는 같은 식을 Module 11에서 active region에만 다시 평가한다.

\subsection*{필수 검증}

\begin{itemize}[leftmargin=1.5em]
  \item Zero relative wind에서 force가 정확히 0인가?
  \item Rigid co-rotation에서 force가 frame-equivariant한가?
  \item Surface velocity를 제거한 baseline과 강한 gust에서 차이가 나는가?
  \item Current normal/area 대신 rest 값을 쓴 baseline 대비 deformed-state force가 정확한가?
\end{itemize}

\section{Module 4: Global aerodynamic load construction}

\subsection*{푸는 질문}

\emph{Current-surface force가 Global coordinate에 하는 virtual work를 어떻게 정확히 구성하며,
full-anchor evaluation이 병목일 때에만 이를 얼마나 줄일 수 있는가?}

\subsection*{Solver 분류}

\physics V0는 full-anchor force evaluation과 Galerkin projection이다.
\optional Profiling gate를 통과한 경우에만 empirical cubature sample/weight optimization과
weighted runtime backend를 추가한다.
Cubature 자체는 고전적이며 \cite{an2008cubature}; main contribution으로 가정하지 않는다.

\subsection*{Full reference}

V0 physical anchor 또는 fixed linear sample lifting을
\(T_s^{A\rightarrow S}\in\R^{3\times3K}\)라 하면

\begin{equation}
\begin{aligned}
  x_s^t&=T_s^{A\rightarrow S}x^t,\\
  v_s^t&=T_s^{A\rightarrow S}v^t,\\
  \Phi_s&:=T_s^{A\rightarrow S}\Phi,\\
  J_{s,g}&:=\frac{\partial x_s}{\partial q}=\Phi_s.
  \label{eq:fixed-global-sample-jacobian}
\end{aligned}
\end{equation}

모든 physical anchor/sample을 평가하는 V0 Global generalized force는

\begin{equation}
  Q_{g,\mathrm{aero}}^t
  =
  \sum_{s=1}^{N_S}
  \Phi_s^\top f_{a,s}^t,
  \label{eq:full-generalized-aero}
\end{equation}

Current geometry는 normal, area, relative wind와 결과 force를 바꾸지만,
fixed anchor/sample lifting의 virtual-work Jacobian을 state-dependent하게 만들지 않는다.
Structural corotation만으로 current sample Jacobian이 필요한 것도 아니다.
Nonlinear sample-position reconstruction을 시험할 때에만 \(J_{s,g}^t\)를 갖는
별도 versioned optional backend로 분리하며 canonical V0나 default corrector에 섞지 않는다.

Attached object의 deformation에는 net wrench보다 \(Q_g\) 정확도가 직접적이므로
식~\eqref{eq:full-generalized-aero}를 primary target으로 둔다.
Net force/torque는 force-transfer sanity와 보조 constraint/metric이다.

\subsection*{조건부 hyper-reduced evaluation}

\begin{equation}
  \widehat Q_{g,\mathrm{aero}}^t
  =
  \sum_{s\in\mathcal S_c}
  \omega_s
  \Phi_s^\top f_{a,s}^t,
  \qquad
  |\mathcal S_c|\ll N_S.
  \label{eq:hyper-aero}
\end{equation}

Full-anchor aero와 projection이 target GPU의 전체 dynamics에서
유의미한 병목일 때만 sample set \(\mathcal S_c\)와 nonnegative weights \(\omega_s\)를 fit한다.
다양한 deformation/wind training snapshots에서 다음 목표를 함께 최소화한다.

\begin{equation}
\begin{aligned}
  \min_{\mathcal S_c,\omega\ge0}
  \sum_m
  &\left\|
  Q_{g,m}^{\mathrm{full}}
  -\widehat Q_{g,m}
  \right\|_{W_Q}^2\\
  +\lambda_F&
  \left\|
  F_m^{\mathrm{full}}
  -\widehat F_m
  \right\|_2^2
  +\lambda_\tau
  \left\|
  \tau_m^{\mathrm{full}}
  -\widehat\tau_m
  \right\|_2^2.
  \label{eq:cubature-fit}
\end{aligned}
\end{equation}

Net force와 torque는

\begin{equation}
  F=\sum_s f_{a,s},
  \qquad
  \tau=\sum_s(x_s-c)\times f_{a,s}
\end{equation}

로 계산한다.
Centroid \(c\) convention은 teacher, training, runtime에서 동일해야 한다.

\subsection*{입력}

Module 3의 causal half-step/frozen-velocity full anchor forces와
식~\eqref{eq:fixed-global-sample-jacobian}의 fixed \(\Phi_s\)가 기본 입력이다.
조건부 backend만 offline sample indices/weights와 object center를 추가로 받는다.

\subsection*{출력과 다음 단계}

\begin{equation}
  Q_{g,\mathrm{aero}}^t,\quad
  F_{\mathrm{aero}}^t,\quad
  \tau_{\mathrm{aero}}^t,\quad
  \epsilon_{\mathrm{red}}^t.
\end{equation}

V0에서는 식~\eqref{eq:full-generalized-aero}의 exact discrete sum을 출력하고
\(\epsilon_{\mathrm{red}}=0\)으로 둔다.
조건부 backend에서는 hat quantity를 위 공통 output schema로 반환하고 full reference 대비 error estimate를 기록한다.
Generalized force는 Module 5로 넘어간다.

\subsection*{왜 network가 아닌가}

한 frame에 현재 wind query 하나만 필요하므로
모든 방향 coefficient를 생성하는 state-conditioned RSH보다
현재 sample force를 직접 평가하고 합산하는 편이 자연스럽다.
RSH, directional MLP, LUT는 E2 공력 표현 비교에만 남긴다.

\subsection*{필수 검증}

\begin{itemize}[leftmargin=1.5em]
  \item 512/1,024/2,048 anchors에서 full aero와 fixed projection의 p50/p95/p99 및 dynamics 비율.
  \item Full samples 대비 generalized-force를 primary, net-force/net-torque를 auxiliary error로 보고.
  \item Random/FPS/nonconservative learned sampling과의 비교.
  \item Sample count 대 actual GPU latency의 Pareto curve.
  \item Held-out deformation, traveling/localized spatial gust에서의 error.
  \item Full backend가 병목이 아니면 cubature를 구현하지 않거나 no-go로 종료.
\end{itemize}

\section{Module 5: Always-on Global reduced physics predictor}

\subsection*{푸는 질문}

\emph{현재 전체 힘을 받았을 때 물체의 큰 굽힘, sway, twist, 저주파 진동이
다음 시간에 어떻게 변하는가?}

\subsection*{Solver 분류}

\physics Reduced second-order structural dynamics와 implicit time integration이다.
기본형에는 direct next-state network를 쓰지 않는다.

\subsection*{입력}

\begin{itemize}[leftmargin=1.5em]
  \item 이전 \(q^{t-1},\dot q^{t-1}\),
  \item Module 4의 \(Q_{g,\mathrm{aero}}^t\),
  \item reduced \(M_{gg},C_{gg},K_{gg}\),
  \item gravity와 explicit non-attachment external force,
  \item 식~\eqref{eq:canonical-structural-force-ledger}의 operator target와
  식~\eqref{eq:hard-support-lift-force}의 support lift,
  \item 이전 fixed-superset Local state가 만드는 predictor-consumed cross load.
\end{itemize}

\subsection*{Global equation}

Module 5가 이번 frame에 정확히 한 번 소비하는 structural cross load를

\begin{equation}
\begin{aligned}
  Q_{g\leftarrow\ell}^{\mathrm{pred},t}
  &:={}
  -C_{g,\mathrm{sup}}\dot z_{\mathrm{sup}}^{\mathrm{pred},t-1/2}
  -K_{g,\mathrm{sup}}^{0}z_{\mathrm{sup}}^{\mathrm{pred},t-1/2},\\
  z_{\mathrm{sup}}^{\mathrm{pred},t-1/2}
  &:={}
  z_{\mathrm{sup}}^{t-1}
  +\frac{\Delta t}{2}\dot z_{\mathrm{sup}}^{t-1},
  \qquad
  \dot z_{\mathrm{sup}}^{\mathrm{pred},t-1/2}
  :=\dot z_{\mathrm{sup}}^{t-1}
\end{aligned}
  \label{eq:predictor-cross-load}
\end{equation}

로 이름 붙이고 force ledger에 저장한다. Dormant coordinate는 0이므로
fixed-superset 표기와 active-block 표기는 같은 virtual work를 만든다.
Canonical projective backend의 Hessian은 fixed \(K^0\)이고 rotation/compliant target은
아래 operator-owned generalized load로 들어가므로 이 cross load에 섞지 않는다.

\begin{equation}
  Q_{g,\mathrm{op}}^t
  :=\Phi^\top(h_K^t+h_C^t+f_{\mathrm p}^t).
  \label{eq:global-operator-target-load}
\end{equation}

\begin{equation}
\begin{aligned}
  M_{gg}\ddot q
  +C_{gg}\dot q
  +K_{gg}q
  =&\;
  Q_{g,\mathrm{aero}}
  +Q_{g,\mathrm{gravity}}
  +Q_{g,\mathrm{other}}
  +Q_{g,\mathrm{op}}\\
  &+Q_{g\leftarrow\ell}^{\mathrm{pred},t}.
  \label{eq:global-dynamics}
\end{aligned}
\end{equation}

\(Q_{g,\mathrm{op}}\)의 relation target, compliant target와 hard lift는
각각 정확히 한 번만 투영한다. 이를 \(K_{gg}q\) 안에 다시 넣거나
patch force proposal로 재조립하지 않는다.
별도 learned Global correction은 first-paper method에 넣지 않고 direct/hybrid neural baseline으로만 비교한다.

\subsection*{Time integration}

식~\eqref{eq:v0-structural-freeze-signature}의 implicit midpoint만 사용한다.
Half-step geometry/frozen-velocity predictor에서 계산한 exogenous/operator 항과
식~\eqref{eq:predictor-cross-load}의 causal structural cross approximation을

\begin{equation}
  F_{g,\mathrm{held}}^t
  :=Q_{g,\mathrm{aero}}^t+Q_{g,\mathrm{gravity}}^t
  +Q_{g,\mathrm{other}}^t+Q_{g,\mathrm{op}}^t
  +Q_{g\leftarrow\ell}^{\mathrm{pred},t}
\end{equation}

로 한 frame 동안 고정한다. 그러면 식~\eqref{eq:midpoint-effective-update}와 동치인

\begin{equation}
\begin{gathered}
  q^{t-\frac12}=\frac12(q^t+q^{t-1}),
  \qquad
  \dot q^{t-\frac12}=\frac12(\dot q^t+\dot q^{t-1}),\\
  q^t-q^{t-1}=\Delta t\,\dot q^{t-\frac12},\\
  M_{gg}\frac{\dot q^t-\dot q^{t-1}}{\Delta t}
  +C_{gg}\dot q^{t-\frac12}
  +K_{gg}q^{t-\frac12}
  =F_{g,\mathrm{held}}^t.
  \label{eq:global-implicit-midpoint}
\end{gathered}
\end{equation}

Backward Euler는 식~\eqref{eq:backward-euler-reference-update}의 debug/robustness ablation으로만
보고하고 canonical rollout과 섞지 않는다.
Global dimension \(r_g\)가 작으므로 same-frame corrector를 한 번 추가해도 비용이 제한적이다.

\subsection*{출력과 다음 단계}

Local을 아직 새로 적용하지 않은

\begin{equation}
  \widetilde q^t,\quad
  \dot{\widetilde q}^t,\quad
  \widetilde x_G^t
  =x^0+u_{\mathrm p}^t+\Phi\widetilde q^t
  +\Psi_{\mathrm{sup}}z_{\mathrm{sup}}^{t-1}
\end{equation}

를 Global predictor라 한다.
이 값과 patch features가 Module 6 selector로 넘어간다.

\subsection*{왜 Global을 physics로 두는가}

Global system은 이미 저차원이므로 direct network로 대체했을 때의 속도 이득이 제한적이다.
반면 physics는 다음을 명시적으로 유지한다.

\begin{itemize}[leftmargin=1.5em]
  \item wind-stop recovery,
  \item resonance와 phase,
  \item mass/stiffness/damping control,
  \item unseen wind와 선언된 V0 scalar material parameter에 대한 구조적 외삽,
  \item long rollout energy behavior.
\end{itemize}

Direct neural Global은 중요한 baseline이지만 default method가 아니다.

\subsection*{필수 검증}

\begin{itemize}[leftmargin=1.5em]
  \item Zero wind에서 rest로 수렴하는가?
  \item Mode frequency와 damping envelope가 teacher와 일치하는가?
  \item Direct MLP/GNN/GRU Global보다 OOD와 장시간 rollout이 안정적인가?
  \item 동일 rest geometry/anchor, mass, user material/attachment, load, timestep,
  precision와 hardware에서 sparse PD/XPBD보다 실제로 유리한가?
  각 backend의 native energy/operator privilege와 simulated DOF/rank를 별도 보고하고,
  동일 operator는 두 backend가 실제로 공유할 수 있을 때만 주장한다.
\end{itemize}

\section{Module 6: Fixed-cardinality Local selector와 조건부 learned gate}

\subsection*{푸는 질문}

\emph{고정된 수의 atomic gating unit을 어느 영역에 써야 Global-only의 미래 오차를 가장 많이 줄이는가?}

\subsection*{Solver 분류}

V0 development에는 \trainonly Oracle future-benefit Top-\(K\),
V0 runtime에는 \physics cached-force dynamic-compliance Top-\(K\)를 사용한다.
현재 geometry와 force를 다시 평가하는 full Global-only counter-state residual은
teacher-free이지만 비싸므로 every-frame mainline이 아니라 offline 또는 저주기 analysis baseline으로 둔다.
\learned Lightweight future-benefit network는 Oracle과 cheap runtime selector 사이에
유의미한 gap이 있을 때만 추가한다.

\subsection*{입력}

\begin{itemize}[leftmargin=1.5em]
  \item Global predictor \(q,\dot q,\widetilde x_G,\widetilde v_G\),
  \item patch strain, curvature, deformation rate,
  \item current relative wind와 wind gradient,
  \item optional aero-reduction/topology uncertainty,
  \item past Local state energy와 recent activation history,
  \item basis truncation indicators.
\end{itemize}

\subsection*{V0 cached-force dynamic-compliance score}

Module 3--4에서 이미 계산한 owner-resolved anchor force를 재사용한다.
\(P_{\mathrm{free}}\)는 prescribed support DOF를 제거하는 고정 free-DOF mask이고,
\(f_{\mathrm{aero},A}^{\mathrm{base},t}\)는 causal half-step/frozen-velocity surface sample force를
보존적으로 transfer/assemble한 base anchor force다.
Hard prescribed support와 compliant moving target은 각각
식~\eqref{eq:hard-support-lift-force}와 식~\eqref{eq:canonical-structural-force-ledger}가 소유한다.
이를 별도 attachment force로 다시 넣지 않는다.

\begin{equation}
\begin{aligned}
  f_{\mathrm{cache,ext}}^t
  =P_{\mathrm{free}}
  \left(
    f_{\mathrm{aero},A}^{\mathrm{base},t}
    +f_{\mathrm{gravity}}^t+f_{\mathrm{other,free}}^t
  \right),\\
  f_{\mathrm{cache,sel}}^t
  =f_{\mathrm{cache,ext}}^t
  +P_{\mathrm{free}}(h_K^t+h_C^t+f_{\mathrm p}^t).
\end{aligned}
\end{equation}

로 둔다. Structural backend가 이미 만든 operator target/lift를 재사용하므로
이 score를 위해 auxiliary time integration, predictor geometry 재구성,
두 번째 aerodynamic/internal-force evaluation을 수행하지 않는다.
강한 overlap 또는 cross block을 가진 patch는 미리 정한 gating group \(\mathcal G\)로 묶고,
해당 complement에 작용하는 cached driving load를

\begin{equation}
  g_{\mathcal G}^t
  =
  \Psi_{\mathcal G}^{\top}f_{\mathrm{cache,sel}}^t
  -C_{\mathcal Gg}\dot{\widetilde q}^{t}
  -K_{\mathcal Gg}^{0}\widetilde q^{t}
  \label{eq:cached-local-drive}
\end{equation}

로 정의한다. Group 밖의 weak Local cross term은 V0 score에서 생략하고,
강한 cross term은 group 내부에 포함한다. 이미 active/decay인 state의 energy와
minimum dwell time은 별도 retention/hysteresis 신호로 유지한다.

선택한 implicit integrator의 group effective operator를

\begin{equation}
  A_{\mathcal G}^{\mathrm{eff}}
  =a_M(\Delta t)M_{\mathcal G\mathcal G}
  +a_C(\Delta t)C_{\mathcal G\mathcal G}
  +K_{\mathcal G\mathcal G}^{0}
\end{equation}

로 둔다. Canonical implicit midpoint에서
\(a_M=4\Delta t^{-2}\), \(a_C=2\Delta t^{-1}\)이다.
식~\eqref{eq:canonical-structural-force-ledger}의 projective target은 RHS만 바꾸므로
이 fixed-Hessian operator를 setup에서 prefactorize해 cheap ranking에 그대로 사용한다.
Backward Euler diagnostic만 \(a_M=\Delta t^{-2}\), \(a_C=\Delta t^{-1}\)로 바꾼다.
Anchor-scalar owner weight는 세 vector DOF에서 반복되도록
\(W_r=\operatorname{diag}(w_{rk}I_3)_{k\in r}\)로 정의한다.
Owner-weighted group mass는 overlap anchor를 중복 계수하지 않도록
\(m_{\mathcal G}^{\mathrm{own}}:=\sum_{r\in\mathcal G}\sum_{k\in r}w_{rk}m_k\)로 두고,
singleton이면 \(m_r^{\mathrm{own}}=m_{\{r\}}^{\mathrm{own}}\)이다.
\(\eps_m>0\)은 mass 단위를 가진 numerical floor다.
every-frame score는

\begin{equation}
  \eta_{\mathcal G,\mathrm{comp}}^t
  =
  \left[
  \frac{
    (g_{\mathcal G}^t)^\top
    (A_{\mathcal G}^{\mathrm{eff}})^{-1}g_{\mathcal G}^t
  }{
    m_{\mathcal G}^{\mathrm{own}}+\eps_m
  }
  \right]^{1/2}.
  \label{eq:cheap-dynamic-compliance}
\end{equation}

이는 raw aero norm이 아니라 current load에 대한 mass-normalized compliance work의 제곱근,
즉 velocity-scale의 one-step response indicator다.
따라서 stiff patch의 과대 선택을 줄이면서도 full counter-state residual보다 싸다.
Predictor 이후 geometry 변화와 group 밖의 weak coupling을 생략한 first-order indicator라는 한계는
Oracle regret과 다음 analysis baseline으로 측정한다.

\subsection*{조건부 deterministic hybrid 진단}

\reviewpartial{Cached compliance가 이미 큰 strain/curvature 또는 package truncation risk를
늦게 감지할 수 있다는 진단은 수용한다. 그러나 compliance, strain, curvature와 truncation indicator를
raw 값으로 더하는 리뷰 식은 단위가 서로 달라 asset scale과 단위계에 따라 ranking이 바뀌므로 채택하지 않는다.
Pure compliance를 V0 mainline으로 유지하고, 무차원화한 deterministic hybrid만 blind-spot ablation으로 둔다.
또한 reference-compliance selector에 state-dependent stiffness detector를 필수화하지 않는다.
Canonical projective backend는 fixed Hessian과 frame target RHS를 사용하며,
그 one-shot target 근사 위험은 Step 0 rollout과 full-counter analysis에서 별도로 측정한다.}

Group support의 owner-weighted strain과 rest 대비 curvature magnitude를

\begin{equation}
\begin{aligned}
  \overline\epsilon_{\mathcal G}^t
  &=
  \left[
    \frac{\sum_{k\in\mathcal I_{\mathcal G}}
      \omega_{\mathcal Gk}\|\mathsf E_k^t\|_F^2}
    {\sum_{k\in\mathcal I_{\mathcal G}}\omega_{\mathcal Gk}+\eps}
  \right]^{1/2},\\
  \overline\kappa_{\mathcal G}^t
  &=
  \left[
    \frac{\sum_{k\in\mathcal I_{\mathcal G}}
      \omega_{\mathcal Gk}\|\mathsf B_k^t-\mathsf B_k^0\|_F^2}
    {\sum_{k\in\mathcal I_{\mathcal G}}\omega_{\mathcal Gk}+\eps}
  \right]^{1/2}
\end{aligned}
\end{equation}

로 둔다. 여기서 \(\omega_{\mathcal Gk}=\sum_{r\in\mathcal G}w_{rk}\)이고,
\(\ell_{\mathcal G}\)는 owner-area에서 얻은 characteristic length다.
Runtime에서 teacher 없이 정의 가능한 discarded-energy/confidence ratio만
\(\overline e_{\mathcal G,\mathrm{trunc}}\in[0,1]\)로 허용한다. Feature vector는

\begin{equation}
  a_{\mathcal G}^t
  =
  \left(
    \eta_{\mathcal G,\mathrm{comp}}^t,
    \overline\epsilon_{\mathcal G}^t,
    \ell_{\mathcal G}\overline\kappa_{\mathcal G}^t,
    \overline e_{\mathcal G,\mathrm{trunc}}
  \right)
\end{equation}

이며, 각 항과 같은 단위를 가진 positive scale \(s_j\)를 development split에서 고정하여

\begin{equation}
  \xi_{j,\mathcal G}^t
  =
  \frac{a_{j,\mathcal G}^t}{a_{j,\mathcal G}^t+s_j},
  \qquad
  s_{\mathcal G,\mathrm{hyb}}^t
  =
  \sum_j\lambda_j\xi_{j,\mathcal G}^t,
  \qquad
  \lambda_j\ge0,
  \quad
  \sum_j\lambda_j=1
  \label{eq:normalized-hybrid-selector}
\end{equation}

로 무차원화한다. \(s_j,\lambda_j\)는 held-out test와 test Oracle을 보지 않고 고정한다.
Teacher 또는 새 internal-force evaluation이 필요한 residual은 cheap hybrid에 넣지 않고
analysis backend로 보낸다. Strain/curvature derivative가 다른 module에서 이미 계산되지 않았다면
그 gather와 geometry-derivative 비용을 selector-only latency에 포함한다.
Pure compliance보다 Oracle regret, high-benefit miss와 activation delay를 낮추면서
동일 measured-p95 quality--latency Pareto를 개선할 때만 runtime 후보로 승격한다.

\subsection*{Pure-compliance selector의 mandatory blind-spot fixture}

\reviewvfouraccept{Cached compliance를 cheap V0 selector로 유지하되, 그 score가 직접 관측하지
않는 상태를 mandatory stress test로 분리한다. 이는 learned gate를 미리 정당화하기 위한 것이 아니라
pure compliance의 적용 경계와 더 비싼 indicator의 실제 Pareto를 측정하기 위한 것이다.}

다음 네 fixture를 고정한다.

\begin{enumerate}[leftmargin=1.7em]
  \item \textbf{Cross-group excitation:} group 밖의 load 또는 retained Local state가 cross block을 통해
  target group에 큰 Oracle benefit을 만드는 경우.
  \item \textbf{High-curvature, low-current-load release:} preloaded/released state처럼 cached external
  drive는 작지만 curvature, strain 또는 retained Local energy가 큰 경우.
  \item \textbf{Topology/operator-error negative control:} controlled close-layer false edge,
  missing relation 또는 local stiffness perturbation으로 internal residual은 크지만 cached external
  drive는 작은 경우. Selector가 invalid package를 복구해야 한다는 요구가 아니라 package validation이
  먼저 실패하는지를 확인한다.
  \item \textbf{Fast traveling/reversing gust:} benefit support가 group 사이를 빠르게 이동하여
  delayed activation과 active-set churn이 생기는 경우.
\end{enumerate}

각 fixture에서 pure compliance, normalized deterministic hybrid, low-rate/full counter와 Oracle benefit을
같은 \(K_A\) 및 measured-p95로 비교한다. Oracle regret, high-benefit miss, activation delay,
trajectory/velocity error, target-band PSD와 selector-only p50/p95/p99를 함께 기록한다.
Cross-group coupling이 구조적으로 크면 selector가 이를 학습해 숨기지 않고 gating group을 병합한다.
Topology-error negative control이 package validation을 통과하면 MC1 package gate를 고친다.
Hybrid/full counter가 정확도를 높여도 추가 비용을 포함한 quality--latency Pareto가 좋아질 때만
runtime 후보로 승격한다.

\subsection*{Analysis/low-rate full counter-state residual}

Analysis용 auxiliary \(\widetilde q_G^t\)는 Module 5와 같은 reduced solve에서
이전 Local cross load를 0으로 둔 Global-only counter-predictor다.
Prescribed hard lift를 \(u_{\mathrm p}^t\)라 할 때
\(u_G^t=u_{\mathrm p}^t+\Phi\widetilde q_G^t\)와 그 시간 미분을 만들고,
\(f_{\mathrm{ext},G}^t,h_{K,G}^t,h_{C,G}^t\)를 같은 counter-state와
식~\eqref{eq:canonical-structural-force-ledger}의 sampling convention에서 다시 평가한다.
Group owner restriction을
\(\mathcal R_{\mathcal G}^{\mathrm{own}}
:=[W_r^{1/2}R_r]_{r\in\mathcal G}\)로 stack하면

\begin{equation}
  \eta_{\mathcal G,\mathrm{counter}}^t
  =
  \frac{
  \left\|
  \mathcal R_{\mathcal G}^{\mathrm{own}}P_{\mathrm{free}}
  \left(
    f_{\mathrm{ext},G}^t+h_{K,G}^t+h_{C,G}^t
    -M\ddot u_G^t-C^0\dot u_G^t-K^0u_G^t
  \right)
  \right\|_2
  }{
  \sqrt{m_{\mathcal G}^{\mathrm{own}}+\eps_m}
  }.
  \label{eq:full-counter-state-residual}
\end{equation}

이 backend는 offline 또는 \(J\) frame마다 실행하고 중간 score를 hold하는 저주기 진단으로만 사용한다.
\(J\in\{2,4,8\}\)과 실제 비용을 sweep한다. Support reaction은 \(P_{\mathrm{free}}\) 밖에서 별도 진단하며,
teacher residual과 이름과 구현을 분리한다. 두 indicator의 normalization은 owner-mass convention으로 통일한다.

\subsection*{조건부 learned 출력}

\begin{equation}
  (\widehat b_{\mathcal G}^t,\widehat\sigma_{\mathcal G}^t)
  =
  G_\varphi(\chi_{\mathcal G}^t),
\end{equation}

여기서 \(\widehat b\)는 raw future error reduction,
\(\widehat\sigma\)는 uncertainty다.
Risk-aware score는 예를 들어

\begin{equation}
  s_{\mathcal G}^t
  =
  \widehat b_{\mathcal G}^t
  -\lambda_\sigma\widehat\sigma_{\mathcal G}^t.
\end{equation}

\subsection*{Active set와 fixed-cardinality contract}

\begin{equation}
  \mathcal A_t^{\mathrm{grp}}
  =
  \operatorname{TopK}
  \left(\{s_{\mathcal G}^t\},K_A\right),
  \qquad
  s_{\mathcal G}^t\in
  \{b_{\mathcal G}^t,\eta_{\mathcal G,\mathrm{comp}}^t,
  s_{\mathcal G,\mathrm{hyb}}^t,
  \eta_{\mathcal G,\mathrm{counter}}^t,
  \widehat b_{\mathcal G}^t-\lambda_\sigma\widehat\sigma_{\mathcal G}^t\}.
\end{equation}

Patch active set은

\begin{equation}
  \mathcal A_t
  =\bigcup_{\mathcal G\in\mathcal A_t^{\mathrm{grp}}}\mathcal G
\end{equation}

로 내려보낸다. 같은 group의 patch는 selection status를 공유하지만
각 \(\gamma_r\)와 persistent Local state는 patch slot에 남는다.

\reviewaccept{\(K_A\)를 fixed compute budget이라고 부르는 표현은 철회한다.
\(K_A\)는 동일한 \(\mathfrak G\) 위에서 선택하는 atomic gating-unit의 cardinality일 뿐이다.
Group마다 Local rank, unique-anchor support, retained cross block과 factor-cache 상태가 다르므로
같은 \(K_A\)가 같은 FLOP, GPU time 또는 memory traffic을 뜻하지 않는다.}

Patch \(r\)의 fixed coordinate rank를 \(r_r\), core--halo unique-anchor support를
\(\mathcal I_r\)라 하고 \(\mathcal S_t=\mathcal A_t\cup\mathcal D_t\)라 두면

\begin{equation}
\begin{aligned}
  r_A^t&=\sum_{r\in\mathcal A_t}r_r,
  &n_A^t&=\left|\bigcup_{r\in\mathcal A_t}\mathcal I_r\right|,\\
  r_D^t&=\sum_{r\in\mathcal D_t}r_r,
  &n_D^t&=\left|\bigcup_{r\in\mathcal D_t}\mathcal I_r\right|,\\
  r_{\mathcal S}^t&=\sum_{r\in\mathcal S_t}r_r,
  &n_{\mathcal S}^t&=\left|\bigcup_{r\in\mathcal S_t}\mathcal I_r\right|.
\end{aligned}
\label{eq:active-solve-size-diagnostics}
\end{equation}

모든 selector에 같은 grouping과 \(K_A\)를 사용하되,
\(|\mathcal A_t|,|\mathcal D_t|,r_A^t,n_A^t,r_{\mathcal S}^t,n_{\mathcal S}^t\)와
actual p50/p95/p99 latency를 함께 보고하며 \(K_A\) sweep으로 quality--latency curve를 만든다.
Decay와 shared assembly/corrector 때문에 하나의 active-only cost로 전체 비용을 귀속하지 않고,
selector, proposal, assembly/Schur, Local solve, corrected-aero reevaluation과 corrector 시간을 분리한다.
\(\eta_{\mathrm{counter}}\)는 analysis/low-rate backend에서만 후보이며 every-frame V0에는 포함하지 않는다.
\(s_{\mathcal G,\mathrm{hyb}}\)도 식~\eqref{eq:normalized-hybrid-selector}의 조건부 deterministic ablation이다.
Oracle benefit, compliance, counter residual과 learned score는 단위가 다르므로
서로 합산하거나 공통 threshold를 공유하지 않고 각 selector 내부 ranking에만 사용한다.
동점, incumbent bias와 minimum dwell time을 deterministic하게 정의하고,
Top-\(K\)에서 빠진 unit의 patch는 state를 reset하지 않고 decay로 이동한다.
Unit별 benefit이나 비용을 더해 global-optimal set이라고 주장하지 않는다.
Per-patch cost prediction과 hardware-aware TopBudget은 measured profiler가 필요성을 보인 뒤의 후속 후보다.

\subsection*{연속 activation}

새 patch는 local coordinate를 0에서 시작하고,
force injection coefficient를 몇 frame에 걸쳐 증가시킨다.

\begin{equation}
  \gamma_r^{t+1}
  =
  \clamp(\gamma_r^t+\Delta\gamma_{\mathrm{on}},0,1)
\end{equation}

비활성 전환에서는

\begin{equation}
  \gamma_r^{t+1}
  =
  \clamp(\gamma_r^t-\Delta\gamma_{\mathrm{off}},0,1)
\end{equation}

로 줄인다.
조건부 F가 유지된 경우에는 expert evaluation 시작 threshold와 full blend threshold를 분리하여
expert가 처음 실행되는 frame에 full-scale force가 갑자기 나타나지 않게 한다.

\subsection*{출력과 다음 단계}

\(\mathcal A_t\), decay set \(\mathcal D_t\), \(\gamma_r^t\),
selector 종류, score와 fixed-cardinality/actual-cost diagnostics가 Module 7로 넘어간다.
Dormant/decay patch의 missing-force kernel은 실제로 launch하지 않는다.

\subsection*{필수 검증}

\begin{itemize}[leftmargin=1.5em]
  \item V0: 같은 \(K_A\)의 random/curvature/strain/cached-compliance selector와
  Oracle Top-\(K\)의 error/regret 비교.
  \item Pure cached-compliance 대 normalized deterministic hybrid의 regret,
  high-benefit miss, activation delay, feature latency와 measured-p95 Pareto 비교.
  \item Cached-compliance, sparse probe, low-rate/full counter-state의 selector 비용과 activation delay 비교.
  \item Activation jerk, patch toggle count, hysteresis ablation.
  \item Kernel profiler로 inactive Local block이 실제 skip되는지 확인.
  \item 조건부 G: Oracle benefit에 대한 recall, AUPRC와 calibration error.
  \item Oracle 대 cached-compliance gap이 작으면 learned gate를 구현하지 않는 go/no-go.
  \item Always-on Local과의 quality--latency Pareto.
\end{itemize}

\section{Module 7: Active-patch Local force proposal}

\subsection*{푸는 질문}

\emph{선택된 patch의 known external force를 Local complement에 어떻게 제안하며,
physics-only Local이 부족할 때 어떤 추가 force만 학습할 것인가?}

\subsection*{Solver 분류}

\physics V0는 cached full-anchor known external force의 active/decay restriction과 owner-aware proposal이다.
\learned Shared patch-local equivariant GNN/MLP와 recurrent memory는
physics-only Local gap이 확인된 경우의 조건부 channel이다.
최종 displacement나 next position을 직접 출력하지 않는다.

\subsection*{입력}

Active patch \(r\)마다 다음을 canonical/co-rotated local frame에 넣는다.

\begin{equation}
\begin{aligned}
  \xi_r^t=\{&
  R_r\widetilde x_G^t,\;
  R_r\widetilde v_G^t,\;
  R_r\Psi_{\mathrm{sup}} z_{\mathrm{sup}}^{t-1},\;
  R_r\Psi_{\mathrm{sup}}\dot z_{\mathrm{sup}}^{t-1},\\
  &\epsilon_r^t,\kappa_r^t,\dot\epsilon_r^t,\;
  W_r^t,\nabla W_r^t,\;
  \vartheta_r,\mathcal C_r,\;
  h_r^{t-1}\}.
\end{aligned}
\end{equation}

\subsection*{출력}

Module 3--4와 explicit-force ledger에서 이미 anchor에 조립한
\(f_{\mathrm{cache,ext}}^t\)를 재평가하지 않고

\begin{equation}
  f_{r,\mathrm{base}}^t
  =R_r f_{\mathrm{cache,ext}}^t,
  \qquad r\in\mathcal S_t=\mathcal A_t\cup\mathcal D_t
\end{equation}

로 gather한다. Owner weight는 여기서 다시 곱하지 않고 Module 10 scatter에서 한 번만 적용한다.
따라서 V0 base proposal은 추가 wind query나 analytic-law evaluation이 아니라 restriction 비용만 든다.
이는 network output이 아니며 \texttt{analytic aero}가 physical owner다.
Projective/compliant/support target는 식~\eqref{eq:local-dynamics}의
\(Q_{\mathcal S_t,\mathrm{op}}\)가 직접 소유하므로 patch base proposal에 넣지 않는다.

V0에서는 \(\widehat{\Delta f}_{r,c}=0\)으로 두고 physics-only Local을 먼저 검증한다.
Full-resolution Local과의 반복 가능한 gap, 식~\eqref{eq:patch-force-label}의 유일한/검증된 label,
그리고 예상 accuracy--cost 이득이 모두 존재할 때만

\begin{equation}
  (\widehat{\Delta f}_{r,\mathrm{aero}}^t,
   \widehat{\Delta f}_{r,\mathrm{struct}}^t,
   h_r^t)
  =
  N_\theta(\xi_r^t)
  \label{eq:missing-force-net}
\end{equation}

을 활성화한다.

Teacher force breakdown이 신뢰 가능하면 aerodynamic residual과 structural/model residual을 분리한다.
그렇지 않으면 conditional F를 활성화하지 않는다. 두 residual은 Module 10에서 서로 다른
constraint/reaction type을 소비하므로, source가 불명확한 single total-defect head로 합치지 않는다.

Stage 7 orchestration의 최종 \texttt{PatchForceBatch}는
\begin{equation}
  \left\{
    f_{r,\mathrm{base}}^t,
    \widehat{\Delta f}_{r,\mathrm{aero}}^t,
    \widehat{\Delta f}_{r,\mathrm{struct}}^t
  \right\}_{r\in\mathcal S_t}
\end{equation}
의 세 channel을 분리해 담는다.
첫 channel은 analytic aero, gravity와 explicit non-attachment external force를 포함한
known base external force이고 뒤의 두 channel만 조건부 learned missing force다.
Compliant target와 hard-support lift는 operator-owned
\(Q_{\mathcal S_t,\mathrm{op}}^t\)에만 들어가며 이 channel에 중복 포함하지 않는다.
조건부 F network는 \(r\in\mathcal A_t\)에서만 실행한다.
Patch가 decay로 전환되면 전환 frame의 마지막 learned proposal을 저장하고
\(\gamma_r\)와 함께 0으로 fade하며 새 expert call은 하지 않는다.
Base proposal은 \(r\in\mathcal D_t\)에서도 current cached force를 gather하되
같은 \(\gamma_r\)로 줄여 abrupt force removal을 피한다.
\(\gamma_r=0\) 뒤에는 두 excitation 모두 0이고 기존 Local state만 물리적으로 감쇠한다.
Module 10과 9는 이 channel들을 합쳐 owner를 잃지 않고 각각 조립·투영한다.

\begin{itemize}[leftmargin=1.5em]
  \item Aerodynamic head는 external missing force이며 nonzero net wrench를 가질 수 있다.
  \item Structural head는 internal correction이므로 free object 부분에서는 net internal force/torque가 0이어야 한다.
  \item Attachment-related correction은 explicit reaction channel로 분리한다.
\end{itemize}

\subsection*{Network를 유지하는 조건}

Network는 알려진 inertia, elasticity, damping을 대체하지 않는다.
다음처럼 analytic하게 정확히 풀면 runtime 비용이 커지거나 target 정보가 부족한 항만 근사한다.

\begin{itemize}[leftmargin=1.5em]
  \item Global/Local basis truncation,
  \item large-deformation thin-shell nonlinearity,
  \item topology/material distillation error,
  \item Teacher-FSI가 있을 때의 unsteady aerodynamic defect.
\end{itemize}

정확한 topology와 계산 여유가 있다면 active patch에 fine FEM/shell solve를 수행할 수 있다.
이는 반드시 포함해야 할 physics-only Local baseline이다.
Reduced physics-only가 같은 quality를 내면 missing-force network는 no-go로 제거하며,
selector/dispatcher 연구는 physics-only Local을 대상으로 독립적으로 계속할 수 있다.

\subsection*{출력과 다음 단계}

Patch force proposals는 즉시 적분하지 않는다.
Module 10의 overlap-aware constrained assembly와 Module 9의 complement decomposition을 거친 뒤
Module 8의 단일 Local physical system에 들어간다.

\subsection*{필수 검증}

\begin{itemize}[leftmargin=1.5em]
  \item V0 analytic Local proposal의 frame/owner/virtual-work와 finite-difference force test.
  \item Full-resolution Local 대비 physics-only reduced Local의 spectrum/rollout error.
  \item 조건부 F: force target error와 multi-step rollout error를 함께 평가한다.
  \item 조건부 F: direct displacement network 대비 wind-stop decay/phase 안정성과
  physics-only Local 대비 learned defect의 순수 이득을 측정한다.
  \item Teacher-FSI가 없을 때 unsteady-aero claim이 결과표에 섞이지 않게 한다.
\end{itemize}

\section{Module 8: Single Local complementary physics update}

\subsection*{푸는 질문}

\emph{조립·투영된 Local force channel을 받았을 때 Global이 표현하지 못하는 Local state가
관성, 복원, 감쇠를 따라 어떻게 움직이는가?}

\subsection*{Solver 분류}

\physics Active/decay Local basis 위의 reduced second-order dynamics다.

\subsection*{입력}

\begin{itemize}[leftmargin=1.5em]
  \item Module 10과 9를 channel별로 거친
  \(Q_{\ell,\mathrm{base}}^{\perp,t}\),
  \(\Delta Q_{\ell,\mathrm{aero}}^{\perp,t}\),
  \(\Delta Q_{\ell,\mathrm{struct}}^{\perp,t}\),
  \item active/decay basis \(\Psi_{\mathcal S_t}\),
  \item 식~\eqref{eq:patch-cross-blocks}에서 gather한
  \(M_{\mathcal S_t\mathcal S_t},C_{\mathcal S_t\mathcal S_t}\)와
  fixed projective Hessian \(K_{\mathcal S_t\mathcal S_t}^{0}\),
  \item operator-owned target/lift
  \(Q_{\mathcal S_t,\mathrm{op}}^t
  :=\Psi_{\mathcal S_t}^\top(h_K^t+h_C^t+f_{\mathrm p}^t)\),
  \item predictor Global state \(\widetilde q,\dot{\widetilde q}\),
  \item 이전 \(z,\dot z\).
\end{itemize}

\subsection*{Local equation}

\begin{equation}
\begin{aligned}
M_{\mathcal S_t\mathcal S_t}\ddot z_{\mathcal S_t}
  +C_{\mathcal S_t\mathcal S_t}\dot z_{\mathcal S_t}
  +K_{\mathcal S_t\mathcal S_t}^{0}z_{\mathcal S_t}
  =&\;
  Q_{\mathcal S_t,\mathrm{base}}^{\perp}
  +\Delta Q_{\mathcal S_t,\mathrm{aero}}^{\perp}
  +\Delta Q_{\mathcal S_t,\mathrm{struct}}^{\perp}
  +Q_{\mathcal S_t,\mathrm{op}}
  -C_{\mathcal S_t g}\dot{\widetilde q}
  -K_{\mathcal S_t g}^{0}\widetilde q.
  \label{eq:local-dynamics}
\end{aligned}
\end{equation}

\subsection*{Discrete solve backend와 factorization 계약}

Global predictor의 midpoint state를
\(\widetilde q^{t-1/2}:=(q^{t-1}+\widetilde q^t)/2\),
\(\dot{\widetilde q}^{t-1/2}:=(\dot q^{t-1}+\dot{\widetilde q}^t)/2\)라 하고,
한 Local step 동안의 held RHS를

\begin{equation}
\begin{aligned}
  F_{\mathcal S_t,\mathrm{held}}^t
  :=&\;Q_{\mathcal S_t,\mathrm{base}}^{\perp,t}
  +\Delta Q_{\mathcal S_t,\mathrm{aero}}^{\perp,t}
  +\Delta Q_{\mathcal S_t,\mathrm{struct}}^{\perp,t}
  +Q_{\mathcal S_t,\mathrm{op}}^t\\
  &-C_{\mathcal S_tg}\dot{\widetilde q}^{t-1/2}
  -K_{\mathcal S_tg}^{0}\widetilde q^{t-1/2}
\end{aligned}
\end{equation}

로 둔다. Canonical implicit midpoint의 active/decay effective system은

\begin{equation}
  A_{\mathcal S_t}^{\mathrm{eff}}
  =\frac{4}{\Delta t^2}M_{\mathcal S_t\mathcal S_t}
  +\frac{2}{\Delta t}C_{\mathcal S_t\mathcal S_t}
  +K_{\mathcal S_t\mathcal S_t}^{0},
  \qquad
  A_{\mathcal S_t}^{\mathrm{eff}}\Delta z_{\mathcal S_t}
  =2F_{\mathcal S_t,\mathrm{held}}^t
  -2K_{\mathcal S_t\mathcal S_t}^{0}z_{\mathcal S_t}^{t-1}
  +\frac{4}{\Delta t}M_{\mathcal S_t\mathcal S_t}
  \dot z_{\mathcal S_t}^{t-1}.
  \label{eq:local-effective-system}
\end{equation}

여기서 \(\Delta z_{\mathcal S_t}:=z_{\mathcal S_t}^t-z_{\mathcal S_t}^{t-1}\)이고
\(\dot z_{\mathcal S_t}^{t}=2\Delta z_{\mathcal S_t}/\Delta t
-\dot z_{\mathcal S_t}^{t-1}\)다.
\(\mathcal S_t\)가 바뀌면 principal matrix와 factor도 바뀌므로
모든 frame을 cached triangular solve라고 가정하지 않는다.
V0 correctness backend는 exact principal block을 gather해 Cholesky/LDL factorization하고,
이 same-\(\mathcal S_t\) exact direct solve를 수치 근사의 oracle로 둔다.
Always-on full-superset solve는 full-Local 물리 upper bound와 latency reference이며,
selected solve의 수치 oracle과 동일시하지 않는다.
반복 active-set signature에는 bounded factor cache를 사용할 수 있다.
그 cache key는 package/operator version, integrator, \(\Delta t\)와
\(\operatorname{sig}(\mathcal S_t)\)를 모두 포함하고 하나라도 바뀌면 invalidate한다.
Projective rotation은 RHS만 바꾸므로 위 key가 같을 때 factor reuse는 exact다.
후속 consistent-tangent corotation을 시험한다면 tangent version을 key에 넣고 stale reuse를 금지한다.
Incremental Cholesky update/downdate는 same-set exact direct solve와 matrix residual 및 rollout을 맞춘 뒤에만 허용한다.

\(\mathfrak G(\mathcal S)\)를 fixed atomic group을 \(\mathcal S\)에 제한한 block partition이라 하자.
\(A_{\mathcal S}^{\mathrm{eff}}\succ0\)인 V0에서
strong-overlap group block으로 만든
\(D_{\mathcal S}=\operatorname{blkdiag}_{\mathcal G\in\mathfrak G(\mathcal S)}
(A_{\mathcal G\mathcal G}^{\mathrm{eff}})\succ0\)와

\begin{equation}
  \chi_{\mathcal S}^{\mathrm{off}}
  =
  \left\|
  D_{\mathcal S}^{-1/2}
  \left(A_{\mathcal S}^{\mathrm{eff}}-D_{\mathcal S}\right)
  D_{\mathcal S}^{-1/2}
  \right\|_2
  \label{eq:off-group-coupling}
\end{equation}

를 cross-group coupling 지표로 기록한다.
Group별 독립 solve는 \(A_{\mathcal S}^{\mathrm{eff}}\)를
\(D_{\mathcal S}\)로 바꾸는 block-diagonal approximation이며,
\(-\sum_{\mathcal H\ne\mathcal G}A_{\mathcal G\mathcal H}^{\mathrm{eff}}
\Delta z_{\mathcal H}\)를 누락한다.
이 항은 same-step effective Local group-to-group coupling load다.
Module 11은 Local solve 뒤에 Global만 고치고 same-frame Local state를 다시 풀지 않는다.
설령 재결합하더라도 Global correction이 Local RHS에 만들 수 있는 방향은
\(B_{\mathcal S}^{\mathrm{eff}}
=K_{\mathcal Sg}^{0}+2\Delta t^{-1}C_{\mathcal Sg}\)의 range에 한정되지만,
누락 load는 일반적인 \(\R^{r_\ell^{\mathcal S}}\) vector다.
따라서 누락 load가 \(\operatorname{range}(B_{\mathcal S}^{\mathrm{eff}})\)에 든다는
특수 조건 없이는 Global corrector가 이를 복원할 수 없다.
따라서 independent group solve는 mainline이 아니며,
식~\eqref{eq:off-group-coupling}과 cross-group gust rollout/energy tolerance가 모두 작을 때의 ablation으로만 허용한다.
Pre-factorized group block은 cross-block matvec를 유지하는 block-preconditioned iterative solve의
preconditioner로는 사용할 수 있다. Backend 선택은 direct factorization, factor cache,
incremental update와 block-PCG를 실제 rank와 active-set trace에서 profile한 뒤 결정한다.

Smooth activation \(\gamma_r^t\)는 식~\eqref{eq:raw-overlap-assembly}의
patch force에 정확히 한 번 적용한다. Module 8 RHS에 다시 곱하지 않는다.
또한 activation은 \(M_{\mathcal S_t\mathcal S_t},C_{\mathcal S_t\mathcal S_t},
K_{\mathcal S_t\mathcal S_t}^{0}\)에 절대 곱하지 않는다.
즉 selector가 off되어도 이미 존재하는 Local state의 inertia, restoring, damping은 계속 작동한다.

\(Q_{\ell,\mathrm{base}}^{\perp}\)는 cached known external force를 overlap-aware하게 조립하고
Global force span을 제거한 Local generalized load다.
두 \(\Delta Q^\perp\)만 learned missing-force channel이다.
조건부 F가 no-go인 V0에서는 두 channel을 생성하지 않으며 0으로 해석한다.
세 channel은 서로 다른 owner/lineage를 유지한 채 같은 Module 10--9 경로를 거치고,
Module 8에서만 합산된다.
Decay patch에서는 새 excitation과 learned defect를 서서히 줄이되,
기존 \(z,\dot z\)는 물리적으로 감쇠한다.

\subsection*{출력}

\begin{equation}
  z_{\mathcal S_t}^t,\quad\dot z_{\mathcal S_t}^t,\quad
  E_{\mathcal S_t}^t
  =
  \frac12\dot z_{\mathcal S_t}^{t\top}
  M_{\mathcal S_t\mathcal S_t}\dot z_{\mathcal S_t}^t
  +E_{\mathrm{str},\mathcal S_t}^t
  (\widetilde q^t,z_{\mathcal S_t}^t).
\end{equation}

두 번째 항은 식~\eqref{eq:projective-structural-energy}를 현재
\(u_{\mathrm p}+\Phi\widetilde q+\Psi z\)에서 owner-consistent하게 평가한 값이며,
fixed Hessian만으로 target affine term을 누락한 \(\tfrac12z^\top K^0z\)로 대체하지 않는다.
State 종료와 decay threshold에는 kinetic/potential energy의 동일한 backend 정의를 사용한다.

결과를 fixed-superset slot에 scatter하고 dormant slot은 0으로 유지한다.
Inter-patch cross energy 때문에 patch self-energy의 합은 \(E_{\mathcal S_t}\)와 일반적으로 같지 않다.
Energy 기반 deactivation은 overlap/coupling group 전체 energy와 해당 patch에 작용하는 cross-block load를 함께 본다.

Local state는 Module 11의 corrected geometry와 Module 14의 decay 판단으로 넘어간다.

\subsection*{왜 patch별 integrate-then-average가 아닌가}

각 patch가 독립된 mass를 갖고 적분한 뒤 위치를 평균하면
overlap anchor의 실제 mass, energy, reaction이 정의되지 않는다.
본 방법은 patch force를 먼저 하나의 global anchor force로 조립하고,
하나의 complement basis system에서 한 번만 적분한다.

\subsection*{필수 검증}

\begin{itemize}[leftmargin=1.5em]
  \item Gate-off/zero-wind에서 Local energy가 증가하지 않는가?
  \item Patch activation/deactivation에서 position, velocity, acceleration이 연속적인가?
  \item Local physics 없이 direct displacement만 더한 방법보다 장시간 안정적인가?
  \item Timestep 변화에 대한 민감도가 허용 범위인가?
\end{itemize}

\section{Module 9: Global-span contract와 Local generalized projection}

\subsection*{푸는 질문}

\emph{Patch별 Local force channel 중 Global이 이미 담당하는 성분과
진짜 Local complement 성분을 어떻게 구분할 것인가?}

\subsection*{Solver 분류}

\deterministic Support-routed 또는 joint-constrained force의 contract 검사와 generalized projection이다.

\subsection*{Force-space complement}

Global virtual work가 0인 force를 만들기 위한 projector는

\begin{equation}
  \Pi_f^\perp
  =
  I-M\Phi(\Phi^\top M\Phi)^{-1}\Phi^\top.
  \label{eq:force-complement}
\end{equation}

이 projector는 unit test와 label construction의 기준이다.
Runtime에서는 Module 10의 external restricted route 또는 structural augmented KKT가
보존 제약과 \(\Phi^\top f=0\)를 동시에 만족하므로 여기서 \(\Pi_f^\perp\)를 다시 적용하지 않는다.
순차 projection은 앞서 만족한 wrench constraint를 깨뜨릴 수 있다.

Channel 집합을

\begin{equation}
  \mathcal C_\ell
  =\{\mathrm{base},\mathrm{miss\mbox{-}aero},\mathrm{miss\mbox{-}struct}\}
\end{equation}

로 두고 route/constraint-verified force \(f_{\mathrm{asm},c}\)마다

\begin{equation}
  f_{\ell,c}^\perp
  =
  f_{\mathrm{asm},c},
  \qquad
  \Phi^\top f_{\ell,c}^\perp=0,
  \quad c\in\mathcal C_\ell.
\end{equation}

Local generalized force는

\begin{equation}
  Q_{\ell,c}^\perp
  =
  \Psi_{\mathcal S_t}^\top f_{\ell,c}^\perp,
  \quad c\in\mathcal C_\ell
\end{equation}

로 만든다.

\subsection*{제거하지 말아야 할 힘}

다음 두 항은 중복이 아니므로 버리지 않는다.

\begin{enumerate}[leftmargin=1.7em]
  \item \(K_{g\ell}z+C_{g\ell}\dot z\)가 나타내는 실제 structural coupling.
  \item Local correction이 표면 normal/area/velocity를 바꾸어 새로 생긴 aerodynamic generalized-force delta.
\end{enumerate}

First-paper method는 별도의 Global missing-force head를 두지 않는다.
\(\Phi^\top\widetilde f\)는 diagnostics에만 기록하고,
조건부 Local head의 Global leakage는 loss와 projection으로 억제한다.

\subsection*{출력과 다음 단계}

세 channel의 \(Q_{\ell,c}^\perp\)는 owner/lineage를 보존한 채 Module 8로,
Global-span leakage norm

\begin{equation}
  \epsilon_{\mathrm{leak}}
  =
  \max_{c\in\mathcal C_\ell}
  \frac{\|\Phi^\top f_{\mathrm{asm},c}\|_2}
  {\|f_{\mathrm{asm},c}\|_2+\eps}
\end{equation}

은 diagnostics와 training loss로 넘어간다.

\section{Module 10: Overlap-aware conservative patch assembly}

\subsection*{푸는 질문}

\emph{여러 active patch가 같은 anchor에서 서로 다른 force를 제안할 때
seam, double counting, 가짜 translation/rotation 없이 하나의 force를 만드는가?}

\subsection*{Solver 분류}

\constrained Sparse scatter, support-restricted projection과 필요한 channel의
equality-constrained weighted projection이다.
Anchor-force primal 자체는 작지 않다. External default는 explicit restricted remainder를 쓰고,
structural/reaction 또는 추가 constraint가 있는 channel만 unique-anchor support의 small dual Schur를 푼다.

\subsection*{초기 조립}

각 channel \(c\in\mathcal C_\ell\)을 합치기 전에 따로 조립한다.

\begin{equation}
  \widetilde f_c^t
  =
  P_{\mathrm{free}}
  \sum_{r\in\mathcal S_t}
  R_r^\top W_r
  \gamma_r^t
  f_{r,c}^t,
  \qquad c\in\mathcal C_\ell.
  \label{eq:raw-overlap-assembly}
\end{equation}

여기서 \(\mathcal S_t=\mathcal A_t\cup\mathcal D_t\)이며,
Decay patch는 expert를 실행하지 않고 current known-base proposal과
전환 frame에 저장한 마지막 learned proposal을 \(\gamma_r^t\)로 fade-out한다.
이 식이 activation weight의 유일한 적용 지점이다.
\(P_{\mathrm{free}}\)도 Local excitation에 정확히 한 번 적용하며,
hard-row force/reaction은 아래 별도 support ledger만 소유한다.

Partition weights는 displacement continuity에 유용하지만
힘과 torque 보존을 자동으로 보장하지 않는다.
따라서 식~\eqref{eq:raw-overlap-assembly}만으로 끝내지 않는다.
\(\mathcal I_t\)를 \(\mathcal S_t\)의 core/halo가 만지는 unique anchor support라 한다.

\subsection*{Support-restricted Global routing과 feasibility certificate}

\reviewvfourpartial{Full-space Global representative가 active support 밖으로 퍼진 뒤
support-limited Local force에 complement와 full-remainder wrench를 동시에 요구하면 joint system이
불가능해질 수 있다는 Critical B의 진단은 수용한다. Global representative도 같은 unique-anchor
support에서 minimum-norm으로 route한다. 다만 \(\operatorname{rank}(\Phi_t)<r_g\)라는 사실만으로
support를 invalid 처리하지 않는다. Raw proposal이 support-contained이면 현재 RHS는 원래부터
\(\operatorname{range}(\Phi_t^\top)\)에 있으며, 필요한 판정은 per-RHS consistency,
nonzero spectrum conditioning과 force amplification이다.}

\(n=3K\)라 하고 \(\mathcal I_t\) 안의 free Cartesian DOF 수를 \(n_t\)라 둔다.
Local excitation injection과 restricted operator는

\begin{equation}
\begin{gathered}
  E_{f,t}\in\R^{n\times n_t},\qquad E_{f,t}^\top E_{f,t}=I,
  \qquad A_{\mathrm h}E_{f,t}=0,\\
  \bar f_c^t=E_{f,t}^\top\widetilde f_c^t,\qquad
  \Phi_t=E_{f,t}^\top\Phi,\qquad
  A_{I,c}^t=A_{\mathrm{cons},c}^tE_{f,t}
  \label{eq:support-restricted-operators}
\end{gathered}
\end{equation}

로 둔다. Raw proposal support와 external channel의 restricted Global route는

\begin{equation}
\begin{aligned}
  \epsilon_{\mathrm{supp},c}^t
  &:={}\frac{\|(I-E_{f,t}E_{f,t}^\top)\widetilde f_c^t\|_2}
              {\|\widetilde f_c^t\|_2+\eps},\\
  q_{g,c}^{\mathrm{full},t}
  &:={}\Phi^\top\widetilde f_c^t,
  &q_{g,c}^{I,t}&:=\Phi_t^\top\bar f_c^t,\\
  \epsilon_{\mathrm{work},c}^t
  &:={}\frac{\|q_{g,c}^{\mathrm{full},t}-q_{g,c}^{I,t}\|_2}
              {\|q_{g,c}^{\mathrm{full},t}\|_2+\eps},\\
  H_{\Phi,t}&:=\Phi_t^\top W_{f,t}^{-1}\Phi_t,\\
  \widehat f_{g,I,c}^t
  &:={}W_{f,t}^{-1}\Phi_tH_{\Phi,t}^{\dagger}q_{g,c}^{I,t},\\
  r_{I,c}^t&:=\bar f_c^t-\widehat f_{g,I,c}^t,
  \qquad c\in\mathcal C_{\mathrm{ext}}
  \label{eq:support-restricted-global-routing}
\end{aligned}
\end{equation}

로 정의한다. 여기서 \(\mathcal C_{\mathrm{ext}}=\{\mathrm{base},\mathrm{miss\mbox{-}aero}\}\)이고
\(\widehat f_{g,I,c}\)는
\(\min_g\tfrac12\|g\|_{W_{f,t}}^2\) subject to
\(\Phi_t^\top g=q_{g,c}^{I}\)의 rank-aware minimum-norm 해다.
Pseudoinverse로 불일치를 숨기지 않고

\begin{equation}
\begin{aligned}
  \epsilon_{\mathrm{route},c}^t
  &:={}\frac{\|(I-H_{\Phi,t}H_{\Phi,t}^{\dagger})q_{g,c}^{I,t}\|_2}
              {\|q_{g,c}^{I,t}\|_2+\eps},\\
  \kappa_{\mathrm{route},t}^{+}
  &:={}\frac{\sigma_{\max}(W_{f,t}^{-1/2}\Phi_t)}
              {\sigma_{\min}^{+}(W_{f,t}^{-1/2}\Phi_t)},\\
  \alpha_{\mathrm{route},c}^t
  &:={}\frac{\|\widehat f_{g,I,c}^t\|_2+\|r_{I,c}^t\|_2}
              {\|\bar f_c^t\|_2+\eps}
  \label{eq:support-routing-diagnostics}
\end{aligned}
\end{equation}

를 기록한다. \(\epsilon_{\mathrm{supp}}\)와 \(\epsilon_{\mathrm{work}}\)가 tolerance 안일 때만
support-contained route로 인정한다. \(\operatorname{rank}(\Phi_t)=0\)이고
\(q_{g,c}^{I}=0\)이면 \(\kappa_{\mathrm{route}}^+\)는 N/A로 기록하고 유효한 zero route로 처리한다.
반대로 attainable singular direction이 없는데 \(q_{g,c}^{I}\ne0\)이면 실패다.
External channel은
\(b_{\mathrm{cons},c}^t:=A_{I,c}^tr_{I,c}^t\)로 두며,
\(y_{\mathrm{wit},c}^t=r_{I,c}^t\)가 route tolerance 안에서

\begin{equation}
  \begin{bmatrix}A_{I,c}^t\\\Phi_t^\top\end{bmatrix}
  y_{\mathrm{wit},c}^t
  =
  \begin{bmatrix}b_{\mathrm{cons},c}^t\\0\end{bmatrix}
  \label{eq:support-feasible-witness}
\end{equation}

를 만족하는 explicit feasible witness다. \(\widehat f_{g,I,c}\)는 Global bookkeeping
representative다. Base channel의 \(q_{g,\mathrm{base}}\)는 Module 4--5가 이미 소비했으므로
다시 더하지 않는다. First-paper missing-aero head는 Local-only label/leakage contract를 사용하므로
\(q_{g,\mathrm{miss\mbox{-}aero}}\)도 새 Global load로 route하지 않는다.
그 leakage가 tolerance를 넘으면 conditional F를 거부하거나 별도 versioned Global-defect owner를
후속 branch로 정의하며, 성분을 조용히 버리거나 double count하지 않는다.

\subsection*{Joint conservation--complement correction}

보정 force가 Dormant 영역에 새 excitation을 만들지 않도록 최종 force도 이 support에 제한한다.
External channel은 \(y_{\mathrm{asm},c}=r_{I,c}\),
\(f_{\mathrm{asm},c}=E_{f,t}r_{I,c}\)가 이미 complement와 chosen remainder resultant를
정확히 만족하므로 같은 equality를 다시 Schur로 풀지 않는다.
Additional low-row constraint가 실제로 붙을 때만 아래 generic backend로 보낸다.

Structural/internal channel을 위해 current geometry의 world-wrench map을

\begin{equation}
  \mathcal W_t f
  =
  \begin{bmatrix}
    \sum_k f_k\\[0.2em]
    \sum_k(x_k^t-x_{\mathrm{ctr}}^t)\times f_k
  \end{bmatrix}
  \label{eq:current-world-wrench-map}
\end{equation}

로 둔다. Free object와 hard attachment의 augmented primal/metric/reference/constraint는

\begin{equation}
\begin{array}{c|c|c|c}
 &\zeta_c&\zeta_c^{\mathrm{ref}}&B_c^t\zeta_c=d_c^t\\
\hline
\mathrm{free}
&y_c&\bar f_c
&\begin{bmatrix}\mathcal W_tE_{f,t}\\\Phi_t^\top\end{bmatrix}y_c=0\\[0.8em]
\mathrm{hard}
&\begin{bmatrix}y_c\\\rho_{h,c}\end{bmatrix}
&\begin{bmatrix}\bar f_c\\0\end{bmatrix}
&\begin{bmatrix}
  \mathcal W_tE_{f,t}&\mathcal W_tA_{\mathrm h}^\top\\
  \Phi_t^\top&0
 \end{bmatrix}
 \begin{bmatrix}y_c\\\rho_{h,c}\end{bmatrix}=0
\end{array}
  \label{eq:structural-augmented-constraints}
\end{equation}

로 고정하고, 각각
\(\overline W_c=W_{f,t}\)와
\(\overline W_c=\operatorname{blkdiag}(W_{f,t},W_h)\)를 쓴다.
Compliant attachment force는 operator-owned attachment channel이므로 이 structural patch primal에 넣지 않는다.
Additional independent low-row constraint가 있으면 \(B_c,d_c\)에 명시적으로 append한다.
Generic joint correction은

\begin{equation}
  \zeta_{\mathrm{asm},c}
  =\argmin_\zeta
  \|\zeta-\zeta_c^{\mathrm{ref}}\|_{\overline W_c}^2
  \quad\mathrm{s.t.}\quad B_c^t\zeta=d_c^t,
  \qquad
  f_{\mathrm{asm},c}=E_{f,t}S_y\zeta_{\mathrm{asm},c}.
  \label{eq:conservative-assembly}
\end{equation}

\(S_y\)는 augmented primal의 support-force coordinate만 고른다.
Hard branch의 object-to-support companion은
\(r_{o\rightarrow h,c}=A_{\mathrm h}^\top S_\rho\zeta_{\mathrm{asm},c}\)이다.

\subsection*{Runtime small-Schur backend와 constraint inventory}

식~\eqref{eq:conservative-assembly}의 full anchor force는 \(3K\) DOF이므로
이를 ``작다''고 부르는 근거는 full saddle matrix가 아니라 constraint 수여야 한다.
Support/augmented coordinate에서 runtime은

\begin{equation}
\begin{aligned}
  G_c^t&=B_c^t\overline W_c^{-1}(B_c^t)^\top,\\
  G_c^t\lambda_c^t&=B_c^t\zeta_c^{\mathrm{ref},t}-d_c^t,\\
  \zeta_{\mathrm{asm},c}^t
  &=\zeta_c^{\mathrm{ref},t}
  -\overline W_c^{-1}(B_c^t)^\top\lambda_c^t
\end{aligned}
  \label{eq:runtime-joint-schur}
\end{equation}

로 푼다. \(W_{f,t}\succ0,W_h\succ0\)는 setup-defined block-diagonal metric이다.
External base/missing-aero의 default route에는 이 factor/solve 비용을 청구하지 않고
restricted projection 비용만 보고한다.

Additional low-row constraint가 붙는 channel은 explicit witness만으로 blanket feasibility를 주장하지 않고

\begin{equation}
  \epsilon_{\mathrm{feas},c}^t
  =
  \frac{\|(I-B_c^t(B_c^t)^\dagger)d_c^t\|_2}
       {\|d_c^t\|_2+\eps}
  \label{eq:joint-feasibility-residual}
\end{equation}

와 independent-row rank를 먼저 검사한다.
Free/hard structural 기본식은 \(d_c=0\)이라 zero field가 항상 feasible할 수 있으므로
\(\epsilon_{\mathrm{feas}}\)만으로 support quality를 통과시키지 않는다. 실제 raw correction 보존량은

\begin{equation}
\begin{aligned}
  \epsilon_{\mathrm{adjust},c}^t
  &:={}\frac{\|S_y\zeta_{\mathrm{asm},c}^t-\bar f_c^t\|_{W_{f,t}}}
              {\|\bar f_c^t\|_{W_{f,t}}+\eps},\\
  \alpha_{h,c}^t
  &:={}\frac{\|S_\rho\zeta_{\mathrm{asm},c}^t\|_2}
              {\|\bar f_c^t\|_2+\eps},\\
  \epsilon_{\Psi,c}^t
  &:={}\frac{\|
  \Psi_{\mathcal S_t}^\top E_{f,t}
  (S_y\zeta_{\mathrm{asm},c}^t-\bar f_c^t)\|_2}
  {\|\Psi_{\mathcal S_t}^\top E_{f,t}\bar f_c^t\|_2+\eps}
  \label{eq:structural-schur-adjustment-diagnostics}
\end{aligned}
\end{equation}

로 별도 기록한다. \(\alpha_h\)는 free branch에서 0으로 둔다.
큰 adjustment, reaction amplification 또는 intended Local generalized-work loss가 반복되면
halo/support, \(W_h\), channel parameterization을 고치거나 package를 거부한다.
Feasible rate 100\%만으로 통과시키지 않는다.

\reviewvfiveaccept{Zero RHS와 explicit witness 때문에 feasibility rate가 높아도
sparse execution의 locality와 비용이 보장되지 않는다는 지적을 수용한다.
Primary diagnostic은 \(\alpha_{\mathrm{route}}\), \(\epsilon_{\mathrm{adjust}}\),
\(\alpha_h\), \(\epsilon_\Psi\), halo-expansion rate, unique support size \(n_{\mathcal S}\)와
\(T_{\mathrm{Schur}}\)이며, feasible rate는 이들을 대신하지 않는다.}

Package/operator version, free/hard attachment class, active/decay group과 unique-anchor support,
independent constraint-row set으로 support signature \(\sigma_{\mathrm{sup}}^t\)를 만든다.
Exact active-set ID의 조합 수는 폭발할 수 있으므로 development split에서
attachment class, support-size/rank bin, constraint-row inventory와 conditioning bin으로 정의한
signature 등가 class별 primary-diagnostic cap과 허용 halo 횟수를 fit/freeze한다.
Test의 새로운 exact active-set identity는 검증된 class에 속하고 runtime diagnostic을 통과하면 허용한다.
Unseen class가 나타나거나 frozen cap을 넘으면
halo를 무제한 확대하거나 값을 clip해 adaptive 결과로 계상하지 않는다.
Setup validation에서 발견되면 package를 reject하고, runtime trace에서 발견되면 해당 asset/run을
사전 선언한 full-superset always-on Local backend로 전환한다. Fallback, 최초 초과 시점과
원인 diagnostic을 결과에 함께 보고한다.

즉 ``small''은 anchor primal이 아니라
\(m_c^t=\operatorname{rank}(B_c^t)\)인 dual Schur system만 뜻한다.
중복 row를 제거한 일반 channel의 row 수는 대략 \(r_g+6\)이며 overlap patch 수에 비례해 증가하지 않아야 한다.
Generic \((n_t+m_c^t)\)-dimensional saddle solve는 작은 fixture의 correctness reference일 수 있지만
runtime mainline이 아니다. Full-row-rank SPD Schur에는 Cholesky를 쓰고,
near-rank-deficient system에는 RHS consistency 검사와 pivoted LDL/QR 또는 RRQR/SVD를 사용한다.
Infeasible하면 support halo를 넓히거나 redundant constraint, Global/attachment routing과
channel definition을 먼저 고치며,
regularization이나 순차 projection으로 feasibility failure를 숨기지 않는다.
Halo 확대는 current RHS consistency, conditioning/amplification 또는 additional-constraint failure에
근거하며 \(\operatorname{rank}(\Phi_t)<r_g\) 자체를 확대 조건으로 쓰지 않는다.

\(A_{\mathrm{cons},c}^t\)의 row는 channel별 resultant force/torque와 필요한 low-dimensional
support/reaction constraint로 제한한다. Shared anchor는 scatter 뒤 하나의 force/state만 가지며,
velocity continuity는 식~\eqref{eq:local-dynamics}의 single coupled Local state가 보장한다.
따라서 patch pair마다 velocity-continuity row를 추가하지 않는다.
Internal action--reaction은 가능하면 antisymmetric edge-force parameterization으로 만족시키고,
남은 low-row resultant만 joint solve에 둔다.
Torque row는 current anchor geometry와 center에 따라 변할 수 있으므로 \(B_c^t\) 전체를 setup 상수로 간주하지 않는다.
\(\overline W_c^{-1}\), \(\Phi\) block, force-sum row와 sparsity/index structure만 가능한 범위에서 precompute한다.
Active unique-anchor support, metric version, independent row set, current torque geometry/center와
package/operator version이 모두 같은 signature일 때만 Schur factor를 channel/frame 사이에 재사용한다.
같은 \(B_c^t,W_{f,t}\)에서 RHS만 다른 channel은 같은 frame의 factor를 공유할 수 있지만,
internal/attachment constraint set이 다르면 별도로 갱신한다.

Constraint는 channel별로 다르다.

\begin{itemize}[leftmargin=1.5em]
  \item Structural/internal residual:
  free object에서는 intrinsic full force의 net wrench를 0으로 둔다.
  Hard attachment가 있으면 free-node force 자체에 zero wrench를 강제하지 않고,
  별도 object-to-support reaction과 합친 wrench를 닫는다.
  Compliant attachment force는 internal residual이 아니라 attachment owner로 유지한다.
  \item Aerodynamic/external residual:
  식~\eqref{eq:support-restricted-global-routing}으로 같은 support에서 Global channel을 route하고,
  \(b_{\mathrm{cons},c}=A_{I,c}r_{I,c}\)인 Local remainder의 resultant force/torque를 유지.
  \item Shared anchor:
  owner-weighted scatter 뒤 하나의 anchor force/state를 사용한다.
  Patch별 velocity constraint를 중복 생성하지 않는다.
\end{itemize}

가능하면 structural residual을 antisymmetric edge force로 parameterize하고,
최종 assembled field에서 net force/torque와 Global virtual work를 함께 검사한다.
Joint system이 rank-deficient하거나 infeasible하면 channel/global routing을 고치며,
conservation solve 뒤 별도 complement projection으로 우회하지 않는다.

\subsection*{Free, hard와 compliant reaction convention}

\reviewvfouraccept{Hard attachment가 있는 물체의 free-node structural force에 blanket zero wrench를
강제하는 것은 잘못이다. Free object, hard support와 compliant handle을 서로 다른 reaction ledger로
분리하고, support reaction은 active Local excitation support 밖의 별도 channel로 허용한다.}

식~\eqref{eq:current-world-wrench-map}에 대해 free object의 intrinsic structural correction은
\(\mathcal W_tf_{\mathrm{int}}^t=0\)을 만족해야 한다.
Hard attachment에서는 object가 support에 전달하는 companion reaction을
식~\eqref{eq:structural-augmented-constraints}의
\(r_{o\rightarrow h}^t=A_{\mathrm h}^\top\rho_h^t\)라 정의하고

\begin{equation}
  \mathcal W_t(E_{f,t}y_{\mathrm{int}}^t+r_{o\rightarrow h}^t)=0,
  \qquad
  f_{h\rightarrow o}^t=-r_{o\rightarrow h}^t
  \label{eq:hard-reaction-wrench-balance}
\end{equation}

로 부호를 고정한다. Fixture/reference의 minimum-norm companion은

\begin{equation}
  \rho_h^t
  =\argmin_\rho\frac12\|\rho\|_{W_h}^2
  \quad\mathrm{s.t.}\quad
  \mathcal W_tA_{\mathrm h}^\top\rho
  =-\mathcal W_tE_{f,t}y_{\mathrm{int}}^t
  \label{eq:hard-reaction-companion}
\end{equation}

로 구한다. \(A_{\mathrm h}\Phi=0\)이므로
\(\Phi^\top r_{o\rightarrow h}=(A_{\mathrm h}\Phi)^\top\rho_h=0\)이고,
이 reaction은 complement excitation을 새로 만드는 Local patch force가 아니다.
Channel companion과 달리 runtime total support-on-object reaction은 Module 11 corrector와
최종 \(q,z\)가 확정된 뒤의 \emph{post-step diagnostic}으로만 복원한다.
그 source of truth는 final-state discrete residual이다.
Implicit midpoint timing과 맞추어

\begin{equation}
  u^{\mathrm{mid},t}:=\frac12(u^t+u^{t-1}),\qquad
  v^{\mathrm{mid},t}:=\frac12(v^t+v^{t-1}),\qquad
  a^{\mathrm{disc},t}:=\frac{v^t-v^{t-1}}{\Delta t}
  \label{eq:post-step-midpoint-kinematics}
\end{equation}

를 사용한다.
이 진단에서 사용하는 full-anchor applied external field는

\begin{equation}
  f_{\mathrm{ext,applied}}^t
  :=f_{\mathrm{aero,base}}^t+\Delta f_{\mathrm{aero,corr}}^t
  +f_{\mathrm{gravity}}^t
  +f_{\mathrm{other,nonattach}}^t
  \label{eq:full-applied-external-force-ledger}
\end{equation}

로 정의한다. Prescribed hard support, compliant attachment, projective target와
조건부 missing-force channel은 이 항에 넣지 않고 각각의 owner에서만 한 번 더한다.
V0의 hard row가 orthonormal coordinate selector라는 계약 아래 hard lift
\(f_{\mathrm p}\)를 제외한 full-space owned force와 조건부 defect field를

\begin{equation}
\begin{aligned}
  f_{\mathrm{owned,full}}^t
  &:={}f_{\mathrm{ext,applied}}^t+h_K^t+h_C^t
  -C^0v^{\mathrm{mid},t}-K^0u^{\mathrm{mid},t}\\
  &\quad+\mathbf1_F f_{\mathrm{asm,miss\mbox{-}aero}}^t
  +\mathbf1_F f_{\mathrm{miss\mbox{-}struct,full}}^t,\\
  f_{\mathrm{miss\mbox{-}struct,full}}^t
  &:={}E_{f,t}S_y\zeta_{\mathrm{asm,miss\mbox{-}struct}}^t
  +r_{o\rightarrow h,\mathrm{miss\mbox{-}struct}}^t,\\
  \rho_{h\rightarrow o}^{\mathrm{tot},t}
  &:={}A_{\mathrm h}(Ma^{\mathrm{disc},t}-f_{\mathrm{owned,full}}^t),\\
  f_{h\rightarrow o}^{\mathrm{tot},t}
  &:={}A_{\mathrm h}^\top\rho_{h\rightarrow o}^{\mathrm{tot},t},
  &\qquad
  r_{o\rightarrow h}^{\mathrm{tot},t}:=-f_{h\rightarrow o}^{\mathrm{tot},t},\\
  r_{\mathrm{ROM}}^t
  &:={}P_{\mathrm{free}}(Ma^{\mathrm{disc},t}-f_{\mathrm{owned,full}}^t)
\end{aligned}
  \label{eq:hard-total-reaction-residual}
\end{equation}

로 복원한다. 여기서 \(f_{\mathrm{ext,applied}}\)는 Module 5가 소비한 full-anchor base
external field와 Module 11이 실제 적용한 corrected-aero anchor delta를 각각 한 번만 포함하며,
그 support-routed \(f_{\mathrm{asm,base}}\) 또는 corrected generalized delta를 다시 더하지 않는다.
Missing-aero는 별도 Global-defect owner가 없는 Local complement field만 한 번 더하고,
missing-struct는 free support force와 hard-row intrinsic companion을 합친 full field로 기록한다.
\(f_{\mathrm p}\)는 prescribed lift를 free-coordinate equation으로 옮긴 항일 뿐
full-space physical reaction residual에 다시 넣지 않는다.
Conditional F가 켜졌을 때 teacher missing-aero의 hard-row residual을 별도 reaction owner가
모델링하지 않는다면, 이 값은 \emph{전체 teacher reaction}이 아니라
\emph{현재 modeled force ledger에 대한 reaction}으로만 보고한다.
Full teacher-reaction accuracy는 그 hard-row residual이 tolerance 아래이거나
별도 reaction-only channel을 명시한 경우에만 주장한다.

Reduced trajectory는 일반적으로 full free-DOF equation을 pointwise 만족하지 않으므로 정확한 분해는

\begin{equation}
  Ma^{\mathrm{disc},t}
  =
  f_{\mathrm{owned,full}}^t
  +f_{h\rightarrow o}^{\mathrm{tot},t}
  +r_{\mathrm{ROM}}^t
  \label{eq:hard-reaction-rom-residual-decomposition}
\end{equation}

이다. Synchronized monolithic reference에서 같은 시점/force를 모두 소비했을 때에만
\(\Phi^\top r_{\mathrm{ROM}}=0\)과
\(\Psi_{\mathcal S_t}^\top r_{\mathrm{ROM}}=0\)을 기대한다.
One-corrector partitioned runtime은 이를 항등식으로 가정하지 않고 두 virtual-work residual,
\(\|r_{\mathrm{ROM}}\|\), reaction error와 teacher/full-free fixture 차이를 보고한다.
즉 hard reaction 하나로 reduced truncation residual까지 닫았다고 주장하지 않는다.
Environment-side reaction은
\(r_{o\rightarrow h}^{\mathrm{tot}}\)로 부호가 반대다.

Compliant attachment는

\begin{equation}
\begin{aligned}
  s_{\mathrm c}^t
  &:={}D_{\mathrm c}(A_{\mathrm c}u^t-d_{\mathrm c}^t)
  +V_{\mathrm c}(A_{\mathrm c}\dot u^t-\dot d_{\mathrm c}^t),\\
  f_{\mathrm c,o}^t&:=-A_{\mathrm c}^\top s_{\mathrm c}^t,
  &r_{o\rightarrow c}^t&:=s_{\mathrm c}^t
  \label{eq:compliant-reaction-ledger}
\end{aligned}
\end{equation}

로 object force와 handle reaction을 분리한다. Common admissible motion
\(\delta d_{\mathrm c}=A_{\mathrm c}\delta u\)에서는
\((f_{\mathrm c,o})^\top\delta u+(r_{o\rightarrow c})^\top\delta d_{\mathrm c}=0\)이어야 한다.
Generic \(\mathcal U_{\mathrm c}\)를 허용하는 package는 handle-control reaction adjoint와
world application point도 저장하며, 그렇지 않으면 compliant handle의 world torque를 주장하지 않는다.

\subsection*{실제 계산 순서}

\begin{enumerate}[leftmargin=1.7em]
  \item Module 7 patch proposals 생성.
  \item 본 Module 10에서 overlap 조립, support-restricted Global route와 필요한 channel의
  conservation--complement augmented joint solve.
  \item Module 9에서 Global-span contract 검사와 Local generalized projection.
  \item Module 8에서 단일 Local system 적분.
\end{enumerate}

\subsection*{출력과 다음 단계}

Channel별 conservative \(f_{\mathrm{asm},c}\), hard structural companion,
overlap constraint residual이 Module 9와 Module 11로 넘어간다.
식~\eqref{eq:hard-total-reaction-residual}의 total reaction과 \(r_{\mathrm{ROM}}\)은
Module 11 corrector 뒤 최종-state diagnostics에서만 계산한다.

\subsection*{필수 검증}

\begin{itemize}[leftmargin=1.5em]
  \item 단순 averaging 대비 patch boundary displacement/velocity jump.
  \item 최종 assembled net force/torque error.
  \item Free object에서 artificial translation/rotation drift.
  \item Patch partition이 달라져도 결과가 크게 달라지지 않는가?
  \item Signature별 \(\epsilon_{\mathrm{supp}},\epsilon_{\mathrm{work}},\epsilon_{\mathrm{route}},
  \kappa_{\mathrm{route}}^+,\alpha_{\mathrm{route}},\epsilon_{\mathrm{feas}},
  \epsilon_{\mathrm{adjust}},\alpha_h,\epsilon_\Psi\),
  independent row rank, feasible rate와 force amplification.
  \item Halo expansion 횟수/final support ratio, factor rebuild/cache hit와 Schur p50/p95/p99.
  \item Free/hard/compliant fixture의 sign-resolved reaction wrench, post-step \(r_{\mathrm{ROM}}\),
  Global/Local virtual-work residual과 teacher/full-free reaction error.
\end{itemize}

\section{Module 11: Local-to-Global aeroelastic feedback와 corrector}

\subsection*{푸는 질문}

\emph{Local flutter/fold가 표면 방향과 구조 반작용을 바꾸었을 때
그 변화가 물체 전체의 다음 움직임에 어떻게 영향을 주는가?}

\subsection*{Solver 분류}

\physics Active-region aerodynamic reevaluation과 작은 reduced Global corrector다.

\reviewaccept{Corrector 비용을 작은 \(r_g\)-dimensional solve 하나로 요약하면 안 된다는 지적을 수용한다.
Active/halo geometry, normal, area와 aerodynamic force 재평가가
reduced solve보다 비쌀 수 있으므로 corrected-aero reevaluation, delta projection,
structural-delta construction과 Global corrector solve를 각각 분리 측정한다.}

\subsection*{Predictor와 corrected geometry}

\begin{equation}
  x_{\mathrm{pred}}^t
  =
  x^0+u_{\mathrm p}^t+\Phi\widetilde q^t
  +\Psi_{\mathrm{sup}}z_{\mathrm{sup}}^{t-1},
  \qquad
  v_{\mathrm{pred}}^t
  =
  \dot u_{\mathrm p}^t+\Phi\dot{\widetilde q}^t
  +\Psi_{\mathrm{sup}}\dot z_{\mathrm{sup}}^{t-1},
\end{equation}

\begin{equation}
  x_{\mathrm{corr}}^t
  =
  x^0+u_{\mathrm p}^t+\Phi\widetilde q^t
  +\Psi_{\mathrm{sup}}z_{\mathrm{sup}}^t,
  \qquad
  v_{\mathrm{corr}}^t
  =
  \dot u_{\mathrm p}^t+\Phi\dot{\widetilde q}^t
  +\Psi_{\mathrm{sup}}\dot z_{\mathrm{sup}}^t.
\end{equation}

Active/halo region에서 Module 2와 3의 normal/area/relative-wind force를 다시 평가한다.

\subsection*{Aerodynamic delta}

\begin{equation}
\begin{aligned}
  \Delta f_{\mathrm{aero,corr}}^t
  &:={}
  \sum_{s\in\mathcal S_t^{\mathrm{halo}}}
  (T_s^{A\rightarrow S})^\top
  \left(
    f_{a,s}^{\mathrm{corr},t}-f_{a,s}^{\mathrm{pred},t}
  \right),\\
  \Delta Q_g^{\mathrm{aero},t}
  &:={}\Phi^\top\Delta f_{\mathrm{aero,corr}}^t
  =
  \sum_{s\in\mathcal S_t^{\mathrm{halo}}}
  \Phi_s^\top
  \left(
    f_{a,s}^{\mathrm{corr},t}-f_{a,s}^{\mathrm{pred},t}
  \right).
\end{aligned}
  \label{eq:corrected-aero-delta}
\end{equation}

식~\eqref{eq:fixed-global-sample-jacobian}의 V0에서는 두 state가 같은 \(\Phi_s\)를 쓴다.
첫째 줄의 conservative transpose scatter를 reaction/force ledger용 anchor field로 저장하고,
둘째 줄만 reduced corrector RHS에 사용한다.
Optional nonlinear sample-map backend만 force와 Jacobian을 predictor/corrected state에서 각각 평가하고,
그 reconstruction/update 비용을 별도 profiler 항목으로 공개한다.
Same-frame corrector에는 식~\eqref{eq:corrected-aero-delta}의 차이만 추가한다.
Module 5에서 이미 사용한 base aerodynamic force 전체를 다시 넣으면 double counting이다.

\subsection*{Structural coupling}

\begin{equation}
\begin{aligned}
  Q_{g\leftarrow\ell}^{\mathrm{mid},t}
  &:={}
  -C_{g,\mathrm{sup}}\dot z_{\mathrm{sup}}^{t-1/2}
  -K_{g,\mathrm{sup}}^{0}z_{\mathrm{sup}}^{t-1/2},\\
  z_{\mathrm{sup}}^{t-1/2}
  &:={}
  \frac12(z_{\mathrm{sup}}^t+z_{\mathrm{sup}}^{t-1}),
  \qquad
  \dot z_{\mathrm{sup}}^{t-1/2}
  :=\frac12(\dot z_{\mathrm{sup}}^t+\dot z_{\mathrm{sup}}^{t-1}).
\end{aligned}
  \label{eq:structural-feedback}
\end{equation}

이 항은 reduced block derivation에서 실제로 nonzero일 때만 사용한다.
Cross blocks가 0인 basis라면 식~\eqref{eq:structural-feedback}를 억지로 만들어내지 않는다.
식~\eqref{eq:canonical-structural-force-ledger}의 projective target는 같은 frame의
Global predictor와 Local solve에서 동일하게 held되므로 state-dependent cross reaction에 다시 넣지 않는다.
Module 5는 식~\eqref{eq:predictor-cross-load}의
\(Q_{g\leftarrow\ell}^{\mathrm{pred},t}\)를 이미 소비하므로,
same-frame corrector는 전체 endpoint cross load가 아니라

\begin{equation}
\begin{aligned}
  \Delta Q_{g\leftarrow\ell}^{\mathrm{struct},t}
  =&\;
  Q_{g\leftarrow\ell}^{\mathrm{mid},t}
  -Q_{g\leftarrow\ell}^{\mathrm{pred},t}\\
  =&\;
  -C_{g,\mathrm{sup}}
  \left[
    \frac12(\dot z_{\mathrm{sup}}^t+\dot z_{\mathrm{sup}}^{t-1})
    -\dot z_{\mathrm{sup}}^{\mathrm{pred},t-1/2}
  \right]\\
  &-K_{g,\mathrm{sup}}^{0}
  \left[
    \frac12(z_{\mathrm{sup}}^t+z_{\mathrm{sup}}^{t-1})
    -z_{\mathrm{sup}}^{\mathrm{pred},t-1/2}
  \right].
\end{aligned}
\end{equation}

만 소비한다. Actual midpoint가 predictor midpoint와 같으면 이 delta는 정확히 0이어야 한다.
이는 canonical fixed-Hessian backend에서 actual-midpoint coupling과
predictor-consumed coupling 사이의 exact updated-minus-consumed 식이다.

\subsection*{Corrector solve}

Module 5와 같은 implicit-midpoint operator를 사용한다.

\begin{equation}
  A_{gg}^{\mathrm{MP}}\delta q^t
  =
  2
  \left(
    \Delta Q_g^{\mathrm{aero},t}
    +\Delta Q_{g\leftarrow\ell}^{\mathrm{struct},t}
  \right),
  \qquad
  \delta\dot q^t=2\Delta t^{-1}\delta q^t.
  \label{eq:global-incremental-corrector}
\end{equation}

여기서
\(A_{gg}^{\mathrm{MP}}
=4\Delta t^{-2}M_{gg}+2\Delta t^{-1}C_{gg}+K_{gg}\)이다.
Endpoint에서 얻은 aerodynamic delta를 held midpoint-load delta로 한 번 적용하는
partitioned approximation이므로, 식~\eqref{eq:midpoint-effective-update}의 monolithic/iterated
reference와 오차를 별도로 보고한다.

\begin{equation}
  q^t=\widetilde q^t+\delta q^t,
  \qquad
  \dot q^t=\dot{\widetilde q}^t+\delta\dot q^t.
\end{equation}

한 번의 corrector로 충분한지 실험하고,
추가 fixed-point iteration은 quality--latency ablation으로 둔다.

\subsection*{출력과 다음 단계}

Corrected \(q^t,\dot q^t,z^t,\dot z^t\)가 Module 12로 넘어간다.
One-way baseline은 aerodynamic delta와 structural coupling을 모두 끈다.

\subsection*{필수 검증}

\begin{itemize}[leftmargin=1.5em]
  \item Feedback을 껐을 때 Global modal amplitude/phase가 실제로 달라지는가?
  \item Corrector가 base force를 두 번 적분하지 않는가?
  \item Same-frame corrector와 next-frame-only feedback의 비용/품질 비교.
  \item Corrected active/halo geometry/normal/area reconstruction,
  aerodynamic force reevaluation와 delta projection, structural-cross-delta construction,
  reduced Global corrector solve의 p50/p95/p99를 각각 보고하는가?
  Optional nonlinear sample-map backend를 켰을 때만 Jacobian reconstruction/update를 별도 보고한다.
  \item Feedback이 미미하면 closed-loop Global--Local/geometry-to-load feedback을 main claim에서 내릴 것.
\end{itemize}

\section{Module 12: 최종 anchor state와 deformation field}

\subsection*{푸는 질문}

\emph{Corrected Global/Local coordinate를 최종 anchor 위치, 속도,
deformation gradient로 어떻게 한 번만 복원하고, total hard reaction과 ROM residual을
어떤 discrete timing에서 진단하는가?}

\subsection*{Solver 분류}

\deterministic Basis lifting과 surface-aware affine reconstruction이다.

\subsection*{최종 상태}

\begin{equation}
  x^t=x^0+u_{\mathrm p}^t+\Phi q^t+\Psi_{\mathcal S_t}z^t,
  \qquad
  v^t=\dot u_{\mathrm p}^t+\Phi\dot q^t+\Psi_{\mathcal S_t}\dot z^t.
\end{equation}

이전/최종 state와 frame 동안 실제 소비한 held force ledger로
식~\eqref{eq:post-step-midpoint-kinematics}를 만든 뒤,
식~\eqref{eq:hard-total-reaction-residual}--\eqref{eq:hard-reaction-rom-residual-decomposition}의
\(f_{h\rightarrow o}^{\mathrm{tot},t}\), \(r_{\mathrm{ROM}}^t\)와
\(\Phi^\top r_{\mathrm{ROM}}^t,\Psi_{\mathcal S_t}^\top r_{\mathrm{ROM}}^t\)를 계산한다.
이 계산은 Local update와 Global corrector 전에 실행하지 않는다.

각 anchor/sample의 final tangential Jacobian \(J_k^t\)와
full-rank shell transport map \(F_{k,\mathrm{shell}}^t\)는
rest/current neighborhoods에서 한 번만 계산한다.
Global map과 Local map을 Gaussian에 순차 적용하여 double deformation을 만들지 않는다.

\subsection*{출력과 다음 단계}

\begin{equation}
  \mathcal X^t
  =
  \{x_k^t,v_k^t,J_k^t,F_{k,\mathrm{shell}}^t,n_k^t,A_k^t\}_{k=1}^K,
  \qquad
  \mathfrak r^t
  =
  (f_{h\rightarrow o}^{\mathrm{tot},t},r_{\mathrm{ROM}}^t,
   \Phi^\top r_{\mathrm{ROM}}^t,\Psi_{\mathcal S_t}^\top r_{\mathrm{ROM}}^t).
\end{equation}

이 결과는 Module 13의 Gaussian transport와
다음 frame Module 2의 current surface state에 사용되며,
\(\mathfrak r^t\)는 diagnostics/profiler에만 저장한다.

\subsection*{필수 검증}

\begin{itemize}[leftmargin=1.5em]
  \item Global-only, Local-only, combined reconstruction의 linearity/consistency.
  \item Pin/soft attachment 오차.
  \item \(J\) rank/orientation과 \(F_{\mathrm{shell}}\) determinant degeneracy.
  \item Global-span leakage와 patch seam diagnostics.
  \item Full-free/teacher fixture 대비 hard reaction error와
  \(\|r_{\mathrm{ROM}}\|\), Global/Local virtual-work residual.
\end{itemize}

\section{Module 13: Gaussian affine transport와 rendering}

\subsection*{푸는 질문}

\emph{Sparse anchor의 최종 deformation을 수많은 render Gaussians에
center-only artifact 없이 어떻게 전달하는가?}

\subsection*{Solver 분류}

\deterministic Affine MLS, covariance transport, co-rotation,
standard 3DGS rendering이다.
이 연산 자체는 기존 Gaussian physics 계열의 선행이 있으므로 contribution으로 주장하지 않는다
\cite{xie2024physgaussian,yan2026gaussianswaying,shao2025gausim}.

\subsection*{Gaussian mean}

Gaussian \(j\)에 대해 full-rank local shell map \((F_{j,\mathrm{shell}}^t,t_j^t)\)를
binding support \(\mathcal B_j\)에서 맞춘다.

\begin{equation}
  \mu_j^t
  =
  F_{j,\mathrm{shell}}^t\mu_j^0+t_j^t.
\end{equation}

\subsection*{Covariance}

\begin{equation}
  \Sigma_j^t
  =
  F_{j,\mathrm{shell}}^t
  \Sigma_j^0
  F_{j,\mathrm{shell}}^{t\top}.
  \label{eq:covariance-transport}
\end{equation}

Eigenvalue clamp 또는 polar decomposition으로
\(\Sigma_j^t\)가 symmetric positive definite를 유지하도록 한다.
Center-only transport는 anisotropic splat orientation과 scale을 보존하지 못하므로 baseline으로 둔다.

\subsection*{Appearance frame}

Canonical appearance SH coefficient를 유지하고
rigid/co-rotational frame change가 필요하면 analytic SH rotation 또는
local view direction transform을 사용한다.
새 appearance residual network는 본 논문의 main contribution에 포함하지 않는다.

\subsection*{출력}

\begin{equation}
  G^t
  =
  \{(\mu_j^t,\Sigma_j^t,\alpha_j,c_j^t)\}_{j=1}^{N_G},
  \qquad
  I_t=\mathcal R_{\mathrm{3DGS}}(G^t,\Pi_t).
\end{equation}

\subsection*{필수 검증}

\begin{itemize}[leftmargin=1.5em]
  \item Identity, translation, rigid rotation, affine reproduction.
  \item Invalid covariance와 scale explosion count.
  \item Rendered coverage hole/overlap count와 temporal popping.
  \item Center-only, rigid-only, full affine transport 비교.
  \item Novel-view PSNR/SSIM/LPIPS와 temporal flicker.
\end{itemize}

\section{Module 14: State, hysteresis, decay, diagnostics 갱신}

\subsection*{푸는 질문}

\emph{어떤 Local state를 다음 frame까지 유지하고,
언제 Local block과 조건부 expert를 끄며, 언제 완전히 dormant로 만들 것인가?}

\subsection*{Solver 분류}

\deterministic State machine과 energy-based cleanup이다.

\subsection*{세 가지 patch 상태}

\begin{longtable}{L{0.18\linewidth}L{0.30\linewidth}L{0.44\linewidth}}
\toprule
상태 & 실행 & 전환 규칙 \\
\midrule
\endfirsthead
\toprule
상태 & 실행 & 전환 규칙 \\
\midrule
\endhead
Dormant & Cheap selector만 평가, Local solve/expert allocation 없음 & Fixed-cardinality selector에 들면 Active \\
\addlinespace
Active & Local excitation/physical solve; 조건부 expert & Fixed-cardinality selector에서 빠지고 dwell 조건을 만족하면 Decay \\
\addlinespace
Decay & Expert off, excitation fade-out, restoring/damping solve 유지 & Group energy/cross load가 작으면 Dormant; selector가 다시 고르면 Active \\
\bottomrule
\end{longtable}

Actual-skip 보고는 세 비용을 분리한다. Dormant는 Local force/solve와 expert를 모두 skip하고,
Decay는 새 expert와 새 excitation을 skip하되 남은 state의 관성, restoring과 damping solve는 유지한다.
따라서 ``inactive/decay skip''을 Decay physical solve까지 끈다는 뜻으로 사용하지 않는다.

\subsection*{Energy criterion}

Inter-patch cross block 때문에 patch self-energy는 가산적이지 않다.
강하게 coupling된 overlap group \(\mathcal G\)에 대해

\begin{equation}
  E_{\mathcal G}^t
  =
  \frac12\dot z_{\mathcal G}^{t\top}
  M_{\mathcal G\mathcal G}\dot z_{\mathcal G}^t
  +E_{\mathrm{str},\mathcal G}^t
  (\widetilde q^t,z_{\mathcal G}^t)
\end{equation}

를 추적한다. Group 밖에서 patch \(r\in\mathcal G\)로 들어오는 dynamic cross load

\begin{equation}
  \ell_{r,\mathrm{cross}}^t
  =
  \left\|
  \sum_{s\notin\mathcal G}
  (M_{rs}\ddot z_s^t+C_{rs}\dot z_s^t+K_{rs}^{0}z_s^t)
  \right\|_2
\end{equation}

도 기록한다. \(E_{\mathcal G}^t<\tau_E\),
\(\ell_{r,\mathrm{cross}}^t<\tau_{\mathrm{cross}}\)이고 selector score도 낮을 때만
해당 group/patch를 dormant로 만든다.
Dormant 전환 시 이미 수치적으로 작은 fixed slot만 0으로 만들므로 popping을 제한한다.

\subsection*{저장 항목}

\begin{itemize}[leftmargin=1.5em]
  \item \(q^t,\dot q^t,z_{\mathrm{sup}}^t,\dot z_{\mathrm{sup}}^t\),
  \item active/decay/dormant state와 \(\gamma_r^t\),
  \item recurrent aerodynamic memory \(h_r^t\)가 있다면 그 state,
  \item force/torque/energy/conservation residual,
  \item active ratio, Local block solves, optional expert calls, per-module p50/p95/p99 latency.
\end{itemize}

\subsection*{필수 검증}

\begin{itemize}[leftmargin=1.5em]
  \item Threshold-crossing wind에서 chattering과 jerk.
  \item Wind-stop 후 Global/Local energy monotonic decay.
  \item Patch off 순간 position/velocity discontinuity.
  \item Memory leak와 active-state 누적.
\end{itemize}

\section{통합 runtime algorithm}

\begin{lstlisting}
Input:
  versioned target object package P(G0; metric, material, attachment)
  prescribed spatial wind W(x,t)
  optional motion of the fixed support layout

State:
  global q, qdot
  fixed-superset local z, zdot with active/decay block indices
  selector state, activation blend, optional expert memory

for frame t:
    # Known state and causal half-step load geometry
    load and validate previous state
    sample prescribed controls and wind at t-1/2
    extrapolate q,z positions to half-step; hold endpoint structural velocity
    fit predicted surface deformation maps, normals, and areas

    # Global physical predictor
    interpolate sampled wind and compute predicted relative wind
    evaluate half-step-geometry/frozen-velocity quasi-steady aerodynamic force
    fit one-shot relation rotations and assemble h_K,h_C plus hard-support lift
    project all anchor forces to Global load
    use cubature only if its profiling gate passed
    use exactly the frozen structural_backend_id and integrator_id
    reuse the fixed projective Hessian; relation targets change RHS only
    solve always-on Global reduced structural dynamics -> q_tilde

    # Adaptive Local decision and proposal
    reuse the owner-resolved anchor forces already evaluated above
    project cached force to each pre-factorized candidate group
    score by dynamic compliance (V0), oracle benefit (development),
        full counter-state residual (offline/low-rate analysis),
        or learned benefit (conditional)
    select a fixed TopK under hysteresis
    restrict cached full-anchor known external force to active/decay patches
    evaluate missing-force network only on active patches if F passed
    hold and fade the last missing-force proposal on decay patches

    # Assemble before integrating
    keep known base, missing aero, missing structural as separate channels
    apply activation blend once and condense patch copies to unique anchors
    route external Global work on the same support and take its explicit Local remainder
    solve structural/reaction or additional constraints by one augmented small dual Schur system
    project each assembled channel to the Local generalized RHS
    gather full cross blocks for one active/decay complementary Local system
    integrate Local mass-stiffness-damping system -> z

    # Mechanical feedback
    reconstruct corrected surface in active/halo regions
    recompute only corrected-minus-predictor aerodynamic force
    add structural_cross_delta = actual-midpoint coupling - predictor-consumed coupling
    solve one small Global corrector -> q

    # Output and state
    lift final q,z once to anchors
    recover total hard reaction from the final full residual
    report free-space ROM residual and its Global/Local virtual work
    transport Gaussian means and covariances by affine MLS
    render standard 3DGS
    update active/decay/dormant state and diagnostics
\end{lstlisting}

\subsection{One-frame data dependency}

\begin{equation}
\boxed{
\begin{aligned}
(s^{t-1},u_{\mathcal C}^{t-1/2},W^{t-1/2})
&\rightarrow \mathcal G_{\mathrm{surf}}^{\mathrm{pred},t-1/2}\\
&\rightarrow (f_{\mathrm{cache,ext}}^t,f_{\mathrm{cache,sel}}^t)
\rightarrow \widetilde q^t\\
&\rightarrow \text{cached-compliance Top-}K
\rightarrow \text{active Local force proposal}\\
&\rightarrow \text{support route/conditional constrained assembly}\\
&\rightarrow \{Q_{\ell,c}^t=\Psi_{\mathcal S_t}^\top f_{\ell,c}^t\}_c
\rightarrow z^t\\
&\rightarrow
(\Delta Q_g^{\mathrm{aero}},\Delta Q_g^{\mathrm{struct}})
\rightarrow q^t
\rightarrow (\mathcal X^t,\mathfrak r^t)
\rightarrow (G^t,s^t).
\end{aligned}
}
\end{equation}

여기서
\(\mathcal G_{\mathrm{surf}}^{\mathrm{pred},t-1/2}
:=(x^{\mathrm{pred},t-1/2},v^{\mathrm{pred},t-1/2},
n^{\mathrm{pred},t-1/2},A^{\mathrm{pred},t-1/2})\)이고,
superscript \(t\)의 force는 이 predictor sample을 frame 동안 hold한 load label이다.

\section{학습 objective}

\subsection{Topology와 physical package loss}

Candidate material edge classification은 class imbalance를 고려한다.

\begin{equation}
  \mathcal L_{\mathrm{edge}}
  =
  \operatorname{BCE}_{\mathrm{focal}}
  (\widehat y_{ij}^{\mathrm{mat}},y_{ij}^{*,\mathrm{mat}}).
\end{equation}

Rest length, tangent direction, oriented normal, area, mass는

\begin{equation}
\begin{aligned}
  \mathcal L_{\mathrm{geom}}
  =&\;
  \lambda_\ell\|\widehat\ell^0-\ell^{0,*}\|_1
  +\lambda_t(1-|\widehat t^\top t^*|)\\
  &+\lambda_n(1-\widehat n^\top n^*)
  +\lambda_A\|\widehat A-A^*\|_1
  +\lambda_m\|\widehat m-m^*\|_1.
\end{aligned}
\end{equation}

Total mass/area conservation은 별도로 검사한다.

\begin{equation}
  \mathcal L_{\mathrm{total}}
  =
  \left|
  \sum_k\widehat m_k-M^*
  \right|
  +
  \left|
  \sum_k\widehat A_k-A^*
  \right|.
\end{equation}

Attachment layout/compliance는 user setup metadata를 anchor scaffold에 transfer하고
그 consistency를 검사한다. Appearance-only attachment classifier를 mainline loss로 두지 않는다.
Mass는 inferred area와 setup density/thickness에서 유도하며 raw Gaussian density로 회귀하지 않는다.
Disconnected close layers를 잘못 연결하는 edge에 큰 penalty를 둔다.
Relation logit은 physical modulus를 회귀하지 않고
식~\eqref{eq:relation-area-partition}의 normalized owner share만 만든다.
Relation validity는 식~\eqref{eq:relation-validity-confidence}의 별도 calibrated head로

\begin{equation}
  \mathcal L_{\mathrm{rel\mbox{-}valid}}
  =
  \sum_{\alpha,k,c}
  \left[
    \operatorname{BCE}(q_{\alpha k}^{c,\mathrm{val}},y_{\alpha k}^{c,\mathrm{val}})
    +\lambda_{\mathrm{Brier}}
    (q_{\alpha k}^{c,\mathrm{val}}-y_{\alpha k}^{c,\mathrm{val}})^2
  \right]
  \label{eq:relation-validity-loss}
\end{equation}

를 사용해 식별하고, calibration error도 held-out package metric으로 보고한다.
Teacher deformation probe \(n\)에서 relation-local membrane/curvature target을
\(\epsilon_{\alpha,c}^{*,n}\), teacher owner area를 \(a_\alpha^{c,*}\)라 하면
teacher relation rotation \(R_\alpha^{*,n}\)과 co-rotated local displacement를

\begin{equation}
  \bar u_{\alpha i}^{*,n}
  :=
  (R_\alpha^{*,n})^\top
  (x_{\alpha i}^{*,n}-\bar x_\alpha^{*,n})
  -(x_{\alpha i}^0-\bar x_\alpha^0)
\end{equation}

로 만든다. Large rigid rotation을 physical strain으로 벌점 주지 않도록 scalar small-strain
supervision은 \(\widehat B_{\alpha,c}
[\bar u_{\alpha i}^{*,n}]_i\)를 teacher의 동일 material-frame target과 비교한다.
Canonical projective channel의 \(R_\alpha^{*,n}\)는 teacher가 임의로 제공한 frame이 아니라
식~\eqref{eq:projective-rotation-fit}과 같은 membrane descriptor fit으로 계산하고,
\(n_\alpha^{*,n}:=R_\alpha^{*,n}n_\alpha^0\)로 강제한다.
Teacher geometric normal과의 차이는 orientation diagnostic으로만 기록한다.
Canonical projective channel은

\begin{equation}
\begin{aligned}
  \mathsf P_{\alpha,n}^{*,n}
  &:={}I_{m_{\mathrm b}}\otimes
  (n_\alpha^{*,n}(n_\alpha^{*,n})^\top),\\
  \mathsf P_{\alpha,t}^{*,n}
  &:={}I_{3m_{\mathrm b}}-\mathsf P_{\alpha,n}^{*,n},\\
  y_{\alpha,\mathrm m}^{*,n}
  &:={}(I_{m_{\mathrm m}}\otimes R_\alpha^{*,n})d_{\alpha,\mathrm m}^0,\\
  y_{\alpha,\mathrm b}^{*,n}
  &:={}\mathsf P_{\alpha,t}^{*,n}\mathsf A_{\alpha,\mathrm b}x^{*,n}
  +(I_{m_{\mathrm b}}\otimes R_\alpha^{*,n})d_{\alpha,\mathrm b}^{0,n},\\
  r_{\alpha,c}^{*,n}
  &:={}\mathsf A_{\alpha,c}x^{*,n}-y_{\alpha,c}^{*,n},\\
  f_{\mathrm{proj}}^{*,n}
  &:={}-\sum_{\alpha,c}\mathsf A_{\alpha,c}^\top
  \Omega_{\alpha,c}r_{\alpha,c}^{*,n}
\end{aligned}
\end{equation}

의 assembled internal-force action을 conservative transfer한 teacher anchor internal force
\(f_{\mathrm{int},A}^{*,n}\)와 비교한다. 따라서 operator supervision은

\begin{equation}
\begin{aligned}
  \mathcal L_{\mathrm{op\mbox{-}action}}
  =&\;
  \sum_{n,\alpha,c}
  \left\|
  \widehat B_{\alpha,c}[\bar u_{\alpha i}^{*,n}]_i
  -\epsilon_{\alpha,c}^{*,n}
  \right\|_2^2\\
  &+\lambda_f\sum_n
  \left\|
  P_{\mathrm{free}}
  (f_{\mathrm{proj}}^{*,n}-f_{\mathrm{int},A}^{*,n})
  \right\|_{M^{-1}}^2
  +\lambda_a\sum_{\alpha,c}
  \left|a_\alpha^c-a_\alpha^{c,*}\right|
  \label{eq:operator-geometry-loss}
\end{aligned}
\end{equation}

로 둔다. \(\mathsf G_\alpha\mathsf V_\alpha=I_6\)와
\(\sum_\alpha a_\alpha^c=\sum_kA_k^0\)는 full-rank/normalization construction에서
자동 성립하므로 학습 loss로 가장하지 않고 numerical hard check로만 기록한다.
Rank/conditioning failure, nonpositive weight/length, denominator 0과 metric/material type mismatch는 loss로 완화하지 않고
package build를 중단하는 hard error다.

\subsection{Reduced operator와 modal loss}

Predicted physical package에서 얻은 eigenfrequency와 mode subspace를 teacher와 비교한다.
Mode의 sign과 repeated eigenvalue ambiguity 때문에 개별 vector L2보다 subspace projector를 사용한다.

\begin{equation}
  \mathcal L_{\mathrm{subspace}}
  =
  \left\|
  \Phi\Phi^\top M
  -
  \Phi^*\Phi^{*\top}M^*
  \right\|_F^2.
\end{equation}

\begin{equation}
  \mathcal L_{\mathrm{freq}}
  =
  \sum_{j=1}^{r_g}
  \left|
  \log(\widehat\omega_j+\eps)
  -
  \log(\omega_j^*+\eps)
  \right|.
\end{equation}

Operator construction은 학습 loss와 별도로
식~\eqref{eq:operator-construction-hard-checks}와 다음 reduced reciprocity를 통과해야 한다.

\begin{equation}
\begin{gathered}
  M\succ0,\qquad K^0\succeq0,\qquad C^0\succeq0,\\
  K_{gr}^0=(K_{rg}^0)^\top,\qquad C_{gr}=C_{rg}^\top,\\
  A_{rs}=A_{sr}^\top,\qquad A\in\{M,C,K^0\}.
\end{gathered}
\end{equation}

\subsection{조건부 aerodynamic cubature loss}

식~\eqref{eq:cubature-fit}의 generalized force와 wrench loss에
held-out deformation regularization을 추가한다.

\begin{equation}
  \mathcal L_{\mathrm{cub}}
  =
  \lambda_Q\mathcal L_Q
  +\lambda_F\mathcal L_F
  +\lambda_\tau\mathcal L_\tau
  +\lambda_{\mathrm{sparse}}|\mathcal S_c|.
\end{equation}

단순히 training force snapshot을 맞추고
current normal/area 또는 traveling gust에 실패하면 hyper-reduction claim을 지지하지 못한다.

Full-anchor profiling gate가 실패하면 이 loss와 model은 구성하지 않는다.

\subsection{조건부 learned gate loss}

Continuous benefit regression은

\begin{equation}
  \mathcal L_{\mathrm{benefit}}
  =
  \operatorname{Huber}
  (\widehat b_{\mathcal G}^t,b_{\mathcal G}^{*,t}).
\end{equation}

같은 frame의 atomic gating-unit 순위를 맞추기 위한 ranking loss는

\begin{equation}
  \mathcal L_{\mathrm{rank}}
  =
  \sum_{(\mathcal G,\mathcal H):b_{\mathcal G}^*>b_{\mathcal H}^*}
  \max
  \left(
    0,\;
    m-\widehat b_{\mathcal G}+\widehat b_{\mathcal H}
  \right).
\end{equation}

Uncertainty calibration을 더한다.

\begin{equation}
  \mathcal L_{\mathrm{gate}}
  =
  \lambda_b\mathcal L_{\mathrm{benefit}}
  +\lambda_r\mathcal L_{\mathrm{rank}}
  +\lambda_{\mathrm{cal}}\mathcal L_{\mathrm{calibration}}.
\end{equation}

V0는 predicted cost나 compute regularizer를 사용하지 않고 모든 selector를 같은 \(K_A\)에서 비교한다.
Learned gate는 Oracle Top-\(K\)가 cached-compliance Top-\(K\)보다 충분히 나은 경우에만 학습한다.

\subsection{조건부 missing-force loss}

Force-space supervision은

\begin{equation}
  \mathcal L_{\Delta f}
  =
  \sum_{t,r,c}
  w_r^t
  \operatorname{Huber}
  \left(
    \widehat{\Delta f}_{r,c}^t,
    \Delta f_{r,c}^{*,t}
  \right).
\end{equation}

식~\eqref{eq:patch-force-label}의 \(\Delta f_{r,c}^{*,t}\)는 raw-head target이다.
Hard structural layout에서는 projected output \(E_{f,t}y_{\mathrm{struct}}^*\)가 아니라
unprojected free-support remainder를 가르치고, runtime과 같은 augmented projection 뒤의
\((y^*,\rho_h^*)\)는 아래 assembled constraint/reaction loss가 감독한다.
Hard-row raw teacher force를 Local head target에 섞지 않는다.

Global leakage penalty는

\begin{equation}
  \mathcal L_{\mathrm{leak}}
  =
  \sum_{t,c}
  \left\|
  \Phi^\top \widehat{\Delta f}_{\ell,c}^{\mathrm{raw},t}
  \right\|_2^2.
\end{equation}

Channel별 leakage를 따로 벌점화하므로 aero와 structural leakage가 서로 상쇄될 수 없다.
여기서 \(\widehat{\Delta f}_{\ell,c}^{\mathrm{raw}}\)는 channel별 owner-weighted
patch proposal을 route 전 anchor space에 조립한 field다. Post-assembly leakage는 exact solver
constraint residual이라 학습 신호로 재사용하지 않고 아래 KKT diagnostic으로 따로 보고한다.

Internal structural channel의 net force/torque loss는

\begin{equation}
\begin{aligned}
  f_{\ell,\theta,\mathrm{struct}}^{\mathrm{full},t}
  &:={}E_{f,t}y_{\ell,\theta,\mathrm{struct}}^t
  +r_{o\rightarrow h,\theta,\mathrm{struct}}^t,\\
  \mathcal L_{\mathrm{int\mbox{-}wrench}}
  &:={}\left\|
  \mathcal W_t f_{\ell,\theta,\mathrm{struct}}^{\mathrm{full},t}
  \right\|_2^2 .
\end{aligned}
\end{equation}

Free object에서는 \(r_{o\rightarrow h}=0\)이다.
Hard structural channel은 free force만 zero-wrench로 만들지 않고 augmented companion과 함께 검사한다.
External aerodynamic channel은 0으로 만들지 않고 teacher external wrench를 맞춘다.

\(\widehat{\Delta f}_{\ell,c}^{\mathrm{raw},t}\)는 patch prediction
\(\widehat{\Delta f}_{r,c}^t\)를 식~\eqref{eq:raw-overlap-assembly}처럼
owner weight로 조립한 anchor-space Local prediction이다.
External/free shortcut은 식~\eqref{eq:teacher-local-force-label}, hard structural channel은
식~\eqref{eq:teacher-hard-structural-augmented-label}과 같은 route/constraint backend를 적용한다.
그 post-assembly field를 \(f_{\ell,\theta,c}^{\mathrm{asm},t}\)로 쓴다.
다음 \(A_{\mathrm{lift},c}\) shorthand는 external/free shortcut에만 사용한다.
\(\mathcal C_{\Delta,\mathrm{ext}}:=\{\mathrm{miss\mbox{-}aero}\}\)라 둔다.

\begin{equation}
  A_{\mathrm{lift},c}
  =
  \begin{bmatrix}\Phi^\top\\A_{\mathrm{cons},c}^t\end{bmatrix},
  \qquad
  b_{\mathrm{lift},c}^t
  =
  \begin{bmatrix}0\\b_{\ell,c}^{*,t}\end{bmatrix}.
\end{equation}

Nodal prediction이 intended generalized force와 joint constraints에 일치하도록

\begin{equation}
\begin{aligned}
  \mathcal L_{Q\mathrm{-cons}}
  ={}&
  \sum_{c\in\mathcal C_{\Delta,\mathrm{ext}}}\left\|
  \Psi_{\mathrm{sup}}^\top f_{\ell,\theta,c}^{\mathrm{asm},t}
  -\Delta Q_{\ell,c}^{*,t}
  \right\|_1
  +\left\|
  \Psi_{\mathrm{sup}}^\top E_{f,t}y_{\ell,\theta,\mathrm{struct}}^t
  -\Delta Q_{\ell,\mathrm{struct}}^{*,t}
  \right\|_1\\
  &+\lambda_{\mathrm{KKT}}
  \sum_{c\in\mathcal C_{\Delta,\mathrm{ext}}}\left\|
  A_{\mathrm{lift},c}
  \widehat{\Delta f}_{\ell,c}^{\mathrm{raw},t}-b_{\mathrm{lift},c}^t
  \right\|_2^2\\
  &+\lambda_{\mathrm{KKT}}
  \left\|
  B_{\mathrm{struct}}^t
  \begin{bmatrix}
    y_{\ell,\theta,\mathrm{struct}}^t\\
    \rho_{h,\theta,\mathrm{struct}}^t
  \end{bmatrix}
  \right\|_2^2
  +\lambda_\rho
  \left\|
  \rho_{h,\theta,\mathrm{struct}}^t
  -\rho_{h,\mathrm{struct}}^{*,t}
  \right\|_{W_h}^2 .
\end{aligned}
\end{equation}

를 추가한다. 식~\eqref{eq:generalized-force-lifting}을 사용한 sample은
KKT condition, feasibility와 label confidence를 함께 기록한다.

\subsection{One-step과 rollout state loss}

\begin{equation}
  \mathcal L_{\mathrm{state}}
  =
  \lambda_q\|q-q^*\|_1
  +\lambda_{\dot q}\|\dot q-\dot q^*\|_1
  +\lambda_z\|z-z^*\|_1
  +\lambda_{\dot z}\|\dot z-\dot z^*\|_1.
\end{equation}

Multi-step rollout은

\begin{equation}
  \mathcal L_{\mathrm{roll}}
  =
  \sum_{h=1}^{H}
  \beta^{h-1}
  \left[
  \lambda_x\|x^{t+h}-x^{*,t+h}\|_1
  +\lambda_v\|v^{t+h}-v^{*,t+h}\|_1
  +\lambda_n\mathcal L_n^{t+h}
  \right].
\end{equation}

Flutter의 amplitude만 맞고 frequency/phase가 틀리는 것을 막기 위해
window spectrum loss를 사용한다.

\begin{equation}
  \mathcal L_{\mathrm{spec}}
  =
  \left\|
  \log(\operatorname{PSD}(\widehat x)+\eps)
  -
  \log(\operatorname{PSD}(x^*)+\eps)
  \right\|_1.
\end{equation}

\subsection{Overlap, feedback, rendering loss}

\begin{equation}
  \mathcal L_{\mathrm{overlap}}
  =
  \lambda_c\|\Delta x_{\mathrm{boundary}}\|_2^2
  +\lambda_v\|\Delta v_{\mathrm{boundary}}\|_2^2
  +\lambda_w\|A_{\mathrm{cons}}f-b_{\mathrm{cons}}\|_2^2.
\end{equation}

Corrected-aero feedback은 teacher Global generalized force delta와 비교한다.

\begin{equation}
  \mathcal L_{\mathrm{feedback}}
  =
  \left\|
  \Delta Q_g^{\mathrm{aero}}
  -
  \Delta Q_g^{\mathrm{aero},*}
  \right\|_1.
\end{equation}

Gaussian transport는 mean/covariance와 rendered frame에서 감독할 수 있다.

\begin{equation}
  \mathcal L_{\Sigma}
  =
  \left\|
  \log\Sigma_j^t-\log\Sigma_j^{*,t}
  \right\|_F^2,
\end{equation}

\begin{equation}
  \mathcal L_{\mathrm{render}}
  =
  \lambda_1\|\widehat I-I^*\|_1
  +\lambda_{\mathrm{SSIM}}(1-\operatorname{SSIM}(\widehat I,I^*))
  +\lambda_p\operatorname{LPIPS}(\widehat I,I^*).
\end{equation}

\subsection{Total objective와 staged activation}

\begin{equation}
\begin{aligned}
  \mathcal L
  =&\;
  \mathcal L_{\mathrm{topology/geometry}}
  +\mathcal L_{\mathrm{rel\mbox{-}valid}}
  +\mathcal L_{\mathrm{operator/modal}}
  +\mathbf 1_{H}\mathcal L_{\mathrm{cub}}\\
  &+\mathbf 1_{G}\mathcal L_{\mathrm{gate}}
  +\mathbf 1_{F}
  (\mathcal L_{\Delta f}+\mathcal L_{Q\mathrm{-cons}}
  +\mathcal L_{\mathrm{leak}}+\mathcal L_{\mathrm{int\mbox{-}wrench}})\\
  &+\mathcal L_{\mathrm{state}}
  +\mathcal L_{\mathrm{roll}}
  +\mathcal L_{\mathrm{spec}}
  +\mathcal L_{\mathrm{overlap}}
  +\mathcal L_{\mathrm{feedback}}
  +\mathcal L_{\mathrm{render}}.
\end{aligned}
\end{equation}

여기서 \(\mathbf 1_H,\mathbf 1_F,\mathbf 1_G\)는 각각
hyper-reduction, missing-force, learned-gate의 독립 go/no-go가 통과했을 때만 1이다.
모든 항을 처음부터 동시에 켜지 않는다.
Stage별로 학습하고 각 stage의 oracle upper bound를 먼저 확인한다.

\section{학습 curriculum}

\subsection{Stage 0: Deterministic physics/transport unit tests}

Network training 전에 다음을 통과한다.

\begin{itemize}[leftmargin=1.5em]
  \item 식~\eqref{eq:v0-structural-freeze-signature}의
  objective-normal-curvature-projective-frozen-target/implicit-midpoint backend,
  \(\Delta t\)와 operator/target version 동결,
  \item E0와 MVP Step 0의 pure-rigid/strained-state superposed-rotation,
  planar bending-pollution, relative-rate damping,
  rest-linear/BE 및 one-shot/two-step/converged local--global projective reference 검증,
  \item teacher coordinate와 independent GS coordinate alignment,
  \item conservative area/mass partition,
  \item synthetic confidence를 사용한 absolute-evidence aggregator와 low-evidence fallback/reject
  deterministic fixture; learned head calibration은 Track B0/E1에서 별도 검증,
  \item attachment force/projection ownership,
  \item support-restricted route/Schur feasibility와 free/hard/compliant reaction ledger,
  \item \(M,K,C\) positivity와 symmetry,
  \item patch별 Global/Local \(M\)-orthogonality와 inter-patch block conditioning,
  \item tangent-cross area-normal update와 full-shell extension,
  \item affine MLS identity/rigid/affine reproduction,
  \item covariance SPD.
  \item 식~\eqref{eq:fixed-global-sample-jacobian}의 fixed sample lifting.
\end{itemize}

\subsection{Track A0: Oracle topology Global physics와 same-anchor kill test}

Teacher source mesh의 topology/relation label과 attachment mapping을 privileged Oracle로 사용하되,
main fairness path의 anchor mass와 structure는 같은 user \(\nu=0\) material metadata 및
식~\eqref{eq:v0-structural-freeze-signature}의 canonical assembly/backend로 만든다.
Teacher shell의 native stiffness/operator를 그대로 투영한 결과는 별도의 privileged physics ceiling으로만 보고,
canonical Track A0 speed/quality 결과와 섞지 않는다. 이 상태에서 Global reduced physics만 구현한다.
이 단계에서 current-state analytic aero와 full-sample projection을 먼저 사용한다.
Global-only가 teacher의 큰 motion과 wind-stop recovery를 학습 없이 재현해야 한다.
이어 동일 rest geometry/anchor, mass, user material/attachment, wind, timestep,
precision과 hardware에서
sparse corotational/PD/XPBD solver와 dynamics-only p50/p95/p99를 비교한다.
각 baseline의 native constitutive energy/operator privilege와 실제 DOF를 별도 보고하며,
동일 full-space \(K\)를 모든 solver에 강제했다고 쓰지 않는다.
512/1,024/2,048 anchors, batch 1과 multi-object에서
개발용 기준인 same-quality dynamics speedup 약 \(2\times\) 또는
same-p95 velocity/trajectory/spectrum 개선 또는 multi-object throughput 이득이 없으면
SYS quality--latency headline을 철회하고 MC2의 selective-compute 부분을 재평가한다.

\subsection{Track B0: 조기 representation feasibility preflight}

Step 0 동안에는 fixture에 필요한 최소 independent-GS coordinate/relation schema probe만 병렬 허용한다.
이는 coordinate alignment, candidate valence/rank, area/mass와 close-layer false edge를 조기에
반증하기 위한 것이며 learned calibration, full Track B0 acceptance 또는 대규모 data generation이 아니다.
Step-0 gate가 통과한 뒤 Track A0와 병렬로 5--10개 simple thin-surface object와
각 object의 독립 GS seed/view/density variant에서
anchor coverage, oriented graph/local relation, total area/mass,
첫 8--16개 eigenfrequency, modal subspace principal angle과 close-layer false edge를 측정한다.
Material, metric scale과 attachment는 setup metadata로 제공한다.
이 preflight가 실패하면 전체 adaptive solver를 만들기 전에
``without explicit target triangle-mesh reconstruction'' 주장을 철회하거나 proxy baseline으로 축소한다.

\subsection{Track A1: Oracle topology Local physics Gate B}

Oracle Local basis와 oracle active set을 사용하여
physics-only Local과 full-resolution Local upper bound를 비교한다.
공식 Gate B rollout 전에 Module 10--9--8의 unique-anchor joint assembly,
Global complement와 single coupled Local update를 먼저 고정한다.
One-corrector는 Local necessity 판정 뒤 별도 SYS ablation으로 검증한다.
Global rank를 늘리는 것만으로 local flutter가 해결되는지도 이 단계에서 반증한다.
Patch-indexed fixed-superset basis, full cross-block solve와 Oracle Top-\(K\)로
same-p95 Global rank sweep 대비 Gate B를 먼저 판정한다.

\subsection{Track A2: Cached-compliance selector와 decay Gate C}

Gate B를 통과한 Oracle Local system을 고정한 뒤
cached-force dynamic-compliance Top-\(K\), Oracle Top-\(K\)와 heuristic을 비교한다.
Active/Decay/Dormant state, actual active/solve-set skip과 gust-stratified decay profiler를 연결해
Gate C의 actual-skip overhead를 판정한다.
Full counter-state residual은 offline 또는 저주기 analysis baseline으로 추가 비용까지 함께 비교한다.

\subsection{조건부 H: Aerodynamic hyper-reduction}

Track A0 profile에서 full-anchor aero/fixed projection이 dynamics의 유의미한 병목일 때만
cubature를 구현한다. Generalized load fidelity를 primary target으로 두며,
localized/traveling gust Pareto가 개선되지 않으면 no-go로 종료한다.

\subsection{조건부 F: Always-on missing-force network}

Full-resolution Local과 reduced physics-only 사이에 반복 가능한 defect가 있고
teacher nodal breakdown 또는 식~\eqref{eq:generalized-force-lifting}의 검증된 lift가 있을 때만 진입한다.
Selector 없이 모든 oracle-relevant patch의 missing-force network를 실행한다.
Network capacity와 label correctness를 먼저 검증한다.
Always-on model이 physics-only Local을 개선하지 못하면 network만 제거하고
physics-only selector 실험은 계속한다.

\subsection{조건부 G: Learned gate와 adaptive execution}

Oracle Top-\(K\)가 cached-compliance Top-\(K\)보다 충분히 나을 때만
raw future-benefit label로 learned gate를 학습한다.
모든 selector는 같은 \(K_A\)에서 비교한다.
Hysteresis와 off-policy noisy rollout state를 학습/검증한다.
Learned selector가 cached-compliance/curvature/strain을 이기지 못하면 no-go로 종료한다.
Per-patch cost head는 profiler가 추가 필요성을 보인 경우의 후속 연구다.

\subsection{Track B1: Predicted topology와 operator}

Representation preflight와 Oracle physics 요구가 둘 다 통과한 뒤
Oracle topology를 learned topology/distilled operator로 교체한다.
Topology error와 dynamics error를 동시에 처음부터 디버깅하지 않는다.
Oracle 대비 degradation을 분리해 보고한다.

\subsection{Track B2: Independent GS reconstruction robustness}

다른 view 수, random seed, density, holes, floaters를 가진 independent GS variants로
object package 추론과 dynamics를 평가한다.
식~\eqref{eq:independent-gs-package-consistency}로 area/mass, eigenspace,
tail distribution과 common-probe rollout의 pairwise median/IQR/95th percentile을 보고한다.
Source mesh에 Gaussian을 직접 bind한 clean variant만으로 결과를 주장하지 않는다.

\subsection{통합: End-to-end Gaussian rendering}

Dynamics가 안정된 뒤 affine Gaussian transport와 rendering loss를 활성화한다.
Appearance correction은 별도 supporting experiment로 제한한다.

\section{계산 복잡도와 runtime contract}

\subsection{One-time target setup}

\begin{equation}
\begin{aligned}
  T_{\mathrm{setup}}
  =&\;
  T_{\mathrm{GS\mbox{-}preprocess}}
  +T_{\mathrm{anchor}}
  +T_{\mathrm{topology}}\\
  &+T_{\mathrm{operator}}
  +T_{\mathrm{basis}}
  +T_{\mathrm{group\mbox{-}compliance\mbox{-}factor}}\\
  &+T_{\mathrm{superset\mbox{-}reference\mbox{-}factor}}
  +T_{\mathrm{joint\mbox{-}constraint\mbox{-}fixed\mbox{-}precompute}}
  +\mathbf 1_H T_{\mathrm{cubature}}
  +T_{\mathrm{binding}}.
\end{aligned}
\end{equation}

Mainline setup은 explicit target triangle-mesh reconstruction을 호출하지 않아야 한다.
Target에서 SuGaR/GOF mesh를 만드는 variant는 end-to-end baseline이지 mainline이 아니다
\cite{guedon2024sugar,yu2024gof}.

\subsection{Per-frame dynamics}

Global mode 수 \(r_g\), active/decay Local coordinate 수 \(r_\ell^{\mathcal S}\),
patch 수 \(P\), atomic group 수 \(P_G\), active patch 수 \(P_A\),
solve patch 수 \(P_S=|\mathcal A_t\cup\mathcal D_t|\), full anchor 수 \(K\),
active/decay unique-anchor 수 \(n_{\mathcal S}^t=|\mathcal I_t|\),
hard row를 제외한 injection dimension
\(n_t=\dim\operatorname{range}(E_{f,t})\le3n_{\mathcal S}^t\),
corrected active/halo aerodynamic sample 수 \(K_{\mathrm{corr}}^t\),
hard constraint row 수 \(m_h=\operatorname{rows}(A_{\mathrm h})\),
조건부 cubature sample 수 \(S_c\)라 하면

\begin{equation}
\begin{aligned}
  T_{\mathrm{dyn}}
  \approx&
  T_{\mathrm{surface/aero}}(K\;\text{or}\;S_c)
  +T_{\mathrm{projection}}(K\;\text{or}\;S_c)\\
  &+T_{\mathrm{struct\mbox{-}rotation/target}}(|\mathcal Q_{\mathrm{op}}|)\\
  &+T_{\mathrm{global}}(r_g)
  +T_{\mathrm{selector}}^{\mathrm{V0}}(P_G)\\
  &\quad
  +\mathbf 1_G T_{\mathrm{gate\mbox{-}net}}(P_G)
  +\mathbf 1_F T_{\mathrm{missing\mbox{-}force}}(P_A)\\
  &\quad
  +T_{\mathrm{restriction/scatter}}(P_S,n_t)
  +T_{\mathrm{support\mbox{-}route}}(n_t,r_g,|\mathcal C_{\mathrm{ext}}|)
  +T_{\mathrm{joint\mbox{-}Schur}}(\{p_c\},\{m_c\})\\
  &\quad
  +T_{\mathrm{local}}(r_\ell^{\mathcal S})
  +T_{\mathrm{corr\mbox{-}geom}}(K_{\mathrm{corr}}^t)\\
  &\quad
  +T_{\mathrm{corr\mbox{-}aero}}(K_{\mathrm{corr}}^t,r_g)
  +T_{\mathrm{corr\mbox{-}struct}}(r_\ell^{\mathcal S},r_g)\\
  &\quad
  +T_{\mathrm{corr\mbox{-}solve}}(r_g)
  +T_{\mathrm{post\mbox{-}reaction/ROM\mbox{-}residual}}(K,m_h).
\end{aligned}
\end{equation}

Relation rotation fit과 \(h_K^t,h_C^t\) target assembly를
\(T_{\mathrm{global}}\) 또는 \(T_{\mathrm{local}}\)에 숨겨 중복/누락 계상하지 않는다.
Fixed Hessian의 reduced block projection과 factorization은 setup/cache 비용으로 별도 기록한다.

여기서 \(\mathcal C_{\mathrm{Schur}}\)는 structural/reaction 또는 additional independent row 때문에
실제로 augmented solve가 필요한 channel 집합이다. External default
\(\mathcal C_{\mathrm{ext}}\)는 \(T_{\mathrm{support\mbox{-}route}}\)에만 계상한다.
Dependent row를 제거한 뒤 \(m_c:=\operatorname{rank}(B_c^t)\)이고,
augmented primal dimension은
\(p_c:=n_t+\mathbf1_{\mathrm{hard},c}m_h\)이며,
generic full saddle solve를 숨기지 않고

\begin{equation}
\begin{aligned}
T_{\mathrm{joint\mbox{-}Schur}}
=\sum_{c\in\mathcal C_{\mathrm{Schur}}}
\bigl(&T_{B\mathrm{\mbox{-}update}}(p_c,m_c)
+T_{\mathrm{Schur\mbox{-}build}}(p_c,m_c)\\
&+T_{\mathrm{Schur\mbox{-}factor}}(m_c)
+T_{\mathrm{Schur\mbox{-}apply}}(p_c,m_c)\bigr)
\end{aligned}
\end{equation}

로 channel별 비용을 분리한다. Fixed block precompute와 factor reuse가 가능한 frame도
torque geometry update가 필요한 frame과 따로 집계한다.
Hard reaction DOF를 6D wrench image로 사전 응축하는 최적화는 그 image rank/conditioning,
recovery error와 precompute/update 비용을 별도 보고할 때만 \(p_c\)를 줄여 계상한다.

Setup에서 factorize한 작은 group compliance operator를 쓰면

\begin{equation}
  T_{\mathrm{selector}}^{\mathrm{V0}}(P_G)
  =T_{\mathrm{cached\mbox{-}projection}}(P_G)
  +\sum_{\mathcal G\in\mathfrak G}O(r_{\mathcal G}^2),
\end{equation}

이며 두 번째 surface/force evaluation은 없다. 반면 full counter-state analysis backend는

\begin{equation}
\begin{aligned}
  T_{\mathrm{selector}}^{\mathrm{counter}}
  =&\;T_{\mathrm{global}}(r_g)+T_{\mathrm{lift}}(K)
  +T_{\mathrm{surface/aero}}(K)\\
  &+T_{\mathrm{internal}}(K)+T_{\mathrm{group\mbox{-}residual}}(P_G)
\end{aligned}
\end{equation}

를 추가한다. 이는 every-frame V0 식에 숨기지 않으며, \(J\)-frame cadence에서는
amortized \(T_{\mathrm{selector}}^{\mathrm{counter}}/J\)와 activation delay를 함께 보고한다.

Canonical Global matrix는 fixed이므로 setup factorization 뒤 triangular solve가 된다.
Projective target assembly는 위의 structural-target 비용에 포함한다.
Local direct backend의 비용은 factor cache \(\mathfrak C_{\mathrm{fact}}\)에 대해

\begin{equation}
\begin{aligned}
  T_{\mathrm{local}}^{\mathrm{direct}}(t)
  =&\;T_{\mathrm{gather}}(\mathcal S_t)
  +\mathbf 1[k_t^{\mathrm{fact}}\notin\mathfrak C_{\mathrm{fact}}]
   T_{\mathrm{factor}}(r_\ell^{\mathcal S})\\
  &+T_{\mathrm{tri}}(r_\ell^{\mathcal S})
  +T_{\mathrm{scatter}}(\mathcal S_t)
\end{aligned}
\end{equation}

로 쓰며
\(k_t^{\mathrm{fact}}
=(v_{\mathrm{package/operator}},\mathrm{backend},\mathrm{integrator},\Delta t,
\operatorname{sig}(\mathcal S_t))\)다.
Projective target은 RHS만 바꾸므로 이 signature가 같으면 exact cross-frame hit가 가능하다.
즉 factorization miss를 포함하고 Local을 항상 cached triangular solve라고 간주하지 않는다.
Block-preconditioned PCG backend는

\begin{equation}
  T_{\mathrm{local}}^{\mathrm{PCG}}(t)
  =T_{\mathrm{gather}}(\mathcal S_t)
  +n_{\mathrm{it}}(t)
  \left(T_{\mathrm{cross\mbox{-}matvec}}+T_{\mathrm{block\mbox{-}prec}}\right)
  +T_{\mathrm{scatter}}(\mathcal S_t)
\end{equation}

로 분리하고, 모든 inter-group cross block을 matvec에 유지한다.
Iteration cap, 최종 linear residual과 exact-direct rollout 차이를 함께 기록한다.
Material/attachment layout 변화는 first-paper per-frame runtime control이 아니라
새 setup/package event이므로 위 per-frame 식에 포함하지 않는다.

\subsection{Gaussian transport와 resolution-scaling claim}

\begin{equation}
  T_{\mathrm{GS\mbox{-}transport}}
  =
  O(N_GK_{\mathrm{MLS}}).
\end{equation}

Dynamics resolution \(K,P,r_g,r_\ell^{\mathcal S}\)는 render Gaussian 수 \(N_G\)와 분리되어야 한다.
고정된 physical budget에서 분리되는 것은 structural dynamics solve 비용이다. 전체 frame 비용은

\begin{equation}
  T_{\mathrm{frame}}(N_G)
  =T_{\mathrm{dyn}}(K,P,r_g,r_\ell^{\mathcal S})
  +T_{\mathrm{GS\mbox{-}transport}}(N_G)
  +T_{\mathrm{render}}(N_G).
\end{equation}

따라서 end-to-end 비용이 \(N_G\)와 독립이라고 주장하지 않는다.
Gaussian 수 증가는 transport와 rendering에 반영되며,
동일 scaffold와 physical budget에서 \(T_{\mathrm{dyn}}\)이 거의 일정한지를 system evidence로 측정한다.
이 scaling은 독립 algorithmic novelty가 아니다.

\subsection{Measured break-even}

Adaptive method의 이론적 sparsity가 실제 GPU speedup을 보장하지 않는다.
작은 patch가 많으면 kernel launch와 gather/scatter가 병목이 될 수 있다.
따라서 target GPU에서

\begin{equation}
  T_{\mathrm{always}},\quad
  T_{\mathrm{adaptive}}(\rho_A),\quad
  T_{\mathrm{physics\mbox{-}only}},\quad
  T_{\mathrm{direct}}
\end{equation}

를 실제 측정한다.
식~\eqref{eq:active-solve-size-diagnostics}의 frame superscript를 생략하고,
Patch \(r\)의 fixed coordinate rank를 \(r_r\),
\(r_A=\sum_{r\in\mathcal A_t}r_r\),
\(r_D=\sum_{r\in\mathcal D_t}r_r\),
\(r_{\mathcal S}=r_A+r_D\),
\(r_{\mathrm{sup}}=\sum_{r=1}^P r_r\)라 두고
\(\rho_A=r_A/r_{\mathrm{sup}}\),
\(\rho_{\mathcal S}=r_{\mathcal S}/r_{\mathrm{sup}}\)로 정의한다.
Fixed \(K_A\)라도 \(\mathcal S_t=\mathcal A_t\cup\mathcal D_t\)이므로
traveling gust와 긴 decay에서 실제 solve dimension이 증가할 수 있다.
이 효과는 stationary, slow-traveling, fast-traveling과 reversing gust로 층화해 측정한다.
따라서 \(|\mathcal A_t|,|\mathcal D_t|,r_A,n_A,r_D,n_D,r_{\mathcal S},n_{\mathcal S}\),
\(|\mathcal S_t\triangle\mathcal S_{t-1}|\)와 toggle coordinate/unique-anchor count,
decay residence distribution, \(\max_t|\mathcal D_t|\)와 \(\max_t r_{\mathcal S}\),
unique signature/eviction/cold--warm cache hit, factor/update/downdate-fallback 횟수,
cache key별 factor memory와 lookup/transfer 비용,
exact active matrix condition estimate와 LDL pivot/fallback,
retained cross-block density/bytes와 PCG matvec bandwidth/iteration/residual 및
gather/factor/solve/scatter p50/p95/p99를 기록한다.
Joint assembly는 channel별 constraint row 수와 rank/condition,
dynamic torque-row update, Schur build/factor/apply, fixed-factor reuse와 fallback 횟수 및
assembly 대비 전체 dynamics p50/p95/p99 비율을 별도로 기록한다.
Corrector는 \(K_{\mathrm{corr}}^t\), corrected geometry update,
aerodynamic reevaluation/projection, structural delta와 reduced solve의 p50/p95/p99를 분리한다.
Optional nonlinear sample-map backend만 Jacobian update를 별도 항목으로 추가한다.
Incremental backend는 factor reconstruction residual과 downdate failure를,
모든 approximate backend는 same-set exact direct 대비 state/energy error를 보고한다.
Fixed \(K_A\), fixed \(r_A\) 또는 fixed \(n_A\) 중 어느 하나도 동일 compute를 보장하지 않으므로
이 진단량과 measured p50/p95/p99를 공동 break-even 축으로 사용한다.
Active ratio만이 아니라 solve-set ratio \(\rho_{\mathcal S}\)와 set churn을 break-even 축으로 사용한다.
Empirical break-even \(\rho_A^{\mathrm{BE}}\) 또는 \(\rho_{\mathcal S}^{\mathrm{BE}}\) 이상에서는
always-on batched Local 또는 direct baseline이 더 빠를 수 있다.
이 경우 runtime dispatcher를 두거나 adaptive claim을 축소한다.

\subsection{Explicit target triangle-mesh reconstruction 비사용 계약}

\begin{longtable}{L{0.28\linewidth}L{0.25\linewidth}L{0.39\linewidth}}
\toprule
항목 & 허용 여부 & 설명 \\
\midrule
\endfirsthead
\toprule
항목 & 허용 여부 & 설명 \\
\midrule
\endhead
Training source mesh & 허용 & Privileged topology/physics teacher \\
\addlinespace
Explicit target triangle-mesh reconstruction & Mainline 금지 & 별도 end-to-end baseline/fallback에서만 허용 \\
\addlinespace
Sparse point/edge/local-relation scaffold & 허용 & Local tangent frame와 area weight를 가지되 global face connectivity, watertight/manifold triangle mesh를 요구하지 않음 \\
\addlinespace
Runtime MPM/FEM grid & 금지 & Teacher와 full-physics baseline에서만 사용 \\
\addlinespace
Prescribed wind grid/query & 허용 & Fluid dynamics state가 아니라 control input \\
\addlinespace
Runtime CFD/LBM & 금지 & Teacher-FSI 생성 또는 accuracy upper bound \\
\addlinespace
Static GS preprocessing & 허용 & Reliability, anchors, topology, binding inference \\
\bottomrule
\end{longtable}

\section{관련 연구: 유사점, 차이점, 비교 용도}

\subsection{직접적인 target-domain 경계}

\begin{longtable}{L{0.16\linewidth}L{0.23\linewidth}L{0.30\linewidth}L{0.23\linewidth}}
\toprule
방법 & 유사점 & 핵심 차이 & 실험에서의 역할 \\
\midrule
\endfirsthead
\toprule
방법 & 유사점 & 핵심 차이 & 실험에서의 역할 \\
\midrule
\endhead
Gaussian Swaying \cite{yan2026gaussianswaying}
& Current Gaussian surface normal, area, relative wind로 drag/lift 계열 force 계산
& Full MPM dynamics, 수동 simulation/material setup; object ROM과 error-gated Local 없음
& 동일 prescribed wind에서 Gaussian-surface aero front-end를 맞추고 MPM 대 ROM 정확도/latency 비교 \\
\addlinespace
DiffWind \cite{lei2026diffwind}
& Spatial wind, 3DGS-derived particles, wind-driven forward dynamics
& Grid/MPM/LBM 계열; surface-force-to-ROM과 sparse Local complement가 아님
& Spatial-wind subset에서 조건부 end-to-end 비교; teacher/upper-bound 후보 \\
\addlinespace
PhysGaussian \cite{xie2024physgaussian}
& Gaussian/particle deformation gradient와 covariance transport
& MPM 기반 generic material dynamics; wind-specific topology/gate 없음
& 동일 external load subset의 physics cost와 GS transport sanity comparison \\
\addlinespace
DynamicTree \cite{li2026dynamictree}
& Static real 3DGS, compact spectral/modal-like state, external-force response
& Tree/mesh/sparse voxel spectrum; current-surface aero와 Local corrector 없음
& Tree/leaf subset의 Global response와 latency 조건부 비교 \\
\addlinespace
FreeForm \cite{xiang2026freeform}
& Point/GS representation에서 mesh-free reduced deformable simulation
& Wind/thin-surface topology/current aero/error-gated local 없음
& 동일 anchor/외력에서 Global ROM 정확도와 basis construction 비교 \\
\addlinespace
PhysSkin \cite{lei2026physskin}
& Mesh-free/discretization-agnostic reduced skinning field, GS 적용
& Wind-specific force/operator와 local adaptive correction 없음
& Reduced lifting/geometry generalization subset 비교; 주로 구조 비교 \\
\addlinespace
GausSim \cite{shao2025gausim}
& Gaussian kernels의 hierarchical physics-aware dynamics와 계산 절감
& Fixed hierarchy/coarse-to-fine; runtime fixed-cardinality complement-response indicator가 아님
& 공통 force/input을 만들 수 있을 때만 conditional quantitative comparison \\
\addlinespace
NGFF \cite{li2026ngff}
& Gaussian object graph의 global/local neural force와 ODE dynamics
& Global은 rigid 6-DoF 중심이며 local stress; elastic modal complement와 wind-specific gate가 아님
& Global/local Gaussian force-field capability 비교; task 불일치 시 정성/인용 \\
\addlinespace
PGRD \cite{patel2026pgrd}
& Physics backbone + learned residual, paper/flag, closed-loop rollout, GS rendering
& Full particle graph의 all-particle residual; spatial error gate와 modal complement 없음
& PGRD-style always-on residual 대 adaptive Local의 핵심 정량 비교 \\
\addlinespace
FastPhysGS \cite{ma2026fastphysgs}
& MPM 기반 physics-aware dynamic 3DGS의 setup/optimization 가속
& Interior particle completion과 material optimization이 중심이며,
runtime reduced Global/selected Local wind solver가 아님
& Broad ``fast physics-based 3DGS'' claim의 경계;
same-scaffold runtime solver baseline으로 간주하지 않음 \\
\addlinespace
BendTwin \cite{jing2026bendtwin}
& Sparse physical graph에 local surface-triplet bending stiffness와 damping 사용
& Sparse-view RGB-D video 기반 digital-twin reconstruction이며,
independent static GS의 prescribed-wind reduced/adaptive animation이 아님
& ``Sparse graph + bending relation'' 자체의 신규성 경계;
입력/task 불일치로 related-work comparison에 사용 \\
\bottomrule
\end{longtable}

\subsection{Global aerodynamic/reduced dynamics 경계}

\begin{longtable}{L{0.16\linewidth}L{0.23\linewidth}L{0.30\linewidth}L{0.23\linewidth}}
\toprule
방법 & 유사점 & 핵심 차이 & 실험에서의 역할 \\
\midrule
\endfirsthead
\toprule
방법 & 유사점 & 핵심 차이 & 실험에서의 역할 \\
\midrule
\endhead
Wind Projection Basis \cite{diener2009windbasis}
& Wind load를 object modes에 투영하고 reduced dynamics 적분
& Rest geometry, uniform/directional wind 가정이 중심; current-deformed GS surface 없음
& Matched WPB-style baseline으로 current normal/area, relative/spatial wind의 필요성 검증 \\
\addlinespace
OmniAD \cite{martin2015omniad}
& Directional aerodynamic force/torque를 spherical harmonics로 표현
& Rigid object 전체, deformation state 없음
& RSH novelty 경계 인용; analytic/LUT/MLP/RSH 공력 표현 비교 \\
\addlinespace
MeshGraph\-Nets \cite{pfaff2021meshgraphnets}
& Current mesh state와 wind에서 flag-like dynamics를 direct neural rollout
& Mesh 입력/all-node graph network; explicit reduced physics base 없음
& 동일 topology/data로 direct neural Global 일반화와 안정성 비교 \\
\addlinespace
Nonlinear subspace integration \cite{barbic2005stvk}
& Reduced second-order nonlinear structural dynamics
& Generic volumetric elasticity; GS/wind/local gate 없음
& 수식의 선행 경계와 Global solver 구현 참고 \\
\addlinespace
Optimized cubature \cite{an2008cubature}
& Subspace force integration을 sparse samples로 가속
& GS surface, spatial aero, net wrench 목적이 아님
& Cubature 선행 인용; 본 연구는 force/moment-preserving aero adaptation을 검증 \\
\bottomrule
\end{longtable}

\subsection{Local correction/adaptivity 경계}

\begin{longtable}{L{0.16\linewidth}L{0.23\linewidth}L{0.30\linewidth}L{0.23\linewidth}}
\toprule
방법 & 유사점 & 핵심 차이 & 실험에서의 역할 \\
\midrule
\endfirsthead
\toprule
방법 & 유사점 & 핵심 차이 & 실험에서의 역할 \\
\midrule
\endhead
Subspace Condensation \cite{teng2015condensation}
& Novel event 주변 full-space region을 runtime 활성화, Global/Local coupling
& Volumetric contact, analytic adaptive solve; learned benefit/missing force/GS wind 없음
& Adaptive-local 선행 경계 인용; 동일 구현 비교보다 conceptual comparison \\
\addlinespace
Trading Spaces \cite{trusty2024tradingspaces}
& Solution quality oracle, local enrichment, quality--cost adaptivity
& Full-space contact elastodynamics; static GS/wind/neural defect 없음
& Gate/oracle 설계의 선행 경계; oracle benefit 및 tolerance curve 비교 개념 \\
\addlinespace
Learning Contact Corrections \cite{romero2021contact}
& Global reduced deformation 위 learned high-frequency local correction
& Contact-specific, always evaluated correction; runtime complement-response selector와 wind feedback 없음
& Direct local displacement correction ablation과 선행 인용 \\
\addlinespace
Fast Complementary Dynamics \cite{benchekroun2023complementary}
& Primary motion과 겹치지 않는 complementary secondary dynamics
& Rig-driven secondary motion; error-gated patch missing force와 wind 없음
& \(M\)-orthogonal complement 선행 경계와 no-complement ablation의 근거 \\
\addlinespace
Physics-Inspired Cloth Upsampling \cite{kavan2011upsampling}
& Coarse cloth 위 high-frequency learned/oscillatory detail, wind flag 사례
& Offline/detail upsampling 중심; current aero와 closed-loop local force 아님
& Global motion+detail 개념 경계 인용; 직접 quantitative baseline은 선택 \\
\bottomrule
\end{longtable}

\subsection{비교 구현의 명명 규칙}

공식 코드와 원 설정을 재현하지 못한 경우 논문명을 그대로 baseline 이름으로 쓰지 않는다.
다음처럼 구현 범위를 명시한다.

\begin{itemize}[leftmargin=1.5em]
  \item \emph{Wind Projection Basis-style rest-state modal load}.
  \item \emph{Gaussian Swaying-style surface aero + matched MPM}.
  \item \emph{PGRD-style all-node physics residual}.
  \item \emph{MeshGraphNets-style direct graph rollout}.
\end{itemize}

입력 privilege, physical scaffold/topology, timestep, wind control 또는
accuracy--latency operating point를 맞추지 못한 방법은
main quantitative leaderboard가 아니라 capability/failure-mode comparison으로 분리한다.
이하 same-anchor 또는 same-scaffold 비교는 동일 anchor graph, material, attachment,
wind, timestep, precision과 hardware를 뜻하며 동일 DOF를 뜻하지 않는다.
Reduced solver와 full-space solver의 실제 DOF, reduced rank, memory와 active coordinate 수를 별도 보고한다.

\clearpage
\section{비교 실험 대장}

이 절은 방법 설명에서 분리된 실험 계획이다.
모든 실험은 먼저 검증하려는 claim을 정하고,
통제 변수, 비교군, 지표, 실패 시 결정을 명시한다.
모든 조합의 Cartesian product를 실행하기보다
단계별 controlled test와 최종 모델의 one-factor ablation을 병행한다.

\subsection{검증할 핵심 claim}

\begin{longtable}{L{0.08\linewidth}L{0.36\linewidth}L{0.48\linewidth}}
\toprule
ID & Claim & 반드시 필요한 증거 \\
\midrule
\endfirsthead
\toprule
ID & Claim & 반드시 필요한 증거 \\
\midrule
\endhead
MC1 & Training-only mesh privilege로 학습하되 independently reconstructed static GS와 fixed setup metadata에서
distill한 scaffold가 explicit target triangle-mesh reconstruction 없이 thin-surface reduced physics를 유지한다.
& Oracle/proposed/kNN/FPS/target-mesh baseline의 graph, area/mass, eigenstructure와 downstream dynamics 차이 \\
\addlinespace
K0 & Global reduced physics가 object-wide 저주파 response를 안정적이고 controllable하게 계산한다.
& Full physics, sparse physics, WPB-style, direct/recurrent neural Global, hybrid Global과의 trajectory/frequency/damping/latency 비교 \\
\addlinespace
MC2 & Selective complementary Local physics가 Global ROM이 놓친 high-frequency response를 필요한 patch에서 복원한다.
& Global/POD/wind-aware basis와 rank sweep, full/reduced physics Local,
never/always/oracle/cached-compliance selector의 trajectory/spectrum/latency 비교 \\
\addlinespace
SYS & Patch force 선조립, complement, decay와 corrected-minus-predictor feedback이 coherent한 closed-loop state와 directly renderable transport를 만든다.
& Integrate-average/no-complement/no-feedback ablation의 seam, force/torque, energy, Global response와 Gaussian validity 비교 \\
\addlinespace
O1 & 조건부 aero hyper-reduction이 full-anchor backend보다 generalized-load quality--latency를 개선한다.
& 512/1,024/2,048 anchors의 full/random/FPS/cubature와 localized gust profile; 실패 시 module 제거 \\
\addlinespace
O2 & 조건부 learned missing force가 physics-only Local의 반복 가능한 defect를 줄인다.
& Full/reduced/hybrid Local, nodal lifting consistency와 rollout 비교; 실패 시 network 제거 \\
\addlinespace
O3 & 조건부 learned benefit gate가 cached-compliance Top-\(K\)보다 같은 \(K_A\)/measured p95에서 낫다.
& Fixed-cardinality \(K_A\) error와 actual skip/wall-clock;
random, heuristic, cached compliance, full-counter analysis, oracle, learned 비교 \\
\addlinespace
SYS & 전체 시스템이 wind-driven static 3DGS에서 유용한 quality--latency Pareto와 OOD/장기 안정성을 제공한다.
& Same-anchor kill test, end-to-end static-GS track, real captured GS demo, 30--60초 stress test와
고정 physical budget에서 분리한 \(T_{\mathrm{dyn}}/T_{\mathrm{transport}}/T_{\mathrm{render}}\) 대 \(N_G\) scaling \\
\bottomrule
\end{longtable}

\subsection{평가 Track 분리}

\subsubsection{Track A: Oracle-topology matched solver comparison}

모든 방법에 다음 정보를 동일하게 제공한다.

\begin{itemize}[leftmargin=1.5em]
  \item 동일 teacher surface/topology,
  \item 동일 mass, material, attachment,
  \item 동일 wind field와 initial state,
  \item 동일 timestep,
  \item 동일 rest geometry/anchor, mass, user material/attachment와 load;
  native energy/operator privilege와 실제 simulated DOF/rank는 별도 보고,
  \item 동일 rendering transport.
\end{itemize}

이 Track은 topology reconstruction 품질을 배제하고
aero operator, Global solver, Local correction, gate, coupling, runtime만 비교한다.
Method accuracy의 주장은 Track A에서 먼저 성립해야 한다.

\subsubsection{Track B: Static-3DGS end-to-end comparison}

각 방법은 원래 요구하는 입력을 사용하되, 입력 특권을 명시한다.

\begin{itemize}[leftmargin=1.5em]
  \item 제안 방법: independently reconstructed static 3DGS와 setup 시 지정한 metric/material/attachment metadata.
  \item Target mesh baseline: target static GS에서 mesh를 별도 추출한 뒤 simulation.
  \item MPM baseline: 해당 preprocessing과 particle/grid representation.
  \item Neural baseline: 요구하는 topology/mesh/data를 표에 명시.
\end{itemize}

이 Track은 explicit target triangle-mesh reconstruction 비사용 property, preprocessing robustness,
최종 dynamics/rendering, setup/runtime cost를 평가한다.
입력 특권이 다른 결과를 하나의 공정한 solver leaderboard로 해석하지 않는다.

\subsubsection{Track C: Real captured static GS}

실제 multi-view capture에서 재구성한 flag/cloth/leaf asset을 사용한다.
가능하면 marker, silhouette, tracked feature, wind sensor를 수집한다.
정확한 material/force ground truth가 없으면 정성 결과와 제한된 trajectory/silhouette metric으로 보고한다.

\subsection{Test asset와 wind program}

\begin{longtable}{L{0.19\linewidth}L{0.25\linewidth}L{0.25\linewidth}L{0.23\linewidth}}
\toprule
Asset & 핵심 현상 & Attachment & 주요 실험 \\
\midrule
\endfirsthead
\toprule
Asset & 핵심 현상 & Attachment & 주요 실험 \\
\midrule
\endhead
Short cantilever strip & 단일 bending mode, damping, 시작/정지 & 한쪽 edge hard pin & Global unit test, frequency/damping \\
\addlinespace
Rectangular long flag & traveling wave, tip flutter, mode truncation & 한쪽 edge hard/compliant & Global/Local/gate main benchmark \\
\addlinespace
Triangular pennant & 비대칭 geometry, tip motion & 한쪽 short edge & Generalization과 aero torque \\
\addlinespace
Hanging cloth & 여러 attachment와 broad deformation & top edge 또는 two points & Secondary topology/attachment stress test \\
\addlinespace
Curtain strip & spatially varying gust와 long-range response & top edge/track & Spatial wind, active patch 이동 \\
\addlinespace
Paper-like elastic sheet & stiff bending, large-bending local response without permanent creasing & two points/partial clamp & Transfer/qualitative stress test \\
\addlinespace
Single leaf & anisotropy, stem attachment, twist & compliant stem & Material/stem OOD qualitative \\
\bottomrule
\end{longtable}

\reviewaccept{첫 논문의 core quantitative dynamics suite를
\emph{short cantilever strip}, \emph{rectangular long flag},
\emph{triangular pennant}, \emph{curtain strip} 네 종류로 고정한다.
Main dynamics table과 Gate B/C의 primary 통계는 이 네 asset에서 산출한다.
Hanging cloth는 secondary topology/attachment stress test로 두고,
paper-like elastic sheet와 single leaf/petal은 transfer 또는 qualitative 결과로 내린다.
특히 leaf의 anisotropy, compliant stem과 front/back appearance를 동시에 해결한 결과를
첫 논문의 필수 성공 조건으로 사용하지 않는다. 다만 MC1 Gate A의 representation preflight는
목적이 다르므로 5--10개 independent-GS simple object 범위를 그대로 유지한다.}

Wind program은 다음을 포함한다.

\begin{enumerate}[leftmargin=1.7em]
  \item Zero-wind relaxation.
  \item Constant weak/medium/strong wind.
  \item Abrupt wind start와 abrupt stop.
  \item Sinusoidal gust.
  \item Frequency chirp.
  \item Yaw/elevation directional sweep.
  \item Stationary, slow-traveling, fast-traveling과 reversing gust.
  \item Localized spatial gust.
  \item OOD strong wind failure test.
\end{enumerate}

\subsection{공통 baseline registry}

\begin{longtable}{L{0.08\linewidth}L{0.30\linewidth}L{0.29\linewidth}L{0.25\linewidth}}
\toprule
ID & Baseline & 통제/입력 & 검증 목적 \\
\midrule
\endfirsthead
\toprule
ID & Baseline & 통제/입력 & 검증 목적 \\
\midrule
\endhead
B0 & High-resolution nonlinear thin-shell/MPM teacher & Full mesh/full resolution &
Privileged accuracy oracle/label source; runtime efficiency competitor가 아님 \\
\addlinespace
B1 & Same-anchor sparse corotational shell, PD, 또는 XPBD & 동일 rest geometry/anchor, mass, material/attachment/wind/hardware; native energy/operator privilege와 실제 DOF 별도 보고 & MPM 대비 해상도 차이만으로 얻은 speedup인지 확인 \\
\addlinespace
B2 & WPB-style rest-state modal wind load & 동일 \(\Phi,M,C,K\), rest normal/area, uniform wind & Current state/spatial aero contribution \\
\addlinespace
B3 & Current aero + all samples + Global ROM & Oracle topology, no hyper-reduction & Global aero accuracy upper bound \\
\addlinespace
B4 & Direct neural Global MLP/GNN & 동일 state/features/data & Physics Global 필요성 \\
\addlinespace
B5 & Recurrent neural Global GRU/state-space & 동일 state/history/data & Neural memory를 포함한 공정 비교 \\
\addlinespace
B6 & Physics Global + learned \(\Delta Q_g\) & 동일 Global base & Global hybrid correction 필요성 \\
\addlinespace
B7 & Global-only proposed & Local off & Local 순수 기여 \\
\addlinespace
B8 & Global + full-resolution Local physics & Oracle active region/full local solver & Local accuracy upper bound \\
\addlinespace
B9 & Global + reduced Local physics only & 동일 Local basis, network off & Missing-force network 필요성 \\
\addlinespace
B10 & Global + direct local displacement network & 동일 patch features/data & Force prediction+integration 필요성 \\
\addlinespace
B11 & PGRD-style all-node/all-patch residual & 동일 base graph와 training data & Always-on physics residual 대 adaptive Local \\
\addlinespace
B12 & Proposed + always-on Local physics/expert & 모든 patch active; physics-only와 조건부 hybrid 분리 & Local quality upper bound/비용 상한 \\
\addlinespace
B13 & Proposed + oracle-benefit Top-\(K\) & Offline oracle selection & Selector upper bound \\
\addlinespace
B14 & Proposed + cached-compliance Top-\(K\) & Teacher 없이 every-frame runtime 가능 & V0 main selector \\
\addlinespace
B14h & Proposed + normalized deterministic hybrid Top-\(K\) & Runtime-observable strain/curvature/truncation feature와 construction cost 포함 & Cached-compliance blind-spot diagnostic; Pareto 개선 시에만 후보 \\
\addlinespace
B14b & Proposed + full counter-state residual Top-\(K\) & Teacher-free이나 extra solve와 force 재평가 비용 포함 & Offline/low-rate analysis ceiling \\
\addlinespace
B15 & Proposed + learned-benefit Top-\(K\) & 조건부 gate backend & O3가 통과할 때만 최종 후보 \\
\bottomrule
\end{longtable}

\section{E0: Deterministic contract와 단위 테스트}

\subsection*{목적}

Network error와 physics/geometry implementation bug를 분리한다.
한 항목이라도 실패하면 long training을 시작하지 않는다.

\subsection*{테스트}

\begin{longtable}{L{0.26\linewidth}L{0.34\linewidth}L{0.32\linewidth}}
\toprule
대상 & 입력 & 통과 조건 \\
\midrule
\endfirsthead
\toprule
대상 & 입력 & 통과 조건 \\
\midrule
\endhead
Coordinate alignment & Known mesh/GS cameras와 marker & Metric scale/axis/handedness 일치 \\
\addlinespace
Area/mass partition & Dense Gaussian support & Total area/mass 보존, density dependence 없음 \\
\addlinespace
Confidence aggregator & Synthetic low-confidence/variable-valence/ill-conditioned relation candidates & 식~\eqref{eq:relation-absolute-confidence}의 count-aware aggregation, support 확대와 fallback/reject 경로 통과; learned calibration accuracy는 E1에서 별도 판정 \\
\addlinespace
Operator & Random V0-valid \((\rho,h,E,\tau_{\mathrm m},\tau_{\mathrm b})\), \(\nu=0\), scale/refinement fixtures & 식~\eqref{eq:operator-construction-hard-checks}, MLS rank/moment reproduction, area normalization, hard/compliant attachment exclusivity, 식~\eqref{eq:klin-projective-hard-gate}와 식~\eqref{eq:operator-dimensional-scaling} 통과 \\
\addlinespace
Bending projector & \(x_{\mathrm{tan}}=(\xi_1+a\xi_1^2,\xi_2,0)\), \(x_{\mathrm{nor}}=(\xi_1,\xi_2,\tfrac12\kappa\xi_1^2)\), 대응 rate, shallow rest와 superposed rotation & Planar in-plane probe의 projected bending energy/force/dissipation이 tolerance 이하; normal quadratic scaling과 shallow-rest zero energy 통과; one/two/converged projection의 tangential lag, torque, phase/PSD tolerance 통과 \\
\addlinespace
Patch complement & Random \(q,z\), active groups & 모든 \(r\)의 \(\|\Phi^\top M\Psi_r\|\), fixed ID, full \(M_{\mathcal SS}\) rank/conditioning 통과; tail은 식~\eqref{eq:tail-downstream-degradation}과 함께 calibration하며 보편 threshold를 요구하지 않음 \\
\addlinespace
Aero geometry & Rest/rigid/in-plane affine & Tangent-cross area/normal과 full-shell map의 identity/rotation/rank 통과 \\
\addlinespace
Aero force & Zero/rigid/relative wind & Zero relative wind force 0, rotation equivariance \\
\addlinespace
Optional cubature & Held-out/localized-gust force snapshots & Generalized force primary와 auxiliary wrench tolerance \\
\addlinespace
Force lifting & Nodal/generalized teacher labels & KKT feasibility, \(\epsilon_Q\), partition reassembly tolerance \\
\addlinespace
Overlap & Synthetic opposing patch forces & Constraint residual, net force/torque 보존 \\
\addlinespace
Joint constraint backend & Singleton, overlap, boundary와 changing support; rank-deficient but RHS-consistent \(\Phi_t\), support leakage와 hard edge &
External restricted projection과 augmented structural/hard-reaction small-Schur를 각각 dense reference와 대조; 식~\eqref{eq:support-routing-diagnostics}--\eqref{eq:structural-schur-adjustment-diagnostics}, feasible rate뿐 아니라 adjustment/reaction amplification/generalized-work 보존, halo frequency, stale-factor와 sequential fallback 거부 \\
\addlinespace
Reaction ledger & Free sheet, one-edge hard attachment와 coordinate compliant spring/damper & Free intrinsic wrench 0; free correction+object-to-support hard reaction wrench 0; support-on-object 부호 반대; compliant object/handle virtual-work closure와 single ownership; final \(r_{\mathrm{ROM}}\) 및 reaction-reference error 보고 \\
\addlinespace
Structural backend/integrator & Pure rigid 및 strained-state superposed rotation, known oscillator, cantilever/flag & 식~\eqref{eq:v0-structural-freeze-signature}의 canonical 조합 검증; rest-linear/BE와 one-shot/two-step/converged local--global projective reference 대비 energy/force objectivity, torque residual, frequency/damping/phase/PSD, timestep convergence와 p50/p95/p99 통과 \\
\addlinespace
Corrector & Synthetic Local deformation & Corrected-minus-predictor delta만 적용 \\
\addlinespace
GS transport & Identity, translation, rotation, affine & Mean / covariance reproduction, SPD 유지 \\
\bottomrule
\end{longtable}

\section{E1: Topology와 operator distillation}

\subsection*{검증 claim}

MC1: Static GS와 setup metadata에서 만든 runtime package가 Oracle topology에 가까운 dynamics를 제공하는가?

\subsection*{비교군}

\begin{enumerate}[leftmargin=1.7em]
  \item Oracle source mesh topology/relation + canonical \(\nu=0\) assembly;
  native teacher operator projection은 별도 privileged ceiling.
  \item Proposed oriented topology distillation.
  \item Euclidean kNN graph.
  \item FPS anchors + kNN.
  \item Orientation/bending relation 제거.
  \item Target GS에서 SuGaR/GOF-style mesh extraction 후 physics.
  \item 동일 user-provided attachment의 Oracle-anchor mapping 대 proposed anchor transfer.
\end{enumerate}

\subsection*{정적 topology 지표}

여기서 eigenfrequency와 modal subspace는 명시적으로
rest-linear pair \((M_{\mathrm f},K_{\mathrm{lin,f}})\)에서 계산한다.
Canonical projective runtime operator는 같은 eigenspectrum이라고 가정하지 않고,
teacher internal-force action, objectivity probe와 nonlinear rollout 지표로 별도 판정한다.

\begin{itemize}[leftmargin=1.5em]
  \item Material edge precision/recall/F1.
  \item Close disconnected layer false-positive rate.
  \item Normal orientation consistency.
  \item User attachment-to-anchor transfer error.
  \item Area/mass relative error.
  \item Relation-validity Brier/ECE, anchor absolute-confidence coverage와
  low-confidence fallback/rejection precision/recall.
  \item Eigenfrequency relative error.
  \item Modal subspace principal angles.
  \item Patch/group complement tail-ratio distribution,
  asset/load별 calibrated threshold 초과율과 tail--trajectory/PSD conditional-quantile curve.
  \item Anchor-density별 affine/quadratic patch energy와 첫 8--16개 eigenfrequency convergence.
\end{itemize}

\subsection*{Downstream 지표}

Global/Local trajectory, normal/curvature, aerodynamic load,
long-rollout stability를 Oracle topology와 비교한다.
Eigenfrequency/modal-subspace error와 downstream rollout error의 관계를
asset와 wind/load별 scatter로 보고하되 보편적인 단조관계로 과장하지 않는다.
Topology F1만 높고 downstream dynamics가 나쁘면 MC1은 성립하지 않는다.

\subsection*{Independent-GS package consistency}

같은 underlying asset \(a\)에서 view 수, seed, Gaussian density, hole과 floater를 달리해
독립적으로 복원한 variant \(v\)마다 package \(\mathcal P_{a,v}\)를 다시 만든다.
Anchor ID와 matrix dimension이 다를 수 있으므로 raw adjacency나 dense \(K\) entry를 직접 비교하지 않는다.
Modal vector와 rollout은 공통 material probe set으로 transfer한 뒤 비교한다.
Variant pair \(v,v'\)에 대해

\begin{equation}
\begin{aligned}
  d_A^{v,v'}
  &=\frac{|A_v-A_{v'}|}{\tfrac12(A_v+A_{v'})+\eps},
  &d_m^{v,v'}
  &=\frac{|m_v-m_{v'}|}{\tfrac12(m_v+m_{v'})+\eps},\\
  d_\omega^{v,v'}
  &=\|\log\boldsymbol\omega_v-\log\boldsymbol\omega_{v'}\|_2,
  &d_\Phi^{v,v'}
  &=\|\sin\Theta_{M_\star}(\widehat\Phi_v,\widehat\Phi_{v'})\|_2,\\
  d_\rho^{v,v'}
  &=W_1(\mathcal R_{\mathrm{tail},v},\mathcal R_{\mathrm{tail},v'}),
  &d_{\mathrm{roll}}^{v,v'}
  &=\operatorname{NRMSE}(x_{\star,v}^{1:T},x_{\star,v'}^{1:T}).
  \label{eq:independent-gs-package-consistency}
\end{aligned}
\end{equation}

여기서 \(\widehat\Phi_v\)와 \(x_{\star,v}^{1:T}\)는 common probe로 transfer한 mode와 trajectory이고,
\(\mathcal R_{\mathrm{tail},v}\)는 tail-ratio empirical distribution이다.
각 asset에서 pairwise median, IQR, 95th percentile과 worst case를 보고하고,
Oracle-relative package error와 variant 간 consistency를 별도 열로 유지한다.
서로 단위가 다른 위 항을 임의의 단일 scalar score로 합치지 않는다.

\subsection*{통제 변수}

동일 teacher state, wind, material, timestep, Global rank,
Local basis rank, aerodynamic samples를 사용한다.

\subsection*{반증과 결정}

\killgate Predicted topology dynamics error가 Oracle 대비 지속적으로 크게 악화되거나
close-layer false connection이 안정성을 깨면
explicit target triangle-mesh reconstruction 비사용 claim을 보류한다.
첫 paper를 oracle/simple topology asset으로 축소하거나
target mesh extraction을 명시적인 setup baseline으로 채택한다.

\section{E2: Aerodynamic operator와 RSH 대안}

\subsection*{검증 claim}

O1: Full-anchor current-deformed analytic aero의 정확도와 비용은 얼마이며,
profiling을 통과한 경우에만 conservative hyper-reduction이 필요한가?

\subsection*{실험 격리}

Teacher/current surface positions, velocities, normals, areas를 고정 입력으로 제공하고
dynamics integration을 끈다.
이렇게 해야 구조 solver error를 aerodynamic operator error로 오해하지 않는다.

\subsection*{비교군}

\begin{enumerate}[leftmargin=1.7em]
  \item Rest surface + uniform wind modal load: WPB-style.
  \item Current surface + absolute wind.
  \item Current surface + relative wind + all samples.
  \item Current surface + random samples.
  \item Current surface + FPS samples.
  \item Learned samples without wrench constraint.
  \item Proposed conservative hyper-reduction.
  \item Directional LUT + interpolation.
  \item Small directional MLP.
  \item Fitted RSH upper bound.
  \item Predicted/state-conditioned RSH.
  \item Teacher-FSI force, 해당 subset이 있을 때만.
\end{enumerate}

\subsection*{입력 sweep}

Held-out yaw/elevation, wind speed, deformed state,
surface velocity, traveling gust, local gust,
different GS reconstruction을 사용한다.

\subsection*{지표}

\begin{itemize}[leftmargin=1.5em]
  \item Sample force error.
  \item Primary: Global generalized-force relative error.
  \item Auxiliary: net force와 net torque relative error.
  \item Drag/lift coefficient error.
  \item Unsteady subset의 amplitude/phase error.
  \item Sample count, memory, p50/p95/p99 GPU latency.
\end{itemize}

\subsection*{결정 규칙}

\killgate 512/1,024/2,048 anchors에서 full-anchor aero+fixed \(\Phi_s^\top\) projection이
dynamics의 유의미한 병목이 아니면 cubature를 구현하지 않고 full backend를 mainline으로 확정한다.
\killgate Direct current-surface analytic force가 RSH/LUT/MLP보다
정확하고 같거나 더 빠르면 RSH는 method에서 제거하고 comparison 결과로만 남긴다.
Conservative sampling이 full/random/FPS보다 generalized-load quality--latency에서 이득이 없거나
localized gust를 놓치면 hyper-reduction을 제거한다.

\section{E3: Global physics 대 neural Global}

\subsection*{검증 claim}

K0: Global reduced physics가 direct neural state transition보다
안정성, control, generalization에서 이점이 있는가?

\subsection*{비교군}

\begin{enumerate}[leftmargin=1.7em]
  \item Full teacher.
  \item Same-anchor sparse PD/XPBD/corotational shell.
  \item WPB-style Global.
  \item Proposed physics Global.
  \item Physics Global + learned generalized-force correction.
  \item Direct next-state MLP.
  \item MeshGraphNets-style direct GNN rollout.
  \item Recurrent GRU/state-space Global.
\end{enumerate}

Direct neural baselines는 가능한 한 동일 topology,
동일 wind/material/attachment features, 동일 data volume을 사용한다.

\subsection*{조건}

Local을 모두 끄고 Oracle topology를 사용한다.
Steady wind, start/stop, sinusoidal gust, chirp, traveling gust,
unseen direction/speed, declared-V0-family scalar material parameter와 attachment에서 평가한다.

\subsection*{지표}

\begin{itemize}[leftmargin=1.5em]
  \item Object-scale normalized position/velocity error.
  \item Tip amplitude와 phase lag.
  \item Dominant frequency와 mode response.
  \item Damping ratio와 wind-stop envelope.
  \item Work--energy residual과 divergence rate.
  \item 4--10초 및 30--60초 rollout.
  \item Per-step p50/p95/p99 latency와 memory.
\end{itemize}

\subsection*{결정 규칙}

\killgate 동일 full-anchor aero와 accuracy level에서 same-anchor sparse solver보다
dynamics-only 개발용 기준인 약 \(2\times\) 빠르지도 않고
same-p95 velocity/trajectory/spectrum 이득도, multi-object throughput 이득도 없으면
MPM 대비 결과와 무관하게 SYS quality--latency headline을 철회한다.
Direct neural Global이 OOD와 long rollout에서도 동등하거나 우수하고 더 빠르면
Global physics 선택을 재검토한다.
Physics + learned \(\Delta Q_g\)가 pure physics를 유의미하게 개선하지 않으면
Global neural correction을 제거한다.

\section{E4: Local physics와 조건부 missing-force network}

\subsection*{검증 claim}

MC2/O2: Global basis가 놓친 detail이 실제로 존재하며,
reduced complementary Local physics와 조건부 learned missing force가 각각 필요한가?

\subsection*{비교군}

\begin{longtable}{L{0.30\linewidth}L{0.58\linewidth}}
\toprule
Variant & 검증 내용 \\
\midrule
\endfirsthead
\toprule
Variant & 검증 내용 \\
\midrule
\endhead
Global-only & Local contribution 자체의 필요성 \\
\addlinespace
Global rank 증가 & Local 대신 Global capacity 확대로 해결 가능한지 \\
\addlinespace
Full-resolution Local physics & 비용을 무시한 local accuracy upper bound \\
\addlinespace
Reduced Local physics-only & 저렴한 known physics만으로 충분한지 \\
\addlinespace
Reduced Local + learned missing force & 최종 Local expert의 추가 이득 \\
\addlinespace
Learned force, local integrator 없음 & 물리 적분기의 필요성 \\
\addlinespace
Direct local displacement/velocity network & Force-space hybrid의 안정성 이점 \\
\addlinespace
PGRD-style all-patch residual & Sparse Local 선택의 이점 \\
\bottomrule
\end{longtable}

\subsection*{평가 영역}

Flag free tip, high-curvature/fold-like region, attachment transition,
traveling gust front, topology uncertainty가 큰 영역을 따로 보고한다.

\subsection*{지표}

\begin{itemize}[leftmargin=1.5em]
  \item Patch-local position/velocity/normal/curvature error.
  \item High-frequency power spectrum distance.
  \item Flutter amplitude, frequency, phase.
  \item Wind-stop decay와 residual energy.
  \item Local force error와 local compute time.
  \item Global-span leakage.
\end{itemize}

\subsection*{Oracle latency accounting}

Oracle future-benefit과 oracle active set 생성은 privileged offline computation이므로
가상의 online selector latency를 부여하지 않는다.
대신 Oracle 결과는

\begin{equation}
\begin{aligned}
  T_{\mathrm{oracle\mbox{-}exec}}
  =&\;T_{\mathrm{proposal}}
  +T_{\mathrm{unique\mbox{-}assembly/Schur/projection}}
  +T_{\mathrm{factor/solve}}
  +T_{\mathrm{scatter}}
  +T_{\mathrm{corrector}},\\
  T_{\mathrm{oracle\mbox{-}label}}
  =&\;T_{\mathrm{privileged\ future\ rollout}}
  \label{eq:oracle-latency-accounting}
\end{aligned}
\end{equation}

로 분리한다.

Same-p95 Gate B는 증분 비용끼리 비교하지 않고 synchronized total-frame trace

\begin{equation}
\begin{aligned}
  T_{G+\mathrm{Oracle}}
  &=T_{\mathrm{common}}+T_{\mathrm{Global}}(r_g)
    +T_{\mathrm{oracle\mbox{-}exec}},\\
  T_{G\mathrm{\mbox{-}rank}}
  &=T_{\mathrm{common}}+T_{\mathrm{Global}}(r_g')
  \label{eq:oracle-total-latency-fairness}
\end{aligned}
\end{equation}

의 measured p95로 맞춘다. \(T_{\mathrm{common}}\)은 두 경로가 공유하는 surface/aero와
frame bookkeeping이고, Oracle 경로에서 제외하는 것은
\(T_{\mathrm{oracle\mbox{-}label}}\)의 privileged future rollout뿐이다.

\(T_{\mathrm{oracle\mbox{-}exec}}\)만 사용한 curve는
\emph{post-selection Local-capacity upper bound}로 표시하고 deployable runtime speed claim에 사용하지 않는다.
식~\eqref{eq:oracle-total-latency-fairness}의 Global rank sweep 대 Oracle Local 비교는
Gate B에서 Local complement의 필요성과 capacity를 판정할 때만 사용한다.
Gate C와 최종 deployable Pareto에는 cached-compliance, heuristic 또는 learned selector의
feature construction, score, Top-\(K\), gather/assembly, cache miss와 solve 비용을 모두 포함한다.

\subsection*{결정 규칙}

\killgate 동일 p95 dynamics latency의 Global rank sweep이 Oracle Local과
velocity/trajectory/PSD--spectrum에서 동등하거나 우수한 경우
Local contribution을 축소한다.
Mode/rank-matched 비교는 보조 진단으로만 사용한다.
Reduced Local physics-only가 learned missing force와 동등하면 network를 제거하고
pure adaptive physics로 전환한다.
Direct displacement가 같은 데이터에서 더 안정적이고 빠르면
force-space 설계를 다시 검토하되 physical control/OOD 결과를 함께 판단한다.

\section{E5: Fixed-cardinality selector와 조건부 learned gate}

\subsection*{검증 claim}

MC2/O3: Fixed-cardinality selector가 always-on Local의 품질을 유지하면서 실제 실행 비용을 줄이며,
learned gate가 cached-force dynamic-compliance selector보다 추가 이득을 주는가?

\subsection*{동일 Local expert를 사용한 gate 비교}

\begin{enumerate}[leftmargin=1.7em]
  \item Never Local.
  \item Always Local.
  \item Random Top-\(K\).
  \item Curvature Top-\(K\).
  \item Strain/deformation severity Top-\(K\).
  \item 식~\eqref{eq:cheap-dynamic-compliance}의 cached-compliance Top-\(K\).
  \item 식~\eqref{eq:normalized-hybrid-selector}의 normalized deterministic hybrid Top-\(K\):
  feature construction 비용을 포함한 조건부 ablation.
  \item Nominal-coordinate-density diagnostic:
  \(d_{\mathcal G}^t=s_{\mathcal G}^t/(r_{\mathcal G}+\eps_r)\)로 정렬하고
  \(\sum_{\mathcal G\in\mathcal A_t^{\mathrm{grp}}}r_{\mathcal G}\le R_A\)를 만족하도록
  greedy pack하는 coordinate-cap baseline.
  \item 식~\eqref{eq:full-counter-state-residual}의 full counter-state residual Top-\(K\):
  추가 비용을 포함한 offline analysis baseline.
  \item Low-rate full counter-state residual Top-\(K\), \(J\in\{2,4,8\}\): 중간 score hold.
  \item Teacher residual threshold: runtime 불가능한 분석 기준.
  \item Oracle raw future-benefit Top-\(K\).
  \item Learned binary error Top-\(K\).
  \item Conditional learned raw-benefit Top-\(K\).
\end{enumerate}

모든 비교는 Module 6의 cross-group excitation, high-curvature/low-load release,
topology/operator-error negative control과 fast traveling/reversing gust fixture를 포함한다.

\subsection*{Gate 자체 지표}

\begin{itemize}[leftmargin=1.5em]
  \item High-benefit patch recall/precision/AUPRC.
  \item Expected calibration error.
  \item Oracle 대비 benefit regret.
  \item Active/decay patch ratio.
  \item Actual Local block solve와 조건부 missing-force network call 수.
  \item Kernel launch 수와 p50/p95/p99 dynamics latency.
  \item Selector-only p50/p95/p99 latency와 추가 geometry/aero/internal evaluation 수.
  \item Hybrid feature construction latency와 feature별 ablation.
  \item Cached-compliance와 full-counter score의 rank correlation, Oracle regret와 activation delay.
  \item Patch toggle count와 activation duration.
  \item Activation acceleration/jerk.
\end{itemize}

\subsection*{공정한 비교}

세 방식으로 비교한다.

\begin{enumerate}[leftmargin=1.7em]
  \item 동일 atomic gating-unit cardinality \(K_A\)에서 가장 낮은 dynamics error와 실제 비용.
  \item 동일 nominal coordinate cap \(R_A\)에서 coordinate-density diagnostic.
  \item 동일 p95 latency budget에서 가장 낮은 dynamics error.
\end{enumerate}

Gate classification score만 높고 wall-clock가 줄지 않으면 adaptive-compute claim은 성립하지 않는다.

\reviewrebut{V2 리뷰의 ``동일 p95 또는 비슷한 coordinate budget''을 동급 Gate B/C 기준으로
두는 처방은 채택하지 않는다. Global과 Local coordinate는 spatial support, cross block,
joint assembly, factorization/cache miss, gather/scatter와 kernel-launch 비용이 달라 동일 coordinate 수가
동일 계산량을 뜻하지 않는다. Coordinate-density와 mode/rank-matched 결과는 구조적 원인을 해석하는
diagnostic ablation으로만 허용하며, 그 결과만으로 Gate를 통과시키거나 hardware-optimal selection을 주장하지 않는다.
Primary 판정은 동일 measured p95의 object-scale-normalized velocity RMSE이며,
trajectory와 PSD/spectrum 및 전체 quality--latency Pareto를 함께 본다.}

\subsection*{결정 규칙}

\killgate Oracle Top-\(K\)가 cached-compliance Top-\(K\)보다 거의 낫지 않으면 learned gate를 학습하지 않는다.
\killgate Learned gate가 cached-compliance/curvature/strain보다 같은 \(K_A\)에서 동등하면
learned gate를 제거하고 가장 단순한 runtime selector를 쓴다.
Full counter-state residual이 cheap score보다 정확하더라도 추가 비용을 포함한
quality--latency가 나쁘면 every-frame runtime mainline으로 승격하지 않는다.
Always-on Local 대비 의미 있는 speedup 없이 error만 증가하면 adaptive claim을 철회한다.

\section{E6: Complement와 overlap conservation}

\subsection*{검증 claim}

SYS: Local force/state를 Global complement로 제한하고
overlap을 보존적으로 조립해야 double counting과 seam이 줄어드는가?

\subsection*{Ablation}

\begin{enumerate}[leftmargin=1.7em]
  \item No complement.
  \item Runtime post-projection only.
  \item Offline complement basis only.
  \item Full-force target instead of missing-force target.
  \item Patch integrate-then-average.
  \item Partition-of-unity displacement only.
  \item Force assembly without torque constraint.
  \item Small-fixture dense primal KKT reference.
  \item Proposed unique-anchor small-Schur assembly.
  \item Equivalent null-space joint solve.
  \item Frozen rest-torque map approximation.
\end{enumerate}

\subsection*{지표}

\begin{itemize}[leftmargin=1.5em]
  \item \(\|\Phi^\top M\Psi\|\)와 Global-span force leakage.
  \item Global modal amplitude inflation.
  \item Patch boundary position/velocity jump.
  \item Net force/torque error.
  \item Artificial center-of-mass drift/rotation.
  \item Energy jump at activation.
  \item Patch partition 변화에 대한 sensitivity.
  \item Dense reference 대비 force/objective gap과 모든 constraint residual.
  \item Constraint raw-row/rank/condition, factor reuse/rebuild와 assembly latency share.
  \item Support-route consistency/conditioning/amplification,
  signature별 feasible rate, halo expansion frequency와 final support ratio.
  \item Free/hard/compliant reaction의 sign-resolved world-wrench 및 virtual-work residual.
\end{itemize}

\subsection*{필수 시각화}

Active core/halo map, overlap force arrows,
Global component와 Local complement component,
seam acceleration, net wrench curve를 한 sequence에 나란히 표시한다.

\section{E7: Local-to-Global feedback}

\subsection*{검증 claim}

SYS의 feedback 부분:
Local correction이 visual overlay가 아니라 Global aeroelastic state를 실제로 바꾸는가?

\subsection*{비교군}

\begin{enumerate}[leftmargin=1.7em]
  \item No feedback: Local displacement는 rendering에만 추가.
  \item Structural cross-block feedback only.
  \item Corrected-aero delta only.
  \item Next-frame feedback.
  \item Same-frame one corrector.
  \item Same-frame multiple fixed-point iterations.
  \item Proposed minimal valid combination.
\end{enumerate}

\subsection*{지표}

\begin{itemize}[leftmargin=1.5em]
  \item Global modal coordinate와 tip trajectory 차이.
  \item Corrected aerodynamic generalized-force delta.
  \item Phase/amplitude와 teacher error.
  \item Corrected geometry reconstruction latency;
  optional nonlinear sample-map일 때만 Jacobian reconstruction latency를 추가.
  \item Corrected-aero force reevaluation와 delta-projection latency.
  \item Structural-delta construction과 reduced Global corrector solve latency.
  \item 위 항목의 합, 전체 dynamics latency에서의 비율과 iteration 수.
  \item Base-force double-count test.
\end{itemize}

\subsection*{결정 규칙}

\killgate Feedback 제거가 모든 benchmark에서 영향을 거의 주지 않으면
closed-loop geometry-to-load feedback을 contribution으로 주장하지 않고
Local visual/kinematic corrector로 단순화한다.
Multiple corrector가 one-corrector보다 이득이 미미하면 한 번만 사용한다.

\section{E8: End-to-end static-3DGS comparison}

\subsection*{검증 claim}

SYS: Static GS 입력에서 explicit target triangle-mesh reconstruction과
runtime full-order continuum solver 없이
유용한 dynamics와 rendering Pareto를 제공하는가?
여기서 full-order는 full-space shell/FEM/MPM/particle/grid solve를 뜻하며,
reduced Global/Local physics integration이 없다는 뜻이 아니다.

\subsection*{직접 또는 matched 비교 우선순위}

\begin{enumerate}[leftmargin=1.7em]
  \item Same-anchor sparse PD/XPBD \cite{bouaziz2014projective,macklin2016xpbd}.
  \item Gaussian Swaying-style surface aero + matched MPM.
  \item FreeForm-style point/GS ROM.
  \item MeshGraphNets-style direct rollout.
  \item PGRD-style all-node residual.
  \item Proposed Global-only, always-on Local, Oracle Top-\(K\),
  cached-compliance Top-\(K\), full-counter analysis와 조건부 learned Top-\(K\).
\end{enumerate}

\subsection*{조건부 subset 비교}

\begin{itemize}[leftmargin=1.5em]
  \item DynamicTree: tree/leaf/branch-like subset.
  \item DiffWind: 공통 spatial wind forward-simulation subset.
  \item PhysSkin: 동일 reduced deformation/external load subset.
  \item PhysGaussian: 공통 force/material 설정 subset.
  \item GausSim/NGFF: 입력과 action/force를 공통화할 수 있을 때만.
\end{itemize}

\subsection*{Main table 항목}

\begin{itemize}[leftmargin=1.5em]
  \item Runtime explicit target triangle-mesh reconstruction 필요 여부.
  \item Runtime full-order shell/FEM/MPM/particle/grid solver 필요 여부.
  \item Dynamics position/normal/frequency error.
  \item Rendering PSNR/SSIM/LPIPS와 temporal error.
  \item Setup time, dynamics와 end-to-end p50/p95/p99,
  transport/render time 및 cold/warm factor-cache 상태.
  \item Peak GPU memory.
  \item Unseen wind/object와 declared-V0-family material 결과.
\end{itemize}

\section{E9: Real static-GS transfer}

\subsection*{Capture}

최소 3개 real asset에 대해 canonical multi-view capture와 static 3DGS를 만든다.
가능하다면 fan/wind generator의 속도 profile,
고정점, object scale, material mass를 기록한다.
동적 ground truth는 별도 high-speed or multi-view video에서 얻는다.

\subsection*{정량 가능 항목}

\begin{itemize}[leftmargin=1.5em]
  \item Tracked marker/feature trajectory.
  \item Silhouette IoU.
  \item Tip amplitude와 dominant frequency.
  \item Wind-stop damping envelope.
  \item Novel-view temporal consistency.
\end{itemize}

정확한 force/material ground truth가 없으면
synthetic physics table과 real qualitative table을 분리한다.

\section{E10: Quality--latency와 scaling}

\subsection*{조절 변수}

\begin{itemize}[leftmargin=1.5em]
  \item Global rank \(r_g\): 예를 들어 8/16/32/64.
  \item Aerodynamic backend: full anchor가 기본이며, profiling gate를 통과한 경우에만 sample 수.
  \item Local basis rank.
  \item Patch core/halo 크기.
  \item Fixed Local cardinality \(K_A\)와 active/solve-set ratio.
  \item Selector score 및 hysteresis threshold; hardware cost model은 조건부 후속 항목.
  \item Corrector iteration 수.
  \item Gaussian 수와 object batch 수.
\end{itemize}

Gaussian-resolution scaling에서는 \(K,P,r_g,r_\ell^{\mathcal S},K_A\), wind와 timestep을 고정하고
동일 asset의 controlled pruning/densification으로 \(N_G\)만 최소 네 단계 변화시킨다.

\subsection*{보고할 curve}

\begin{itemize}[leftmargin=1.5em]
  \item Object-scale normalized velocity RMSE 대 p95 dynamics latency를 Gate B/C primary curve로 사용.
  \item Position trajectory error와 PSD/spectrum error 대 p95 latency를 공동 필수 curve로 사용.
  \item Quality--p95 Pareto에서 Global rank \(r_g\)와 Local cardinality \(K_A\)를 marker/color로 함께 표시.
  \item Generalized-force/wrench error 대 aerodynamic sample 수.
  \item \(K_A,r_A,r_{\mathcal S},n_A,n_{\mathcal S}\), active/solve-set ratio와
  set churn 대 error와 measured latency.
  \item Topology modal error 대 downstream rollout error를 asset/load별로 층화한 scatter.
  \item LPIPS 대 end-to-end latency.
  \item \(N_G\) 대 분리 측정한 dynamics, Gaussian transport, rendering, end-to-end latency와 memory.
\end{itemize}

\subsection*{Latency breakdown}

\begin{enumerate}[leftmargin=1.7em]
  \item Surface reconstruction and wind query.
  \item Full-anchor aerodynamic force/projection과 조건부 hyper-reduction.
  \item Global predictor.
  \item Cached-compliance selector, low-rate/full counter-state analysis와 조건부 learned gate.
  \item Physics-only Local force와 조건부 missing-force network.
  \item Patch-copy condensation/unique-anchor scatter와 external support-restricted route.
  \item 필요한 structural/reaction channel의 joint constraint operator update,
  Schur factor/dual solve와 primal backprojection.
  \item Local integrator.
  \item Corrected active/halo geometry, normal와 area update;
  optional nonlinear sample-map일 때만 current Jacobian update.
  \item Corrected-aero force reevaluation와 predictor-delta projection.
  \item Structural cross-delta construction.
  \item Reduced Global corrector solve.
  \item Final-state hard reaction과 free-space ROM residual recovery.
  \item Gaussian transport.
  \item Rendering.
\end{enumerate}

Warm-up 후 GPU synchronization을 사용하고 p50/p95/p99를 보고한다.
Batch 1과 multi-object batch, peak memory, kernel 수를 함께 기록한다.
Resolution scaling은 system evidence이며 독립적인 방법 신규성으로 계산하지 않는다.
고정 physical budget에서도 \(T_{\mathrm{dyn}}\)이 \(N_G\)에 따라 증가하면
resolution-decoupling 문구를 철회한다.

\section{전체 ablation registry}

\begin{longtable}{L{0.20\linewidth}L{0.37\linewidth}L{0.12\linewidth}L{0.22\linewidth}}
\toprule
분류 & Ablation & Claim & 예상 실패 신호 \\
\midrule
\endfirsthead
\toprule
분류 & Ablation & Claim & 예상 실패 신호 \\
\midrule
\endhead
Topology & Oracle / proposed / kNN / FPS & MC1 & Layer shortcut, mode error, rollout drift \\
\addlinespace
Topology & Oriented normal 제거 & MC1,K0 & Front/back force error, normal flip \\
\addlinespace
Topology & Bending / stretching 구분 제거 & MC1,K0 & Frequency / material control error \\
\addlinespace
Aero & Current 대신 rest normal/area & K0 & Large deformation force error \\
\addlinespace
Aero & Relative 대신 absolute wind & K0 & Phase/amplitude error \\
\addlinespace
Aero & Spatial 대신 object-average wind & K0 & Traveling gust failure \\
\addlinespace
Sampling & Full, random, FPS, nonconservative, cubature & O1,SYS & Generalized-load/wrench / latency tradeoff \\
\addlinespace
Global & Physics, hybrid, direct, recurrent & K0 & OOD / energy / stability difference \\
\addlinespace
Global & Mode count sweep & K0, MC2, SYS & Local necessity and Pareto \\
\addlinespace
Local & None, full, reduced physics, 조건부 hybrid & MC2,O2 & Flutter / fold recovery \\
\addlinespace
Local & Physics-only 대 force residual 대 displacement residual & O2 & Wind-stop/phase stability \\
\addlinespace
Local & Temporal memory 제거 & O2 & Unsteady/history subset failure \\
\addlinespace
Selector & Never, always, random, cached compliance, full-counter analysis, oracle, 조건부 learned Top-\(K\) & MC2,O3 & Benefit regret and compute \\
\addlinespace
Selector & Cached compliance 대 full-counter/future-benefit score & O3 & Late activation, selector overhead and drift \\
\addlinespace
Selector & Hysteresis/cross-fade 제거 & SYS & Chattering and jerk \\
\addlinespace
Coupling & Complement 제거 & SYS & Amplitude double counting \\
\addlinespace
Coupling & Conservative overlap 제거 & SYS & Seam, net torque drift \\
\addlinespace
Coupling & Structural reaction 제거 & SYS & Global response mismatch if cross block nonzero \\
\addlinespace
Coupling & Corrected aero 제거 & SYS & Local normal change ignored \\
\addlinespace
Coupling & Corrector 제거 & SYS & One-way feedback gap \\
\addlinespace
Transport & Center-only, rigid, affine & SYS & Covariance / orientation artifact \\
\bottomrule
\end{longtable}

\section{Generalization과 stress-test protocol}

\subsection{Generalization matrix}

\begin{longtable}{L{0.22\linewidth}L{0.31\linewidth}L{0.39\linewidth}}
\toprule
Split & 학습에 없는 조건 & 검증 항목 \\
\midrule
\endfirsthead
\toprule
Split & 학습에 없는 조건 & 검증 항목 \\
\midrule
\endhead
Direction & Unseen yaw/elevation & Directional aero and Global response \\
\addlinespace
Magnitude & Weak/strong extrapolation & Force scale, stability, gate recall \\
\addlinespace
Temporal & Unseen gust frequency/chirp & Resonance, phase, history \\
\addlinespace
Spatial & Moving/local gust & Hyper-reduction and active patch migration \\
\addlinespace
Material & V0 \(\nu=0\) family의 stiffness/damping/density & Physical editability; nonzero \(\nu\)/anisotropy 제외 \\
\addlinespace
Attachment & Held-out support layout의 setup-time package rebuild & Topology/operator/control transfer \\
\addlinespace
Geometry & New shape within category & Object generalization \\
\addlinespace
Category & Flag--curtain--paper-like sheet--leaf 사이 bounded stress test & Cross-family limit; arbitrary zero-shot claim 없음 \\
\addlinespace
Reconstruction & New GS fit/view count/seed & Target robustness \\
\addlinespace
Corruption & Floaters/holes/density change & Topology and binding robustness \\
\addlinespace
Combined OOD & New object+material+wind & Practical worst case \\
\bottomrule
\end{longtable}

\subsection{Stability stress tests}

\begin{longtable}{L{0.24\linewidth}L{0.33\linewidth}L{0.35\linewidth}}
\toprule
입력 & 확인할 현상 & 주요 지표 \\
\midrule
\endfirsthead
\toprule
입력 & 확인할 현상 & 주요 지표 \\
\midrule
\endhead
Zero-wind relaxation & Rest 복귀와 energy decay & Energy slope, residual displacement \\
\addlinespace
Abrupt start & 순간 가속/activation & Peak acceleration, jerk \\
\addlinespace
Abrupt stop & 복원과 damping & Envelope/frequency error \\
\addlinespace
Threshold-crossing wind & Gate flicker & Toggle count, jerk \\
\addlinespace
Constant strong wind & Steady drift/divergence & Energy growth, inversion \\
\addlinespace
Sinusoidal/chirp & Resonance/phase & Amplitude, dominant frequency \\
\addlinespace
Traveling gust & Active patch 이동 & Gate recall, spatial continuity \\
\addlinespace
30--60초 rollout & 누적 drift & Divergence rate, state norm \\
\addlinespace
\(\Delta t/2,2\Delta t\) & Integrator sensitivity & Trajectory/frequency change \\
\addlinespace
Strong OOD wind & Graceful failure & Covariance SPD, inversion, gate saturation \\
\bottomrule
\end{longtable}

\section{평가 지표}

\begin{longtable}{L{0.20\linewidth}L{0.31\linewidth}L{0.41\linewidth}}
\toprule
범주 & 지표 & 의미 \\
\midrule
\endfirsthead
\toprule
범주 & 지표 & 의미 \\
\midrule
\endhead
Geometry & Object-scale normalized position RMSE & 전체 deformation error \\
\addlinespace
State & Velocity/acceleration RMSE & Dynamic state error \\
\addlinespace
Surface & Normal angular error & Aero/render direction accuracy \\
\addlinespace
Surface & Strain/curvature error & Thin-shell response \\
\addlinespace
Dynamics & Tip amplitude/phase & Large response \\
\addlinespace
Dynamics & Dominant frequency/damping ratio & Oscillation and recovery \\
\addlinespace
Spectrum & PSD distance & Flutter/high-frequency detail \\
\addlinespace
Aero & Sample/generalized-force error & Aerodynamic operator accuracy \\
\addlinespace
Conservation & Net force/torque error & Cubature/overlap validity \\
\addlinespace
Energy & Work--energy residual/max growth & Nonphysical energy injection \\
\addlinespace
Coupling & Global-span leakage & Double-counted motion/force \\
\addlinespace
Patch & Boundary position/velocity jump & Seam artifact \\
\addlinespace
Selector & Fixed-\(K\) error/oracle regret; 조건부 AUPRC/calibration & Selection quality \\
\addlinespace
Stability & Divergence/inversion rate & Long rollout robustness \\
\addlinespace
Gaussian & Invalid covariance/scale explosion & Render representation stability \\
\addlinespace
Gaussian coverage & Hole/overlap count와 temporal popping & Transport/binding coverage stability \\
\addlinespace
Rendering & PSNR/SSIM/LPIPS & Image fidelity \\
\addlinespace
Temporal render & Flow/feature warp error & Flicker and coherence \\
\addlinespace
Runtime & p50/p95/p99, memory, kernel count, cache hit/miss & Measured cost; interactive status는 조건부 \\
\bottomrule
\end{longtable}

Rendered PSNR만으로 physics claim을 지지하지 않는다.
Main table에는 trajectory, frequency/damping, force/wrench, latency가 함께 있어야 한다.

\section{통계와 reporting}

\begin{itemize}[leftmargin=1.5em]
  \item Frame을 독립 sample로 취급하지 않는다.
  Sequence별 metric을 먼저 구하고 asset 단위로 집계한다.
  \item Asset bootstrap 95\% confidence interval을 보고한다.
  \item Learned baseline은 가능하면 최소 3개 seed를 사용한다.
  \item 모든 method에 동일 test trajectories와 initial conditions를 사용한다.
  \item Failure/divergence sequence를 평균에서 제외하지 않고 별도 비율과 capped metric으로 보고한다.
  \item 공식 구현이 아닌 style baseline은 implementation 차이를 supplementary에 기록한다.
  \item Runtime은 동일 hardware, precision, warm-up, synchronization 조건으로 측정한다.
\end{itemize}

\section{개발용 반증 기준과 go/no-go}
\label{sec:development-go-no-go}

다음 수치는 논문 약속이 아니라 구현 중 방향을 유지할지 판단하는 임시 기준이다.

\begin{longtable}{L{0.31\linewidth}L{0.28\linewidth}L{0.33\linewidth}}
\toprule
질문 & 임시 통과 목표 & 실패 시 결정 \\
\midrule
\endfirsthead
\toprule
질문 & 임시 통과 목표 & 실패 시 결정 \\
\midrule
\endhead
Predicted topology가 충분한가? & Oracle dynamics error 대비 추가 악화 약 10--15\% 이내 & Oracle/simple topology scope 또는 target mesh baseline으로 축소 \\
\addlinespace
Local이 필요한가? & 동일 p95의 Global rank sweep보다 velocity와 local spectrum/trajectory 개선 & Local contribution 제거/축소 \\
\addlinespace
Missing-force network가 필요한가? & Physics-only Local 대비 반복 가능한 개선 & Network 제거, pure adaptive physics \\
\addlinespace
Learned gate가 필요한가? & Cached-compliance보다 같은 \(K_A\)/measured p95에서 낮은 error & 조건부 G module과 O3 claim 제거 \\
\addlinespace
Adaptive Pareto가 있는가? & Always-on 대비 개발용 기준으로 Local cost 약 2배 감소, dynamics error 증가 약 5\% 내외 & Adaptive claim 철회 \\
\addlinespace
End-to-end speedup이 있는가? & Always-on 대비 최소 약 1.25배 또는 same-quality 유의미한 개선 & Dispatcher/always-on/simple solver 전환 \\
\addlinespace
Global dynamics가 맞는가? & Tip frequency 약 5\%, amplitude 약 10\% 이내 & Basis/operator/aero 재설계 \\
\addlinespace
Long rollout이 안정적인가? & 4--10초 normalized position error 약 object scale 5\% 내외, divergence 없음 & Integrator/selector를 수정하고 조건부 learned branch는 별도 재학습 또는 제거 \\
\addlinespace
Conservative sampling이 필요한가? & Full backend보다 generalized-load Pareto 개선 & H 제거, full-anchor backend 확정 \\
\addlinespace
Feedback이 필요한가? & One-way 대비 measurable Global response/teacher error 개선 & Feedback claim 제거 \\
\bottomrule
\end{longtable}

약 \(2\times\), 10--15\%, 5\%와 1.25\(\times\)는 모두 구현 자원 배분을 위한
내부 threshold이지 venue가 요구하는 보편적 수치가 아니다.
최종 판정은 동일 hardware/input privilege에서 confidence interval을 가진
quality--latency Pareto, actual DOF/rank, active/decay set churn과 multi-object throughput으로 내린다.
또한 ``same scaffold''는 graph, material, attachment, load와 timestep을 맞춘다는 뜻이지
ROM과 full-space solver의 simulated DOF가 같다는 뜻이 아니며, 실제 DOF와 rank를 별도로 보고한다.
Global coordinate와 Local coordinate는 support와 실행 비용이 다르므로 동일 coordinate count를
공정성의 충분조건으로 사용하지 않는다. Mode/rank-matched 결과는 진단용이고,
primary 판정은 same-quality speedup, same-p95 velocity/trajectory/spectrum error 또는
multi-object throughput 중 하나가 만드는 Pareto non-dominated region이다.
모든 operating point에서 pointwise strict dominance를 요구하지 않는다.

\reviewvthreerebut{V3 리뷰가 strong Global rank sweep, always-on Local,
same-anchor sparse solver와 direct neural rollout을 모두 포함하라는 요구는 수용한다.
그러나 ``모든 baseline을 이겨야 한다''를 모든 operating point에서 더 빠르고 더 정확해야 한다는
pointwise strict dominance로 해석하지 않는다.
제안 frontier가 이 strong-baseline union의 Pareto envelope에 confidence interval을 고려한
non-dominated region을 추가해야 한다.
반대로 하나의 충분히 최적화된 matched baseline이 제안 frontier 전체를 지배하면 MC2를 기각한다.}

\section{논문 결과표와 figure 계획}

\begin{enumerate}[leftmargin=1.7em]
  \item \textbf{Figure 1:} Static GS \(\rightarrow\) topology/operator package
  \(\rightarrow\) Global predictor \(\rightarrow\) active Local patches
  \(\rightarrow\) corrected Gaussians.
  \item \textbf{Table 1:} Oracle-topology matched solver accuracy/latency.
  \item \textbf{Table 2:} Static-GS end-to-end comparison과 input privilege.
  \item \textbf{Table 3:} Full-anchor aero profile과, profiling gate를 통과한 경우에만
  random/FPS/RSH/cubature 비교.
  \item \textbf{Table 4:} Local physics/selector/coupling 및 조건부 missing-force/learned-gate ablations.
  \item \textbf{Table 5:} OOD와 long-rollout stability.
  \item \textbf{Figure 2:} Current-state aero와 rest-state WPB-style force 차이.
  \item \textbf{Figure 3:} Teacher/Global-only/physics-only Local/hybrid Local의 time response와 PSD.
  \item \textbf{Figure 4:} Oracle benefit, cached-compliance와 full-counter analysis score,
  actual Top-\(K\) set과
  조건부 learned score map.
  \item \textbf{Figure 5:} Complement/overlap conservation 제거 시 seam, drift, torque error.
  \item \textbf{Figure 6:} Quality--latency Pareto, active/solve-set ratio와 set-churn break-even,
  고정 physical budget의 \(N_G\) scaling breakdown.
  \item \textbf{Figure 7:} Oracle topology 대 predicted topology downstream motion과
  modal error--rollout error의 asset/load-stratified 관계.
  \item \textbf{Figure 8:} Setup별 material/attachment metadata를 가진 real static-GS wind response.
\end{enumerate}

\section{주요 위험과 실패 모드}

\subsection{``Without explicit target triangle-mesh reconstruction'' 표현이 사실상 mesh가 되는 위험}

Oriented anchor graph가 watertight triangle connectivity와 full shell elements를 요구하면
reviewer는 target mesh를 다른 이름으로 부른 것이라고 볼 수 있다.
Runtime package는 sparse, non-watertight, locally oriented material graph로 정의하고,
global face connectivity와 watertight/manifold 조건을 요구하지 않으며,
mesh extraction baseline과 representation/storage/solver 차이를 명시해야 한다.
식~\eqref{eq:relation-mls-map}의 quadratic MLS relation은 bounded local support와
oriented material coordinates만 사용하며 face cycle이나 paired triangle hinge를 요구하지 않는다.
그 rank/conditioning을 만족하려고 사실상 전역 face connectivity를 복원한다면 main claim에서 제외한다.
반대로 stable area-normal과 bending operator를 위해 실제 triangles가 필수라면
``without explicit target triangle-mesh reconstruction'' claim을 약화하고
lightweight inferred proxy로 정직하게 표현한다.

\subsection{Topology distillation이 dynamics bottleneck이 되는 위험}

Close disconnected layers, thin tips, missing backside, floaters는
Euclidean graph와 orientation 추론을 쉽게 깨뜨린다.
Topology confidence가 낮을 때에는

\begin{itemize}[leftmargin=1.5em]
  \item uncertain edges 제거,
  \item user hint 또는 category prior 요청,
  \item conservative stiffness,
  \item target mesh extraction fallback
\end{itemize}

을 고려한다.
Fallback을 사용한 sequence는 main no-explicit-target-triangle-reconstruction 결과와 구분한다.

\subsection{Metric scale와 material identifiability}

일반 SfM/3DGS는 절대 scale이 없고,
wind speed, density, thickness, Young's modulus는 서로 보상될 수 있다.
User-supplied scale marker, known object dimension,
또는 nondimensional response preset 중 하나를 runtime contract에 넣는다.
실측 물성 추정과 artist-friendly control은 별도 평가한다.

\subsection{Global basis만으로 충분할 위험}

Global rank를 조금 늘리면 local flutter까지 충분히 재현될 수 있다.
이 경우 Local/gate/coupling 복잡성은 정당화되지 않는다.
E4에서 mode-count Pareto를 가장 먼저 수행하고,
Local이 같은 latency에서 분명한 high-frequency 이득을 보일 때만 유지한다.

\subsection{Local physics-only로 충분할 위험}

Teacher-S가 runtime과 같은 aerodynamic law와 structural law를 쓰면
learned missing force가 학습할 defect가 작을 수 있다.
Network가 이득을 주지 않으면 pure adaptive local physics가 더 깨끗한 논문이 된다.
Unsteady force를 인위적으로 만들기 위해 teacher와 runtime physics를 불필요하게 다르게 하지 않는다.

\subsection{Selector 계산이 Local skip 이득을 먹는 위험}

Every-frame auxiliary Global solve, predictor geometry 재구성과 aero/internal-force 재평가는
skip하려는 Local block보다 비쌀 수 있다는 지적을 수용한다.
따라서 V0 mainline은 이미 계산한 owner-resolved force의 complement projection과
setup에서 factorize한 small-group compliance solve만 사용한다.
Full counter-state residual은 offline 또는 저주기 analysis ceiling으로 내리고,
추가 surface/force evaluation을 포함한 비용을 숨기지 않는다.
반대로 cached score는 predictor 뒤 geometry 변화, group 밖 coupling과 갑자기 나타나는
localized gust를 놓칠 수 있으므로 Oracle regret, full-counter rank correlation,
late-activation과 selector-only p50/p95/p99를 함께 측정한다.
식~\eqref{eq:normalized-hybrid-selector}의 deterministic hybrid는 이 blind spot을 검사하는
조건부 ablation이며, raw heterogeneous-unit 합이나 무료 internal residual을 허용하지 않는다.
Cheap score가 실패하고 full counter-state도 end-to-end Pareto를 개선하지 못하면
learned gate 또는 budgeted probe를 조건부로 검토하되 residual 계산을 무료라고 가정하지 않는다.

\subsection{조건부 learned gate 비용이 절감량을 먹는 위험}

모든 patch에 heavy GNN gate를 실행하면 Local skip 이득이 사라진다.
Gate는 pooled Global state와 cheap local invariants를 사용하는 작은 model이어야 한다.
Gate latency, feature construction, gather/scatter를 별도 profile한다.

\subsection{Active principal-block factorization이 sparse 이득을 먹는 위험}

Fixed-cardinality \(K_A\)라도 실제 solve 집합은 \(\mathcal S_t=\mathcal A_t\cup\mathcal D_t\)이고,
active/decay 조합이 바뀔 때 exact principal matrix의 factorization이 필요할 수 있다는 지적을 수용한다.
Traveling gust, 긴 decay와 set churn에서는 factor-cache miss와 gather/scatter가
Local 계산 절감량을 상쇄할 수 있다.
V0는 exact active direct solve와 always-on full-superset reference를 먼저 profile하고,
반복 signature cache, 검증된 incremental Cholesky, cross-block matvec를 유지하는
block-preconditioned PCG를 순서대로 비교한다.

반면 patch/group을 독립적으로 적분하고 Global corrector가 누락 coupling을 대신하게 하자는 처방은
일반적으로 성립하지 않으므로 수용하지 않는다.
누락되는 항은 식~\eqref{eq:off-group-coupling}의 Local--Local cross interaction이고,
same-step effective Local coupling load는 일반적으로
\(\operatorname{range}(B_{\mathcal S}^{\mathrm{eff}})\)에 있지 않으며,
Local state를 다시 풀지 않는 Global corrector는 이를 복원할 수 없다.
독립 group solve는 off-group coupling과 rollout/energy error가 모두 작은 경우의 근사 ablation일 뿐이며,
group factor는 기본적으로 coupled iterative solve의 preconditioner로 사용한다.
어떤 backend도 p95/p99를 낮추지 못하면 adaptive-compute headline을 축소하거나
exact direct/always-on backend로 돌아간다.

\subsection{Fixed-superset conditioning과 cross-patch energy 귀속}

V0는 patch별 basis slot을 고정한 \(\Psi_{\mathrm{sup}}=[\Psi_1,\ldots,\Psi_P]\)를 사용하므로
activation 때 basis dimension이나 좌표 의미를 바꾸지 않고 state remapping도 하지 않는다.
대신 강하게 겹치는 patch의 cross block이 ill-conditioned할 수 있다.
Offline rank 검사와 overlap group 병합을 수행하고,
energy는 독립 patch coordinate의 단순 합이 아니라 support/owner-weighted energy 또는
group의 full kinetic energy와 식~\eqref{eq:projective-structural-energy}로 보고한다.
Dynamic basis와 energy-consistent remapping은 V0 profiling 이후의 후속 최적화다.

\subsection{조건부 missing-force inverse dynamics noise}

Position trajectory를 두 번 미분한 acceleration은 noisy하다.
가능하면 teacher force breakdown을 직접 저장하고,
그렇지 않으면 smoothing, weak-form/integral loss,
differentiable integrator를 통한 multi-step state supervision을 같이 쓴다.

\subsection{Anchor-primal joint KKT 비용과 conservation--complement 충돌}

Wrench constraint와 Global complement를 순차 projection하면
한 constraint가 다른 constraint를 다시 깨뜨릴 수 있다.
External channel은 식~\eqref{eq:support-feasible-witness}의 support-restricted route가
두 조건을 이미 동시에 만족하므로 동일 constraint를 다시 KKT로 풀지 않는다.
Structural/hard-reaction channel과 추가 independent row가 있는 경우에는
식~\eqref{eq:conservative-assembly}에서 wrench와 \(\Phi^\top f=0\)를 하나의 augmented
KKT system에 동시에 넣는다. Constraint feasibility와 rank를 signature마다 검사한다.
다만 anchor-force primal 전체의 generic saddle KKT를 매 frame 직접 푸는 것은 금지한다.
Production backend는 unique-anchor condensation 뒤 식~\eqref{eq:runtime-joint-schur}의
small Schur, 동등한 null-space solve 또는 intrinsic edge-force parameterization만 허용한다.
Constraint row 삭제, torque 생략과 순차 projection은 최적화가 아니라 method ablation이다.
동등한 joint backend를 최적화해도 assembly latency가 Local skip 이득을 지우면
actual-skip/SYS headline을 축소하며 conservation--complement 계약 자체를 깨지 않는다.

\subsection{Feedback이 수치적으로 불안정할 위험}

Local correction 후 aero를 다시 계산하면 follower-force feedback이 커질 수 있다.
One-corrector, under-relaxation, force-delta clamp,
smaller timestep을 stress test한다.
Full fixed-point convergence를 보장하지 않는다면
method를 predictor--one-corrector approximation으로 명시한다.

\subsection{Teacher와 real wind의 domain gap}

Quasi-steady teacher만으로 real cloth flutter의 phase와 spectrum을 맞추기 어렵다.
Real sequence는 system identification 또는 limited Teacher-FSI fine-tuning이 필요할 수 있다.
Synthetic physics accuracy와 real visual plausibility를 같은 claim으로 섞지 않는다.

\subsection{Gaussian transport와 physical surface 불일치}

Visual Gaussian support가 실제 material surface를 벗어나거나
두 겹 surface를 동시에 표현하면 동일 \(F\) transport가 artifact를 만든다.
Surface reliability, layer separation, covariance clamp,
outlier Gaussian fallback을 둔다.

\section{P0 implementation kill gates}

\begin{longtable}{L{0.06\linewidth}L{0.45\linewidth}L{0.41\linewidth}}
\toprule
\# & 중단 조건 & 필요한 조치 \\
\midrule
\endfirsthead
\toprule
\# & 중단 조건 & 필요한 조치 \\
\midrule
\endhead
1 & \(\|\Phi^\top M\Psi\|\)가 tolerance를 넘는다. & Basis construction/complement를 고치기 전 Local training 금지. \\
\addlinespace
2 & \(M\succ0,K^0,C^0\succeq0\) 또는 block symmetry가 깨진다. & Runtime solve 금지; operator construction 수정. \\
\addlinespace
3 & Gravity, attachment, aero, restoring, residual 중 owner가 둘 이상이다. & Force ledger를 다시 만들고 residual label 재생성. \\
\addlinespace
4 & Gate-off/zero-wind에서 rest로 돌아오지 않거나 Local energy가 증가한다. & Activation equation, damping, state decay 수정. \\
\addlinespace
5 & Patch별 integrate-then-average가 구현되어 있다. & Force assembly 후 단일 Local update로 변경. \\
\addlinespace
6 & 조립 후 Global leakage, overlap seam, support-route/Schur feasibility,
adjustment/reaction amplification/generalized-work 보존 또는
free/hard/compliant reaction-wrench ledger가 tolerance를 넘는다. & Joint constraint solve,
support/route 또는 basis/patch를 수정한다. Ridge, row 삭제나 sequential projection으로 통과 처리하지 않는다. \\
\addlinespace
7 & Global corrector가 corrected-minus-predictor delta 외 base force를 다시 사용한다. & Double-counting 제거. \\
\addlinespace
8 & Inactive 전환 시 \(z,\dot z\)를 즉시 freeze/reset한다. & Decay state와 energy threshold 추가. \\
\addlinespace
9 & Rigid/affine transport, covariance SPD, aero/render normal consistency가 실패한다. & Gaussian 결과 생성 금지; transport 수정. \\
\addlinespace
10 & Oracle Local이 Global-only를 개선하지 못한다. & Local branch 제거 또는 Global basis 재설계. \\
\addlinespace
11 & Always-on missing-force가 physics-only Local을 개선하지 못한다. & Network만 제거하고 pure adaptive Local physics를 유지한다. \\
\addlinespace
12 & Learned gate가 cached-compliance Top-\(K\)보다 같은 \(K_A\)/measured p95에서 개선하지 못한다. & Learned gate만 제거하고 cached-compliance/Oracle selector를 유지한다. \\
\addlinespace
13 & Same-anchor sparse PD/XPBD 대비 개발용 기준인 dynamics-only 약 \(2\times\) speedup도,
동일 p95 latency의 velocity/trajectory/spectrum 개선도, multi-object throughput 이득도 없다.
& SYS selective-compute/quality--latency headline을 철회하고 always-on Local 또는 same-scaffold simple solver를 검토한다.
이 실패만으로 MC1 representation이나 Gate B의 Local accuracy/necessity를 기각하지 않는다. \\
\addlinespace
14 & Full-anchor aero가 target hardware의 유의미한 병목이 아니거나 cubature가
localized/traveling gust Pareto를 개선하지 못한다. & Hyper-reduction branch를 제거하고 full-anchor backend를 확정한다. \\
\addlinespace
15 & 고정 physical budget에서 \(N_G\) 증가와 함께 structural dynamics 비용이 유의미하게 증가한다. & Resolution-decoupling system evidence를 철회하고 transport/render 병목과 binding을 재설계한다. \\
\addlinespace
16 & Selector, joint assembly/Schur, active/decay gather 또는 factorization 때문에 adaptive p95/p99가 always-on Local보다 나쁘다. & Backend/cache를 재검토하고, 개선되지 않으면 dispatcher 또는 always-on을 채택하며 actual-skip headline을 철회한다. \\
\addlinespace
17 & Incremental/PCG/independent backend가 same-set exact direct 대비 linear residual, state 또는 energy tolerance를 넘는다. & 해당 근사 backend만 제거하고 exact active direct를 유지한다. \\
\addlinespace
18 & Stationary/slow/fast/reversing gust에서 decay 누적으로 \(\rho_{\mathcal S}\simeq1\)이 지속되거나 cache churn/memory cap으로 latency 이득이 사라진다. & Physical decay를 끄거나 state를 reset하지 않고 decay/dispatcher를 재설계한다. 개선되지 않으면 always-on으로 후퇴하고 actual-skip headline을 철회한다. \\
\addlinespace
19 & Production trace가 full anchor-primal generic KKT를 풀거나, duplicate overlap/velocity row로 constraint 수가 patch 수에 비례하거나, 잘못된 factor reuse/sequential fallback을 사용하거나, signature별 support/halo budget을 반복 초과한다. & Unique-anchor support-restricted routing과 동등한 small-Schur/null-space backend로 고친다. Joint 계약을 유지해도 latency 이득이 없으면 actual-skip/SYS headline을 축소한다. \\
\addlinespace
20 & 식~\eqref{eq:canonical-structural-force-ledger}의 pure-rigid zero energy,
strained-state superposed-rotation energy/force equivariance, objective relative-rate damping,
planar nonuniform in-plane probe의 bending energy/dissipation,
midpoint phase/PSD 또는 one-shot/two-step/converged projective rollout의 torque/phase tolerance가 실패한다.
& Track A0에 진입하지 않고 projective target/operator contract를 재설계한다.
Asset별 backend 변경이나 backward Euler의 수치 감쇠로 실패를 숨기지 않는다. \\
\addlinespace
21 & Normalized relation share는 존재하지만 식~\eqref{eq:relation-absolute-confidence}의
calibrated absolute confidence가 낮거나 candidate-count/conditioning별 calibration이 없는 package가 통과한다.
& Support 확대 후 재검사하고, 실패하면 user hint/fallback/reject한다. Confidence를 stiffness에 곱해 통과시키지 않는다. \\
\addlinespace
22 & 식~\eqref{eq:klin-projective-hard-gate}의 energy, stiffness/damping force action,
action-tangent symmetry/PSD, modal angle 또는 common-probe tolerance가 실패한다.
& Global basis/projective backend를 재설계한다. Learned residual이나 raw \(K_{\mathrm{proj}}^0\) mode 비교로 우회하지 않는다. \\
\addlinespace
23 & Canonical V0가 식~\eqref{eq:fixed-global-sample-jacobian} 대신 state-dependent Jacobian을
암묵적으로 만들거나 corrector가 Jacobian update 비용을 숨긴다.
& V0를 fixed lifting으로 고치고 nonlinear sample map은 별도 versioned optional backend/profiler로 분리한다. \\
\addlinespace
24 & Mandatory selector blind-spot fixture에서 cached compliance가 반복적으로 큰 Oracle-benefit region을 놓치고,
hybrid/full-counter/조건부 learned 대안도 measured-p95 Pareto를 개선하지 못한다.
& Error-triggered/adaptive selector claim을 축소하고 always-on Local, merged group 또는 단순 dispatcher로 후퇴한다. \\
\bottomrule
\end{longtable}

연구-level Gate A는 representation preflight, Gate B는 Oracle Local necessity,
Gate C는 same-anchor quality--latency와 actual-skip overhead를 함께 뜻한다.
세 연구 gate가 위의 force ownership, joint constraint, energy/decay와 transport invariant를 대체하지 않는다.
Gate A 실패는 MC1 representation claim만 축소하며 Gate B/C는 Oracle 또는
명시적 fallback scaffold에서 계속 판정할 수 있다.
Gate B 실패는 Local/selector claim을 제거하지만 MC1과 Global physics를 기각하지 않는다.
Gate C 실패는 selective-compute/SYS performance headline만 제거하며
Gate A와 Gate B의 accuracy 판정을 뒤집지 않는다.

이 표의 숫자는 보편적인 출판 합격선이 아니라 개발 자원을 어디까지 투입할지 정하는 내부 기준이다.
특히 약 \(2\times\)는 solver/runtime 환경에 종속되며, 더 작은 speedup이라도
동일 p95 latency의 velocity/trajectory/spectrum 개선이나 multi-object throughput 이득이 크면 MC2를 유지할 수 있다.
반대로 exact chain이 새롭다는 이유만으로 위 claim-specific failure를 무시하지 않는다.

\section{문서와 구현 계약 관리}

\reviewvfivepartial{111쪽 canonical sketch를 구현자가 매번 순회하기 어렵다는 진단과
짧은 V0 contract의 필요성은 수용한다. 그러나 이 문서를 rationale archive로 내리고
수식을 복제한 20--30쪽 문서를 새로운 독립 source of truth로 만드는 처방은 채택하지 않는다.
두 문서가 같은 수식을 별도로 소유하면 review revision 때 drift와 owner 불일치가 다시 생긴다.}

Canonical derivation과 normative equation의 유일한 source of truth는 이 문서로 유지한다.
짧은 V0 implementation contract는 여기서 파생한 tensor shape, typed API, force owner,
artifact와 kill gate만 담고, 수식을 복제하지 않고 안정적인 \texttt{eq:} 및 section label로 연결한다.
PDF의 숫자 식 번호는 checklist와 코드 주석의 영구 식별자로 사용하지 않는다.
파생 contract/checklist에는 canonical source revision, label 존재 여부와 schema version을 기록해
stale label/API/owner를 검사한다. 대안 backend와 review verdict는 이 canonical rationale에 남기되,
파생 contract에서는 현재 동결한 V0 branch와 failure policy만 노출한다.
이 파생 artifact와 stale-reference test의 구현 owner는 TD01이며, TD00의 완료 상태를 소급해
바꾸지 않는다.

\section{MVP와 구현 순서}

\subsection{Step 0: Structural/operator-load preflight와 V0 backend freeze}

\reviewvfiveaccept{현재 architecture를 더 확장하지 않고 Step-0-relevant P0 subset에서 deterministic
structural/assembly backend를 학습 없이 먼저 반증하라는 최종 판정을 수용한다.
Architecture freeze는 현 backend가 정답이라는 뜻이 아니다. 아래 fixture를 통과하기 전에는
H/F/G, Teacher-FSI, full predicted-topology training 또는 대규모 data generation을 mainline에
추가하지 않으며, 실패하면 식~\eqref{eq:klin-projective-hard-gate}의 원인별 정책에 따라
backend/contract를 교체한다.}

Track A0 전에 다음을 완료한다.

여기서 Step-0 gate는 K00--K09의 deterministic 해당 항목, K20, K22--K23과
K21의 confidence schema/synthetic aggregator/fallback fixture를 뜻한다.
Object-disjoint learned K21 calibration은 Track B0/E1에서, selector K24와 나머지
runtime P0는 해당 downstream stage에서 판정한다. 따라서 ``Step 0/P0 통과''를
아직 학습하지 않은 confidence calibration까지 선행 완료했다는 뜻으로 사용하지 않는다.

\begin{enumerate}[leftmargin=1.7em]
  \item 식~\eqref{eq:relation-mls-map}--\eqref{eq:hard-attachment-free-operators}의
  relation-level \(M,K,C\) assembly, material/metric scaling, attachment 처리와
  rigid-motion nullspace, absolute relation-confidence schema/aggregator/fallback을 고정한다.
  이 단계는 synthetic score로 deterministic 경로만 검사하며 learned calibration은 Track B0/E1에서 수행한다.
  \item 식~\eqref{eq:normal-curvature-one-shot-residual}의 flat-rest bending-pollution fixture와
  planar in-plane, normal-curvature 및 shallow-rest probe를 실행한다.
  \item 식~\eqref{eq:refit-projective-action-closure}--\eqref{eq:klin-projective-hard-gate}의
  rest-linear/locally refit energy, held-target stiffness와 descriptor-rate damping action,
  action tangent symmetry/PSD와 modal closure를 통과한다.
  \item 식~\eqref{eq:v0-structural-freeze-signature}의 one-shot/two-step/converged projective solve를
  pure rigid, strained-state superposed rotation, oscillator, cantilever와 flag에서 비교하고
  torque residual, phase, target-band PSD와 timestep convergence를 함께 판정한다.
  \item 식~\eqref{eq:fixed-global-sample-jacobian}의 fixed anchor/sample lifting을
  V0 Global load와 corrector의 유일한 Jacobian contract로 동결한다.
  \item Singleton, overlap, boundary, hard attachment와 traveling-support signature에서
  support-restricted route, joint Schur feasibility뿐 아니라 raw-force adjustment,
  reaction amplification, generalized-work 보존, post-step ROM residual/reaction ledger와
  halo-expansion rate를 측정하고, development에서 \(\sigma_{\mathrm{sup}}\)별 primary-diagnostic
  cap과 test-time reject/always-on fallback을 동결한다.
\end{enumerate}

Oracle oscillator, cantilever와 flag preflight에서 canonical objective normal-curvature projective
frozen-target backend와 implicit midpoint를 rest-linear, backward-Euler 및
같은 target-refit held action의 two-step 또는 converged fixed-point accuracy reference와
pure-rigid 및 strained-state superposed-rotation objectivity,
first-mode frequency/damping, phase와 target-band spectrum, energy behavior,
최대 허용 timestep, p50/p95/p99와 factor-cache compatibility로 비교한다.
Held-out 결과를 본 뒤 backend를 고르지 않는다. Canonical 조합이 development tolerance에 실패하면
Track A0로 진행하지 않고 structural contract를 재설계한다.
이 fixture가 모두 통과하기 전에는 Track A0의 speed result, full predicted-topology training,
large teacher/data generation 또는 H/F/G 구현에 진입하지 않는다.
단, fixture 자체에 필요한 최소 Oracle strip/flag와 independent-GS probe 생성은 허용한다.

\subsection{Track A0: Oracle-topology Global strip와 same-anchor kill test}

\begin{itemize}[leftmargin=1.5em]
  \item 한쪽 고정 cantilever strip/flag.
  \item Oracle source topology/relation 및 attachment mapping;
  mass/material/backend는 canonical \(\nu=0\) assembly와 맞추고 native teacher operator는 별도 ceiling.
  \item Causal half-step-geometry/frozen-velocity analytic aero with all samples.
  \item \(r_g=8\)--32 Global modes.
  \item Global physics, Gaussian affine transport.
  \item 512/1,024/2,048 anchors에서 same-anchor sparse PD/XPBD/corotational solver.
\end{itemize}

Kill test는 먼저 target asset에 가장 잘 맞고 최적화된 same-anchor solver 하나로 수행한다.
추가 PD/XPBD/corotational variant는 서로 다른 Pareto 영역이나 failure mode를 보일 때 확장하며,
여러 미최적화 baseline 구현 자체를 선행 조건으로 만들지 않는다.

통과 조건은 start/stop, sinusoidal/chirp의 frequency/damping과
same-anchor sparse solver 대비 개발용 기준인 dynamics-only 약 \(2\times\) speedup 또는
동일 p95 latency에서 유의미하게 낮은 velocity/trajectory/spectrum error,
또는 multi-object throughput 이득이다.
통과하지 못하면 SYS quality--latency headline을 철회한다.

\subsection{Track B0: 조기 representation feasibility}

\begin{itemize}[leftmargin=1.5em]
  \item 5--10개의 simple thin-surface object와 독립 재구성 3DGS variant.
  \item Oriented anchor graph, area/mass 보존, close-layer false edge.
  \item 첫 8--16개 eigenfrequency error와 modal-subspace principal angle.
  \item 식~\eqref{eq:patch-complement-tail}의 patch/group tail distribution과 rollout-weighted tail.
  \item 식~\eqref{eq:independent-gs-package-consistency}의
  area/mass, eigenspace, tail과 common-probe rollout consistency.
  \item User-specified material/attachment metadata에서 package를 다시 빌드하는 setup contract.
\end{itemize}

\reviewvfivepartial{Full Track B0와 learned calibration/training은 Step-0 gate 통과 뒤에
시작하라는 자원 배분 원칙을 수용한다. 다만 coordinate/schema, relation rank/valence,
area/mass와 close-layer false edge를 보는 값싼 independent-GS probe까지 금지하면
representation의 조기 kill signal을 잃는다. 이 최소 readout-only probe는 Step 0과 병렬 허용하고,
그 결과를 MC1 acceptance나 calibrated confidence로 주장하지 않는다.}

Full Track B0는 Step-0 gate 통과 뒤 Track A0와 병렬로 수행하며,
실패하면 adaptive system 전체 구현 전에 MC1 표현 주장을 축소한다.

\subsection{Track A1: Oracle active Local Gate B}

\begin{itemize}[leftmargin=1.5em]
  \item Patch-local fixed-superset complement basis와 oracle active patch.
  \item 공식 Oracle Local 비교 전에 unique-anchor conservative joint assembly,
  Global complement와 single coupled Local update를 correctness backend로 먼저 고정.
  \item Physics-only Local과 full-superset full-Local upper bound/always-on latency;
  canonical fixed-Hessian factor cache와 projective target assembly 비용을 분리.
  \item Complement/overlap ablation. One-corrector는 Gate B 뒤의 별도 SYS ablation이며
  Local necessity 판정을 gate하지 않음.
\end{itemize}

Localized/traveling gust와 여러 excitation에서 Oracle Local이
동일 p95의 충분히 강한 Global/POD/wind-aware basis rank sweep보다
velocity/trajectory/spectrum Pareto를 개선하지 못하면 Local 방향을 중단한다.
이 same-p95 Oracle 비교는 식~\eqref{eq:oracle-total-latency-fairness}의
synchronized total-frame cost를 맞춘 Local-capacity ceiling이며,
deployable selector의 speed claim은 Track A2/Gate C에서 별도로 판정한다.

\subsection{Track A2: Cached-compliance selector와 decay Gate C}

\begin{itemize}[leftmargin=1.5em]
  \item Fixed-cardinality \(K_A\) Oracle benefit 대 cached-compliance Top-\(K\)와
  normalized deterministic hybrid Top-\(K\).
  \item Full counter-state residual은 offline 또는 \(J\)-frame analysis baseline으로 실제 추가 비용까지 비교.
  \item 동일 active/decay trace에서 exact active direct를 numerical oracle로 두고,
  bounded factor cache/incremental Cholesky, block-PCG와 independent-group 근사의
  p50/p95/p99, residual, rollout/energy error 비교.
  \item \(\chi_{\mathcal S}^{\mathrm{off}}\) bin별 cross-group gust error,
  cold/warm cache, eviction/downdate fallback, decay 누적과 max \(\rho_{\mathcal S}\) 비교.
  \item Mandatory blind-spot 네 fixture에서 regret, high-benefit miss, activation delay,
  target-band PSD와 selector-only latency를 같은 \(K_A\)/measured-p95로 비교.
\end{itemize}

\subsection{조건부 H: Aerodynamic hyper-reduction}

\begin{itemize}[leftmargin=1.5em]
  \item Full-anchor profile이 병목일 때만 random/FPS/nonconservative/cubature를 구현한다.
  \item Held-out deformation과 localized/traveling gust에서 generalized load를 primary metric으로 본다.
  \item Net wrench는 auxiliary constraint/metric으로 보고 GPU wall-clock까지 비교한다.
\end{itemize}

\subsection{조건부 F/G: Missing force와 learned gate}

\begin{itemize}[leftmargin=1.5em]
  \item Physics-only Local과 full-resolution Local 사이에 gap이 있을 때만 missing-force를 학습한다.
  \item Teacher nodal breakdown 또는 검증된 constrained lifting으로만 patch force label을 만든다.
  \item Oracle benefit과 cached-compliance selector 사이에 gap이 있을 때만 learned future-benefit gate를 학습한다.
  \item 두 모듈은 독립 go/no-go이며, 한쪽 실패가 다른 쪽이나 V0를 중단시키지 않는다.
\end{itemize}

\subsection{Track B1: Predicted target package와 end-to-end transfer}

Track A의 solver baseline 우선순위는 MC1의 representation baseline을 대체하지 않는다.
Oracle/proposed/kNN/FPS와 explicit target-mesh extraction comparison을 모두 유지한다.

\begin{itemize}[leftmargin=1.5em]
  \item Independent static GS variants와 held-out setup asset.
  \item Oriented topology, area weights, setup-conditioned operator distillation.
  \item Oracle/predicted/kNN/FPS/explicit mesh extraction baseline 비교.
  \item Real static-GS qualitative transfer와 rendering-resolution scaling.
\end{itemize}

\subsection{첫 논문 권장 최소 범위}

\begin{itemize}[leftmargin=1.5em]
  \item Core quantitative dynamics: short cantilever strip, rectangular flag,
  triangular pennant와 curtain strip.
  \item Secondary/transfer: hanging cloth, permanent creasing 없는 paper-like elastic sheet와
  compliant stem의 single leaf/petal; 첫 논문 필수 성공 조건이 아님.
  \item Prescribed spatial wind.
  \item Teacher-S 중심; Teacher-FSI는 작은 optional subset.
  \item Sparse oriented point/edge/local-relation scaffold;
  explicit target triangle-mesh reconstruction과 runtime full-order continuum solve는 없음.
  \item Setup에서 고정한 metric scale, user material과 attachment metadata.
  \item Full-anchor Global physics + fixed-superset Local physics + fixed-cardinality cached-compliance/Oracle selector.
  \item Learned missing force, learned gate, aerodynamic cubature는 통과 시에만 포함.
  \item Affine GS transport와 standard rendering.
\end{itemize}

Full tree, self-contact, tearing, free flight, appearance residual,
arbitrary cross-family material generalization은 후속으로 둔다.

\section{논문 framing}

\subsection{권장 제목}

\begin{quote}
\textbf{Topology-Distilled, Error-Triggered Global--Local Wind Dynamics
for Static 3D Gaussian Thin Surfaces}
\end{quote}

이는 working title이며 앞의 error-relevance evidence gate를 통과할 때만 최종 제목으로 유지한다.
Learned future-benefit gate의 채택 여부 자체는 제목 유지 조건이 아니다.

\reviewpartial{Error relevance를 결과에 종속시키라는 진단은 수용한다.
다만 learned gate가 없다는 이유만으로 지금 즉시 \emph{Selective}로 낮추는 처방은 수용하지 않는다.
Deterministic cached-compliance selector도 held-out Oracle rank correlation, regret,
activation delay/high-benefit miss와 quality--latency evidence를 만족하면
\emph{Error-Triggered}를 유지할 수 있다. 그 증거가 실패할 때만
\emph{Compliance-Guided}, \emph{Selective Complementary}, MC1 중심 제목으로 순차 축소한다.
결과 전 canonical 파일명을 반복 변경하지 않는다.}

파일명은 제목의 핵심을 반영해

\begin{center}
\path{3dgs_topology_distilled_error_triggered_wind_dynamics_2026-08-13.tex}
\end{center}

로 정한다.

\subsection{대안 제목}

\begin{enumerate}[leftmargin=1.7em]
  \item \emph{Topology-Distilled, Compliance-Guided Global--Local Wind Dynamics
  for Static 3D Gaussian Thin Surfaces}.
  \item \emph{Topology-Distilled Selective Complementary Local Wind Dynamics
  for Static 3D Gaussian Thin Surfaces}.
  \item \emph{Aeroelastic Animation of Static Gaussian Thin Surfaces
  without Explicit Target Triangle-Mesh Reconstruction}.
  \item \emph{Physics-Global, Selective-Local Wind Response for Static 3D Gaussian Thin Surfaces}.
  \item \emph{Conservative Reduced Wind Dynamics with Selective Local Refinement for Static 3DGS}.
\end{enumerate}

\subsection{안전한 novelty 문장}

\begin{quote}
We introduce a setup-conditioned, topology-distilled physical scaffold for static
Gaussian thin surfaces and a selective complementary Local dynamics scheme with measured quality--latency validation.
The method is evaluated by operator/eigenspace preservation and matched
quality--latency gains over strong Global and same-scaffold solver baselines;
conservative assembly, delta feedback, and Gaussian transport close the system loop.
\end{quote}

조사 범위에서 동일한 전체 input/output/coupling chain을 찾지 못한 것은 관련 연구의 경계를 정하지만,
그 부재만으로 novelty가 증명되지는 않는다.
따라서 본문 주장은 exact chain의 최초성보다 MC1의 operator/eigenspace 보존과
MC2의 strong-baseline quality--latency 우위에 둔다.
개별 modal equation, drag law, residual, adaptive subspace와 Gaussian transport의 최초성은 주장하지 않는다.

\subsection{권장 contribution 문장}

\begin{enumerate}[leftmargin=1.7em]
  \item A setup-conditioned, topology-distilled reduced physical scaffold for
  independently reconstructed static Gaussian thin surfaces, without explicit
  target triangle-mesh reconstruction.
  \item An error-triggered complementary Local dynamics scheme that selectively
  restores high-frequency wind response beyond the always-on Global ROM.
  The selector may be a validated deterministic cached-compliance proxy;
  a learned future-benefit gate is conditional, not a prerequisite.
  Learned missing force is an independently conditional realization.
  \item A closed-loop geometry-to-load update with conservative patch-force assembly,
  corrected-minus-predictor structural/aerodynamic feedback, and coherent,
  directly renderable Gaussian transport.
\end{enumerate}

Aerodynamic hyper-reduction is reported only as a supporting module if full-anchor
profiling and localized-gust tests justify it.

\subsection{피해야 할 표현}

\begin{itemize}[leftmargin=1.5em]
  \item First global--local Gaussian dynamics.
  \item First modal wind simulation for 3DGS.
  \item First Gaussian aerodynamic force.
  \item First adaptive local physics.
  \item First physics + neural residual 3DGS.
  \item Full fluid--structure interaction.
  \item Runtime cost independent of Gaussian count.
  \item Without runtime physics/continuum simulation.
  \item Interactive 또는 real-time system, target hardware, batch, resolution과
  renderer를 포함한 end-to-end p95가 측정되기 전.
\end{itemize}

필요하면 \emph{without runtime full-order continuum simulation}이라고 한정한다.
이는 runtime full-space shell/FEM/MPM/grid solve를 뜻한다.
본 방법은 매 frame reduced Global/Local physical dynamics를 실제로 적분한다.

\reviewdefer{\emph{Directly renderable}은 최종 Gaussian state가 standard renderer에 직접 연결된다는
I/O 계약일 뿐 interactive 또는 real-time 성능 주장이 아니다.
Renderer를 포함한 지정 hardware/batch/resolution의 end-to-end p95를 측정하기 전에는
\emph{interactive system}을 contribution 문장에 넣지 않는다.
CGF/TVCG 가능성 평가는 개발 목표 후보로만 보존하고 방법의 타당성이나 신규성 증거로 사용하지 않는다.
Gate A/B/C, real static-GS transfer와 end-to-end visual/runtime 결과가 나온 뒤에만 venue scope를 다시 판정한다.}

\reviewvthreedefer{V3 리뷰 \S4.3과 \S12.1의 \emph{interactive system} 문구와
``MC1/MC2 통과 시 CGF/TVCG 근거가 충분하다''는 venue 판정은 같은 이유로 보류한다.
또한 executive summary의 ``without runtime continuum solver''는 reduced physics까지 없는 것으로
오독될 수 있으므로 사용하지 않고, 위에서 정의한
\emph{without runtime full-order continuum simulation}만 사용한다.}

\reviewvfivedefer{V5 리뷰의 ``대응 적절, architecture freeze, Step 0 구현 착수''는
구현 단계 판정으로 수용한다. 그러나 CGF/TVCG 가능성은 외부 문헌 최신성 검증, 실제 코드,
teacher와 profiler 결과 없이 확정하지 않는다. 문서 길이와 structural backend의 복잡성,
또는 조사 범위에서 exact chain을 찾지 못했다는 사실도 venue나 novelty의 증거가 아니다.
MC1의 independent-GS operator/eigenspace/action 보존과 MC2의 strong-baseline
non-dominated quality--latency Pareto가 나온 뒤에만 venue scope와 최종 제목을 다시 판정한다.}

\section{최종 method 요약}

Training에서는 source thin-surface meshes와 high-resolution structural simulations를
privileged teacher로 사용한다.
Multi-view rendering에서 독립적으로 복원한 static 3DGS를 입력으로,
oriented material anchor graph, conservative area/mass,
normalized metric-MLS relation stencil과 setup에서 지정한 material/attachment로
식~\eqref{eq:exact-positive-stiffness-assembly}의 rest-linear membrane/bending reference와
식~\eqref{eq:fixed-vector-projective-hessian}의 objective normal-curvature projective runtime operator,
Global modes, patch-local complement modes와 Gaussian binding을 만든다.
Normalized area share와 별도로 calibrated absolute relation confidence를 저장하며,
low-evidence package는 support 확대/fallback/rejection을 통과해야 한다.
첫 논문의 정량 core는 flat/shallow-rest strip, flag, pennant와 curtain의
asset-uniform isotropic \(\nu=0\) family다. Nonzero Poisson ratio, anisotropy,
pre-curved leaf/shell은 이 V0 operator의 정량 일반화 claim이 아니라 qualitative/후속 범위다.
Runtime scaffold는 sparse point/edge/local-relation graph, tangent frame과 area weight를 가지지만
global face connectivity, watertightness 또는 manifold triangle mesh를 요구하지 않는다.

Target runtime에서는 이전 Global/Local state의 causal half-step predictor에서 anchor 위치, 속도,
normal, area를 복원하고, half-step에 sampling한 prescribed spatial wind와
frozen endpoint structural velocity의 relative velocity로
quasi-steady aerodynamic sample force를 계산한다.
구조 dynamics는 식~\eqref{eq:v0-structural-freeze-signature}의
one-shot objective normal-curvature projective frozen target, projected descriptor-rate damping과
implicit midpoint를 사용하며, Step 0에서 pure-rigid 및 strained-state superposed-rotation
objectivity, spectrum과 cost tolerance를 통과해야 한다.
Homogeneous structural update와 달리 이 one-pass exogenous-load/coupling predictor는
자동으로 2차 정확하다고 주장하지 않으며 true-midpoint/iterated reference와 오차를 보고한다.
V0는 모든 physical anchor에서 힘을 계산해 fixed sample lifting의
\(Q_g=\sum_s\Phi_s^\top f_s\)로 Global generalized load를 직접 투영하고,
항상 켜진 Global mass--stiffness--damping system을 적분하여 큰 sway와 bending을 예측한다.
Cubature는 이 단계가 실제 병목이고 held-out gust에서 generalized-load Pareto를 개선할 때만 사용한다.

V0 selector는 이미 계산한 free-DOF/owner-resolved force를 patch/group complement로 투영하고,
pre-factorized canonical normal-curvature-projective Local Hessian의 owner-mass-normalized dynamic compliance로
fixed-cardinality \(K_A\)개를 선택한다. 이 every-frame 경로는 auxiliary rollout이나 두 번째
aero/internal-force 평가를 요구하지 않는다. Full Global-only counter-state residual은
추가 비용을 포함한 offline/low-rate analysis baseline으로만 사용한다.
Normalized deterministic hybrid는 pure compliance의 blind spot과 feature 비용을 확인하는
조건부 ablation이며, measured-p95 Pareto를 개선할 때만 runtime 후보가 된다.
선택된 patch의 analytic Local excitation을 먼저 anchor force로 조립하고
같은 unique-anchor support에서 Global component를 route한다. External channel은 explicit
feasible Local remainder를 직접 사용하고, structural/hard-reaction 또는 추가 constraint channel만
augmented small-Schur로 conservation/complement를 동시에 만족시킨다. 이어 inter-patch
\(M/K/C\) cross block을 포함한 하나의 fixed-superset
Local complement system에서 관성, restoring, damping과 함께 적분한다.
Hard support reaction은 Local excitation과 분리한 ledger에서, compliant handle reaction은
attachment owner에서 sign-resolved world wrench/virtual work로 검사한다.
선택 해제 뒤에도 남은 Local energy는 같은 물리식으로 decay한다.
Oracle benefit은 selector upper bound로 사용한다.
조건부 learned gate는 같은 fixed cardinality \(K_A\)에서 score만 예측하고,
조건부 missing-force network는 검증된 nodal force defect만 더하며 위치를 직접 출력하지 않는다.

Local correction은 렌더링에만 더해지지 않는다.
수정된 surface normal, area, velocity에서 aerodynamic force를 다시 계산하고,
predictor 대비 aerodynamic delta와
actual-midpoint coupling에서 predictor-consumed coupling을 뺀 structural cross delta만
작은 Global corrector에 되돌린다.
최종 \(q,z\)에서 anchor deformation을 한 번만 복원하고,
affine MLS로 Gaussian mean과 covariance를 운반하여 standard 3DGS로 렌더링한다.
Dormant patch의 Local force/solve block은 실제 skip한다.
조건부 F가 유지된 경우에는 expert call도 함께 skip하지만,
이미 남은 Local energy는 물리적인 restoring/damping으로 decay한 뒤 종료한다.

\section{최종 연구 가치 판정}

이 방향의 강점은 Global/Local이라는 고전적 계층화 자체가 아니다.
각 레벨이 무엇을 입력으로 받고 무엇을 출력하는지,
어떤 force를 known physics가 소유하고 조건부 network가 어떤 defect만 학습할 수 있는지,
patch overlap과 Global complement를 어떻게 보존적으로 묶는지,
Local geometry 변화가 어떻게 aerodynamic load를 통해 다시 Global에 들어가는지를
wind-driven static 3DGS라는 target domain에 맞게 하나의 operator chain으로 정의한다는 점이다.

현재 가장 강한 novelty 후보는 조건부 learned gate가 아니라
selective complementary execution을 포함한 두 algorithmic contribution이다.
세 번째 항목은 이를 완결하는 system contribution과 scaling evidence다.

\begin{equation}
\boxed{
\begin{aligned}
&\text{MC1: setup-conditioned topology-distilled physical scaffold}\\
+\;&\text{MC2: error-indicator-triggered selective complementary Local execution}\\
+\;&\text{SYS: joint assembly, delta feedback, coherent transport, and scaling evidence}.
\end{aligned}
}
\end{equation}

그러나 exact chain이 선행에 없다는 사실만으로 충분하지 않다.
Gaussian Swaying, Wind Projection Basis, FreeForm/DynamicTree,
PGRD, Subspace Condensation/Trading Spaces의 단순 조합처럼 보이지 않으려면
E1--E10에서 topology operator, same-anchor speed/quality kill test,
selective Local physics, complement/feedback 각각의 비자명한 필요성과
quality--latency Pareto를 반증 중심으로 보여야 한다.

\begingroup
\sloppy
\bibliographystyle{plain}
\bibliography{refs_topology_distilled_error_triggered_wind_dynamics}
\endgroup

\end{document}
